\beginnumbering
\eledsection*{Adhyāya 1}
\pstart \ledsidenote{\passref{5.1.1}}\skt{अथा\-तोऽ\-न्न\-पा\-न\-र\-क्षा\-कल्पं }\edtext{\skt{व्याख्यास्यामः}}{\lemma{\skt{व्याख्यास्यामः}}\Afootnote{A H; \emph{om.}\ K}}\skt{~॥}\pend

\pstart \ledsidenote{\passref{5.1.2}}\skt{यथोवाच }\edtext{\skt{भग\-वा\-न्ध\-न्व\-न्तरिः}}{\lemma{\skt{धन्वन्तरिः}}\Afootnote{A K H\textsuperscript{ac}; \skt{धन्वन्तरिः~॥ अथ खलु वत्ससुश्रुतः~। १}\ H\textsuperscript{pc}}}\skt{~॥}\pend

\pstart\versepar \ledsidenote{\passref{5.1.3}}\edtext{\skt{दिवो\-दा\-स\-क्षि\-ति\-प\-ति\-स्त\-पो\-ध\-र्म\-भृतां}}{\lemma{\skt{दिवो\appbrk{}दा\appbrk{}स\appbrk{}क्षि\appbrk{}ति\appbrk{}प\appbrk{}तिस्}}\Afootnote{K H; \skt{धन्वन्तरिः काशिपतिस्}\ A}}\skt{ वरः~। }\pend
\pstart\versepar \skt{}\edtext{\edtext{\skt{सुश्रु\-त\-प्र\-मु\-खा\-ञ्छि\-ष्या\-ञ्श\-शा\-सा\-ह\-त\-शा\-सनः}}{\lemma{\skt{सुश्रु\appbrk{}त\appbrk{}प्र\appbrk{}मु\appbrk{}खाञ्}}\Afootnote{K H; \skt{सुश्रु\appbrk{}त\appbrk{}प्र\appbrk{}भृ\appbrk{}तीञ्}\ A}}}{\lemma{\skt{छिष्या\appbrk{}ञ्श\appbrk{}शा\appbrk{}सा\appbrk{}ह\appbrk{}त\appbrk{}शा\appbrk{}सनः}}\Afootnote{\skt{छिष्या\appbrk{}ञ्छ\appbrk{}शा\appbrk{}सा\appbrk{}ह\appbrk{}त\appbrk{}शा\appbrk{}सनः~॥}\ A; \skt{छिष्यां शशा\appbrk{}सा\appbrk{}ह\appbrk{}त\appbrk{}शा\appbrk{}सनः~।}\ K; \skt{छिष्या\appbrk{}न्श\appbrk{}शा\appbrk{}सा\appbrk{}ह\appbrk{}त\appbrk{}सा\appbrk{}तनः~॥}\ H\textsuperscript{ac}; \skt{छिष्या\appbrk{}न्श\appbrk{}शा\appbrk{}सा\appbrk{}हत शासनः~। २~॥}\ H\textsuperscript{pc}}}\skt{~॥}\pend

\pstart\versepar \ledsidenote{\passref{5.1.4}}\skt{रिपवो }\edtext{\skt{विक्रमाक्रान्ताः स्वे वा स्युः}}{\lemma{\skt{विक्रमाक्रान्ताः स्वे वा स्युः}}\Afootnote{K H; \skt{विक्रमाक्रान्ता ये च स्वे}\ A}}\skt{ कृत्यताङ्गताः~। }\pend
\pstart\versepar \skt{सिसृक्षवः क्रोधविषं विवरं प्राप्य तादृशम्~॥}\pend

\pstart \ledsidenote{\passref{5.1.5ab}}\edtext{\edtext{\skt{विषै\-र्हिं\-स्यु\-र\-कि\-ञ्चिज्ज्ञं}}{\lemma{\skt{हिंस्युर्}}\Afootnote{K H; \skt{निह}\ A}}}{\lemma{\skt{अकिञ्चिज्ज्ञं}}\Afootnote{\skt{न्युर्निपुणं}\ A; \skt{अकिञ्चिज्ञं}\ K H}}\skt{ नृपतिं }\edtext{\skt{दुष्टचेतसः}}{\lemma{\skt{दुष्टचेतसः}}\Afootnote{A K; \skt{दुष्टवेतसः~।}\ H}}\skt{~।}\pend

\pstart \ledsidenote{\passref{5.1.6cd}}\skt{तस्माद्वैद्येन }\edtext{\skt{सततं}}{\lemma{\skt{सततं}}\Afootnote{A K; \skt{शततं}\ H}}\skt{ }\edtext{\skt{विषाद्रक्ष्यो}}{\lemma{\skt{विषाद्रक्ष्यो}}\Afootnote{A K; \skt{विषाद्रक्षेन्}\ H}}\skt{ नराधिपः~॥}\pend

\pstart\versepar \ledsidenote{\passref{5.1.7}}\edtext{\skt{यस्मा\-च्चा\-नि\-त्य\-चि\-त्त\-त्व\-म\-श्व\-व\-त्प्र\-थितं}}{\lemma{\skt{यस्मा\appbrk{}च्चा\appbrk{}नि\appbrk{}त्य\appbrk{}चि\appbrk{}त्त\appbrk{}त्वम्}}\Afootnote{K H\textsuperscript{pc}; \skt{यस्माच्च चेतोऽनित्यत्वम्}\ A; \skt{यस्मा\appbrk{}च्चा\appbrk{}नि\appbrk{}त्य\appbrk{}चि\appbrk{}न्न\appbrk{}त्वम्}\ H\textsuperscript{ac}}}\skt{ }\edtext{\skt{नृषु}}{\lemma{\skt{नृषु}}\Afootnote{K H; \skt{नृ}\ A}}\skt{~। }\pend
\pstart\versepar \skt{}\edtext{\skt{तस्मान्न}}{\lemma{\skt{तस्मान्न}}\Afootnote{K H; \skt{णाम्~। न}\ A}}\skt{ }\edtext{\skt{विश्वसेद्राजा}}{\lemma{\skt{विश्वसेद्राजा}}\Afootnote{K H; \skt{विश्वस्यात्ततो राजा}\ A}}\skt{ कदाचिदपि कस्य चित्~॥}\pend

\pstart \ledsidenote{\passref{5.1.8ab}}\skt{कुलीनं धार्मिकं }\edtext{\skt{स्निग्धमकृशं}}{\lemma{\skt{अकृशं}}\Afootnote{K H; \skt{सुभृतं}\ A}}\skt{ }\edtext{\skt{सततोत्थितम्}}{\lemma{\skt{सततोत्थितम्}}\Afootnote{K H; \skt{संततोत्थितम्~।}\ A}}\skt{~।}\pend

\pstart \ledsidenote{\passref{5.1.11cd}}\skt{महानसे }\edtext{\skt{नियुञ्जीत}}{\lemma{\skt{नियुञ्जीत}}\Afootnote{K H; \skt{प्रयुञ्जीत}\ A}}\skt{ वैद्य\-न्त\-द्वि\-द्य\-पू\-जि\-तम्~॥}\pend

\pstart \ledsidenote{\passref{5.1.12ab}}\skt{प्रश\-स्त\-दि\-ग्दे\-श\-कृतं }\edtext{\edtext{\skt{शुचि\-भा\-ण्ड\-म्म\-ह\-च्छुचिः}}{\lemma{\skt{शुचिभाण्डम्}}\Afootnote{A; \skt{शुचिभाण्डां}\ K; \skt{सूचिभाण्डम्}\ H}}}{\lemma{\skt{महच्छुचिः}}\Afootnote{K H; \skt{महच्छुचि~।}\ A}}\skt{~।}\pend

\pstart \ledsidenote{\passref{5.1.13cd}}\skt{परी\-क्षि\-त\-स्त्री\-पु\-रुषं भवेच्चापि महानसम्~॥}\pend

\pstart \ledsidenote{\passref{5.1.17cd}}\edtext{\skt{महा\-न\-सि\-क\-वो\-ढारः}}{\lemma{\skt{महा\appbrk{}न\appbrk{}सि\appbrk{}क\appbrk{}वो\appbrk{}ढारः}}\Afootnote{H; \skt{माहा\appbrk{}न\appbrk{}सि\appbrk{}क\appbrk{}वो\appbrk{}ढारः}\ A; \skt{महा\appbrk{}न\appbrk{}सि\appbrk{}क\appbrk{}वो\appbrk{}ढ्ढारः}\ K}}\skt{ }\edtext{\skt{सौपो\-द\-नि\-क\-पू\-पिकाः}}{\lemma{\skt{सौपो\appbrk{}द\appbrk{}नि\appbrk{}क\appbrk{}पू\appbrk{}पिकाः}}\Afootnote{K; \skt{सौपौ\appbrk{}द\appbrk{}नि\appbrk{}क\appbrk{}पौ\appbrk{}पिकाः~॥}\ A; \skt{सूपो\appbrk{}द\appbrk{}नि\appbrk{}क\appbrk{}पू\appbrk{}पिकाः~।}\ H}}\skt{~॥}\pend

\pstart\versepar \ledsidenote{\passref{5.1.18}}\skt{भवे\-यु\-र्वै\-द्य\-व\-शगा ये }\edtext{\skt{चाप्यन्येऽत्र}}{\lemma{\skt{ऽत्र}}\Afootnote{A; \skt{त्र}\ K H}}\skt{ केचन~। }\pend
\pstart\versepar \skt{इङ्गितज्ञो मनुष्याणां वाक्चे\-ष्टा\-मु\-ख\-वै\-कृतैः~॥}\pend

\pstart\versepar \ledsidenote{\passref{5.1.19}}\edtext{\skt{जानी\-या\-द्वि\-ष\-दा\-ता\-र\-मे\-भि\-र्लि\-ङ्गैश्च}}{\lemma{\skt{जानी\appbrk{}या\appbrk{}द्वि\appbrk{}ष\appbrk{}दा\appbrk{}ता\appbrk{}रम्}}\Afootnote{K H; \skt{विद्याद्विषस्य दातारम्}\ A}}\skt{ बुद्धिमान्~। }\pend
\pstart\versepar \skt{}\edtext{\skt{विवक्षुर्मुह्यते}}{\lemma{\skt{मुह्यते}}\Afootnote{K H; \emph{om.}\ A}}\skt{ }\edtext{\skt{पृष्टो}}{\lemma{\skt{पृष्टो}}\Afootnote{K H; \emph{om.}\ A}}\skt{ }\edtext{\skt{नोत्तरं}}{\lemma{\skt{नोत्तरं}}\Afootnote{K H; \skt{न्मोहम्}\ A}}\skt{ }\edtext{\skt{प्रतिपद्यते}}{\lemma{\skt{प्रतिपद्यते}}\Afootnote{K H; \skt{एति च~॥}\ A}}\skt{~॥}\edtext{}{\lemma{\skt{बुद्धिमान्~। विवक्षुर्}}\Afootnote{K; \skt{बुद्धिमान्~। न ददात्युत्तरं पृष्टो विवक्ष}\ A; \skt{वुद्धिमान्~। विचक्षुर्}\ H}}\pend

\pstart\versepar \ledsidenote{\passref{5.1.20}}\skt{अपार्थं }\edtext{\skt{बहुसङ्कीर्णं}}{\lemma{\skt{बहुसङ्कीर्णं}}\Afootnote{A K; \skt{वहुसङ्कीर्ण्णम्}\ H}}\skt{ भाषते चापि मूढवत्~। }\pend
\pstart\versepar \skt{}\edtext{\skt{हस\-त्य\-क\-स्मा\-द\-ङ्गुलीः}}{\lemma{\skt{अकस्मादङ्गुलीः}}\Afootnote{K H; \skt{अकस्मा}\ A}}\skt{ }\edtext{\skt{स्फोट\-ये\-द्वि\-लि\-खे\-द्म\-हीम्}}{\lemma{\skt{स्फोट\appbrk{}ये\appbrk{}द्वि\appbrk{}लि\appbrk{}खे\appbrk{}द्म\appbrk{}हीम्}}\Afootnote{\skt{द्विलिखेद्धसेत्~॥}\ A; \skt{स्फोटा\appbrk{}ये\appbrk{}द्वि\appbrk{}लि\appbrk{}खे\appbrk{}द्महीं~।}\ K; \skt{स्फोट\appbrk{}ये\appbrk{}द्वि\appbrk{}लि\appbrk{}खे\appbrk{}न्महीं~।}\ H}}\skt{~॥}\edtext{}{\lemma{\skt{मूढवत्~। हसत्य्}}\Afootnote{K H; \skt{मूढवत्~। स्फोट\appbrk{}य\appbrk{}त्य\appbrk{}ङ्गु\appbrk{}ली\appbrk{}र्भू\appbrk{}मिम्}\ A}}\pend

\pstart\versepar \ledsidenote{\passref{5.1.21}}\edtext{\skt{वेपथुश्चास्य}}{\lemma{\skt{चास्य}}\Afootnote{K H; \skt{जाय}\ A}}\skt{ }\edtext{\skt{भवति}}{\lemma{\skt{भवति}}\Afootnote{K H; \skt{ते तस्य}\ A}}\skt{ }\edtext{\edtext{\skt{त्रस्त\-श्चा\-न्योऽ\-न्य\-मी\-क्षते}}{\lemma{\skt{ऽन्यम्}}\Afootnote{A; \skt{न्यम्}\ K H}}}{\lemma{\skt{ईक्षते}}\Afootnote{A K; \skt{ईक्ष्यते~।}\ H}}\skt{~। }\pend
\pstart\versepar \skt{}\edtext{\skt{विवर्णवक्त्रो}}{\lemma{\skt{विवर्णवक्त्रो}}\Afootnote{K H; \skt{क्षामो विवर्णवक्त्रश्}\ A}}\skt{ }\edtext{\skt{ध्यामश्च}}{\lemma{\skt{ध्यामश्च}}\Afootnote{K H; \skt{च}\ A}}\skt{ नखैः }\edtext{\skt{किञ्चिच्छिनत्यपि}}{\lemma{\skt{किञ्चिच्छिनत्य्}}\Afootnote{K H; \skt{किञ्चिच्छिनत्त्य्}\ A}}\skt{~॥}\pend

\pstart \ledsidenote{\passref{5.1.23}}\skt{वर्तते }\edtext{\skt{विपरीतश्च}}{\lemma{\skt{विपरीतश्च}}\Afootnote{K; \skt{विपरीतं तु}\ A; \skt{विपरीतञ्च}\ H}}\skt{ विषदाता }\edtext{\skt{विचेतनः}}{\lemma{\skt{विचेतनः}}\Afootnote{K H; \skt{विचेतनः~। केचि\appbrk{}द्भ\appbrk{}या\appbrk{}त्पा\appbrk{}र्थि\appbrk{}वस्य त्वरिता वा तदाज्ञया~॥}\ A}}\skt{~।}\pend

\pstart\versepar \ledsidenote{\passref{5.1.25}}\edtext{\edtext{\skt{दन्त\-का\-ष्ठेऽ\-न्न\-पाने}}{\lemma{\skt{दन्तकाष्ठे}}\Afootnote{\skt{अन्ने पाने दन्तकाष्ठे}\ A; \skt{दन्तकाष्ठे न्}\ K H}}\skt{ च}}{\lemma{\skt{ऽन्नपाने च}}\Afootnote{\emph{om.}\ A; \skt{नपाने च}\ K H}}\skt{ }\edtext{\edtext{\skt{तथा\-भ्य\-ङ्गेऽ\-व\-ले\-खने}}{\lemma{\skt{तथाभ्यङ्गे}}\Afootnote{K H; \skt{तथाऽभ्यङ्गे}\ A}}}{\lemma{\skt{ऽवलेखने}}\Afootnote{A; \skt{वलेखने}\ K H}}\skt{~। }\pend
\pstart\versepar \skt{उत्सादने }\edtext{\skt{परीषेके}}{\lemma{\skt{परीषेके}}\Afootnote{K H; \skt{कषाये च परिषेके}\ A}}\skt{ }\edtext{\skt{कषाये}}{\lemma{\skt{कषाये}}\Afootnote{K; \skt{ऽ}\ A; \skt{कषायैः}\ H}}\skt{ }\edtext{\skt{सानुलेपने}}{\lemma{\skt{सानुलेपने}}\Afootnote{K H; \skt{नुलेपने~॥}\ A}}\skt{~॥}\pend

\pstart\versepar \ledsidenote{\passref{5.1.26}}\skt{स्रक्षु वस्त्रेषु शय्यासु कव\-चा\-भ\-र\-णेषु च~। }\pend
\pstart\versepar \skt{पादु\-का\-पा\-द\-पी\-ठेषु पृष्ठेषु गजवाजिनाम्~॥}\pend

\pstart\versepar \ledsidenote{\passref{5.1.27}}\edtext{\skt{विषजुष्टेषु}}{\lemma{\skt{विषजुष्टेषु}}\Afootnote{A K; \skt{विषदुष्टेषु}\ H}}\skt{ चान्येषु }\edtext{\skt{नस्य\-धू\-मा\-ञ्ज\-ना\-दिषु}}{\lemma{\skt{नस्य\appbrk{}धू\appbrk{}मा\appbrk{}ञ्ज\appbrk{}ना\appbrk{}दिषु}}\Afootnote{A; \skt{धूम\appbrk{}न\appbrk{}स्या\appbrk{}ञ्ज\appbrk{}ना\appbrk{}दिषु}\ K H}}\skt{~। }\pend
\pstart\versepar \skt{लक्षणानि प्रवक्ष्यामि }\edtext{\skt{चिकि\-त्सा\-ञ्चा\-प्य\-न\-न्त\-रम्}}{\lemma{\skt{चाप्य्}}\Afootnote{K H; \skt{अप्य्}\ A}}\skt{~॥}\pend

\pstart\versepar \ledsidenote{\passref{5.1.28}}\edtext{\edtext{\skt{नृप\-भ\-क्ता\-द्ब\-लि\-न्दत्तं}}{\lemma{\skt{नृप\appbrk{}भ\appbrk{}क्ता\appbrk{}द्ब\appbrk{}लिन्}}\Afootnote{A K; \skt{नृप\appbrk{}भ\appbrk{}क्ता\appbrk{}द्व\appbrk{}लिंन्}\ H}}}{\lemma{\skt{दत्तं}}\Afootnote{K H; \skt{न्यस्तं}\ A}}\skt{ सविषं भक्षयन्ति ये~। }\pend
\pstart\versepar \skt{तत्रैव ते विनश्यन्ति मक्षि\-का\-वा\-य\-सा\-दयः~।}\pend

\pstart\versepar \ledsidenote{\passref{5.1.29}}\skt{हुतभुक्तेन }\edtext{\skt{चान्नेन}}{\lemma{\skt{चान्नेन}}\Afootnote{A K\textsuperscript{pc} H; \skt{चान्ने}\ K\textsuperscript{ac}}}\skt{ भृश\-ञ्च\-ट\-च\-टा\-यते~। }\pend
\pstart\versepar \skt{मयू\-र\-क\-ण्ठ\-प्र\-तिमो जायते चापि दुःसहः~।}\pend

\pstart \ledsidenote{\passref{5.1.30cd}}\skt{चको\-र\-स्या\-क्षि\-वै\-राग्यं जायते क्षिप्रमेव तु~।}\pend

\pstart\versepar \ledsidenote{\passref{5.1.31}}\edtext{\skt{दृष्ट्वान्नं}}{\lemma{\skt{दृष्ट्वान्नं}}\Afootnote{K H; \skt{दृष्ट्वाऽन्नं}\ A}}\skt{ विषसंसृष्टं }\edtext{\skt{म्रियते}}{\lemma{\skt{म्रियते}}\Afootnote{K H; \skt{म्रियन्ते}\ A}}\skt{ }\edtext{\skt{जीवजीवकः}}{\lemma{\skt{जीवजीवकः}}\Afootnote{K H\textsuperscript{pc}; \skt{जीवजीवकाः~।}\ A; \skt{जीव\appbrk{}जी\appbrk{}व\appbrk{}जी\appbrk{}वकः~।}\ H\textsuperscript{ac}}}\skt{~। }\pend
\pstart\versepar \skt{कोकिलः स्वरवैकृत्यं }\edtext{\skt{क्रोञ्चस्तु}}{\lemma{\skt{क्रोञ्चस्तु}}\Afootnote{K; \skt{क्रौञ्चस्तु}\ A H}}\skt{ }\edtext{\skt{मदमर्च्छति}}{\lemma{\skt{अर्च्छति}}\Afootnote{K H; \skt{ऋच्छति~॥}\ A}}\skt{~।}\pend

\pstart\versepar \ledsidenote{\passref{5.1.32}}\edtext{\edtext{\skt{हृष्ये\-न्म\-यू\-र\-स्तू\-द्विग्ने}}{\lemma{\skt{हृष्येन्मयूरस्}}\Afootnote{K H NE-Bikaner\_RORI\_5157; \skt{हृष्येन्मयूर}\ A}}}{\lemma{\skt{तूद्विग्ने}}\Afootnote{K H NE-Bikaner\_RORI\_5157; \skt{उद्विग्नः}\ A}}\skt{ }\edtext{\skt{क्रोशेते}}{\lemma{\skt{क्रोशेते}}\Afootnote{K NE-Bikaner\_RORI\_5157; \skt{क्रोशतः}\ A; \skt{क्रोशैते}\ H}}\skt{ शुकसारिके~। }\pend
\pstart\versepar \skt{}\edtext{\skt{हंसः}}{\lemma{\skt{हंसः}}\Afootnote{A H NE-Bikaner\_RORI\_5157; \skt{हंस}\ K}}\skt{ }\edtext{\skt{क्ष्वेलति}}{\lemma{\skt{क्ष्वेलति}}\Afootnote{K; \skt{क्ष्वेडति}\ A NE-Bikaner\_RORI\_5157; \skt{खेलति}\ H}}\skt{ क्ष्वेडति }\edtext{\edtext{\skt{चात्यर्थं}}{\lemma{\skt{कूजते}}\Afootnote{H NE-Bikaner\_RORI\_5157; \emph{om.}\ A; \skt{कू…जते}\ K}}\skt{ कूजते भृङ्गराजकः}}{\lemma{\skt{भृङ्गराजकः}}\Afootnote{K H NE-Bikaner\_RORI\_5157; \skt{भृङ्गराजस्तु कूजति~॥}\ A}}\skt{~।}\pend

\pstart \ledsidenote{\passref{5.1.33ab}}\edtext{\skt{वृषभो}}{\lemma{\skt{वृषभो}}\Afootnote{H; \skt{पृषतो}\ A NE-Bikaner\_RORI\_5157; \skt{पृषभो}\ K}}\skt{ पृषतो }\edtext{\edtext{\skt{विसृजत्यस्रं}}{\lemma{\skt{अस्रं}}\Afootnote{H NE-Bikaner\_RORI\_5157; \skt{अश्रुं}\ A; \skt{अश्रं}\ K}}\skt{ }\edtext{\skt{मुञ्चते}}{\lemma{\skt{विट्च}}\Afootnote{K H NE-Bikaner\_RORI\_5157; \skt{ति}\ A}}}{\lemma{\skt{मुञ्चते}}\Afootnote{K H NE-Bikaner\_RORI\_5157; \skt{विष्ठां मुञ्च}\ A}}\skt{ विट्च मर्कटः~।}\pend

\pstart \ledsidenote{\passref{5.1.34cd}}\skt{उपक्षिप्तस्य चान्नस्य }\edtext{\edtext{\skt{बाष्पे\-णो\-र्ध्व\-मु\-दी\-यता}}{\lemma{\skt{बाष्पेणोर्ध्वम्}}\Afootnote{A; \skt{वाष्पेणोर्ध्वम्}\ K H}}}{\lemma{\skt{उदीयता}}\Afootnote{K; \skt{प्रसर्पता~॥}\ A; \skt{उदीर्यतां~।}\ H}}\skt{~।}\pend

\pstart\versepar \ledsidenote{\passref{5.1.35}}\skt{हृत्पीडा भ्रान्तनेत्रत्वं शिरोदुःखञ्च जायते~। }\pend
\pstart\versepar \skt{तत्र }\edtext{\skt{नस्याञ्जने}}{\lemma{\skt{नस्याञ्जने}}\Afootnote{A K; \skt{तस्याञ्जने}\ H}}\skt{ कुष्ठं लामज्जं नलदं }\edtext{\skt{मधुः}}{\lemma{\skt{मधुः}}\Afootnote{K; \skt{मधु~॥}\ A; \skt{मधूः।}\ H\textsuperscript{ac}; \skt{मधू।}\ H\textsuperscript{pc}}}\skt{~।}\pend

\pstart \ledsidenote{\passref{5.1.36cd}}\skt{हृदि }\edtext{\skt{चन्दनलेपश्च}}{\lemma{\skt{चन्दनलेपश्च}}\Afootnote{K H; \skt{चन्दनलेपस्तु}\ A}}\skt{ तथा सुख\-म\-वा\-प्नु\-यात्~।}\pend

\pstart\versepar \ledsidenote{\passref{5.1.37}}\skt{पाणिप्राप्तं पाणिदाहं }\edtext{\skt{नख\-शा\-त\-ङ्क\-रोति}}{\lemma{\skt{नखशातङ्}}\Afootnote{A K; \skt{नखसातङ्}\ H}}\skt{ च~। }\pend
\pstart\versepar \skt{}\edtext{\skt{तत्र}}{\lemma{\skt{तत्र}}\Afootnote{K H; \skt{अत्र}\ A}}\skt{ प्रलेपः }\edtext{\skt{श्यामे\-न्द्र\-गो\-प\-सो\-मो\-त्प\-लानि}}{\lemma{\skt{श्यामे\appbrk{}न्द्र\appbrk{}गो\appbrk{}प\appbrk{}सो\appbrk{}मो\appbrk{}त्प\appbrk{}लानि}}\Afootnote{K H; \skt{श्यामे\appbrk{}न्द्र\appbrk{}गो\appbrk{}पा\appbrk{}सो\appbrk{}मो\appbrk{}त्प\appbrk{}लानि}\ A}}\skt{ च~।}\pend

\pstart\versepar \ledsidenote{\passref{5.1.38}}\skt{स चेत्प्र\-मा\-दा\-न्मो\-हाद्वा }\edtext{\skt{तद्भुंक्ते}}{\lemma{\skt{तद्भुंक्ते}}\Afootnote{K H; \skt{तदन्}\ A}}\skt{ }\edtext{\skt{भोजनं}}{\lemma{\skt{भोजनं}}\Afootnote{K H; \skt{नम्}\ A}}\skt{ }\edtext{\skt{यदि}}{\lemma{\skt{यदि}}\Afootnote{K H; \skt{उपसे}\ A}}\skt{~। }\pend
\pstart\versepar \skt{}\edtext{\edtext{\skt{ततोऽ\-स्या\-ष्ठी\-ल\-व\-ज्जिह्वा}}{\lemma{\skt{ततो}}\Afootnote{K H; \skt{वते~।}\ A}}}{\lemma{\skt{ऽस्या\appbrk{}ष्ठी\appbrk{}ल\appbrk{}व\appbrk{}ज्जिह्वा}}\Afootnote{K H; \skt{अष्ठीलावत्ततो जिह्वा}\ A}}\skt{ }\edtext{\edtext{\skt{जाय\-तेऽ\-र\-स\-वे\-दनी}}{\lemma{\skt{जायते}}\Afootnote{K H; \skt{भवत्य्}\ A}}}{\lemma{\skt{ऽरसवेदनी}}\Afootnote{\skt{अरसवेदिनी~॥}\ A; \skt{रसवेदनी~।}\ K H}}\skt{~।}\pend

\pstart\versepar \ledsidenote{\passref{5.1.39}}\skt{तुद्यते दह्यते चापि श्लेष्मा }\edtext{\skt{चास्य}}{\lemma{\skt{चास्य}}\Afootnote{K H; \skt{चास्या}\ A}}\skt{ }\edtext{\skt{प्रसिच्यते}}{\lemma{\skt{प्रसिच्यते}}\Afootnote{K H; \skt{त्प्रसिच्यते~।}\ A}}\skt{~। }\pend
\pstart\versepar \skt{तत्र }\edtext{\skt{बाष्पेरितं}}{\lemma{\skt{बाष्पेरितं}}\Afootnote{A; \skt{वाष्पेरितं}\ K H}}\skt{ कर्म }\edtext{\skt{यच्च}}{\lemma{\skt{यच्च}}\Afootnote{A K H\textsuperscript{pc}; \skt{यच्चा}\ H\textsuperscript{ac}}}\skt{ }\edtext{\skt{स्याद्द\-न्त\-का\-ष्ठि\-कम्}}{\lemma{\skt{स्याद्द\appbrk{}न्त\appbrk{}का\appbrk{}ष्ठि\appbrk{}कम्}}\Afootnote{K H; \skt{स्याद्दा\appbrk{}न्त\appbrk{}का\appbrk{}ष्ठि\appbrk{}कम्~॥}\ A}}\skt{~।}\pend

\pstart\versepar \ledsidenote{\passref{5.1.40}}\edtext{\skt{मूर्छां}}{\lemma{\skt{मूर्छां}}\Afootnote{A H; \skt{मूर्च्छा}\ K}}\skt{ छर्दिं }\edtext{\skt{रोम\-ह\-र्ष\-मा\-ध्मानं}}{\lemma{\skt{रोमहर्षम्}}\Afootnote{K H; \skt{अतीसारम्}\ A}}\skt{ }\edtext{\edtext{\skt{दाहमेव}}{\lemma{\skt{दाहम्}}\Afootnote{K H; \skt{दाह}\ A}}\skt{ च}}{\lemma{\skt{एव च}}\Afootnote{K H; \skt{वेपथू~।}\ A}}\skt{~। }\pend
\pstart\versepar \skt{इन्द्रियाणाञ्च वैकृत्यं }\edtext{\skt{कुर्यादामाशयं}}{\lemma{\skt{कुर्यादामाशयं}}\Afootnote{A K; \skt{कुर्या\appbrk{}दा\appbrk{}मा\appbrk{}स\appbrk{}यङ्}\ H\textsuperscript{ac}; \skt{द्रवद्रव्ये शये गतं~।३कु\appbrk{}र्या\appbrk{}दा\appbrk{}मा\appbrk{}स\appbrk{}यङ्}\ H\textsuperscript{pc}}}\skt{ गतम्~।}\pend

\pstart\versepar \ledsidenote{\passref{5.1.41}}\skt{तत्राशु }\edtext{\skt{मद\-ना\-ला\-बु\-बि\-म्बी\-को\-शा\-त\-की\-फलैः}}{\lemma{\skt{मद\appbrk{}ना\appbrk{}ला\appbrk{}बु\appbrk{}बि\appbrk{}म्बी\appbrk{}को\appbrk{}शा\appbrk{}त\appbrk{}की\appbrk{}फलैः}}\Afootnote{A; \skt{मद\appbrk{}ना\appbrk{}ला\appbrk{}वु\appbrk{}वि\appbrk{}म्वी\appbrk{}को\appbrk{}शा\appbrk{}त\appbrk{}की\appbrk{}फलैः~।}\ K H}}\skt{~। }\pend
\pstart\versepar \skt{छर्दनं }\edtext{\skt{दध्यु\-द\-श्वि\-द्भ्या\-म\-थवा}}{\lemma{\skt{दध्युदश्विद्भ्याम्}}\Afootnote{A; \skt{दध्युदश्विभ्याम्}\ K H}}\skt{ }\edtext{\skt{तण्डुलाम्बुना}}{\lemma{\skt{तण्डुलाम्बुना}}\Afootnote{A; \skt{ताण्डुलाम्वुना~।}\ K; \skt{तण्डुलाम्वुना~।}\ H}}\skt{~।}\pend

\pstart\versepar \ledsidenote{\passref{5.1.42}}\skt{दाहं मूर्च्छामतीसारं तृष्णा\-मि\-न्द्रि\-य\-वै\-कृ\-तम्~। }\pend
\pstart\versepar \skt{आटोपं पाण्डुतां कार्श्यं }\edtext{\skt{कुर्यात्पक्वाशयं}}{\lemma{\skt{कुर्यात्पक्वाशयं}}\Afootnote{A K; \skt{कुर्यात्पक्वासयं}\ H}}\skt{ गतं}\pend

\pstart\versepar \ledsidenote{\passref{5.1.43}}\edtext{\skt{तत्र}}{\lemma{\skt{तत्र}}\Afootnote{K H; \emph{om.}\ A}}\skt{ }\edtext{\skt{नीलीफलं}}{\lemma{\skt{नीलीफलं}}\Afootnote{K H; \skt{विरेच}\ A}}\skt{ }\edtext{\skt{श्रेष्ठं}}{\lemma{\skt{श्रेष्ठं}}\Afootnote{K H; \skt{नं}\ A}}\skt{ ससर्पिष्कं }\edtext{\skt{विरेचनम्}}{\lemma{\skt{विरेचनम्}}\Afootnote{K H; \skt{तत्रोक्तं नीलिनीफलम्~।}\ A}}\skt{~। }\pend
\pstart\versepar \skt{दध्ना दूषीविषारिश्च }\edtext{\skt{पेयो}}{\lemma{\skt{पेयो}}\Afootnote{K H; \skt{पेयो वा}\ A}}\skt{ }\edtext{\skt{मधुसमायुतः}}{\lemma{\skt{मधुसमायुतः}}\Afootnote{K H; \skt{मधुसंयुतः~॥}\ A}}\skt{~।}\pend

\pstart\versepar \ledsidenote{\passref{5.1.44}}\skt{द्रवद्रव्येषु सर्वेषु क्षीर\-म\-द्यो\-द\-का\-दिषु~। }\pend
\pstart\versepar \skt{भवन्ति विविधा राज्यः }\edtext{\skt{फेन\-बु\-द्बु\-द\-जन्म}}{\lemma{\skt{फेन\appbrk{}बु\appbrk{}द्बु\appbrk{}द\appbrk{}जन्म}}\Afootnote{A K; \skt{फेन\appbrk{}वु\appbrk{}द्वु\appbrk{}द\appbrk{}जन्म}\ H}}\skt{ च~।}\pend

\pstart\versepar \ledsidenote{\passref{5.1.45}}\edtext{\skt{छायाश्चात्र}}{\lemma{\skt{छायाश्}}\Afootnote{A; \skt{छाया}\ K H}}\skt{ न दृश्यन्ते दृश्यन्ते यदि वा पुनः~। }\pend
\pstart\versepar \skt{भवन्ति }\edtext{\skt{विकृ\-ता\-श्छि\-द्रा\-स्तन्व्यो}}{\lemma{\skt{विकृताश्}}\Afootnote{\skt{यमलाश्}\ A; \skt{विकृता}\ K H\textsuperscript{ac}; \skt{विकृता यमला}\ H\textsuperscript{pc}}}\skt{ वा विकृतास्तथा~।}\pend

\pstart\versepar \ledsidenote{\passref{5.1.46}}\edtext{\skt{शाक\-सू\-पा\-न्न\-मां\-सानि}}{\lemma{\skt{शाक\appbrk{}सू\appbrk{}पा\appbrk{}न्न\appbrk{}मां\appbrk{}सानि}}\Afootnote{A H; \skt{शाक\appbrk{}सू\appbrk{}पां\appbrk{}ण\appbrk{}मां\appbrk{}सानि}\ K}}\skt{ क्लिन्नानि विरसानि च~। }\pend
\pstart\versepar \skt{सद्यः }\edtext{\skt{पर्युषितानीव}}{\lemma{\skt{पर्युषितानीव}}\Afootnote{A H; \skt{पर्युशितानीव}\ K}}\skt{ }\edtext{\skt{विगन्धीनि}}{\lemma{\skt{विगन्धीनि}}\Afootnote{K H; \skt{विगन्धानि}\ A}}\skt{ भवन्ति च~।}\pend

\pstart\versepar \ledsidenote{\passref{5.1.47}}\skt{गन्ध\-व\-र्ण्ण\-र\-सै\-र्हीनाः सर्वे भक्ष्याः फलानि च~। }\pend
\pstart\versepar \skt{}\edtext{\skt{पक्वान्याशु}}{\lemma{\skt{आशु}}\Afootnote{A; \skt{आशुः}\ K H}}\skt{ }\edtext{\skt{प्रकुथ्यन्ते}}{\lemma{\skt{प्रकुथ्यन्ते}}\Afootnote{K H; \skt{विशीर्यन्ते}\ A}}\skt{ पाकमामानि यान्ति च~।}\pend

\pstart\versepar \ledsidenote{\passref{5.1.48}}\edtext{\skt{विशीर्यन्ते}}{\lemma{\skt{विशीर्यन्ते}}\Afootnote{K; \skt{विशीर्यते}\ A; \skt{विसीर्यन्ते}\ H}}\skt{ कूर्चकस्तु दन्तकाष्ठगते विषे~। }\pend
\pstart\versepar \skt{}\edtext{\skt{जिह्वा\-द\-न्तौ\-ष्ठ\-मां\-सेषु}}{\lemma{\skt{जिह्वा\appbrk{}द\appbrk{}न्तौ\appbrk{}ष्ठ\appbrk{}मां\appbrk{}सेषु}}\Afootnote{K; \skt{जिह्वा\appbrk{}द\appbrk{}न्तौ\appbrk{}ष्ठ\appbrk{}मां\appbrk{}सानां}\ A; \skt{जिह्वा\appbrk{}द\appbrk{}न्तो\appbrk{}ष्ठ\appbrk{}मां\appbrk{}सेषु}\ H}}\skt{ }\edtext{\skt{श्वय\-थु\-श्चो\-प\-जा\-यते}}{\lemma{\skt{श्वयथुश्}}\Afootnote{A K; \skt{स्वयथुश्}\ H}}\skt{~।}\pend

\pstart\versepar \ledsidenote{\passref{5.1.49}}\skt{अथास्य }\edtext{\skt{धात\-की\-पु\-ष्प\-ज\-म्ब्वा\-म्रा\-स्थि\-ह\-री\-तकैः}}{\lemma{\skt{धात\appbrk{}की\appbrk{}पु\appbrk{}ष्प\appbrk{}ज\appbrk{}म्ब्वा\appbrk{}म्रा\appbrk{}स्थि\appbrk{}ह\appbrk{}री\appbrk{}तकैः}}\Afootnote{\skt{धात\appbrk{}की\appbrk{}पु\appbrk{}ष्प\appbrk{}प\appbrk{}थ्या\appbrk{}ज\appbrk{}म्बू\appbrk{}फ\appbrk{}ला\appbrk{}स्थिभिः~।}\ A; \skt{धात\appbrk{}की\appbrk{}पु\appbrk{}ष्प\appbrk{}ज\appbrk{}म्व्वा\appbrk{}म्रा\appbrk{}स्थि\appbrk{}ह\appbrk{}री\appbrk{}ककैः~।}\ K; \skt{धात\appbrk{}की\appbrk{}पु\appbrk{}ष्प\appbrk{}ज\appbrk{}म्व्वा\appbrk{}म्रा\appbrk{}स्थि\appbrk{}ह\appbrk{}री\appbrk{}तकैः~।}\ H}}\skt{~। }\pend
\pstart\versepar \skt{सक्षौद्रैः प्रच्छिते शोफे कर्तव्यं प्रतिसारणम्~।}\pend

\pstart\versepar \ledsidenote{\passref{5.1.50}}\edtext{\skt{अथ\-वा\-ङ्को\-ल्ल\-मू\-लानि}}{\lemma{\skt{अथ\appbrk{}वा\appbrk{}ङ्को\appbrk{}ल्ल\appbrk{}मू\appbrk{}लानि}}\Afootnote{K H; \skt{अथा\appbrk{}वा\appbrk{}ङ्को\appbrk{}ठ\appbrk{}मू\appbrk{}लानि}\ A}}\skt{ त्वचः सप्तछदस्य वा~। }\pend
\pstart\versepar \skt{शिरीषमाषका }\edtext{\skt{वापि}}{\lemma{\skt{वापि}}\Afootnote{\skt{वाऽपि}\ A; \skt{द्वापि}\ K H}}\skt{ }\edtext{\skt{कर्तव्यं}}{\lemma{\skt{कर्तव्यं}}\Afootnote{K H; \skt{सक्षौद्राः}\ A}}\skt{ प्रतिसारणम्~।}\pend

\pstart\versepar \ledsidenote{\passref{5.1.51}}\skt{जिह्वा\-नि\-र्ले\-ख\-क\-वलौ दन्त\-का\-ष्ठ\-व\-दा\-दि\-शेत्~। }\pend
\pstart\versepar \skt{पिच्छलो }\edtext{\edtext{\skt{बहलोऽभ्यङ्गो}}{\lemma{\skt{बहलो}}\Afootnote{\skt{बहुलो}\ A; \skt{वहलो}\ K H}}}{\lemma{\skt{ऽभ्यङ्गो}}\Afootnote{A; \skt{भ्यङ्गो}\ K H}}\skt{ }\edtext{\skt{विवर्णश्च}}{\lemma{\skt{विवर्णश्च}}\Afootnote{K H; \skt{विवर्णो वा}\ A}}\skt{ विषान्वितः~।}\edtext{}{\lemma{\skt{दन्त\appbrk{}का\appbrk{}ष्ठ\appbrk{}व\appbrk{}दा\appbrk{}दि\appbrk{}शेत्~। पिच्छलो}}\Afootnote{K; \skt{दन्त\appbrk{}का\appbrk{}ष्ठ\appbrk{}व\appbrk{}दा\appbrk{}दि\appbrk{}शेत्~। पिच्छिलो}\ A H}}\pend

\pstart\versepar \ledsidenote{\passref{5.1.52}}\edtext{\skt{स्फोटा\-ज\-न्म\-रु\-जा\-स्रा\-व\-स्त्व\-क्पाकः}}{\lemma{\skt{स्फोटा\appbrk{}ज\appbrk{}न्म\appbrk{}रु\appbrk{}जा\appbrk{}स्रा\appbrk{}वस्}}\Afootnote{K; \skt{स्फोट\appbrk{}ज\appbrk{}न्म\appbrk{}रु\appbrk{}जा\appbrk{}स्राव}\ A; \skt{स्फोटा\appbrk{}ज\appbrk{}न्म\appbrk{}रु\appbrk{}जा\appbrk{}स्रा\appbrk{}वास्}\ H}}\skt{ }\edtext{\skt{स्वेद}}{\lemma{\skt{स्वेद}}\Afootnote{A H; \skt{स्वेदम्}\ K}}\skt{ }\edtext{\skt{एव च}}{\lemma{\skt{एव च}}\Afootnote{K H; \skt{नं ज्वरः~।}\ A}}\skt{~। }\pend
\pstart\versepar \skt{}\edtext{\skt{दारणञ्चापि}}{\lemma{\skt{दारणञ्}}\Afootnote{K; \skt{दरणं}\ A; \skt{दारुणञ्}\ H}}\skt{ मांसानामभ्यङ्गे }\edtext{\skt{विषदूषिते}}{\lemma{\skt{विषदूषिते}}\Afootnote{K H; \skt{विषसंयुते~॥}\ A}}\skt{~।}\pend

\pstart\versepar \ledsidenote{\passref{5.1.53}}\skt{तत्र }\edtext{\skt{शीताम्बुसिक्तस्य}}{\lemma{\skt{शीताम्बुसिक्तस्य}}\Afootnote{A; \skt{सीताम्वुसिक्तस्य}\ K; \skt{शीताम्वुसिक्तस्य}\ H}}\skt{ कर्त\-व्य\-म\-नु\-ले\-प\-नम्~। }\pend
\pstart\versepar \skt{चन्दनन्तगरं कुष्ठमुशीरं वेणुपत्रिका~।}\pend

\pstart\versepar \ledsidenote{\passref{5.1.54}}\edtext{\skt{सोमवल्यमृता}}{\lemma{\skt{सोमवल्यमृता}}\Afootnote{K H; \skt{सोमवल्ल्यमृता}\ A}}\skt{ }\edtext{\skt{श्वेता}}{\lemma{\skt{श्वेता}}\Afootnote{A K; \skt{स्वेता}\ H}}\skt{ पद्मं }\edtext{\skt{कालीयकन्तथा}}{\lemma{\skt{तथा}}\Afootnote{K H; \skt{त्वचम्~।}\ A}}\skt{~। }\pend
\pstart\versepar \skt{}\edtext{\skt{कपि\-त्थ\-र\-स\-प\-त्राभ्यां}}{\lemma{\skt{कपि\appbrk{}त्थ\appbrk{}र\appbrk{}स\appbrk{}प\appbrk{}त्राभ्यां}}\Afootnote{K H; \skt{कपि\appbrk{}त्थ\appbrk{}र\appbrk{}स\appbrk{}मू\appbrk{}त्राभ्यां}\ A}}\skt{ }\edtext{\edtext{\skt{पानमेतच्च}}{\lemma{\skt{पानम्}}\Afootnote{A H; \skt{पाण}\ K\textsuperscript{ac}; \skt{पाणम्}\ K\textsuperscript{pc}}}}{\lemma{\skt{एतच्च}}\Afootnote{A K\textsuperscript{pc} H; \skt{तच्च}\ K\textsuperscript{ac}}}\skt{ }\edtext{\skt{पूज्यते}}{\lemma{\skt{पूज्यते}}\Afootnote{K H; \skt{युज्यते~॥}\ A}}\skt{~।}\pend

\pstart\versepar \ledsidenote{\passref{5.1.55}}\skt{उत्सादने परीषेके कषाये }\edtext{\skt{सानुलेपने}}{\lemma{\skt{सानुलेपने}}\Afootnote{K H; \skt{चानुलेपने~।}\ A}}\skt{~। }\pend
\pstart\versepar \skt{}\edtext{\skt{शय्या\-व\-स्त्र\-त\-नु\-त्रेषु}}{\lemma{\skt{शय्या\appbrk{}व\appbrk{}स्त्र\appbrk{}त\appbrk{}नु\appbrk{}त्रेषु}}\Afootnote{A K; \skt{शय्या\appbrk{}व\appbrk{}स्त्र\appbrk{}त\appbrk{}नु\appbrk{}त्वेषु}\ H}}\skt{ }\edtext{\skt{विद्या\-द\-भ्य\-ङ्ग\-व\-द्भि\-षक्}}{\lemma{\skt{विद्या\appbrk{}द\appbrk{}भ्य\appbrk{}ङ्ग\appbrk{}व\appbrk{}द्भि\appbrk{}षक्}}\Afootnote{K H; \skt{ज्ञेय\appbrk{}म\appbrk{}भ्य\appbrk{}ङ्ग\appbrk{}ल\appbrk{}क्षणैः~॥}\ A}}\skt{~।}\pend

\pstart\versepar \ledsidenote{\passref{5.1.56}}\edtext{\skt{केशशातः}}{\lemma{\skt{केशशातः}}\Afootnote{A K; \skt{केशसातः}\ H}}\skt{ शिरोदुःखं खेभ्यश्च रुधिरागमः~। }\pend
\pstart\versepar \skt{}\edtext{\skt{ग्रन्थि\-ज\-न्मो\-त्त\-माङ्गे च}}{\lemma{\skt{ग्रन्थि\appbrk{}ज\appbrk{}न्मो\appbrk{}त्त\appbrk{}माङ्गे च}}\Afootnote{K H; \skt{ग्रन्थि\appbrk{}ज\appbrk{}न्मो\appbrk{}त्त\appbrk{}मा\appbrk{}ङ्गेषु}\ A}}\skt{ }\edtext{\skt{विष\-जु\-ष्टेऽ\-व\-ले\-खने}}{\lemma{\skt{ऽवलेखने}}\Afootnote{A; \skt{वलेखने~।}\ K H}}\skt{~।}\pend

\pstart\versepar \ledsidenote{\passref{5.1.57}}\edtext{\skt{तत्र}}{\lemma{\skt{तत्र}}\Afootnote{K H; \emph{om.}\ A}}\skt{ प्रलेपो }\edtext{\skt{बहुशो}}{\lemma{\skt{बहुशो}}\Afootnote{A; \skt{वहुशो}\ K H}}\skt{ }\edtext{\skt{भाविता}}{\lemma{\skt{भाविता}}\Afootnote{K H; \skt{तत्र भविताः}\ A}}\skt{ }\edtext{\skt{कृष्णमृत्तिका}}{\lemma{\skt{कृष्णमृत्तिका}}\Afootnote{K H; \skt{कृष्णमृत्तिकाः~।}\ A}}\skt{~। }\pend
\pstart\versepar \skt{}\edtext{\skt{ऋक्ष\-पि\-त्त\-घृ\-त\-श्या\-मा\-पा\-लि\-न्दी\-त\-ण्डु\-ली\-यकैः}}{\lemma{\skt{ऋक्ष\appbrk{}पि\appbrk{}त्त\appbrk{}घृ\appbrk{}त\appbrk{}श्या\appbrk{}मा\appbrk{}पा\appbrk{}लि\appbrk{}न्दी\appbrk{}त\appbrk{}ण्डु\appbrk{}ली\appbrk{}यकैः}}\Afootnote{K H; \skt{ऋष्य\appbrk{}पि\appbrk{}त्त\appbrk{}घृ\appbrk{}त\appbrk{}श्या\appbrk{}मा\appbrk{}पा\appbrk{}लि\appbrk{}न्दी\appbrk{}त\appbrk{}ण्डु\appbrk{}ली\appbrk{}यकैः~॥}\ A}}\skt{~।}\pend

\pstart\versepar \ledsidenote{\passref{5.1.58}}\skt{गोमयस्वरसो }\edtext{\skt{वापि}}{\lemma{\skt{वापि}}\Afootnote{K H; \skt{वाऽपि}\ A}}\skt{ हितो वा मालतीरसः~। }\pend
\pstart\versepar \skt{रसो }\edtext{\skt{मूषिककर्ण्या वा}}{\lemma{\skt{मूषिककर्ण्या वा}}\Afootnote{K H; \skt{मूषिकपर्ण्या वा}\ A}}\skt{ धूमो }\edtext{\skt{वागारसंज्ञितः}}{\lemma{\skt{वागारसंज्ञितः}}\Afootnote{K H; \skt{वाऽगारसंभवः~॥}\ A}}\skt{~॥}\pend

\pstart\versepar \ledsidenote{\passref{5.1.59}}\edtext{\skt{शिरोऽभ्यङ्गः}}{\lemma{\skt{ऽभ्यङ्गः}}\Afootnote{A K; \skt{ऽभ्यङ्ग}\ H}}\skt{ शिरस्त्राणं स्नान\-मु\-ष्णी\-ष\-मेव च~। }\pend
\pstart\versepar \skt{स्रजश्च विषसंसृष्टाः }\edtext{\skt{साध\-ये\-द\-व\-ले\-ख\-वत्}}{\lemma{\skt{साध\appbrk{}ये\appbrk{}द\appbrk{}व\appbrk{}ले\appbrk{}ख\appbrk{}वत्}}\Afootnote{\skt{साध\appbrk{}ये\appbrk{}द\appbrk{}व\appbrk{}ले\appbrk{}ख\appbrk{}नात्}\ A K\textsuperscript{pc} H; \skt{साध\appbrk{}ये\appbrk{}द\appbrk{}व\appbrk{}ले\appbrk{}ख\appbrk{}ना\appbrk{}वत्~।}\ K\textsuperscript{ac}}}\skt{~।}\pend

\pstart\versepar \ledsidenote{\passref{5.1.60}}\skt{मुखालेपे }\edtext{\skt{मुखश्यावं}}{\lemma{\skt{मुखश्यावं}}\Afootnote{K H; \skt{मुखं श्यावं}\ A}}\skt{ युक्त\-म\-भ्य\-ङ्ग\-ल\-क्षणैः~। }\pend
\pstart\versepar \skt{}\edtext{\skt{पद्मि\-नी\-क\-ण्ट\-क\-प्रख्यैः}}{\lemma{\skt{पद्मि\appbrk{}नी\appbrk{}क\appbrk{}ण्ट\appbrk{}क\appbrk{}प्रख्यैः}}\Afootnote{A K; \skt{पद्मि\appbrk{}नी\appbrk{}क\appbrk{}ण्ठ\appbrk{}क\appbrk{}प्रख्यैः}\ H}}\skt{ }\edtext{\skt{कण्ट\-कै\-श्चो\-प\-ची\-यते}}{\lemma{\skt{कण्टकैश्}}\Afootnote{A K; \skt{कण्ठकैश्}\ H}}\skt{~।}\pend

\pstart\versepar \ledsidenote{\passref{5.1.61}}\skt{तत्र क्षौद्रघृतं }\edtext{\skt{पान\-मा\-ले\-प\-श्च\-न्दनं}}{\lemma{\skt{आलेपश्}}\Afootnote{K H; \skt{प्रलेपश्}\ A}}\skt{ घृतम्~। }\pend
\pstart\versepar \skt{पयस्या }\edtext{\skt{मधुका}}{\lemma{\skt{मधुका}}\Afootnote{K H; \skt{मधुकं}\ A}}\skt{ फञ्जी }\edtext{\skt{बन्धुजीवः}}{\lemma{\skt{बन्धुजीवः}}\Afootnote{A; \skt{वन्धुजीव}\ K H}}\skt{ }\edtext{\skt{पुनर्ण्णवा}}{\lemma{\skt{पुनर्ण्णवा}}\Afootnote{K H; \skt{पुनर्नवा~॥}\ A}}\skt{~॥}\pend

\pstart\versepar \ledsidenote{\passref{5.1.62}}\edtext{\skt{अस्वास्थ्यं}}{\lemma{\skt{अस्वास्थ्यं}}\Afootnote{A K; \skt{आस्वास्थ्यं}\ H}}\skt{ कुञ्जरादीनां }\edtext{\edtext{\skt{लाला\-स्र\-व\-ण\-मेव}}{\lemma{\skt{लालास्रवणम्}}\Afootnote{K H; \skt{लालास्रावोऽ}\ A}}}{\lemma{\skt{एव}}\Afootnote{K H; \skt{क्षिर}\ A}}\skt{ }\edtext{\skt{च}}{\lemma{\skt{च}}\Afootnote{K H; \skt{क्तता~। स्फिक्पायुमेढ्र}\ A}}\skt{~। }\pend
\pstart\versepar \skt{}\edtext{\skt{यातुश्च}}{\lemma{\skt{यातुश्च}}\Afootnote{K H; \skt{मुष्केषु यातुश्च}\ A}}\skt{ }\edtext{\skt{स्फोटनास्रावौ}}{\lemma{\skt{स्फोटनास्रावौ}}\Afootnote{\skt{स्फोटसंभवः~॥}\ A; \skt{स्फोटनंमास्रावौ}\ K; \skt{स्फोटनं स्रावौ}\ H}}\skt{ }\edtext{\edtext{\skt{मुष्क\-मे\-ढ्र\-गु\-दे\-ष्वथ}}{\lemma{\skt{मुष्क\appbrk{}मे\appbrk{}ढ्र\appbrk{}गु\appbrk{}देष्व्}}\Afootnote{K H; \emph{om.}\ A}}}{\lemma{\skt{अथ}}\Afootnote{K H; \emph{om.}\ A}}\skt{~।}\pend

\pstart\versepar \ledsidenote{\passref{5.1.63}}\edtext{\skt{तत्रा\-भ्य\-ङ्ग\-व\-दा\-चष्टे}}{\lemma{\skt{तत्रा\appbrk{}भ्य\appbrk{}ङ्ग\appbrk{}व\appbrk{}दा\appbrk{}चष्टे}}\Afootnote{K H; \skt{तत्रा\appbrk{}भ्य\appbrk{}ङ्ग\appbrk{}व\appbrk{}दे\appbrk{}वेष्टा}\ A}}\skt{ }\edtext{\skt{यात्रिवाहनयोः}}{\lemma{\skt{यात्रिवाहनयोः}}\Afootnote{K H; \skt{यातृवाहनयोः}\ A}}\skt{ }\edtext{\skt{क्रियाम्}}{\lemma{\skt{क्रियाम्}}\Afootnote{K H; \skt{क्रिया~।}\ A}}\skt{~। }\pend
\pstart\versepar \skt{शोणितागमनं खेभ्यः शिरो\-रु\-क्क\-फ\-सं\-स्रवः~।}\pend

\pstart\versepar \ledsidenote{\passref{5.1.64}}\skt{नस्यधूमगते }\edtext{\edtext{\skt{लिङ्ग\-मि\-न्द्रि\-या\-णाञ्च}}{\lemma{\skt{लिङ्गम्}}\Afootnote{K; \skt{लिङ्ग}\ A H}}}{\lemma{\skt{इन्द्रियाणाञ्च}}\Afootnote{K; \skt{मिन्दिर्याणां च}\ A; \skt{मिन्द्रियाणाञ्च}\ H}}\skt{ वैकृतम्~। }\pend
\pstart\versepar \skt{तत्र }\edtext{\edtext{\skt{सर्पिर्गवादीनां}}{\lemma{\skt{सर्पिर्}}\Afootnote{\skt{दुग्धैर्}\ A; \skt{सर्पि}\ K H}}}{\lemma{\skt{गवादीनां}}\Afootnote{A; \skt{र्गवादीनां}\ K; \skt{र्गवादिनां}\ H}}\skt{ }\edtext{\skt{दुग्धैः}}{\lemma{\skt{दुग्धैः}}\Afootnote{K H; \skt{सर्पिः}\ A}}\skt{ सातिविषैः }\edtext{\skt{शृतम्}}{\lemma{\skt{शृतम्}}\Afootnote{A; \skt{श्रितं~।}\ K; \skt{स्रितं~।}\ H}}\skt{~।}\pend

\pstart\versepar \ledsidenote{\passref{5.1.65}}\edtext{\skt{नस्यं}}{\lemma{\skt{नस्यं}}\Afootnote{K H; \skt{पाने नस्ये}\ A}}\skt{ }\edtext{\skt{पानञ्च}}{\lemma{\skt{पानञ्च}}\Afootnote{K H; \skt{च}\ A}}\skt{ }\edtext{\skt{विहितं}}{\lemma{\skt{विहितं}}\Afootnote{K H; \skt{सश्वेतं}\ A}}\skt{ }\edtext{\skt{शीतं}}{\lemma{\skt{शीतं}}\Afootnote{K H; \skt{हितं}\ A}}\skt{ सम\-द\-य\-न्ति\-कम्~॥ }\pend
\pstart\versepar \skt{गन्ध\-हा\-नि\-र्वि\-व\-र्ण्णत्वं पुष्पाणां म्लानता }\edtext{\skt{तथा}}{\lemma{\skt{तथा}}\Afootnote{K H; \skt{भवेत्~॥}\ A}}\skt{~।}\pend

\pstart\versepar \ledsidenote{\passref{5.1.66}}\skt{जिघ्रतश्च शिरोदुःखं वारिपूर्ण्णे च लोचने~। }\pend
\pstart\versepar \skt{}\edtext{\skt{तत्रेष्टं}}{\lemma{\skt{तत्रेष्टं}}\Afootnote{K H; \skt{तत्र बा}\ A}}\skt{ }\edtext{\skt{बाष्पिकं}}{\lemma{\skt{बाष्पिकं}}\Afootnote{\skt{ष्पेरितं}\ A; \skt{वाष्पिकं}\ K H}}\skt{ कर्म मुखालेपे च यत्स्मृतम्~॥}\pend

\pstart\versepar \ledsidenote{\passref{5.1.67}}\skt{कर्ण्णतैलगते }\edtext{\skt{श्रोत्रे}}{\lemma{\skt{श्रोत्रे}}\Afootnote{K; \skt{श्रोत्र}\ A; \skt{स्रोत्रे}\ H}}\skt{ }\edtext{\skt{वैकृत्यं}}{\lemma{\skt{वैकृत्यं}}\Afootnote{K H; \skt{वैगुण्यं}\ A}}\skt{ }\edtext{\skt{शोफवेदना}}{\lemma{\skt{शोफवेदना}}\Afootnote{K H; \skt{शोफवेदने~।}\ A}}\skt{~। }\pend
\pstart\versepar \skt{}\edtext{\skt{कर्ण्णास्रावश्च}}{\lemma{\skt{कर्ण्णास्रावश्च}}\Afootnote{K; \skt{कर्णस्रावश्च}\ A H}}\skt{ }\edtext{\skt{तत्राशु}}{\lemma{\skt{तत्राशु}}\Afootnote{A H; \skt{तत्राशुः}\ K}}\skt{ कर्तव्यं प्रतिपूरणम्~।}\pend

\pstart\versepar \ledsidenote{\passref{5.1.68}}\skt{स्वरसो }\edtext{\skt{बहुपुत्रायाः}}{\lemma{\skt{बहुपुत्रायाः}}\Afootnote{A H; \skt{वहुपुत्रायाः}\ K}}\skt{ }\edtext{\skt{सघृतं}}{\lemma{\skt{सघृतं}}\Afootnote{H; \skt{सघृतः}\ A; \emph{om.}\ K}}\skt{ }\edtext{\skt{क्षौद्रसंयुतम्}}{\lemma{\skt{क्षौद्रसंयुतम्}}\Afootnote{H; \skt{क्षौद्रसंयुतः~।}\ A; \skt{…तं~।}\ K}}\skt{~। }\pend
\pstart\versepar \skt{}\edtext{\skt{सोम\-व\-ल्क\-र\-स\-श्चापि}}{\lemma{\skt{सोम\appbrk{}व\appbrk{}ल्क\appbrk{}र\appbrk{}सश्}}\Afootnote{A K; \skt{सोम\appbrk{}क\appbrk{}ल्क\appbrk{}र\appbrk{}सश्}\ H}}\skt{ सुशीतो }\edtext{\skt{हितमिष्यते}}{\lemma{\skt{हितम्}}\Afootnote{K H; \skt{हित}\ A}}\skt{~।}\pend

\pstart\versepar \ledsidenote{\passref{5.1.69}}\edtext{\skt{अस्रोपदेहौ}}{\lemma{\skt{अस्रोपदेहौ}}\Afootnote{H; \skt{अश्रूपदेहो}\ A; \skt{अश्रोपदेहौ}\ K}}\skt{ दाहश्च वेदना दृष्टिविभ्रमः~। }\pend
\pstart\versepar \skt{अञ्जने विषसंसृष्टे भवे\-दा\-न्ध्य\-म\-थापि }\edtext{\skt{वा}}{\lemma{\skt{वा}}\Afootnote{K H; \skt{च~॥}\ A}}\skt{~।}\pend

\pstart\versepar \ledsidenote{\passref{5.1.70}}\skt{तत्र सद्यो घृतं पेयं }\edtext{\skt{तर्पणे च}}{\lemma{\skt{तर्पणे च}}\Afootnote{H; \skt{तर्पणं च}\ A; \skt{तर्पणै च}\ K}}\skt{ समागधम्~। }\pend
\pstart\versepar \skt{अञ्जनं मेषशृङ्गस्य निर्यासो }\edtext{\skt{वरुणस्य}}{\lemma{\skt{वरुणस्य}}\Afootnote{A K; \skt{वरणस्य}\ H}}\skt{ च~।}\pend

\pstart \ledsidenote{\passref{5.1.71cd}}\edtext{\edtext{\edtext{\skt{कपि\-त्थ\-मे\-ष\-शृ\-ङ्गा\-भ्या\-म्पु\-ष्प\-म्भ\-ल्ला\-त\-कस्य}}{\lemma{\skt{कपि\appbrk{}त्थ\appbrk{}मे\appbrk{}ष\appbrk{}शृ\appbrk{}ङ्गा\appbrk{}भ्याम्}}\Afootnote{\skt{कपि\appbrk{}त्थ\appbrk{}मे\appbrk{}ष\appbrk{}शृ\appbrk{}ङ्ग्योश्च}\ A; \skt{कपि\appbrk{}त्थ\appbrk{}मे\appbrk{}ष\appbrk{}शृङ्ग…}\ K; \skt{कपि\appbrk{}त्थ\appbrk{}मे\appbrk{}ष\appbrk{}श्रृ\appbrk{}ङ्गा\appbrk{}भ्याम्}\ H}}}{\lemma{\skt{पुष्पम्}}\Afootnote{A H; \emph{om.}\ K}}}{\lemma{\skt{भल्लातकस्य}}\Afootnote{A H; \emph{om.}\ K}}\skt{ }\edtext{\skt{च}}{\lemma{\skt{च}}\Afootnote{H; \skt{वा~॥}\ A; \emph{om.}\ K}}\pend

\pstart \ledsidenote{\passref{5.1.72cd}}\edtext{\skt{शोफः}}{\lemma{\skt{शोफः}}\Afootnote{A K; \skt{सोफः}\ H}}\skt{ }\edtext{\skt{स्वापस्तथा}}{\lemma{\skt{स्वापस्}}\Afootnote{K H; \skt{स्रावस्}\ A}}\skt{ }\edtext{\skt{स्रावः}}{\lemma{\skt{स्रावः}}\Afootnote{K H; \skt{स्वापः}\ A}}\skt{ पादयोः स्फोटजन्म च~।}\pend

\pstart\versepar \ledsidenote{\passref{5.1.73}}\skt{भवन्ति }\edtext{\skt{विषदुष्टाभ्यां}}{\lemma{\skt{विषदुष्टाभ्यां}}\Afootnote{K H; \skt{विषजुष्टाभ्यां}\ A}}\skt{ पादु\-का\-भ्या\-म\-सं\-श\-यम्~। }\pend
\pstart\versepar \skt{उपा\-न\-त्पा\-द\-पी\-ठानि }\edtext{\skt{पादुकाभ्यां}}{\lemma{\skt{पादुकाभ्यां}}\Afootnote{K H; \skt{पादुकाव}\ A}}\skt{ }\edtext{\skt{प्रसाधयेत्}}{\lemma{\skt{प्रसाधयेत्}}\Afootnote{K H; \skt{त्प्रसाधयेत्~॥}\ A}}\skt{~।}\pend

\pstart\versepar \ledsidenote{\passref{5.1.74}}\edtext{\skt{भूषणानि}}{\lemma{\skt{भूषणानि}}\Afootnote{A K; \skt{भूषनानि}\ H}}\skt{ हतार्चींषि न विभान्ति यथा पुरा~। }\pend
\pstart\versepar \skt{स्वानि स्थानानि हन्युश्च }\edtext{\skt{दाह\-पा\-का\-व\-दा\-रणैः}}{\lemma{\skt{दाह\appbrk{}पा\appbrk{}का\appbrk{}व\appbrk{}दा\appbrk{}रणैः}}\Afootnote{A H\textsuperscript{pc}; \skt{दाह\appbrk{}पा\appbrk{}का\appbrk{}व\appbrk{}दा\appbrk{}रुणैः}\ K H\textsuperscript{ac}}}\skt{~।}\pend

\pstart\versepar \ledsidenote{\passref{5.1.75}}\edtext{\skt{पादु\-का\-भू\-ष\-णे\-षू\-क्त\-म\-भ्य\-ङ्ग\-वि\-धि\-मा\-च\-रेत्}}{\lemma{\skt{पादु\appbrk{}का\appbrk{}भू\appbrk{}ष\appbrk{}णे\appbrk{}षू\appbrk{}क्तम्}}\Afootnote{A K; \skt{पादु\appbrk{}क\appbrk{}भू\appbrk{}ष\appbrk{}णे\appbrk{}क्तम्}\ H\textsuperscript{ac}; \skt{पादु\appbrk{}क\appbrk{}भू\appbrk{}ष\appbrk{}णे\appbrk{}षू\appbrk{}क्तम्}\ H\textsuperscript{pc}}}\skt{~। }\pend
\pstart\versepar \skt{विषोपसर्गो }\edtext{\skt{बाष्पा\-दि\-र्भू\-ष\-णान्तो}}{\lemma{\skt{बाष्पादिर्}}\Afootnote{A; \skt{वाष्पादिर्}\ K H}}\skt{ य ईरितः~।}\pend

\pstart\versepar \ledsidenote{\passref{5.1.76}}\edtext{\edtext{\skt{उपद्रवांस्तत्र}}{\lemma{\skt{उपद्रवांस्}}\Afootnote{K H; \skt{समी\appbrk{}क्ष्यो\appbrk{}प\appbrk{}द्र\appbrk{}वांस्}\ A}}}{\lemma{\skt{तत्र}}\Afootnote{K H; \skt{त}\ A}}\skt{ }\edtext{\skt{वीक्ष्य}}{\lemma{\skt{वीक्ष्य}}\Afootnote{K H; \skt{स्य}\ A}}\skt{ विदधीत चिकित्सितम्~। }\pend
\pstart\versepar \skt{}\edtext{\skt{महा\-सु\-ग\-न्ध\-म\-गदं}}{\lemma{\skt{महासुगन्धम्}}\Afootnote{K H; \skt{महासुगन्धिम्}\ A}}\skt{ यं प्रवक्ष्यामि तम्भिषक्~।}\pend

\pstart\versepar \ledsidenote{\passref{5.1.77}}\skt{पाना\-ले\-प\-न\-न\-स्येषु विद\-धी\-ता\-ञ्ज\-नेषु च~॥ }\pend
\pstart\versepar \skt{विरेचनानि तीक्ष्णानि कुर्या\-त्प्र\-च्छ\-र्द\-नानि च~।}\pend

\pstart\versepar \ledsidenote{\passref{5.1.78}}\edtext{\skt{शिराश्च}}{\lemma{\skt{शिराश्च}}\Afootnote{K H; \skt{सिराश्च}\ A}}\skt{ }\edtext{\skt{व्यधयेत्प्राप्ताः}}{\lemma{\skt{व्यधयेत्प्राप्ताः}}\Afootnote{K; \skt{व्यधयेत्क्षिप्रं}\ A; \skt{वेधयेत्प्राप्ताः}\ H}}\skt{ प्राप्तं विस्रावणं यदि~। }\pend
\pstart\versepar \skt{}\edtext{\skt{मूषिकाजरुहा}}{\lemma{\skt{मूषिकाजरुहा}}\Afootnote{K; \skt{मूसिकाऽजरुहा}\ A; \skt{मूषिकाजरुका}\ H}}\skt{ }\edtext{\skt{वापि}}{\lemma{\skt{वापि}}\Afootnote{K H; \skt{वाऽपि}\ A}}\skt{ हस्ते }\edtext{\skt{बद्धा तु}}{\lemma{\skt{बद्धा तु}}\Afootnote{A; \skt{वद्धा तु}\ K H}}\skt{ }\edtext{\skt{भूपतेः}}{\lemma{\skt{भूपतेः}}\Afootnote{A K; \skt{भूपते~।}\ H}}\skt{~।}\pend

\pstart\versepar \ledsidenote{\passref{5.1.79}}\skt{करोति निर्विषं सर्वमन्नं विष\-स\-मा\-यु\-तम्~। }\pend
\pstart\versepar \skt{हृदयावरणं नित्यं }\edtext{\skt{कुर्या\-च्चा\-मि\-त्र\-म\-ध्यगः}}{\lemma{\skt{कुर्या\appbrk{}च्चा\appbrk{}मि\appbrk{}त्र\appbrk{}म\appbrk{}ध्यगः}}\Afootnote{K H; \skt{कुर्याच्च मित्रमध्यगः~॥}\ A}}\skt{~।}\pend

\pstart\versepar \ledsidenote{\passref{5.1.80}}\edtext{\edtext{\skt{पिबे\-द्घृ\-त\-म\-जे\-या\-ख्य\-म\-मृतं}}{\lemma{\skt{पिबेद्घृतम्}}\Afootnote{A; \skt{पिवेद्घृतम्}\ K H}}}{\lemma{\skt{अमृतं}}\Afootnote{K H; \skt{अमृ}\ A}}\skt{ }\edtext{\skt{चाप्यभुक्तवान्}}{\lemma{\skt{चाप्य्}}\Afootnote{K; \skt{ताख्यं}\ A; \skt{वाप्य्}\ H}}\skt{~। }\pend
\pstart\versepar \skt{सर्पिः }\edtext{\skt{क्षौद्रं}}{\lemma{\skt{क्षौद्रं}}\Afootnote{K H; \emph{om.}\ A}}\skt{ दधि पयः }\edtext{\skt{पिबेद्वा}}{\lemma{\skt{पिबेद्वा}}\Afootnote{\skt{क्षौद्रं पिबेद्वा}\ A; \skt{पिवेद्वा}\ K H}}\skt{ शीतलञ्जलम्~।}\edtext{}{\lemma{\skt{अभुक्तवान्~। सर्पिः}}\Afootnote{K H; \skt{च बुद्धिमान्~। सर्पिर्}\ A}}\pend

\pstart\versepar \ledsidenote{\passref{5.1.81}}\edtext{\skt{गोधा\-म\-यू\-र\-न\-कु\-ला\-न्पृ\-ष\-ता\-न्ह\-रि\-णा\-नपि}}{\lemma{\skt{गोधा\appbrk{}म\appbrk{}यू\appbrk{}र\appbrk{}न\appbrk{}कु\appbrk{}ला\appbrk{}न्पृ\appbrk{}ष\appbrk{}ता\appbrk{}न्ह\appbrk{}रि\appbrk{}णा\appbrk{}नपि}}\Afootnote{H; \skt{मयू\appbrk{}रा\appbrk{}न्न\appbrk{}कु\appbrk{}ला\appbrk{}न्गोधाः पृष\appbrk{}ता\appbrk{}न्ह\appbrk{}रि\appbrk{}णा\appbrk{}नपि~।}\ A; \skt{गोधा\appbrk{}म\appbrk{}यू\appbrk{}र\appbrk{}न\appbrk{}कुलां पृषता हरिणानपि~।}\ K}}\skt{~। }\pend
\pstart\versepar \skt{}\edtext{\skt{विषघ्नानाञ्च}}{\lemma{\skt{विषघ्नानाञ्च}}\Afootnote{K H; \emph{om.}\ A}}\skt{ सततं }\edtext{\skt{रसांस्तेषां}}{\lemma{\skt{रसांस्}}\Afootnote{\skt{भक्षयेच्चापि रसांस्}\ A; \skt{रसास्}\ K H}}\skt{ }\edtext{\skt{पिबेदपि}}{\lemma{\skt{पिबेदपि}}\Afootnote{A; \skt{पिवेदपि}\ K H}}\skt{~।}\pend

\pstart\versepar \ledsidenote{\passref{5.1.82}}\skt{गोधा\-न\-कु\-ल\-मां\-सेषु हरिणस्य च }\edtext{\skt{बुद्धिमान्}}{\lemma{\skt{बुद्धिमान्}}\Afootnote{A; \skt{वुद्धिमां}\ K H}}\skt{~। }\pend
\pstart\versepar \skt{दद्यात्सुपिष्टां पालिन्दीं }\edtext{\skt{मधुकं}}{\lemma{\skt{मधुकं}}\Afootnote{A K; \skt{मधूकं}\ H}}\skt{ }\edtext{\skt{शर्करन्तथा}}{\lemma{\skt{शर्करन्}}\Afootnote{K H; \skt{शर्करां}\ A}}\skt{~।}\pend

\pstart\versepar \ledsidenote{\passref{5.1.83}}\skt{शर्करातिविषे देये मायूरे समहौषधे~। }\pend
\pstart\versepar \skt{पार्षते चापि }\edtext{\skt{देयाः स्युः}}{\lemma{\skt{देयाः स्युः}}\Afootnote{A H; \skt{देया स्युः}\ K}}\skt{ पिप्पल्यः समहौषधाः~।}\pend

\pstart\versepar \ledsidenote{\passref{5.1.84}}\skt{सक्षौद्रः सघृतः }\edtext{\skt{शीतः}}{\lemma{\skt{शीतः}}\Afootnote{K H; \skt{चैव}\ A}}\skt{ }\edtext{\skt{निम्बयूष}}{\lemma{\skt{निम्बयूष}}\Afootnote{\skt{शिम्बीयूषो}\ A; \skt{निम्वयूष}\ K H}}\skt{ }\edtext{\skt{हितस्तथा}}{\lemma{\skt{तथा}}\Afootnote{K H; \skt{सदा~।}\ A}}\skt{~। }\pend
\pstart\versepar \skt{}\edtext{\skt{विषघ्नानि च}}{\lemma{\skt{विषघ्नानि च}}\Afootnote{A H; \skt{विषाघ्नानि च}\ K}}\skt{ सेवेत भक्ष्यभोज्यानि }\edtext{\skt{बुद्धिमान्}}{\lemma{\skt{बुद्धिमान्}}\Afootnote{A; \skt{वुद्धिमान्}\ K H}}\skt{~।}\pend

\pstart\versepar \ledsidenote{\passref{5.1.85}}\edtext{\skt{पिप्प\-ली\-म\-धु\-क\-क्षौ\-द्र\-श\-र्क\-रे\-क्षु\-र\-सा\-म्बुभिः}}{\lemma{\skt{पिप्प\appbrk{}ली\appbrk{}म\appbrk{}धु\appbrk{}क\appbrk{}क्षौ\appbrk{}द्र\appbrk{}श\appbrk{}र्क\appbrk{}रे\appbrk{}क्षु\appbrk{}र\appbrk{}सा\appbrk{}म्बुभिः}}\Afootnote{A; \skt{पिप्प\appbrk{}ली\appbrk{}म\appbrk{}धु\appbrk{}क\appbrk{}क्षौ\appbrk{}द्र\appbrk{}श\appbrk{}र्क\appbrk{}रे\appbrk{}क्षु\appbrk{}र\appbrk{}सा\appbrk{}म्वुभिः}\ K H}}\skt{~। }\pend
\pstart\versepar \skt{छर्द\-ये\-द्गु\-प्त\-हृ\-दयो }\edtext{\skt{यदि}}{\lemma{\skt{यदि}}\Afootnote{K H; \skt{भक्षितं यदि}\ A}}\skt{ }\edtext{\skt{पीतं}}{\lemma{\skt{पीतं}}\Afootnote{K H; \emph{om.}\ A}}\skt{ }\edtext{\skt{भवेद्विषम्}}{\lemma{\skt{भवेद्विषम्}}\Afootnote{\skt{वै विषम्~॥}\ A; \skt{भवेद्विषमिति~॥}\ K H}}\pend

\pstart \ledsidenote{\passref{5.1.colophon}}\edtext{\skt{इति~॥~॥ कल्पेषु}}{\lemma{\skt{इति~॥~॥ कल्पेषु}}\Afootnote{\skt{इति सुश्रु\appbrk{}त\appbrk{}सं\appbrk{}हि\appbrk{}तायां कल्प\appbrk{}स्था\appbrk{}नेऽ\appbrk{}न्न\appbrk{}पा\appbrk{}न\appbrk{}र\appbrk{}क्षा\appbrk{}कल्पो नाम}\ A; \skt{कल्पे १~॥ {\symfont ❈}~॥}\ K; \skt{॥ कल्पेषु}\ H}}\skt{ }\edtext{\edtext{\skt{प्रथमोऽध्यायः}}{\lemma{\skt{प्रथमो}}\Afootnote{A H; \emph{om.}\ K}}}{\lemma{\skt{ऽध्यायः}}\Afootnote{A H; \emph{om.}\ K}}\skt{~॥}\pend

\endnumbering
\beginnumbering
\eledsection*{Adhyāya 4}
\pstart \ledsidenote{\passref{5.4.1}}\skt{अथातः }\edtext{\skt{सर्प\-द\-ष्ट\-वि\-ष\-वि\-ज्ञा\-नीयं}}{\lemma{\skt{सर्प\appbrk{}द\appbrk{}ष्ट\appbrk{}वि\appbrk{}ष\appbrk{}वि\appbrk{}ज्ञा\appbrk{}नीयं}}\Afootnote{A K; \skt{सर्प्प\appbrk{}द\appbrk{}ष्ट\appbrk{}वि\appbrk{}ज्ञा\appbrk{}नी\appbrk{}यङ्}\ H}}\skt{ कल्पं }\edtext{\skt{व्याख्यास्यामः}}{\lemma{\skt{व्याख्यास्यामः}}\Afootnote{A H; \emph{om.}\ K}}\skt{~॥}\pend

\pstart\versepar \ledsidenote{\passref{5.4.3}}\skt{धन्वन्तरिं महाप्राज्ञं सर्व\-शा\-स्त्र\-वि\-शा\-र\-दम्~। }\pend
\pstart\versepar \skt{}\edtext{\skt{चर\-णा\-वु\-प\-स\-ङ्गृह्य}}{\lemma{\skt{चरणाव्}}\Afootnote{K H; \skt{पादयोर्}\ A}}\skt{ सुश्रुतः परिपृच्छति~॥}\pend

\pstart\versepar \ledsidenote{\passref{5.4.4}}\edtext{\skt{सर्पसंख्यां}}{\lemma{\skt{सर्पसंख्यां}}\Afootnote{A K; \skt{सर्पसंख्या}\ H}}\skt{ विभागं च दष्ट\-ल\-क्ष\-ण\-मेव च~। }\pend
\pstart\versepar \skt{ज्ञानञ्च विषवेगानां }\edtext{\skt{भग\-व\-न्व\-क्तु\-म\-र्हसि}}{\lemma{\skt{भगवन्वक्तुम्}}\Afootnote{A H; \skt{भगवं वक्तुम्}\ K}}\skt{~।}\pend

\pstart\versepar \ledsidenote{\passref{5.4.5}}\skt{तस्य }\edtext{\skt{तद्वचनं}}{\lemma{\skt{तद्वचनं}}\Afootnote{A K\textsuperscript{pc} H; \skt{तद्ववचनं}\ K\textsuperscript{ac}}}\skt{ श्रुत्वा }\edtext{\skt{प्राब्र\-वी\-द्भि\-ष\-जा\-म्वरः}}{\lemma{\skt{प्राब्र\appbrk{}वी\appbrk{}द्भि\appbrk{}ष\appbrk{}जाम्}}\Afootnote{A; \skt{प्राव्र\appbrk{}वी\appbrk{}द्भि\appbrk{}ष\appbrk{}जाम्}\ K H}}\skt{~। }\pend
\pstart\versepar \skt{}\edtext{\skt{असंख्येया}}{\lemma{\skt{असंख्येया}}\Afootnote{K H; \skt{असंख्या}\ A}}\skt{ }\edtext{\skt{महात्मानो}}{\lemma{\skt{महात्मानो}}\Afootnote{K H; \emph{om.}\ A}}\skt{ }\edtext{\skt{वासु\-की\-त\-क्ष\-का\-दयः}}{\lemma{\skt{वासु\appbrk{}की\appbrk{}त\appbrk{}क्ष\appbrk{}का\appbrk{}दयः}}\Afootnote{K H; \skt{वासुकिश्रेष्ठा विख्या\appbrk{}ता\appbrk{}स्त\appbrk{}क्ष\appbrk{}का\appbrk{}दयः~॥}\ A}}\skt{~।}\pend

\pstart\versepar \ledsidenote{\passref{5.4.6}}\skt{महीधराश्च नागेन्द्राः हुता\-ग्नि\-स\-म\-व\-र्चसः~। }\pend
\pstart\versepar \skt{ये चाप्यजस्रं गर्जन्ति वर्षन्ति च तपन्ति च~।}\edtext{}{\lemma{\skt{हुता\appbrk{}ग्नि\appbrk{}स\appbrk{}म\appbrk{}व\appbrk{}र्चसः~। ये}}\Afootnote{K H; \skt{हुता\appbrk{}ग्नि\appbrk{}स\appbrk{}म\appbrk{}ते\appbrk{}जसः~। ये}\ A}}\pend

\pstart\versepar \ledsidenote{\passref{5.4.7}}\skt{ससा\-ग\-र\-गि\-रि\-द्वीपा }\edtext{\skt{यैश्च}}{\lemma{\skt{च}}\Afootnote{K H; \skt{इ}\ A}}\skt{ }\edtext{\skt{सन्धार्यते}}{\lemma{\skt{सन्धार्यते}}\Afootnote{K H; \skt{यं धार्यते}\ A}}\skt{ मही~। }\pend
\pstart\versepar \skt{क्रुद्धा }\edtext{\skt{निश्वासदृष्टिभ्यां ये}}{\lemma{\skt{निश्वासदृष्टिभ्यां ये}}\Afootnote{K H; \skt{निःश्वासदृष्टिभ्यां ये}\ A}}\skt{ हन्युरखिलं जगत्~।}\pend

\pstart\versepar \ledsidenote{\passref{5.4.8}}\skt{नमस्तेभ्यो }\edtext{\skt{न}}{\lemma{\skt{न}}\Afootnote{K H; \skt{ऽस्ति नो}\ A}}\skt{ }\edtext{\skt{तैः}}{\lemma{\skt{तैः}}\Afootnote{K H; \skt{तेषां}\ A}}\skt{ }\edtext{\edtext{\skt{किञ्चित्कार्यमत्र}}{\lemma{\skt{किञ्चित्कार्यम्}}\Afootnote{K H; \skt{कार्यं किञ्चि}\ A}}}{\lemma{\skt{अत्र}}\Afootnote{K H; \emph{om.}\ A}}\skt{ चिकित्सया~। }\pend
\pstart\versepar \skt{ये तु दंष्ट्राविषा भौमा ये }\edtext{\skt{दशन्ति च}}{\lemma{\skt{दशन्ति च}}\Afootnote{A H; \skt{दशंते च}\ K}}\skt{ }\edtext{\skt{मानवान्}}{\lemma{\skt{मानवान्}}\Afootnote{K H; \skt{मानुषान्~॥}\ A}}\skt{~॥}\edtext{}{\lemma{\skt{चिकित्सया~। ये तु}}\Afootnote{K H; \skt{च्चिकित्सया~। ये तु}\ A}}\pend

\pstart\versepar \ledsidenote{\passref{5.4.9}}\skt{तेषां }\edtext{\skt{संख्यां}}{\lemma{\skt{संख्यां}}\Afootnote{K H; \skt{सङ्ख्यां}\ A}}\skt{ प्रवक्ष्यामि यथा\-व\-द\-नु\-पू\-र्वशः~। }\pend
\pstart\versepar \skt{}\edtext{\skt{अशीतिरेव}}{\lemma{\skt{एव}}\Afootnote{K H; \skt{त्वेव}\ A}}\skt{ सर्पाणां }\edtext{\skt{भिद्यते ते तु}}{\lemma{\skt{भिद्यते ते तु}}\Afootnote{K H; \skt{भिद्यते}\ A}}\skt{ }\edtext{\skt{पञ्चधा}}{\lemma{\skt{पञ्चधा}}\Afootnote{K H; \skt{पञ्चधा तु सा~॥}\ A}}\skt{~।}\pend

\pstart\versepar \ledsidenote{\passref{5.4.10}}\skt{दर्वीकरा }\edtext{\skt{मण्डलिनो}}{\lemma{\skt{मण्डलिनो}}\Afootnote{A K; \skt{मण्डलिना}\ H}}\skt{ }\edtext{\skt{राजी\-म\-न्त\-स्त\-थैव}}{\lemma{\skt{राजीमन्तस्}}\Afootnote{K H; \skt{राजिमन्तस्}\ A}}\skt{ च~॥ }\pend
\pstart\versepar \skt{निर्विषा वैकरञ्जाश्च त्रिविधास्ते पुनः स्मृताः~।}\pend

\pstart\versepar \ledsidenote{\passref{5.4.11}}\edtext{\skt{विंशतिः}}{\lemma{\skt{विंशतिः}}\Afootnote{K H; \emph{om.}\ A}}\skt{ }\edtext{\edtext{\skt{फणिनस्तेषां}}{\lemma{\skt{फणिनस्}}\Afootnote{H; \emph{om.}\ A; \skt{फनिनस्}\ K}}}{\lemma{\skt{तेषां}}\Afootnote{K H; \skt{दर्वी}\ A}}\skt{ }\edtext{\skt{षट्च}}{\lemma{\skt{षट्च}}\Afootnote{K; \skt{करा}\ A; \skt{षट्वा}\ H}}\skt{ मण्डलिनः }\edtext{\skt{पुनः}}{\lemma{\skt{पुनः}}\Afootnote{K; \emph{om.}\ A; \skt{पुन~।}\ H}}\skt{~। }\pend
\pstart\versepar \skt{}\edtext{\skt{तावन्त}}{\lemma{\skt{तावन्त}}\Afootnote{K H; \emph{om.}\ A}}\skt{ }\edtext{\skt{एव}}{\lemma{\skt{एव}}\Afootnote{K H; \emph{om.}\ A}}\skt{ }\edtext{\skt{विज्ञेया}}{\lemma{\skt{विज्ञेया}}\Afootnote{K H; \emph{om.}\ A}}\skt{ }\edtext{\edtext{\skt{राजी\-म\-न्त\-स्त्र\-यो\-दश}}{\lemma{\skt{राजीमन्तस्}}\Afootnote{K H; \skt{राजिमन्तश्}\ A}}}{\lemma{\skt{त्रयोदश}}\Afootnote{K H; \skt{च पन्नगाः~। तेषु दर्वीकरा ज्ञेया विंशतिः षट्च पन्नगाः~॥}\ A}}\skt{~।}\pend

\pstart \ledsidenote{\passref{5.4.12}}\edtext{\skt{निर्विषा}}{\lemma{\skt{निर्विषा}}\Afootnote{K H; \skt{द्वाविं\appbrk{}श\appbrk{}ति\appbrk{}र्म\appbrk{}ण्ड\appbrk{}लिनो राजिमन्तस्तथा दश~। निर्विषा}\ A}}\skt{ द्वादश }\edtext{\skt{प्रोक्ता}}{\lemma{\skt{प्रोक्ता}}\Afootnote{K H; \skt{ज्ञेया}\ A}}\skt{ }\edtext{\skt{वैकरञ्जास्त्रयः}}{\lemma{\skt{त्रयः}}\Afootnote{A H; \skt{त्रय}\ K}}\skt{ }\edtext{\skt{स्मृताः}}{\lemma{\skt{स्मृताः}}\Afootnote{K H; \skt{तथा~॥}\ A}}\skt{~।}\pend

\pstart \ledsidenote{\passref{5.4.13}}\edtext{\skt{पादाभिमृष्टा}}{\lemma{\skt{पादाभिमृष्टा}}\Afootnote{K H; \skt{वैकरञ्जोद्भवाः सप्त चित्रा मण्डलिराजिलाः~। पादाभिमृष्टा}\ A}}\skt{ दुष्टा वा क्रुद्धा ग्रासार्थिनोऽपि वा~।}\pend

\pstart\versepar \ledsidenote{\passref{5.4.14}}\skt{ते दशन्ति }\edtext{\skt{महाक्रोधास्तच्च}}{\lemma{\skt{तच्च}}\Afootnote{K H; \emph{om.}\ A}}\skt{ }\edtext{\skt{त्रिविधमुच्यते}}{\lemma{\skt{उच्यते}}\Afootnote{K H; \skt{भीमदर्शनाः~।}\ A}}\skt{~॥ }\pend
\pstart\versepar \skt{}\edtext{\edtext{\skt{सर्पि\-त\-न्द\-रि\-त\-म्वापि}}{\lemma{\skt{दरितम्}}\Afootnote{K H; \skt{रदितं}\ A}}}{\lemma{\skt{वापि}}\Afootnote{K H; \skt{चापि}\ A}}\skt{ तृतीयमथ निर्विषम्~। }\pend
\pstart\versepar \skt{}\edtext{\skt{सर्प\-गा\-त्रा\-ह\-त\-ङ्के\-चि\-दि\-च्छन्ति}}{\lemma{\skt{सर्प\appbrk{}गा\appbrk{}त्रा\appbrk{}ह\appbrk{}तङ्}}\Afootnote{H; \skt{सर्पाङ्गाभिहतं}\ A; \skt{सर्प ग्रा स्रा त्राहतङ्}\ K}}\skt{ खलु तद्विदः~।}\pend

\pstart\versepar \ledsidenote{\passref{5.4.15}}\skt{पदानि यत्र }\edtext{\skt{दन्तानामेकं द्वे वा}}{\lemma{\skt{एकं द्वे वा}}\Afootnote{A H; \skt{एक द्वे वा}\ K}}\skt{ }\edtext{\skt{बहूनि वा}}{\lemma{\skt{बहूनि वा}}\Afootnote{A; \skt{वहूनि वा}\ K H}}\skt{~। }\pend
\pstart\versepar \skt{निम\-ग्ना\-न्य\-ल्प\-र\-क्तानि }\edtext{\skt{यान्युद्वृत्तः}}{\lemma{\skt{उद्वृत्तः}}\Afootnote{K H; \skt{उद्वृत्य}\ A}}\skt{ करोति }\edtext{\skt{च}}{\lemma{\skt{च}}\Afootnote{H; \skt{हि~॥}\ A; \skt{चः~।}\ K}}\skt{~।}\pend

\pstart\versepar \ledsidenote{\passref{5.4.16}}\edtext{\skt{चुञ्चु\-मा\-ल\-क\-यु\-क्तानि}}{\lemma{\skt{चुञ्चु\appbrk{}मा\appbrk{}ल\appbrk{}क\appbrk{}यु\appbrk{}क्तानि}}\Afootnote{K H; \skt{चञ्चु\appbrk{}मा\appbrk{}ल\appbrk{}क\appbrk{}यु\appbrk{}क्तानि}\ A}}\skt{ वैकृ\-त्य\-क\-र\-णानि च~। }\pend
\pstart\versepar \skt{संक्षिप्तानि सशोफानि विद्यात्तत्सर्पितं भिषक्~।}\pend

\pstart\versepar \ledsidenote{\passref{5.4.17}}\edtext{\skt{राज्याः}}{\lemma{\skt{राज्याः}}\Afootnote{K H; \skt{राज्यः}\ A}}\skt{ सलोहिता यत्र }\edtext{\skt{नीला}}{\lemma{\skt{नीला}}\Afootnote{K H; \skt{नीलाः}\ A}}\skt{ }\edtext{\skt{वा}}{\lemma{\skt{वा}}\Afootnote{K H; \skt{पी}\ A}}\skt{ }\edtext{\skt{यदि वा}}{\lemma{\skt{यदि वा}}\Afootnote{K H; \skt{ताः सितास्}\ A}}\skt{ }\edtext{\skt{सिता}}{\lemma{\skt{सिता}}\Afootnote{K H; \skt{तथा~।}\ A}}\skt{~। }\pend
\pstart\versepar \skt{}\edtext{\skt{विज्ञे\-य\-न्द\-रि\-त\-न्तत्तु}}{\lemma{\skt{दरितन्}}\Afootnote{K H; \skt{रदितं}\ A}}\skt{ }\edtext{\skt{भिष\-जा\-ल्प\-वि\-षा\-न्वि\-तम्}}{\lemma{\skt{भिष\appbrk{}जा\appbrk{}ल्प\appbrk{}वि\appbrk{}षा\appbrk{}न्वि\appbrk{}तम्}}\Afootnote{K; \skt{ज्ञेयमल्पविषं च तत्~॥}\ A; \skt{विष\appbrk{}जा\appbrk{}ल्प\appbrk{}वि\appbrk{}षा\appbrk{}न्वितं~।}\ H}}\skt{~।}\pend

\pstart\versepar \ledsidenote{\passref{5.4.18}}\skt{अशो\-फ\-म\-ल्प\-दु\-ष्टा\-सृ\-क्प्र\-कृ\-ति\-स्थस्य देहिनः~। }\pend
\pstart\versepar \skt{पदं पदानि वा }\edtext{\skt{विद्यादविषाणि}}{\lemma{\skt{विद्यादविषाणि}}\Afootnote{K H; \skt{विद्यादविषं त}\ A}}\skt{ }\edtext{\skt{चिकित्सकः}}{\lemma{\skt{चिकित्सकः}}\Afootnote{K H; \skt{च्चिकित्सकः~॥}\ A}}\skt{~।}\pend

\pstart\versepar \ledsidenote{\passref{5.4.19}}\skt{सर्पस्पृष्टस्य भीरोर्हि भयेन कुपितोऽनिलः~। }\pend
\pstart\versepar \skt{कस्यचित्कुरुते शोफं }\edtext{\skt{सर्प\-गा\-त्रा\-ह\-तन्तु}}{\lemma{\skt{सर्प\appbrk{}गा\appbrk{}त्रा\appbrk{}ह\appbrk{}तन्तु}}\Afootnote{K H; \skt{सर्पाङ्गाभिहतं तु}\ A}}\skt{ तत्~।}\pend

\pstart\versepar \ledsidenote{\passref{5.4.20}}\skt{व्याधि\-तो\-द्वि\-ग्न\-द\-ष्टानि ज्ञेया\-न्य\-ल्प\-वि\-षाणि तु~॥ }\pend
\pstart\versepar \skt{}\edtext{\skt{तथा\-ति\-बा\-ल\-वृ\-द्धानां}}{\lemma{\skt{तथा\appbrk{}ति\appbrk{}बा\appbrk{}ल\appbrk{}वृ\appbrk{}द्धानां}}\Afootnote{\skt{तथा\appbrk{}ति\appbrk{}वृ\appbrk{}द्ध\appbrk{}बा\appbrk{}लाभि}\ A; \skt{तथा\appbrk{}ति\appbrk{}वा\appbrk{}ल\appbrk{}वृद्ध}\ K; \skt{तथा\appbrk{}ति\appbrk{}वा\appbrk{}ल\appbrk{}वृ\appbrk{}द्धानां}\ H}}\skt{ दष्टमल्पविषं स्मृतम्~॥}\pend

\pstart\versepar \ledsidenote{\passref{5.4.21}}\edtext{\skt{सुप\-र्ण्ण\-दे\-व\-ब्र\-ह्म\-र्षि\-भू\-त\-सि\-द्ध\-नि\-षे\-विते}}{\lemma{\skt{सुप\appbrk{}र्ण्ण\appbrk{}दे\appbrk{}व\appbrk{}ब्र\appbrk{}ह्म\appbrk{}र्षि\appbrk{}भू\appbrk{}त\appbrk{}सि\appbrk{}द्ध\appbrk{}नि\appbrk{}षे\appbrk{}विते}}\Afootnote{\skt{सुप\appbrk{}र्ण\appbrk{}दे\appbrk{}व\appbrk{}ब्र\appbrk{}ह्म\appbrk{}र्षि\appbrk{}य\appbrk{}क्ष\appbrk{}सि\appbrk{}द्ध\appbrk{}नि\appbrk{}षे\appbrk{}विते~।}\ A; \skt{सुप\appbrk{}र्ण्ण\appbrk{}दे\appbrk{}व\appbrk{}व्र\appbrk{}ह्म\appbrk{}र्षि\appbrk{}भू\appbrk{}त\appbrk{}सि\appbrk{}द्ध\appbrk{}नि\appbrk{}षे\appbrk{}विते}\ K H}}\skt{~। }\pend
\pstart\versepar \skt{}\edtext{\skt{विष\-घ्नौ\-ष\-ध\-जुष्टे च}}{\lemma{\skt{विष\appbrk{}घ्नौ\appbrk{}ष\appbrk{}ध\appbrk{}जुष्टे च}}\Afootnote{K H; \skt{विष\appbrk{}घ्नौ\appbrk{}ष\appbrk{}धि\appbrk{}युक्ते च}\ A}}\skt{ देशे न क्रमते विषम्~।}\pend

\pstart\versepar \ledsidenote{\passref{5.4.22}}\skt{रथा\-ङ्ग\-ला\-ङ्ग\-ल\-च्छ\-त्र\-स्व\-स्ति\-कां\-कु\-श\-धा\-रिणः~। }\pend
\pstart\versepar \skt{ज्ञेया दर्वीकराः सर्पाः फणिनः शीघ्रगामिनः~।}\pend

\pstart\versepar \ledsidenote{\passref{5.4.23}}\edtext{\skt{मण्ड\-लै\-र्वि\-वि\-धै\-श्चित्त्राः}}{\lemma{\skt{चित्त्राः}}\Afootnote{A K; \skt{चित्रा}\ H}}\skt{ पृथवो मन्दगामिनः~। }\pend
\pstart\versepar \skt{ज्ञेया }\edtext{\skt{मण्डलिनश्चापि}}{\lemma{\skt{चापि}}\Afootnote{K H; \skt{सर्पा}\ A}}\skt{ }\edtext{\skt{ज्वलिताग्निसमा}}{\lemma{\skt{ज्वलिताग्निसमा}}\Afootnote{K H; \skt{ज्वलनार्कसम}\ A}}\skt{ }\edtext{\skt{विषैः}}{\lemma{\skt{विषैः}}\Afootnote{K H; \skt{प्रभाः~॥}\ A}}\skt{~।}\pend

\pstart\versepar \ledsidenote{\passref{5.4.24}}\skt{स्निग्धा }\edtext{\skt{विवि\-ध\-व\-र्ण्णा\-भि\-स्ति\-र्य\-गू\-र्ध्वञ्च}}{\lemma{\skt{तिर्यगूर्ध्वञ्च}}\Afootnote{A K; \skt{तिर्यगूर्द्धञ्च}\ H}}\skt{ राजिभिः~। }\pend
\pstart\versepar \skt{}\edtext{\skt{विचित्रा}}{\lemma{\skt{विचित्रा}}\Afootnote{K H; \skt{चित्रिता}\ A}}\skt{ इव ये भान्ति }\edtext{\skt{राजीमन्तस्तु ते}}{\lemma{\skt{राजीमन्तस्तु ते}}\Afootnote{K H; \skt{राजिमन्तस्तु ते}\ A}}\skt{ स्मृताः~।}\pend

\pstart\versepar \ledsidenote{\passref{5.4.25}}\skt{मुक्तारूप्यप्रभा ये च कपिला ये च पन्नगाः~। }\pend
\pstart\versepar \skt{}\edtext{\skt{सुवर्ण्णाभाः}}{\lemma{\skt{सुवर्ण्णाभाः}}\Afootnote{K H; \skt{सुगन्धयः}\ A}}\skt{ }\edtext{\skt{सुगन्धाश्च ते}}{\lemma{\skt{सुगन्धाश्च ते}}\Afootnote{K H; \skt{सुवर्णाभास्ते}\ A}}\skt{ जात्या }\edtext{\skt{ब्राह्मणाः}}{\lemma{\skt{ब्राह्मणाः}}\Afootnote{A; \skt{व्राह्मणाः}\ K H}}\skt{ स्मृताः}\pend

\pstart\versepar \ledsidenote{\passref{5.4.26}}\skt{क्षत्रियाः स्निग्धवर्ण्णास्तु पन्नगा भृशकोपनाः~। }\pend
\pstart\versepar \skt{}\edtext{\edtext{\skt{सूर्यश्चन्द्रः}}{\lemma{\skt{सूर्यश्}}\Afootnote{K H; \skt{सूर्य}\ A}}}{\lemma{\skt{चन्द्रः}}\Afootnote{K H; \skt{चन्द्रा}\ A}}\skt{ }\edtext{\edtext{\skt{क्षितिश्छत्रं}}{\lemma{\skt{क्षितिश्}}\Afootnote{K H; \skt{कृति}\ A}}}{\lemma{\skt{छत्रं}}\Afootnote{K H; \skt{च्छत्र}\ A}}\skt{ }\edtext{\skt{लक्ष्यं}}{\lemma{\skt{लक्ष्यं}}\Afootnote{K H; \skt{लक्ष्म}\ A}}\skt{ }\edtext{\skt{तेषा\-न्त\-था\-द्रि\-जम्}}{\lemma{\skt{तथाद्रिजम्}}\Afootnote{K H; \skt{तथाऽम्बुजम्~॥}\ A}}\skt{~॥}\pend

\pstart\versepar \ledsidenote{\passref{5.4.27}}\skt{कृष्णा }\edtext{\skt{वज्रप्रभा ये च}}{\lemma{\skt{वज्रप्रभा ये च}}\Afootnote{K H; \skt{वज्रनिभा ये च}\ A}}\skt{ लोहिता वर्ण्णतस्तथा~। }\pend
\pstart\versepar \skt{धूम्राः पारावताभाश्च वैश्यास्ते पन्नगाः स्मृताः~॥}\pend

\pstart\versepar \ledsidenote{\passref{5.4.28}}\edtext{\skt{महि\-ष\-द्वी\-पि\-व\-र्ण्णा\-भा\-स्त\-थैव}}{\lemma{\skt{महि\appbrk{}ष\appbrk{}द्वी\appbrk{}पि\appbrk{}व\appbrk{}र्ण्णा\appbrk{}भास्}}\Afootnote{A K\textsuperscript{pc} H\textsuperscript{ac}; \skt{महि\appbrk{}ष\appbrk{}द्वि\appbrk{}पि\appbrk{}व\appbrk{}र्ण्णा\appbrk{}भास्}\ K\textsuperscript{ac}; \skt{महि\appbrk{}ष\appbrk{}द्वी\appbrk{}पि\appbrk{}द्वि\appbrk{}ज\appbrk{}व\appbrk{}र्ण्णा\appbrk{}भास्}\ H\textsuperscript{pc}}}\skt{ परुषत्वचः~।}\skt{ }\pend
\pstart\versepar \skt{भिन्नवर्ण्णाश्च ये }\edtext{\skt{केचिच्छूद्रास्ते}}{\lemma{\skt{केचिच्छूद्रास्ते}}\Afootnote{K H; \skt{केचिच्छूद्रा ते}\ A}}\skt{ परिकीर्त्तिताः~॥}\pend

\pstart\versepar \ledsidenote{\passref{5.4.31}}\skt{रजन्याः }\edtext{\skt{प्रथमे}}{\lemma{\skt{प्रथमे}}\Afootnote{K H; \skt{पश्चिमे}\ A}}\skt{ यामे }\edtext{\edtext{\skt{सर्वा\-श्चि\-त्रा\-श्च\-रन्ति}}{\lemma{\skt{सर्वाश्}}\Afootnote{K H; \skt{सर्पाश्}\ A}}\skt{ ह}}{\lemma{\skt{चरन्ति ह}}\Afootnote{K H; \skt{चरन्ति हि~।}\ A}}\skt{~।}\skt{ }\pend
\pstart\versepar \skt{}\edtext{\edtext{\skt{शेषास्त्वथो}}{\lemma{\skt{शेषास्}}\Afootnote{K H; \skt{शेषे}\ A}}}{\lemma{\skt{त्वथो}}\Afootnote{K H; \skt{षूक्ता}\ A}}\skt{ मण्डलिनो दिवा }\edtext{\skt{दर्वीकरास्तथा}}{\lemma{\skt{तथा}}\Afootnote{K H; \skt{स्मृताः~॥}\ A}}\skt{~।}\pend

\pstart\versepar \ledsidenote{\passref{5.4.29}}\skt{कोप\-य\-न्त्य\-नि\-ल\-ञ्जन्तोः फणिनः सर्व एव तु~। }\pend
\pstart\versepar \skt{पित्तं मण्डलिनश्चापि कफं चानेकराजयः~।}\pend

\pstart\versepar \ledsidenote{\passref{5.4.30}}\edtext{\skt{अत्य\-ल्प\-स\-म\-व\-र्ण्णाभ्यां}}{\lemma{\skt{अत्य\appbrk{}ल्प\appbrk{}स\appbrk{}म\appbrk{}व\appbrk{}र्ण्णाभ्यां}}\Afootnote{K H; \skt{अप\appbrk{}त्य\appbrk{}म\appbrk{}स\appbrk{}व\appbrk{}र्णाभ्यां}\ A}}\skt{ द्विदो\-ष\-क\-र\-ल\-क्ष\-णम्~। }\pend
\pstart\versepar \skt{}\edtext{\skt{दम्प\-त्य\-यो\-गा\-द्वि\-ज्ञेयं}}{\lemma{\skt{दम्प\appbrk{}त्य\appbrk{}यो\appbrk{}गा\appbrk{}द्वि\appbrk{}ज्ञेयं}}\Afootnote{K H; \skt{ज्ञेयौ दोषैश्च दम्पत्योर्वि}\ A}}\skt{ }\edtext{\skt{परवादञ्च}}{\lemma{\skt{परवादञ्च}}\Afootnote{K H; \skt{शेषश्चात्र}\ A}}\skt{ }\edtext{\skt{वक्ष्यति}}{\lemma{\skt{वक्ष्यति}}\Afootnote{K H; \skt{वक्ष्यते~॥}\ A}}\skt{~।}\pend

\pstart \ledsidenote{\passref{5.4.34.add-1}}\skt{तत्र दर्वीकराः कृष्णसर्पो महाकृष्णः कृष्णोदरः~। }\edtext{\skt{सर्वकृष्णः}}{\lemma{\skt{सर्वकृष्णः}}\Afootnote{K H; \skt{श्वेतकपो}\ A}}\skt{ }\edtext{\skt{श्वेतः}}{\lemma{\skt{श्वेतः}}\Afootnote{K; \skt{तो, म}\ A; \skt{स्वेतः}\ H}}\skt{ }\edtext{\skt{कपोतो}}{\lemma{\skt{कपोतो}}\Afootnote{K H; \skt{हाकपोतो,}\ A}}\skt{ }\edtext{\skt{वलाहको}}{\lemma{\skt{वलाहको}}\Afootnote{K H; \skt{बलाहको,}\ A}}\skt{ महासर्पः }\edtext{\skt{शंखपालो}}{\lemma{\skt{शंखपालो}}\Afootnote{K H; \skt{शङ्खकपालो,}\ A}}\skt{ लोहिताक्षो गवेधुकः परिसर्पः खण्डफणः }\edtext{\skt{कूकुटः}}{\lemma{\skt{कूकुटः}}\Afootnote{K; \skt{ककुदः,}\ A; \skt{कुकूटः}\ H\textsuperscript{ac}; \skt{कुक्कुटः}\ H\textsuperscript{pc}}}\skt{ पद्मो महापद्मः }\edtext{\skt{दर्भपुष्पो}}{\lemma{\skt{दर्भपुष्पो}}\Afootnote{A K; \skt{दर्व्भपुष्पो}\ H}}\skt{ दधिमुखः }\edtext{\skt{पुण्डरीकमुखो}}{\lemma{\skt{पुण्डरीकमुखो}}\Afootnote{K H; \skt{पुण्डरीको,}\ A}}\skt{ }\edtext{\skt{बभ्रूकुटीमुखो}}{\lemma{\skt{बभ्रूकुटीमुखो}}\Afootnote{\skt{भ्रूकुटीमुखो,}\ A; \skt{वभ्रूकुटीमुखो}\ K H}}\skt{ }\edtext{\skt{विचित्रः}}{\lemma{\skt{विचित्रः}}\Afootnote{\skt{विष्किरः,}\ A; \skt{विचित्र}\ K H}}\skt{ }\edtext{\skt{पुष्पाभिकीर्ण्णाभो}}{\lemma{\skt{पुष्पाभिकीर्ण्णाभो}}\Afootnote{K H; \skt{पुष्पाभिकीर्णो,}\ A}}\skt{ }\edtext{\skt{गिरिसर्पो}}{\lemma{\skt{गिरिसर्पो}}\Afootnote{\skt{गिरिसर्पः,}\ A; \skt{गिरिसर्पो}\ K H}}\skt{ }\edtext{\skt{ऋजुसर्पः}}{\lemma{\skt{ऋजुसर्पः}}\Afootnote{A; \skt{रिजुसर्पः}\ K H}}\skt{ }\edtext{\skt{श्वेतदरो}}{\lemma{\skt{श्वेतदरो}}\Afootnote{K; \skt{श्वेतोदरो,}\ A; \skt{स्वेतोदरो}\ H}}\skt{ }\edtext{\edtext{\edtext{\skt{महा\-शी\-र्षोऽ\-ल\-ग\-र्द\-श्चेति}}{\lemma{\skt{महाशीर्षो}}\Afootnote{K; \skt{महाशिरा,}\ A; \skt{महाशीर्षा}\ H}}}{\lemma{\skt{ऽलगर्दश्}}\Afootnote{\skt{अलगर्द,}\ A; \skt{लगर्दश्}\ K H}}}{\lemma{\skt{चेति}}\Afootnote{K H; \skt{अशीविष इति (१);}\ A}}\skt{~॥}\pend

\pstart \ledsidenote{\passref{5.4.34.add-2}}\skt{मण्डलिनस्तु आदर्शमण्डलः }\edtext{\skt{श्वेतमण्डलो}}{\lemma{\skt{श्वेतमण्डलो}}\Afootnote{A K; \skt{स्वेतमण्डलो}\ H}}\skt{ }\edtext{\skt{रक्तमण्डलः}}{\lemma{\skt{रक्तमण्डलः}}\Afootnote{K H; \skt{रक्तमण्डलः, चित्रमण्डलः,}\ A}}\skt{ पृषतो }\edtext{\skt{देवदिन्नः}}{\lemma{\skt{देवदिन्नः}}\Afootnote{K H; \skt{रोध्रपुष्पो,}\ A}}\skt{ }\edtext{\skt{पिलिन्दको}}{\lemma{\skt{पिलिन्दको}}\Afootnote{K H; \skt{मिलिन्दको,}\ A}}\skt{ }\edtext{\skt{वृद्धगोनसः}}{\lemma{\skt{वृद्धगोनसः}}\Afootnote{K H; \skt{गोनसो, वृद्धगोनसः,}\ A}}\skt{ }\edtext{\skt{पनसको}}{\lemma{\skt{पनसको}}\Afootnote{K H; \skt{पनसो,}\ A}}\skt{ }\edtext{\skt{महापनसकः}}{\lemma{\skt{महापनसकः}}\Afootnote{K H; \skt{महापनसो,}\ A}}\skt{ वेणुपत्रकः शिशुको }\edtext{\skt{मदनकः}}{\lemma{\skt{मदनकः}}\Afootnote{K H; \skt{मदनः,}\ A}}\skt{ }\edtext{\skt{पालिन्दकः}}{\lemma{\skt{पालिन्दकः}}\Afootnote{K H; \skt{पालिन्दिरः,}\ A}}\skt{ }\edtext{\skt{तन्तुकः}}{\lemma{\skt{तन्तुकः}}\Afootnote{K H; \skt{पिङ्गलः, तन्तुकः}\ A}}\skt{ पुष्पपाण्डुः }\edtext{\edtext{\skt{षडङ्गोऽग्निको}}{\lemma{\skt{षडङ्गो}}\Afootnote{H; \skt{षडङ्गो,}\ A; \skt{षडङ्गा}\ K}}}{\lemma{\skt{ऽग्निको}}\Afootnote{A; \skt{ग्निको}\ K H}}\skt{ }\edtext{\skt{बभ्रु}}{\lemma{\skt{बभ्रु}}\Afootnote{\skt{बभ्रुः,}\ A; \skt{वभ्रु}\ K H}}\skt{ कषायः }\edtext{\skt{खलुषः}}{\lemma{\skt{खलुषः}}\Afootnote{K H; \skt{कलुषः}\ A}}\skt{ पारावतो }\edtext{\skt{हस्ताभरणकः}}{\lemma{\skt{हस्ताभरणकः}}\Afootnote{K H; \skt{हस्ताभर}\ A}}\skt{ }\edtext{\skt{तत्रश्चित्रकः}}{\lemma{\skt{तत्रश्}}\Afootnote{K; \skt{णः,}\ A; \skt{तत्रकश्}\ H}}\skt{ }\edtext{\edtext{\skt{एणीपदश्चेति}}{\lemma{\skt{एणीपदश्}}\Afootnote{K H; \skt{एणीपद}\ A}}}{\lemma{\skt{चेति}}\Afootnote{K H; \skt{इति (२);}\ A}}\skt{~॥}\pend

\pstart \ledsidenote{\passref{5.4.34.add-3}}\edtext{\skt{राजीमन्तस्तु}}{\lemma{\skt{राजीमन्तस्तु}}\Afootnote{K H; \skt{राजिमन्तस्तु}\ A}}\skt{ }\edtext{\skt{पुण्डरीको}}{\lemma{\skt{पुण्डरीको}}\Afootnote{A K; \skt{पुण्डरीको हि}\ H}}\skt{ }\edtext{\skt{राजिचित्रो}}{\lemma{\skt{राजिचित्रो}}\Afootnote{A K; \skt{राजोराजिचित्रो}\ H}}\skt{ }\edtext{\skt{अङ्गुलिराजिः}}{\lemma{\skt{अङ्गुलिराजिः}}\Afootnote{K H; \skt{अङ्}\ A}}\skt{ }\edtext{\skt{द्व्यङ्गुलिराजिः}}{\lemma{\skt{द्व्यङ्गुलिराजिः}}\Afootnote{K H; \skt{गुलराजिः,}\ A}}\skt{~। }\edtext{\skt{बिन्दुराजिः}}{\lemma{\skt{बिन्दुराजिः}}\Afootnote{A; \skt{विन्दुराजिः}\ K H}}\skt{ }\edtext{\edtext{\skt{कर्द\-म\-स्तृ\-ण\-शो\-षकः}}{\lemma{\skt{कर्दमस्}}\Afootnote{K H; \skt{कर्दमकः,}\ A}}}{\lemma{\skt{तृणशोषकः}}\Afootnote{K H; \skt{तृणशोषकः, सर्षपकः}\ A}}\skt{ }\edtext{\edtext{\skt{श्वेत\-ह\-नु\-र्द\-र्भ\-पुष्पो}}{\lemma{\skt{श्वेतहनुर्}}\Afootnote{A; \skt{स्वेतहनुर्}\ K H}}}{\lemma{\skt{दर्भपुष्पो}}\Afootnote{A K; \skt{द्दर्व्भपुष्पो}\ H}}\skt{ }\edtext{\edtext{\skt{लोहि\-ता\-क्ष\-श्च\-क्रकः}}{\lemma{\skt{लोहिताक्षश्}}\Afootnote{K H; \emph{om.}\ A}}}{\lemma{\skt{चक्रकः}}\Afootnote{K H\textsuperscript{ac}; \skt{चक्रको, गोधूमकः,}\ A; \skt{चन्द्रकः}\ H\textsuperscript{pc}}}\skt{ }\edtext{\edtext{\skt{किक्किसादश्चेति}}{\lemma{\skt{किक्किसादश्}}\Afootnote{K H; \skt{किक्किसाद}\ A}}}{\lemma{\skt{चेति}}\Afootnote{K H; \skt{इति (३);}\ A}}\skt{~॥}\pend

\pstart \ledsidenote{\passref{5.4.34.add-4}}\skt{निर्विषास्तु }\edtext{\edtext{\skt{वला\-ह\-कोऽ\-हि\-प\-ताकः}}{\lemma{\skt{वलाहको}}\Afootnote{K H; \skt{गल}\ A}}}{\lemma{\skt{ऽहिपताकः}}\Afootnote{\skt{गोली, शू}\ A; \skt{हि पताकः}\ K H}}\skt{ }\edtext{\edtext{\skt{शुक\-प\-त्रोऽ\-ज\-गरो}}{\lemma{\skt{शुकपत्रो}}\Afootnote{K H; \skt{कपत्रो,}\ A}}}{\lemma{\skt{ऽजगरो}}\Afootnote{A H; \skt{जगरो}\ K}}\skt{ }\edtext{\skt{दीप्यकः}}{\lemma{\skt{दीप्यकः}}\Afootnote{K H; \skt{दिव्य}\ A}}\skt{~। }\edtext{\skt{इलिकिनी}}{\lemma{\skt{इलिकिनी}}\Afootnote{K H; \emph{om.}\ A}}\skt{~। }\edtext{\skt{वर्षाहीको}}{\lemma{\skt{वर्षाहीको}}\Afootnote{K H; \skt{को,}\ A}}\skt{ }\edtext{\skt{द्व्याहिकः}}{\lemma{\skt{द्व्याहिकः}}\Afootnote{K; \skt{वर्षाहिकः,}\ A; \skt{प्याहिकः~।}\ H}}\skt{~। }\edtext{\skt{क्षीरिकापुष्पः}}{\lemma{\skt{क्षीरिकापुष्पः}}\Afootnote{K; \emph{om.}\ A; \skt{क्षीरिकापुष्पकः~।}\ H}}\skt{ }\edtext{\skt{पुष्पसकली}}{\lemma{\skt{पुष्पसकली}}\Afootnote{K H; \skt{पुष्पशकली,}\ A}}\skt{ }\edtext{\skt{ज्योतीरथो}}{\lemma{\skt{ज्योतीरथो}}\Afootnote{A K\textsuperscript{pc} H; \skt{ज्योतिरथो}\ K\textsuperscript{ac}}}\skt{ }\edtext{\edtext{\skt{वृक्षकश्चेति}}{\lemma{\skt{वृक्षकश्}}\Afootnote{K H; \skt{क्षीरिकापुष्पको, अहिपताको, अन्धाहिको, गौराहिको, वृक्षेश}\ A}}}{\lemma{\skt{चेति}}\Afootnote{K H; \skt{य इति (४);}\ A}}\skt{~॥}\pend

\pstart \ledsidenote{\passref{5.4.34.add-5}}\skt{वैकरञ्जास्तु त्रयाणां }\edtext{\skt{वर्ण्णानां}}{\lemma{\skt{वर्ण्णानां}}\Afootnote{K H; \skt{दर्वीकरादीनां}\ A}}\skt{ }\edtext{\skt{व्यति\-रे\-क\-जा\-स्त\-द्यथा}}{\lemma{\skt{व्यतिरेकजास्}}\Afootnote{K H; \skt{व्यतिकराज्जाताः,}\ A}}\skt{~। माकुलिः पोटगलः }\edtext{\skt{स्निग्धराजिश्चेति}}{\lemma{\skt{चेति}}\Afootnote{K H; \skt{इति~।}\ A}}\skt{~॥ तत्र कृष्णसर्पेण गोनस्यां }\edtext{\skt{वैपरीत्येन वा}}{\lemma{\skt{वैपरीत्येन वा}}\Afootnote{A H; \skt{वैपरीतेन वा}\ K}}\skt{ जातो माकुलिः~। राजिलेन गोनस्यां }\edtext{\skt{वैपरीत्येन वा}}{\lemma{\skt{वैपरीत्येन वा}}\Afootnote{A H; \skt{वैपरीत्येन}\ K}}\skt{ जातः पोटगलः~॥ कृष्णसर्पेण राजिमत्यां वैपरीत्येन वा जातः स्निग्धराजिरिति~॥ तेषां }\edtext{\edtext{\edtext{\skt{पितृ\-व\-द्वि\-ष\-मु\-त्क\-र्षा\-द्द्व\-यो\-र्मा\-तृ\-व\-दि\-त्येके}}{\lemma{\skt{पितृवद्विषम्}}\Afootnote{K H; \skt{आद्यस्य पितृवद्विषो}\ A}}}{\lemma{\skt{उत्कर्षाद्द्वयोर्}}\Afootnote{H; \skt{त्कर्षो, द्वयोर्}\ A; \skt{उत्कर्षाद्द्व ꣸ ꣸ ꣸ ꣸ ꣸ ꣸ योर्}\ K}}}{\lemma{\skt{मातृवदित्य्}}\Afootnote{A K\textsuperscript{pc} H; \skt{मातृदित्य्}\ K\textsuperscript{ac}}}\skt{~।}\pend

\pstart \ledsidenote{\passref{5.4.34.5a}}\edtext{\skt{एवमेषां}}{\lemma{\skt{एषां}}\Afootnote{K H; \skt{एतेषां}\ A}}\skt{ }\edtext{\edtext{\skt{सर्पा\-णा\-म\-शी\-ति\-र्व्या\-ख्याता}}{\lemma{\skt{अशीतिर्}}\Afootnote{A K; \skt{अशीति}\ H}}}{\lemma{\skt{व्याख्याता}}\Afootnote{A; \skt{व्याख्याताः~॥}\ K; \skt{व्या𑑛 ख्याताः~।}\ H}}\skt{~॥}\pend

\pstart \ledsidenote{\passref{5.4.35}}\skt{तत्र }\edtext{\skt{महा\-ने\-त्र\-जि\-ह्वा\-शि\-रसः}}{\lemma{\skt{महा\appbrk{}ने\appbrk{}त्र\appbrk{}जि\appbrk{}ह्वा\appbrk{}शि\appbrk{}रसः}}\Afootnote{K H; \skt{महा\appbrk{}ने\appbrk{}त्र\appbrk{}जि\appbrk{}ह्वा\appbrk{}स्य\appbrk{}शि\appbrk{}रसः}\ A}}\skt{ पुमांसः~। }\edtext{\skt{सूक्ष्म\-ने\-त्र\-जि\-ह्वा\-शि\-रसः}}{\lemma{\skt{सूक्ष्म\appbrk{}ने\appbrk{}त्र\appbrk{}जि\appbrk{}ह्वा\appbrk{}शि\appbrk{}रसः}}\Afootnote{K H; \skt{सूक्ष्म\appbrk{}ने\appbrk{}त्र\appbrk{}जि\appbrk{}ह्वा\appbrk{}स्य\appbrk{}शि\appbrk{}रसः}\ A}}\skt{ स्त्रियः~। उभयलक्षणा }\edtext{\skt{मन्दचेष्टाक्रोधा}}{\lemma{\skt{मन्दचेष्टाक्रोधा}}\Afootnote{K H; \skt{मन्दविषा अक्रोधा}\ A}}\skt{ नपुंसका इति~॥}\pend

\pstart \ledsidenote{\passref{5.4.36}}\skt{तत्र }\edtext{\skt{सर्वेषामेव}}{\lemma{\skt{एव}}\Afootnote{K H; \emph{om.}\ A}}\skt{ सर्पाणां सामान्यत एव }\edtext{\skt{दष्ट\-ल\-क्ष\-ण\-मु\-प\-दे\-क्ष्यामः~। किं}}{\lemma{\skt{उपदेक्ष्यामः~। किं}}\Afootnote{K H; \skt{वक्ष्यामः~। किं}\ A}}\skt{ कारणम्~। विषं हि }\edtext{\skt{हुत\-हु\-त\-व\-ह\-नि\-शि\-त\-नि\-स्त्रिं\-शा\-श\-नि\-क\-ल्प\-मा\-शु\-कारि}}{\lemma{\skt{हुत\appbrk{}हु\appbrk{}त\appbrk{}व\appbrk{}ह\appbrk{}नि\appbrk{}शि\appbrk{}त\appbrk{}नि\appbrk{}स्त्रिं\appbrk{}शा\appbrk{}श\appbrk{}नि\appbrk{}क\appbrk{}ल्पम्}}\Afootnote{K H; \skt{निशि\appbrk{}त\appbrk{}नि\appbrk{}स्त्रिं\appbrk{}शा\appbrk{}श\appbrk{}नि\appbrk{}हु\appbrk{}त\appbrk{}व\appbrk{}ह\appbrk{}दे\appbrk{}श्यम्}\ A}}\skt{ मुहू\-र्त्त\-म\-प्यु\-पे\-क्षि\-त\-मा\-तु\-र\-म\-ति\-पा\-त\-यति~। न }\edtext{\edtext{\skt{चावकाशोऽस्ति}}{\lemma{\skt{चावकाशो}}\Afootnote{A K; \skt{चावकासो}\ H}}}{\lemma{\skt{ऽस्ति}}\Afootnote{A H; \skt{स्तिः}\ K}}\skt{ }\edtext{\skt{वाक्स\-मू\-ह\-म\-नु\-स\-र्तुम्}}{\lemma{\skt{अनुसर्तुम्}}\Afootnote{K H; \skt{उपसर्तुं}\ A}}\skt{~। }\edtext{\skt{प्रत्येकमपि च}}{\lemma{\skt{अपि च}}\Afootnote{K H; \skt{अपि}\ A}}\skt{ }\edtext{\edtext{\skt{दष्ट\-ल\-क्ष\-णेऽ\-भि\-हिते}}{\lemma{\skt{दष्टलक्षणे}}\Afootnote{A; \skt{दुष्टलक्षणे}\ K H}}}{\lemma{\skt{ऽभिहिते}}\Afootnote{A H; \skt{भिहिते}\ K}}\skt{ }\edtext{\skt{सर्प\-त्रै\-वि\-ध्या\-त्क्रि\-या\-त्रै\-विध्यं}}{\lemma{\skt{सर्प\appbrk{}त्रै\appbrk{}वि\appbrk{}ध्या\appbrk{}त्क्रि\appbrk{}या\appbrk{}त्रै\appbrk{}विध्यं}}\Afootnote{K\textsuperscript{pc}; \skt{सर्वत्र त्रैविध्यं}\ A; \skt{सर्व\appbrk{}र्प\appbrk{}त्रै\appbrk{}वि\appbrk{}ध्या\appbrk{}त्क्रि\appbrk{}या\appbrk{}त्रै\appbrk{}विध्यं}\ K\textsuperscript{ac}; \skt{सर्व्व\appbrk{}त्रै\appbrk{}वि\appbrk{}ध्या\appbrk{}त्क्रि\appbrk{}या\appbrk{}त्रै\appbrk{}वि\appbrk{}ध्यम्}\ H}}\skt{ भवति~। }\edtext{\skt{तस्मात्त्रैविध्येन}}{\lemma{\skt{तस्मात्त्रैविध्येन}}\Afootnote{K H; \skt{तस्मा\appbrk{}त्त्रै\appbrk{}वि\appbrk{}ध्य\appbrk{}मेव}\ A}}\skt{ वक्ष्यामः~। }\edtext{\edtext{\skt{एत\-द्ध्या\-तु\-र\-हि\-त\-म\-सं\-मो\-ह\-क\-र\-ञ्चा\-स्मि\-न्नेव}}{\lemma{\skt{असं\appbrk{}मो\appbrk{}ह\appbrk{}क\appbrk{}रञ्}}\Afootnote{K; \skt{असंमोहकरं च,}\ A; \skt{असम्मोहे करञ्}\ H}}}{\lemma{\skt{चास्मिन्नेव}}\Afootnote{K H; \skt{अपि चात्रैव स}\ A}}\skt{ }\edtext{\skt{च}}{\lemma{\skt{च}}\Afootnote{K H; \skt{र्व}\ A}}\skt{ }\edtext{\skt{सर्व\-व्य\-ञ्ज\-ना\-व\-रोध}}{\lemma{\skt{सर्व\appbrk{}व्य\appbrk{}ञ्ज\appbrk{}ना\appbrk{}व\appbrk{}रोध}}\Afootnote{K H; \skt{सर्प\appbrk{}व्य\appbrk{}ञ्ज\appbrk{}ना\appbrk{}व\appbrk{}रोधः~॥}\ A}}\skt{ }\edtext{\skt{इति}}{\lemma{\skt{इति}}\Afootnote{K H; \emph{om.}\ A}}\skt{~।}\pend

\pstart \ledsidenote{\passref{5.4.37}}\skt{तत्र }\edtext{\skt{दर्वी\-क\-र\-वि\-षेण}}{\lemma{\skt{दर्वी\appbrk{}क\appbrk{}र\appbrk{}वि\appbrk{}षेण}}\Afootnote{A K\textsuperscript{pc} H; \skt{दर्व\appbrk{}क\appbrk{}र\appbrk{}वि\appbrk{}षेण}\ K\textsuperscript{ac}}}\skt{ }\edtext{\skt{त्वङ्न\-ख\-न\-य\-न\-व\-द\-न\-मू\-त्र\-पु\-री\-ष\-दं\-श\-कृ\-ष्णत्वं}}{\lemma{\skt{त्वङ्न\appbrk{}ख\appbrk{}न\appbrk{}य\appbrk{}न\appbrk{}व\appbrk{}द\appbrk{}न\appbrk{}मू\appbrk{}त्र\appbrk{}पु\appbrk{}री\appbrk{}ष\appbrk{}दं\appbrk{}श\appbrk{}कृ\appbrk{}ष्णत्वं}}\Afootnote{H; \skt{त्वङ्न\appbrk{}य\appbrk{}न\appbrk{}न\appbrk{}ख\appbrk{}द\appbrk{}श\appbrk{}न\appbrk{}व\appbrk{}द\appbrk{}न\appbrk{}मू\appbrk{}त्र\appbrk{}पु\appbrk{}री\appbrk{}ष\appbrk{}दं\appbrk{}श\appbrk{}कृ\appbrk{}ष्णत्वं}\ A; \skt{त्वङ्न\appbrk{}ख\appbrk{}न\appbrk{}य\appbrk{}न\appbrk{}व\appbrk{}द\appbrk{}न\appbrk{}मू\appbrk{}त्र\appbrk{}पु\appbrk{}री\appbrk{}श\appbrk{}दं\appbrk{}श\appbrk{}कृ\appbrk{}ष्णत्वं}\ K}}\skt{ रौक्ष्यं }\edtext{\skt{सन्धिवेदना}}{\lemma{\skt{सन्धिवेदना}}\Afootnote{K H; \skt{शिरसो गौरवं सन्धिवेदना}\ A}}\skt{ }\edtext{\skt{शिरोगौरवं}}{\lemma{\skt{शिरोगौरवं}}\Afootnote{K H; \emph{om.}\ A}}\skt{ }\edtext{\skt{कटी\-पृ\-ष्ठ\-ग्री\-वा\-दौ\-र्बल्यं}}{\lemma{\skt{कटी\appbrk{}पृ\appbrk{}ष्ठ\appbrk{}ग्री\appbrk{}वा\appbrk{}दौ\appbrk{}र्बल्यं}}\Afootnote{A; \skt{कटी\appbrk{}पृ\appbrk{}ष्ठ\appbrk{}ग्री\appbrk{}वा\appbrk{}दौ\appbrk{}र्वल्यं}\ K H}}\skt{ जृम्भणं }\edtext{\skt{स्वरावसादः}}{\lemma{\skt{स्वरावसादः}}\Afootnote{K H; \skt{वेपथुः स्वरावसादो}\ A}}\skt{ }\edtext{\skt{खुरखुरको}}{\lemma{\skt{खुरखुरको}}\Afootnote{\skt{घुर्घुरको}\ A; \skt{खरुखरुको}\ K H}}\skt{ जडता शुष्कोद्गारः }\edtext{\skt{कासः}}{\lemma{\skt{कासः}}\Afootnote{K H; \skt{कास}\ A}}\skt{ }\edtext{\skt{श्वासो}}{\lemma{\skt{श्वासो}}\Afootnote{K H; \skt{श्वासौ}\ A}}\skt{ हिक्का }\edtext{\skt{वायो\-रु\-र्ध्व\-ग\-मनं}}{\lemma{\skt{उर्ध्वगमनं}}\Afootnote{K; \skt{ऊर्ध्वगमनं}\ A; \skt{ऊर्द्धगमनं~।}\ H}}\skt{ शूलोद्वेष्टनं }\edtext{\skt{कृष्ण\-ला\-ला\-स्र\-वणं}}{\lemma{\skt{कृष्ण\appbrk{}ला\appbrk{}ला\appbrk{}स्र\appbrk{}वणं}}\Afootnote{K H; \skt{तृष्णा लालास्रावः}\ A}}\skt{ }\edtext{\skt{फेनागमनं}}{\lemma{\skt{फेनागमनं}}\Afootnote{A K; \skt{फेणागमनं}\ H}}\skt{ }\edtext{\skt{स्रोतो\-व\-रो\-ध\-स्ता\-स्ताश्च}}{\lemma{\skt{स्रोतो\appbrk{}व\appbrk{}रो\appbrk{}ध\appbrk{}स्ता\appbrk{}स्ताश्च}}\Afootnote{K; \skt{स्रोतोऽ\appbrk{}व\appbrk{}रो\appbrk{}ध\appbrk{}स्ता\appbrk{}स्ताश्च}\ A; \skt{स्रोत्रो\appbrk{}व\appbrk{}रो\appbrk{}ध\appbrk{}स्ता\appbrk{}स्ताश्च}\ H}}\skt{ वातवेदना भवन्ति~॥ }\edtext{\skt{मण्डलिविषेण तु}}{\lemma{\skt{मण्डलिविषेण तु}}\Afootnote{H; \skt{मण्डलिविषेण}\ A; \skt{मण्डलविषेण तु}\ K}}\skt{ }\edtext{\skt{त्वङ्न\-ख\-न\-य\-न\-द\-श\-न\-व\-द\-न\-मू\-त्र\-पु\-री\-ष\-दं\-श\-पी\-तत्वं}}{\lemma{\skt{त्वङ्न\appbrk{}ख\appbrk{}न\appbrk{}य\appbrk{}न\appbrk{}द\appbrk{}श\appbrk{}न\appbrk{}व\appbrk{}द\appbrk{}न\appbrk{}मू\appbrk{}त्र\appbrk{}पु\appbrk{}री\appbrk{}ष\appbrk{}दं\appbrk{}श\appbrk{}पी\appbrk{}तत्वं}}\Afootnote{\skt{त्वगादीनां पीतत्वं}\ A; \skt{त्वङ्न\appbrk{}ख\appbrk{}न\appbrk{}य\appbrk{}न\appbrk{}द\appbrk{}श\appbrk{}न\appbrk{}व\appbrk{}द\appbrk{}न\appbrk{}मू\appbrk{}त्र\appbrk{}पु\appbrk{}री\appbrk{}श\appbrk{}दं\appbrk{}पी\appbrk{}तत्वं}\ K\textsuperscript{ac}; \skt{त्वङ्न\appbrk{}ख\appbrk{}न\appbrk{}य\appbrk{}न\appbrk{}द\appbrk{}श\appbrk{}न\appbrk{}व\appbrk{}द\appbrk{}न\appbrk{}मू\appbrk{}त्र\appbrk{}पु\appbrk{}री\appbrk{}श\appbrk{}दं\appbrk{}श\appbrk{}पी\appbrk{}तत्वं}\ K\textsuperscript{pc}; \skt{त्वङ्न\appbrk{}ख\appbrk{}न\appbrk{}य\appbrk{}न\appbrk{}व\appbrk{}द\appbrk{}न\appbrk{}मू\appbrk{}त्र\appbrk{}पु\appbrk{}री\appbrk{}ष\appbrk{}दं\appbrk{}श\appbrk{}पी\appbrk{}तत्वं}\ H}}\skt{ शीताभिलाषः }\edtext{\skt{परिधूपायनं}}{\lemma{\skt{परिधूपायनं}}\Afootnote{K H; \skt{परिधूपनं}\ A}}\skt{ }\edtext{\edtext{\skt{दाहस्तृष्णा}}{\lemma{\skt{दाहस्}}\Afootnote{K H; \skt{दाह}\ A}}}{\lemma{\skt{तृष्णा}}\Afootnote{K H; \skt{स्तृष्णा}\ A}}\skt{ मदो मूर्च्छा ज्वरः }\edtext{\skt{शोणि\-ता\-ग\-म\-न\-मू\-र्ध्व\-म\-धश्च}}{\lemma{\skt{ऊर्ध्वम्}}\Afootnote{A K; \skt{ऊर्द्धम्}\ H}}\skt{ }\edtext{\skt{मांस\-व\-सा\-व\-सादः}}{\lemma{\skt{मांस\appbrk{}व\appbrk{}सा\appbrk{}व\appbrk{}सादः}}\Afootnote{K; \skt{मांसा\appbrk{}ना\appbrk{}म\appbrk{}व\appbrk{}शा\appbrk{}तनं}\ A; \skt{मासावसादः~।}\ H}}\skt{ श्वय\-थु\-र्दं\-श\-कोथो }\edtext{\edtext{\skt{विप\-री\-त\-द\-र्श\-न\-मा\-तु\-र\-को\-प\-स्ता\-स्ताश्च}}{\lemma{\skt{विप\appbrk{}री\appbrk{}त\appbrk{}द\appbrk{}र्श\appbrk{}नम्}}\Afootnote{K H; \skt{पीत\appbrk{}रू\appbrk{}प\appbrk{}द\appbrk{}र्श\appbrk{}नम्}\ A}}}{\lemma{\skt{आतु\appbrk{}र\appbrk{}को\appbrk{}प\appbrk{}स्ता\appbrk{}स्ताश्}}\Afootnote{K H; \skt{आशु\appbrk{}को\appbrk{}प\appbrk{}स्ता\appbrk{}स्ताश्}\ A}}\skt{ पित्तवेदना भवन्ति~॥ }\edtext{\skt{राजीमद्विषेण}\skt{ तु}}{\lemma{\skt{राजीमद्विषेण तु}}\Afootnote{K H; \skt{राजिमद्विषेण}\ A}}\skt{ }\edtext{\skt{त्वङ्न\-ख\-न\-य\-न\-द\-श\-न\-व\-द\-न\-मू\-त्र\-पु\-री\-ष\-दं\-श\-पा\-ण्डुत्वं}}{\lemma{\skt{त्वङ्न\appbrk{}ख\appbrk{}न\appbrk{}य\appbrk{}न\appbrk{}द\appbrk{}श\appbrk{}न\appbrk{}व\appbrk{}द\appbrk{}न\appbrk{}मू\appbrk{}त्र\appbrk{}पु\appbrk{}री\appbrk{}ष\appbrk{}दं\appbrk{}श\appbrk{}पा\appbrk{}ण्डुत्वं}}\Afootnote{\skt{शुक्लत्वं त्व}\ A; \skt{त्वङ्न\appbrk{}ख\appbrk{}न\appbrk{}य\appbrk{}न\appbrk{}दं\appbrk{}श\appbrk{}न\appbrk{}व\appbrk{}द\appbrk{}न\appbrk{}मू\appbrk{}त्र\appbrk{}पु\appbrk{}री\appbrk{}शा\appbrk{}दं\appbrk{}श\appbrk{}पा\appbrk{}ण्डुत्वं}\ K\textsuperscript{ac}; \skt{त्वङ्न\appbrk{}ख\appbrk{}न\appbrk{}य\appbrk{}न\appbrk{}द\appbrk{}श\appbrk{}न\appbrk{}व\appbrk{}द\appbrk{}न\appbrk{}मू\appbrk{}त्र\appbrk{}पु\appbrk{}री\appbrk{}शा\appbrk{}दं\appbrk{}श\appbrk{}पा\appbrk{}ण्डुत्वं}\ K\textsuperscript{pc}; \skt{त्वङ्न\appbrk{}ख\appbrk{}न\appbrk{}य\appbrk{}न\appbrk{}द\appbrk{}श\appbrk{}न\appbrk{}व\appbrk{}द\appbrk{}न\appbrk{}मू\appbrk{}त्र\appbrk{}पु\appbrk{}री\appbrk{}षा\appbrk{}दं\appbrk{}श\appbrk{}पा\appbrk{}ण्डुत्वं}\ H}}\skt{ }\edtext{\skt{शीतज्वरो}}{\lemma{\skt{शीतज्वरो}}\Afootnote{K H; \skt{गादीनां शीतज्वरो}\ A}}\skt{ }\edtext{\skt{रोमहर्षः}}{\lemma{\skt{रोमहर्षः}}\Afootnote{A; \skt{रोमहर्ष}\ K H}}\skt{ }\edtext{\skt{स्तब्धत्वं}}{\lemma{\skt{स्तब्धत्वं}}\Afootnote{A; \skt{स्तव्धत्वं}\ K H}}\skt{ गात्रा\-णा\-मा\-दं\-श\-शोफः }\edtext{\skt{सान्द्र\-क\-फ\-प्र\-से\-क\-श्छ\-र्दि\-रक्ष्णोः}}{\lemma{\skt{अक्ष्णोः}}\Afootnote{K; \skt{अभीक्ष्णमक्ष्णोः}\ A; \skt{अ}\ H}}\skt{ }\edtext{\skt{कण्डू}}{\lemma{\skt{कण्डू}}\Afootnote{K; \skt{कण्डूः}\ A; \emph{om.}\ H}}\skt{ }\edtext{\skt{खुरखुरकः}}{\lemma{\skt{खुरखुरकः}}\Afootnote{K; \skt{कण्ठे श्वय\appbrk{}थु\appbrk{}र्घु\appbrk{}र्घु\appbrk{}रक}\ A; \skt{क्ष्णोः खुरुखुरुकः~।}\ H}}\skt{ }\edtext{\skt{उच्छ्वा\-स\-नि\-रो\-ध\-स्ता\-स्ताश्च}}{\lemma{\skt{तास्ताश्च}}\Afootnote{K H; \skt{तमः प्रवेशस्तास्ताश्च}\ A}}\skt{ कफवेदना भवन्ति~॥}\pend

\pstart \ledsidenote{\passref{5.4.38}}\edtext{\skt{तत्र}}{\lemma{\skt{तत्र}}\Afootnote{K H; \emph{om.}\ A}}\skt{ }\edtext{\skt{पुरुषेण}}{\lemma{\skt{पुरुषेण}}\Afootnote{K H; \skt{पुरुषाभि}\ A}}\skt{ दष्ट }\edtext{\skt{ऊर्ध्वं}}{\lemma{\skt{ऊर्ध्वं}}\Afootnote{A K; \skt{उर्द्ध्वम्}\ H}}\skt{ प्रेक्षते~। }\edtext{\skt{स्त्रिया}}{\lemma{\skt{स्त्रिया}}\Afootnote{\skt{अधस्तात्स्त्रिया}\ A; \emph{om.}\ K H}}\skt{ }\edtext{\skt{तिर्य\-ङ्न\-पुं\-स\-के\-नाध}}{\lemma{\skt{तिर्य\appbrk{}ङ्न\appbrk{}पुं\appbrk{}स\appbrk{}के\appbrk{}नाध}}\Afootnote{\skt{सिराश्चोत्तिष्ठन्ति ललाटे, नपुं\appbrk{}स\appbrk{}का\appbrk{}भि\appbrk{}द\appbrk{}ष्टस्}\ A; \skt{स्त्रिया\appbrk{}स्ति\appbrk{}र्य\appbrk{}ङ्न\appbrk{}पुं\appbrk{}स\appbrk{}के\appbrk{}नाध}\ K H}}\skt{ }\edtext{\skt{इति}}{\lemma{\skt{इति}}\Afootnote{K H; \skt{तिर्यक्प्रेक्षी भवति,}\ A}}\skt{~। }\edtext{\skt{गर्भिण्या}}{\lemma{\skt{गर्भिण्या}}\Afootnote{A K; \skt{गर्ब्भिण्याः}\ H}}\skt{ }\edtext{\skt{पाण्डुमुखो}}{\lemma{\skt{पाण्डुमुखो}}\Afootnote{K H; \skt{पाण्डुमुखो}\ A}}\skt{ }\edtext{\skt{आध्मातश्च}}{\lemma{\skt{आध्मातश्च}}\Afootnote{K H; \skt{ध्मातश्च,}\ A}}\skt{ }\edtext{\skt{भवति}}{\lemma{\skt{भवति}}\Afootnote{K H; \emph{om.}\ A}}\skt{~। सूतिकया कुक्षिशूलार्त्तः सरुधिरं }\edtext{\skt{मेहति}}{\lemma{\skt{मेहति}}\Afootnote{K H; \skt{मेह\appbrk{}त्यु\appbrk{}प\appbrk{}जि\appbrk{}ह्विका चास्य भवति,}\ A}}\skt{~। }\edtext{\edtext{\skt{ग्रासा\-र्थि\-ना\-न्न\-मा\-कां\-क्षति}}{\lemma{\skt{ग्रासार्थिनान्नम्}}\Afootnote{\skt{ग्रासार्थिनाऽन्}\ A; \skt{ग्रासार्थिमान्नम्}\ K; \skt{ग्रासार्थिनोऽन्नम्}\ H}}}{\lemma{\skt{आकांक्षति}}\Afootnote{K H; \skt{नं काङ्क्षति,}\ A}}\skt{~। वृद्धेन }\edtext{\skt{चिरान्मन्दाश्च}}{\lemma{\skt{चिरान्मन्दाश्च}}\Afootnote{A K; \skt{चिरात्मन्दाश्च}\ H}}\skt{ वेगा }\edtext{\skt{भवन्ति}}{\lemma{\skt{भवन्ति}}\Afootnote{K H; \emph{om.}\ A}}\skt{~। }\edtext{\edtext{\skt{बाले\-ना\-शु\-स्ती\-क्ष्णश्च}}{\lemma{\skt{बालेनाशुस्}}\Afootnote{\skt{बालेनाशु}\ A; \skt{वालेनाशुस्}\ K H}}}{\lemma{\skt{तीक्ष्णश्च}}\Afootnote{K H; \skt{मृदवश्च,}\ A}}\skt{~। निर्वि\-षे\-णा\-वि\-ष\-लि\-ङ्गम्~। }\edtext{\skt{अन्धा\-हि\-के\-ना\-न्ध\-त्व\-मेके}}{\lemma{\skt{एके}}\Afootnote{K H; \skt{इत्येके,}\ A}}\skt{~। ग्रस\-ना\-द\-ज\-गरः }\edtext{\skt{प्राणहरो न}}{\lemma{\skt{प्राणहरो न}}\Afootnote{K H; \skt{शरी\appbrk{}र\appbrk{}प्रा\appbrk{}ण\appbrk{}हरो न}\ A}}\skt{ }\edtext{\skt{विषादिति}}{\lemma{\skt{विषादिति}}\Afootnote{K H; \skt{विसात्~। तत्र सद्यः प्राण\appbrk{}ह\appbrk{}रा\appbrk{}हि\appbrk{}दष्टः पतति शास्त्राशनिहत एव भूमौ, स्रस्ताङ्गः स्वपिति~॥}\ A}}\skt{~॥}\pend

\pstart \ledsidenote{\passref{5.4.39}}\skt{तत्र }\edtext{\skt{सर्व\-स\-र्प\-वि\-षाणां}}{\lemma{\skt{सर्व\appbrk{}स\appbrk{}र्प\appbrk{}वि\appbrk{}षाणां}}\Afootnote{K H; \skt{सर्वेषां सर्पाणां विषस्य}\ A}}\skt{ }\edtext{\skt{सप्त विषवेगा}}{\lemma{\skt{सप्त विषवेगा}}\Afootnote{K H; \skt{सप्तवेगा}\ A}}\skt{ भवन्ति~। तत्र दर्वीकराणां प्रथमे }\edtext{\skt{वेगे}}{\lemma{\skt{वेगे}}\Afootnote{A H; \skt{वेग}\ K}}\skt{ विषं शोणितं दूषयति~। तत्प्रदुष्टं कृष्णतामुपैति~। तेन }\edtext{\skt{कार्ष्ण्यं}}{\lemma{\skt{कार्ष्ण्यं}}\Afootnote{A; \skt{कार्ष्ण्य}\ K H}}\skt{ }\edtext{\edtext{\skt{पिपी\-लि\-का\-प\-रि\-स\-र्प\-ण\-मिव}}{\lemma{\skt{पिपी\appbrk{}लि\appbrk{}का\appbrk{}प\appbrk{}रि\appbrk{}स\appbrk{}र्प\appbrk{}णम्}}\Afootnote{A H; \skt{पिपी\appbrk{}लि\appbrk{}का\appbrk{}प\appbrk{}रि\appbrk{}सर्प ञ्च ण म्}\ K}}}{\lemma{\skt{इव}}\Afootnote{A K H\textsuperscript{pc}; \skt{इवा}\ H\textsuperscript{ac}}}\skt{ चाङ्गे }\edtext{\skt{भवन्ति}}{\lemma{\skt{भवन्ति}}\Afootnote{K H; \skt{भवति;}\ A}}\skt{~॥ द्वितीये मांसं दूषयति~। }\edtext{\skt{तेना\-त्य\-र्थ\-कृ\-ष्णता}}{\lemma{\skt{तेना\appbrk{}त्य\appbrk{}र्थ\appbrk{}कृ\appbrk{}ष्णता}}\Afootnote{K H; \skt{तेनात्यर्थं कृष्णता}\ A}}\skt{ }\edtext{\edtext{\skt{ग्रन्थयश्च}}{\lemma{\skt{ग्रन्थयश्}}\Afootnote{K H; \skt{शोफो ग्रन्थयश्चाङ्}\ A}}}{\lemma{\skt{च}}\Afootnote{K H; \skt{गे}\ A}}\skt{ भवन्ति~॥ तृतीये }\edtext{\skt{मेदो}}{\lemma{\skt{मेदो}}\Afootnote{A K\textsuperscript{pc} H; \skt{मे}\ K\textsuperscript{ac}}}\skt{ दूषयति~। तेन }\edtext{\skt{दंशक्लेदः}}{\lemma{\skt{दंशक्लेदः}}\Afootnote{A K H\textsuperscript{pc}; \skt{दंशःक्लेदः}\ H\textsuperscript{ac}}}\skt{ शिरोगौरवं }\edtext{\skt{चक्षुर्ग्रहणञ्च}}{\lemma{\skt{चक्षुर्ग्रहणञ्च}}\Afootnote{K; \skt{स्वेद\appbrk{}श्च\appbrk{}क्षु\appbrk{}र्ग्र\appbrk{}हणं च;}\ A; \skt{चक्षुग्रहणञ्च}\ H}}\skt{ }\edtext{\skt{भवति}}{\lemma{\skt{भवति}}\Afootnote{K H; \emph{om.}\ A}}\skt{~। चतुर्थे }\edtext{\skt{कोष्ठ\-म\-नु\-प्र\-वि\-शति}}{\lemma{\skt{अनुप्रविशति}}\Afootnote{K H; \skt{अनुप्रविश्य}\ A}}\skt{~। }\edtext{\skt{ततः}}{\lemma{\skt{ततः}}\Afootnote{K H; \emph{om.}\ A}}\skt{ }\edtext{\edtext{\skt{कफ\-प्र\-भ\-वा\-न्दो\-षा\-न्को\-प\-यति}}{\lemma{\skt{कफप्रभवान्}}\Afootnote{K H; \skt{कफप्रधानान्}\ A}}}{\lemma{\skt{दोषान्कोपयति}}\Afootnote{H; \skt{दोषान्दूषयति,}\ A; \skt{दोषां कोपयति}\ K}}\skt{ तेन }\edtext{\skt{तन्द्री\-क\-फ\-प्र\-सेकः}}{\lemma{\skt{तन्द्री\appbrk{}क\appbrk{}फ\appbrk{}प्र\appbrk{}सेकः}}\Afootnote{K H; \skt{तन्द्राप्रसेक}\ A}}\skt{ }\edtext{\skt{सन्धिविश्लेषश्च}}{\lemma{\skt{सन्धिविश्लेषश्च}}\Afootnote{K H; \skt{सन्धिविश्लेषा}\ A}}\skt{ }\edtext{\skt{भवति}}{\lemma{\skt{भवति}}\Afootnote{K H; \skt{भवन्ति;}\ A}}\skt{~॥ }\edtext{\skt{पञ्च\-मेऽ\-स्थी\-न्य\-नु\-प्र\-वि\-शति}}{\lemma{\skt{ऽस्थीन्य्}}\Afootnote{A H; \skt{स्थीन्य्}\ K}}\skt{ }\edtext{\skt{तेन}}{\lemma{\skt{तेन}}\Afootnote{K H\textsuperscript{pc}; \skt{प्राणमग्निं च दूषयति, तेन}\ A; \skt{तेन सर्प्प}\ H\textsuperscript{ac}}}\skt{ पर्वभेदो हिक्का दाहश्च भवति~॥ षष्ठे }\edtext{\skt{मज्जा\-म\-नु\-प्र\-वि\-शति}}{\lemma{\skt{मज्जाम्}}\Afootnote{K; \skt{मज्जानम्}\ A H}}\skt{~। }\edtext{\skt{तेन}}{\lemma{\skt{तेन}}\Afootnote{K H; \emph{om.}\ A}}\skt{ }\edtext{\skt{ग्रहणीदोषा}}{\lemma{\skt{ग्रहणीदोषा}}\Afootnote{K H; \skt{ग्रहणीं चात्यर्थं दूषयति, तेन}\ A}}\skt{ }\edtext{\skt{गात्र\-गौ\-र\-व\-म\-ती\-सारो}}{\lemma{\skt{गात्रगौरवम्}}\Afootnote{K H; \skt{गात्राणां गौरवम्}\ A}}\skt{ हृत्पीडा मूर्च्छा च भवति~॥ सप्तमे शुक्र\-म\-नु\-प्र\-वि\-शति व्यानञ्चात्यर्थं कोपयति कफञ्च }\edtext{\skt{सूक्ष्मं}}{\lemma{\skt{सूक्ष्मं}}\Afootnote{K H; \skt{सूक्ष्म}\ A}}\skt{ स्रोतोभ्यः प्रच्यावयति~। तेन }\edtext{\skt{श्लेष्मप्रादुर्भावः}}{\lemma{\skt{श्लेष्मप्रादुर्भावः}}\Afootnote{K; \skt{श्लेष्म\appbrk{}व\appbrk{}र्ति\appbrk{}प्रा\appbrk{}दु\appbrk{}र्भावः}\ A; \skt{श्लेष्मप्रादुर्व्भावः}\ H}}\skt{ }\edtext{\skt{कटी\-पृ\-ष्ठ\-स्क\-न्द\-भङ्गः}}{\lemma{\skt{कटी\appbrk{}पृ\appbrk{}ष्ठ\appbrk{}स्क\appbrk{}न्द\appbrk{}भङ्गः}}\Afootnote{K; \skt{कटीपृष्ठभङ्गः}\ A; \skt{कटीपृष्ठस्कन्दे भङ्गाः}\ H\textsuperscript{ac}; \skt{कटी\appbrk{}स्त\appbrk{}म्भ\appbrk{}पृ\appbrk{}ष्ठ\appbrk{}स्कन्दे भङ्गाः}\ H\textsuperscript{pc}}}\skt{ }\edtext{\skt{सर्व\-चे\-ष्ठा\-वि\-घातः}}{\lemma{\skt{सर्व\appbrk{}चे\appbrk{}ष्ठा\appbrk{}वि\appbrk{}घातः}}\Afootnote{K; \skt{सर्व\appbrk{}चे\appbrk{}ष्टा\appbrk{}वि\appbrk{}घातो लाला\appbrk{}स्वे\appbrk{}द\appbrk{}यो\appbrk{}रति}\ A; \skt{सर्व्व\appbrk{}चे\appbrk{}ष्टा\appbrk{}वि\appbrk{}घातः~।}\ H}}\skt{ }\edtext{\skt{उच्छ्वासविरोधो}}{\lemma{\skt{उच्छ्वासविरोधो}}\Afootnote{K H; \skt{प्रवृ\appbrk{}त्ति\appbrk{}रु\appbrk{}च्छ्वा\appbrk{}स\appbrk{}नि\appbrk{}रो\appbrk{}धश्}\ A}}\skt{ }\edtext{\skt{भवतीति}}{\lemma{\skt{भवतीति}}\Afootnote{K; \skt{च भवति~।}\ A; \skt{भवति~॥}\ H}}\skt{~॥ }\edtext{\skt{मण्डलिनान्तु}}{\lemma{\skt{मण्डलिनान्तु}}\Afootnote{K H; \skt{मण्डलिनां}\ A}}\skt{ प्रथमे वेगे }\edtext{\skt{विषः}}{\lemma{\skt{विषः}}\Afootnote{K H; \skt{विषं}\ A}}\skt{ शोणितं दूषयति~। तत्प्रदुष्टं पीततामुपैति~। }\edtext{\skt{तेन}}{\lemma{\skt{तेन}}\Afootnote{K H; \skt{तत्र परिदाहः}\ A}}\skt{ पीतावभासता }\edtext{\skt{परिदाहश्च}}{\lemma{\skt{परिदाहश्च}}\Afootnote{K H; \skt{चाङ्गानां}\ A}}\skt{ भवति~॥ द्वितीये मांसं दूषयति~। }\edtext{\skt{तेन}}{\lemma{\skt{तेन}}\Afootnote{K; \skt{तेनात्यर्थं पी}\ A; \skt{तेना\appbrk{}त्य\appbrk{}र्थ\appbrk{}पी\appbrk{}ताङ्ग}\ H}}\skt{ }\edtext{\skt{चात्य\-र्थ\-पी\-ता\-ङ्गता}}{\lemma{\skt{चात्य\appbrk{}र्थ\appbrk{}पी\appbrk{}ता\appbrk{}ङ्गता}}\Afootnote{K; \skt{त}\ A; \skt{ता चात्य\appbrk{}र्थ\appbrk{}पी\appbrk{}ता\appbrk{}ङ्गता}\ H}}\skt{ }\edtext{\skt{चात्य\-र्थ\-प\-रि\-दाहो}}{\lemma{\skt{चात्य\appbrk{}र्थ\appbrk{}प\appbrk{}रि\appbrk{}दाहो}}\Afootnote{H; \skt{ता परिदाहो}\ A; \skt{श्वर्थपरिदाहो}\ K}}\skt{ }\edtext{\skt{दंश\-श्व\-य\-थु\-र्भ\-वति}}{\lemma{\skt{दंशश्वयथुर्}}\Afootnote{K\textsuperscript{pc}; \skt{दंशे श्वयथुश्च}\ A; \skt{दंशश्च श्वयथुर्}\ K\textsuperscript{ac}; \skt{दंशश्वयथुश्च}\ H}}\skt{~॥ तृतीये मेदो दूषयति तेन }\edtext{\skt{कृष्णादंशक्लेदः}}{\lemma{\skt{कृष्णादंशक्लेदः}}\Afootnote{K H; \skt{पूर्व\appbrk{}व\appbrk{}च्च\appbrk{}क्षु\appbrk{}र्ग्र\appbrk{}हणं तृष्णा दंशक्लेदः}\ A}}\skt{ स्वेदश्च }\edtext{\skt{भवति}}{\lemma{\skt{भवति}}\Afootnote{K H; \emph{om.}\ A}}\skt{~॥ चतुर्थे }\edtext{\skt{पूर्व\-व\-द\-नु\-प्र\-विश्य}}{\lemma{\skt{पूर्व\appbrk{}व\appbrk{}द\appbrk{}नु\appbrk{}प्र\appbrk{}विश्य}}\Afootnote{\skt{कोष्ठ\appbrk{}म\appbrk{}नु\appbrk{}प्र\appbrk{}विश्य}\ A; \skt{पूर्व\appbrk{}व\appbrk{}द\appbrk{}नु\appbrk{}प्र\appbrk{}विश्य}\ K H}}\skt{ ज्वर\-मा\-पा\-द\-यति~॥ पञ्चमे }\edtext{\skt{दाहं}}{\lemma{\skt{दाहं}}\Afootnote{K H; \skt{परिदाहं}\ A}}\skt{ सर्वगात्रेषु करोति~। षष्ठसप्तमयोः }\edtext{\skt{पूर्ववदिति}}{\lemma{\skt{पूर्ववदिति}}\Afootnote{K H; \skt{पूर्वव}\ A}}\skt{~। }\edtext{\skt{राजीमतां तु}}{\lemma{\skt{राजीमतां तु}}\Afootnote{K H; \skt{त्~। राजिमतां}\ A}}\skt{ प्रथमे वेगे }\edtext{\skt{शोणि\-त\-न्दू\-ष\-यति}}{\lemma{\skt{शोणितन्}}\Afootnote{K H; \skt{विषं शोणितं}\ A}}\skt{~॥ तत्प्रदुष्टं पाण्डुतामुपैति तेन रोमहर्षः }\edtext{\skt{पाण्ड्वावभासश्च}}{\lemma{\skt{पाण्ड्वावभासश्च}}\Afootnote{K H; \skt{शुक्लावभासश्च}\ A}}\skt{ पुरुषो भवति~॥ द्वितीये मांसं दूषयति तेन }\edtext{\edtext{\edtext{\skt{पाण्डु\-र\-त्य\-र्थ\-जा\-ड्यञ्च}}{\lemma{\skt{पाण्डुर्}}\Afootnote{K H; \skt{पाण्डुता}\ A}}}{\lemma{\skt{अत्यर्थजाड्यञ्}}\Afootnote{K H; \skt{ऽत्यर्थं जाड्यं}\ A}}}{\lemma{\skt{च}}\Afootnote{K H; \skt{शिरःशोफश्च}\ A}}\skt{ भवति~। तृतीये मेदो }\edtext{\skt{दूषयति}}{\lemma{\skt{दूषयति}}\Afootnote{A K\textsuperscript{pc} H; \skt{षयति}\ K\textsuperscript{ac}}}\skt{ तेन }\edtext{\edtext{\skt{दंश\-क्ले\-दोऽ\-क्षि\-ना\-सा\-स्रा\-वश्च}}{\lemma{\skt{दंशक्लेदो}}\Afootnote{K H; \skt{चक्षुर्ग्रहणं दंशक्लेदः}\ A}}}{\lemma{\skt{ऽक्षि\appbrk{}ना\appbrk{}सा\appbrk{}स्रा\appbrk{}वश्च}}\Afootnote{K H; \skt{स्वेदो घ्राणाक्षिस्रावश्च}\ A}}\skt{ भवति~॥ चतुर्थे }\edtext{\skt{पूर्व\-व\-द\-नु\-प्र\-विश्य}}{\lemma{\skt{पूर्व\appbrk{}व\appbrk{}द\appbrk{}नु\appbrk{}प्र\appbrk{}विश्य}}\Afootnote{K H; \skt{कोष्ठ\appbrk{}म\appbrk{}नु\appbrk{}प्र\appbrk{}विश्य}\ A}}\skt{ }\edtext{\skt{मन्या\-स्त\-म्भ\-शि\-रो\-गौ\-र\-व\-ञ्चा\-पा\-द\-यति}}{\lemma{\skt{मन्या\appbrk{}स्त\appbrk{}म्भ\appbrk{}शि\appbrk{}रो\appbrk{}गौ\appbrk{}र\appbrk{}वञ्}}\Afootnote{K H; \skt{मन्यास्तम्भं शिरोगौरवं}\ A}}\skt{~॥ पञ्चमे }\edtext{\skt{वाक्सङ्गः}}{\lemma{\skt{वाक्सङ्गः}}\Afootnote{K H; \skt{वाक्सङ्गं}\ A}}\skt{ }\edtext{\skt{शीतज्वरश्च}}{\lemma{\skt{शीतज्वरश्च}}\Afootnote{\skt{शीतज्वरं च}\ A K H}}\skt{~॥ }\edtext{\skt{षष्ठसप्तमयोः}}{\lemma{\skt{षष्ठसप्तमयोः}}\Afootnote{K H; \skt{करोति; षष्ठसप्तमयोः}\ A}}\skt{ }\edtext{\skt{पूर्ववदिति}}{\lemma{\skt{पूर्ववदिति}}\Afootnote{A H; \skt{पूवदिति~॥}\ K}}\skt{~॥}\pend

\pstart\versepar \ledsidenote{\passref{5.4.40}}\edtext{\skt{भवन्ति}}{\lemma{\skt{भवन्ति}}\Afootnote{A H; \skt{भवति~॥}\ K}}\skt{ }\edtext{\skt{चात्र}}{\lemma{\skt{चात्र}}\Afootnote{A H; \emph{om.}\ K}}\skt{ }\edtext{\skt{श्लोकाः}}{\lemma{\skt{श्लोकाः}}\Afootnote{H; \emph{om.}\ A K}}\skt{~॥ }\pend
\pstart\versepar \skt{धात्वन्तरेषु याः सप्त कलाः संपरिकीर्त्तिताः }\pend
\pstart\versepar \skt{}\edtext{\edtext{\skt{तास्वे\-कै\-क\-म\-ति\-क्रम्य}}{\lemma{\skt{तास्व्}}\Afootnote{A; \skt{ताःस्व्}\ K H}}}{\lemma{\skt{एकैकम्}}\Afootnote{K H; \skt{एकैकाम्}\ A}}\skt{ वेगं प्रकुरुते }\edtext{\skt{विषः}}{\lemma{\skt{विषः}}\Afootnote{K H; \skt{विषम्~॥}\ A}}\skt{~॥}\pend

\pstart\versepar \ledsidenote{\passref{5.4.41}}\skt{येनान्तरेण तु }\edtext{\skt{कलाः}}{\lemma{\skt{कलाः}}\Afootnote{K H; \skt{कलां}\ A}}\skt{ कालकल्पं भिनत्ति }\edtext{\skt{ह}}{\lemma{\skt{ह}}\Afootnote{K H\textsuperscript{ac}; \skt{हि~।}\ A; \skt{वदन्ति ह~।}\ H\textsuperscript{pc}}}\skt{~। }\pend
\pstart\versepar \skt{समी\-र\-णे\-नो\-ह्य\-मानं तत्तु वेगान्तरं }\edtext{\skt{मतम्}}{\lemma{\skt{मतम्}}\Afootnote{K H; \skt{स्मृतम्~॥}\ A}}\skt{~॥}\pend

\pstart\versepar \ledsidenote{\passref{5.4.42}}\skt{शूनाङ्गः प्रथमे वेगे पशुः }\edtext{\skt{प्रध्याति}}{\lemma{\skt{प्रध्याति}}\Afootnote{K H; \skt{ध्यायति}\ A}}\skt{ दुःखितः~॥ }\pend
\pstart\versepar \skt{}\edtext{\skt{द्वितीये}}{\lemma{\skt{द्वितीये}}\Afootnote{K H; \skt{लालास्रावो द्वितीये}\ A}}\skt{ }\edtext{\skt{लालि\-मा\-न्कि\-ञ्चि\-द्धृ\-ष्टाङ्गः}}{\lemma{\skt{लालि\appbrk{}मा\appbrk{}न्कि\appbrk{}ञ्चि\appbrk{}द्धृ\appbrk{}ष्टाङ्गः}}\Afootnote{\skt{तु कृष्णाङ्गः}\ A; \skt{लालि\appbrk{}मा\appbrk{}न्कि\appbrk{}ञ्चि\appbrk{}द्धृ\appbrk{}ष्टाङ्गः}\ K; \skt{लालि\appbrk{}मा\appbrk{}न्कि\appbrk{}ञ्चि\appbrk{}द्धृ\appbrk{}ष्टाङ्गाः}\ H\textsuperscript{ac}; \skt{लालि\appbrk{}मा\appbrk{}न्कि\appbrk{}ञ्चि\appbrk{}द्दृ\appbrk{}ष्टाङ्गाः}\ H\textsuperscript{pc}}}\skt{ पीड्यते हृदि~।}\pend

\pstart\versepar \ledsidenote{\passref{5.4.43}}\edtext{\skt{तृतीयस्य}}{\lemma{\skt{तृतीयस्य}}\Afootnote{K H; \skt{तृतीये च}\ A}}\skt{ शिरोदुःखं }\edtext{\skt{कर्ण्णग्रीवाञ्च}}{\lemma{\skt{कर्ण्णग्रीवाञ्च}}\Afootnote{K H\textsuperscript{ac}; \skt{कण्ठग्रीवं च}\ A; \skt{कण्ठग्रीवाञ्च}\ H\textsuperscript{pc}}}\skt{ भज्यते~। }\pend
\pstart\versepar \skt{चतुर्थे वेपते मूढः }\edtext{\skt{खाद\-न्द\-न्ता\-ञ्ज\-हा\-त्य\-सून्}}{\lemma{\skt{दन्ताञ्जहात्य्}}\Afootnote{K; \skt{दन्तान्जहात्य्}\ A; \skt{दन्तान्जहत्य्}\ H}}\skt{~।}\pend

\pstart\versepar \ledsidenote{\passref{5.4.44}}\edtext{\skt{केचिद्वेगत्रयं}}{\lemma{\skt{केचिद्वेगत्रयं}}\Afootnote{A K H\textsuperscript{pc}; \skt{केचि\appbrk{}द्वे\appbrk{}द\appbrk{}त्र\appbrk{}यम्}\ H\textsuperscript{ac}}}\skt{ }\edtext{\skt{प्राहु\-र\-न्तः\-स्वे\-देषु}}{\lemma{\skt{अन्तःस्वेदेषु}}\Afootnote{H; \skt{अन्तं चैतेषु}\ A; \skt{अन्तस्वेदेषु}\ K}}\skt{ }\edtext{\skt{तद्विदः}}{\lemma{\skt{तद्विदः}}\Afootnote{A K; \skt{तद्विदुः~।}\ H}}\skt{~॥ }\pend
\pstart\versepar \skt{}\edtext{\skt{वेगे तु}}{\lemma{\skt{वेगे तु}}\Afootnote{H; \skt{ध्यायति}\ A; \skt{वेदे तु}\ K}}\skt{ प्रथमे }\edtext{\skt{पक्षी}}{\lemma{\skt{पक्षी}}\Afootnote{K H; \skt{वेगे पक्षी}\ A}}\skt{ }\edtext{\skt{ध्याति}}{\lemma{\skt{ध्याति}}\Afootnote{K H; \emph{om.}\ A}}\skt{ मुह्यत्यतः परम्~॥}\pend

\pstart\versepar \ledsidenote{\passref{5.4.45}}\skt{द्वितीये विह्वलः }\edtext{\skt{कूजन्पक्षी}}{\lemma{\skt{कूजन्पक्षी}}\Afootnote{K H; \emph{om.}\ A}}\skt{ }\edtext{\edtext{\skt{मरणमर्च्छति}}{\lemma{\skt{मरणम्}}\Afootnote{K H; \skt{प्रोक्तस्}\ A}}}{\lemma{\skt{अर्च्छति}}\Afootnote{K H; \skt{तृतीये मृत्युमृच्छति~।}\ A}}\skt{~। }\pend
\pstart\versepar \skt{केचिदेकं विहङ्गेषु }\edtext{\edtext{\skt{विष\-वे\-ग\-मु\-शन्ति}}{\lemma{\skt{विषवेगम्}}\Afootnote{A; \skt{विषवेषम्}\ K; \skt{विषवेगेम्}\ H}}\skt{ वै}}{\lemma{\skt{उशन्ति वै}}\Afootnote{\skt{उशन्ति हि~।}\ A; \skt{उषन्ति वै~॥}\ K H}}\skt{~॥ }\pend
\pstart\versepar \skt{मार्जा\-र\-न\-कु\-ला\-दीनां विषं }\edtext{\skt{नातिप्रवर्तत}}{\lemma{\skt{नातिप्रवर्तत}}\Afootnote{K H; \skt{नातिप्रवर्तते~॥}\ A}}\skt{ }\edtext{\skt{इति}}{\lemma{\skt{इति}}\Afootnote{K H; \emph{om.}\ A}}\skt{~॥}\pend

\endnumbering
\beginnumbering
\eledsection*{Adhyāya 5}
\pstart \ledsidenote{\passref{5.5.1}}\skt{अथातः }\edtext{\skt{सर्प\-द\-ष्ट\-चि\-कि\-त्सितं}}{\lemma{\skt{सर्प\appbrk{}द\appbrk{}ष्ट\appbrk{}चि\appbrk{}कि\appbrk{}त्सितं}}\Afootnote{K H; \skt{सर्प\appbrk{}द\appbrk{}ष्ट\appbrk{}वि\appbrk{}ष\appbrk{}चि\appbrk{}कि\appbrk{}त्सितं}\ A}}\skt{ }\edtext{\skt{कल्पं}}{\lemma{\skt{कल्पं}}\Afootnote{A K; \emph{om.}\ H}}\skt{ }\edtext{\skt{व्याख्यास्यामः}}{\lemma{\skt{व्याख्यास्यामः}}\Afootnote{A H; \emph{om.}\ K}}\skt{~॥}\pend

\pstart\versepar \ledsidenote{\passref{5.5.3}}\skt{सर्वैरेवादितः सर्पैः शाखादष्टस्य देहिनः~। }\pend
\pstart\versepar \skt{}\edtext{\edtext{\skt{बध्नी\-या\-द्गा\-ढ\-मु\-परि}}{\lemma{\skt{बध्नीयाद्गाढम्}}\Afootnote{\skt{दंशस्योपरि बध्नीया}\ A; \skt{वध्नीयाद्गाढम्}\ K H}}}{\lemma{\skt{उपरि}}\Afootnote{K H; \skt{दरि}\ A}}\skt{ }\edtext{\skt{दंशात्तु}}{\lemma{\skt{दंशात्तु}}\Afootnote{K H; \skt{ष्टाश्}\ A}}\skt{ }\edtext{\skt{चतुरङ्गुलम्}}{\lemma{\skt{चतुरङ्गुलम्}}\Afootnote{K H; \skt{चतुरङ्गुले~॥}\ A}}\skt{~।}\pend

\pstart\versepar \ledsidenote{\passref{5.5.4}}\skt{प्लोत\-च\-र्मा\-न्त\-व\-ल्कानां मृदु\-ना\-न्य\-त\-मेन वा~। }\pend
\pstart\versepar \skt{न }\edtext{\skt{पर्येति}}{\lemma{\skt{पर्येति}}\Afootnote{K H; \skt{गच्छति}\ A}}\skt{ }\edtext{\skt{विषं}}{\lemma{\skt{विषं}}\Afootnote{A K\textsuperscript{pc}; \skt{चाविषं}\ K\textsuperscript{ac}; \skt{च विषन्}\ H}}\skt{ }\edtext{\skt{देह\-म\-रि\-ष्टा\-भि\-र्नि\-वा\-रि\-तम्}}{\lemma{\skt{निवारितम्}}\Afootnote{A K; \skt{न्निवारतः~॥}\ H}}\skt{~।}\edtext{}{\lemma{\skt{मृदु\appbrk{}ना\appbrk{}न्य\appbrk{}त\appbrk{}मेन वा~। न}}\Afootnote{K; \skt{मृदु\appbrk{}नाऽ\appbrk{}न्य\appbrk{}त\appbrk{}मेन वै~। न}\ A; \skt{मृदु\appbrk{}ना\appbrk{}न्य\appbrk{}त\appbrk{}मेन वा~।}\ H}}\pend

\pstart\versepar \ledsidenote{\passref{5.5.5}}\edtext{\skt{दहे\-द्दं\-श\-म\-थो\-द्धृत्य}}{\lemma{\skt{अथोद्धृत्य}}\Afootnote{K H; \skt{अथोत्कृत्य}\ A}}\skt{ यत्र }\edtext{\skt{बन्धो न}}{\lemma{\skt{बन्धो न}}\Afootnote{A; \skt{वन्धो न}\ K H}}\skt{ जायते~। }\pend
\pstart\versepar \skt{}\edtext{\skt{आचू\-ष\-ण\-च्छे\-द\-दाहाः}}{\lemma{\skt{आचू\appbrk{}ष\appbrk{}ण\appbrk{}च्छे\appbrk{}द\appbrk{}दाहाः}}\Afootnote{A K; \skt{आचू\appbrk{}ष\appbrk{}ण\appbrk{}च्छे\appbrk{}द\appbrk{}दाहः}\ H}}\skt{ सर्वत्रैव तु पूजिताः~।}\pend

\pstart\versepar \ledsidenote{\passref{5.5.6}}\skt{प्रतिपूर्य मुखं }\edtext{\skt{पांशो\-र्हि\-त\-मा\-चू\-षणं}}{\lemma{\skt{पांशोर्}}\Afootnote{K H; \skt{वस्त्रैर्}\ A}}\skt{ भवेत्~। }\pend
\pstart\versepar \skt{}\edtext{\skt{सन्दष्टव्योऽथवा}}{\lemma{\skt{ऽथवा}}\Afootnote{A H; \skt{थवा}\ K}}\skt{ सर्पो }\edtext{\skt{दष्टमात्रेण}}{\lemma{\skt{दष्टमात्रेण}}\Afootnote{K H; \skt{लोष्टो वाऽपि}\ A}}\skt{ }\edtext{\skt{जानता}}{\lemma{\skt{जानता}}\Afootnote{K H; \skt{हि तत्क्षणम्~॥}\ A}}\skt{~॥}\edtext{}{\lemma{\skt{भवेत्~। सन्दष्टव्यो}}\Afootnote{H; \skt{भवेत्~। स दष्टव्यो}\ A; \skt{भवेत्~। सन्दष्टव्यो}\ K}}\pend

\pstart\versepar \ledsidenote{\passref{5.5.7}}\skt{अथ }\edtext{\skt{मण्डलिदष्टन्तु न}}{\lemma{\skt{मण्डलिदष्टन्तु न}}\Afootnote{K H; \skt{मण्डलिना दष्टं न}\ A}}\skt{ }\edtext{\skt{कथञ्चित्तु}}{\lemma{\skt{चित्तु}}\Afootnote{H; \skt{चन}\ A; \skt{च}\ K}}\skt{ दाहयेत्~। }\pend
\pstart\versepar \skt{स }\edtext{\skt{पित्त\-वि\-ष\-बा\-हु\-ल्या\-द्दंशो}}{\lemma{\skt{पित्त\appbrk{}वि\appbrk{}ष\appbrk{}बा\appbrk{}हु\appbrk{}ल्या\appbrk{}द्दंशो}}\Afootnote{\skt{पित्त\appbrk{}बा\appbrk{}हु\appbrk{}ल्य\appbrk{}वि\appbrk{}षा\appbrk{}द्दंशो}\ A; \skt{पित्त\appbrk{}वि\appbrk{}ष\appbrk{}वा\appbrk{}हु\appbrk{}ल्या\appbrk{}द्दंशो}\ K H}}\skt{ }\edtext{\skt{दाहा\-द्वि\-ना\-श\-येत्}}{\lemma{\skt{दाहा\appbrk{}द्वि\appbrk{}ना\appbrk{}श\appbrk{}येत्}}\Afootnote{K H; \skt{दाहाद्विसर्पते~॥}\ A}}\skt{~।}\pend

\pstart\versepar \ledsidenote{\passref{5.5.8}}\edtext{\skt{अरिष्टामपि}}{\lemma{\skt{अरिष्टाम्}}\Afootnote{A; \skt{अरिष्टाव्}\ K H}}\skt{ }\edtext{\skt{मन्त्रैस्तु}}{\lemma{\skt{मन्त्रैस्तु}}\Afootnote{K H; \skt{मन्त्रैश्च}\ A}}\skt{ बध्नी\-या\-त्म\-न्त्र\-को\-विदः~। }\pend
\pstart\versepar \skt{सा तु }\edtext{\edtext{\skt{रज्ज्वादिभिर्बद्धा}}{\lemma{\skt{रज्ज्वादिभिर्}}\Afootnote{A; \skt{रज्वादिभिर्}\ K H}}}{\lemma{\skt{बद्धा}}\Afootnote{A; \skt{वद्धा}\ K H}}\skt{ }\edtext{\skt{विषपूतिकरी}}{\lemma{\skt{विषपूतिकरी}}\Afootnote{K; \skt{विषप्रतिकरी}\ A; \skt{विषपूतीकरी}\ H}}\skt{ मता~।}\edtext{}{\lemma{\skt{बध्नी\appbrk{}या\appbrk{}त्म\appbrk{}न्त्र\appbrk{}को\appbrk{}विदः~। सा तु}}\Afootnote{\skt{बध्नी\appbrk{}या\appbrk{}न्म\appbrk{}न्त्र\appbrk{}को\appbrk{}विदः~। सा तु}\ A; \skt{वध्नी\appbrk{}या\appbrk{}न्म\appbrk{}न्त्र\appbrk{}को\appbrk{}विदः~। सा तु}\ K; \skt{वध्नी\appbrk{}या\appbrk{}त्म\appbrk{}न्त्र\appbrk{}को\appbrk{}विदः~। सा तु}\ H}}\pend

\pstart\versepar \ledsidenote{\passref{5.5.9}}\edtext{\skt{देव\-ब्र\-ह्म\-र्षि\-वि\-हिता}}{\lemma{\skt{देव\appbrk{}ब्र\appbrk{}ह्म\appbrk{}र्षि\appbrk{}वि\appbrk{}हिता}}\Afootnote{\skt{देवब्रह्मर्षिभिः प्रोक्ता}\ A; \skt{देव\appbrk{}व्र\appbrk{}ह्म\appbrk{}र्षि\appbrk{}वि\appbrk{}हिताः}\ K H}}\skt{ मन्त्राः सत्यतपोमयाः~। }\pend
\pstart\versepar \skt{}\edtext{\skt{भवन्त्यनत्ययाः}}{\lemma{\skt{अनत्ययाः}}\Afootnote{K H; \skt{नान्यथा}\ A}}\skt{ क्षिप्रं विषं }\edtext{\skt{हन्युश्च}}{\lemma{\skt{च}}\Afootnote{K H; \skt{सु}\ A}}\skt{ दुस्तरम्~।}\pend

\pstart\versepar \ledsidenote{\passref{5.5.10}}\skt{विषं तेजोमयैर्मन्त्रैः }\edtext{\skt{सत्य\-ब्र\-ह्म\-त\-पो\-मयैः}}{\lemma{\skt{सत्य\appbrk{}ब्र\appbrk{}ह्म\appbrk{}त\appbrk{}पो\appbrk{}मयैः}}\Afootnote{A; \skt{सत्य\appbrk{}व्र\appbrk{}ह्म\appbrk{}त\appbrk{}पो\appbrk{}मयैः}\ K H}}\skt{~। }\pend
\pstart\versepar \skt{यथा निवार्यते क्षिप्रं }\edtext{\skt{प्रयुक्तैर्न}}{\lemma{\skt{प्रयुक्तैर्न}}\Afootnote{A K; \skt{प्रयुक्तन्न}\ H}}\skt{ तथौषधैः~।}\pend

\pstart\versepar \ledsidenote{\passref{5.5.11}}\skt{मन्त्राणां ग्रहणं कार्यं स्त्रीमां\-स\-म\-धु\-व\-र्जिना~। }\pend
\pstart\versepar \skt{}\edtext{\skt{यताहारेण}}{\lemma{\skt{यताहारेण}}\Afootnote{K H; \skt{मिताहारेण}\ A}}\skt{ शुचिना कुशा\-स्त\-र\-ण\-शा\-यिना~।}\pend

\pstart\versepar \ledsidenote{\passref{5.5.12}}\skt{गन्ध\-मा\-ल्यो\-प\-हा\-रैश्च }\edtext{\skt{बलिभिश्चापि}}{\lemma{\skt{बलिभिश्}}\Afootnote{A; \skt{वलिभिश्}\ K H}}\skt{ }\edtext{\skt{देवताम्}}{\lemma{\skt{देवताम्}}\Afootnote{H; \skt{देवताः~।}\ A; \skt{देवता~।}\ K}}\skt{~। }\pend
\pstart\versepar \skt{}\edtext{\skt{पूज\-ये\-त्म\-न्त्र\-सि\-द्ध्यर्थं}}{\lemma{\skt{पूज\appbrk{}ये\appbrk{}त्म\appbrk{}न्त्र\appbrk{}सि\appbrk{}द्ध्यर्थं}}\Afootnote{K H; \skt{पूज\appbrk{}ये\appbrk{}न्म\appbrk{}न्त्र\appbrk{}सि\appbrk{}द्ध्यर्थं}\ A}}\skt{ }\edtext{\skt{जपहोमैश्च}}{\lemma{\skt{जपहोमैश्च}}\Afootnote{A K; \skt{जापहोमैश्च}\ H}}\skt{ यत्नतः~॥}\pend

\pstart\versepar \ledsidenote{\passref{5.5.13}}\skt{मन्त्रास्त्वविधिना प्रोक्ता हीना वा }\edtext{\skt{स्वरवर्ण्णतः}}{\lemma{\skt{स्वरवर्ण्णतः}}\Afootnote{A K; \skt{स्वरवर्ण्णितः~।}\ H}}\skt{~। }\pend
\pstart\versepar \skt{यस्मान्न सिद्धिमायान्ति }\edtext{\skt{तस्मा\-द्यो\-ज्योऽ\-ग\-द\-क्रमः}}{\lemma{\skt{ऽगदक्रमः}}\Afootnote{A; \skt{गदक्रमः}\ K H}}\skt{~।}\pend

\pstart\versepar \ledsidenote{\passref{5.5.14}}\edtext{\skt{दंशात्समन्ताच्च}}{\lemma{\skt{दंशात्समन्ताच्च}}\Afootnote{K H; \skt{समन्ततः}\ A}}\skt{ }\edtext{\skt{सिरां}}{\lemma{\skt{सिरां}}\Afootnote{K H; \skt{सिरा}\ A}}\skt{ }\edtext{\skt{व्यधयेत्कुशलो}}{\lemma{\skt{व्यधयेत्कुशलो}}\Afootnote{K H; \skt{दंशाद्विध्येत्तु कुशलो}\ A}}\skt{ भिषक्~। }\pend
\pstart\versepar \skt{शाखाश्रयां }\edtext{\skt{ललाटे च}}{\lemma{\skt{ललाटे च}}\Afootnote{H; \skt{ललाटे वा}\ A; \skt{ललाटो च}\ K}}\skt{ }\edtext{\skt{वेद्धव्या}}{\lemma{\skt{वेद्धव्या}}\Afootnote{K H; \skt{व्यध्यास्ता}\ A}}\skt{ }\edtext{\skt{विसृते}}{\lemma{\skt{विसृते}}\Afootnote{K H; \skt{विऋते}\ A}}\skt{ विषे~।}\edtext{}{\lemma{\skt{भिषक्~। शाखाश्रयां}}\Afootnote{K H; \skt{भिषक्~। शाखाग्रे वा}\ A}}\pend

\pstart\versepar \ledsidenote{\passref{5.5.15}}\edtext{\edtext{\skt{रक्त\-न्नि\-र्ह्रि\-य\-मा\-नन्तु}}{\lemma{\skt{रक्तन्}}\Afootnote{H; \skt{रक्ते}\ A; \skt{रक्त}\ K}}}{\lemma{\skt{निर्ह्रियमानन्तु}}\Afootnote{K; \skt{निर्ह्रियमाणे तु}\ A; \skt{निर्ह्रियमाणन्तु}\ H}}\skt{ कृत्स्नं }\edtext{\skt{निर्हरते}}{\lemma{\skt{निर्हरते}}\Afootnote{K H; \skt{निर्ह्रियते}\ A}}\skt{ विषम्~। }\pend
\pstart\versepar \skt{तस्मा\-द्वि\-स्रा\-व\-ये\-द्रक्तं सा ह्यस्य }\edtext{\skt{परमा}}{\lemma{\skt{परमा}}\Afootnote{A K; \skt{परम}\ H}}\skt{ क्रिया~।}\pend

\pstart\versepar \ledsidenote{\passref{5.5.16}}\edtext{\skt{दंशं}}{\lemma{\skt{दंशं}}\Afootnote{K H; \emph{om.}\ A}}\skt{ समन्तादगदैः }\edtext{\skt{प्रच्छयित्वा}}{\lemma{\skt{प्रच्छयित्वा}}\Afootnote{K H; \skt{दंशं प्रच्छयित्वा}\ A}}\skt{ }\edtext{\skt{च}}{\lemma{\skt{च}}\Afootnote{K H; \skt{प्र}\ A}}\skt{ लेपयेत्~। }\pend
\pstart\versepar \skt{चन्द\-नो\-शी\-र\-सि\-क्तेन वारिणा }\edtext{\skt{चापि}}{\lemma{\skt{चापि}}\Afootnote{K H; \skt{परि}\ A}}\skt{ }\edtext{\skt{सेचयेत्}}{\lemma{\skt{सेचयेत्}}\Afootnote{K H; \skt{षेचयेत्~॥}\ A}}\skt{~।}\edtext{}{\lemma{\skt{लेपयेत्~। चन्द\appbrk{}नो\appbrk{}शी\appbrk{}र\appbrk{}सि\appbrk{}क्तेन}}\Afootnote{K; \skt{लेपयेत्~। चन्द\appbrk{}नो\appbrk{}शी\appbrk{}र\appbrk{}यु\appbrk{}क्तेन}\ A; \skt{लेपयेत्~। चन्द\appbrk{}नो\appbrk{}शी\appbrk{}र\appbrk{}शि\appbrk{}क्तेन}\ H}}\pend

\pstart\versepar \ledsidenote{\passref{5.5.17}}\edtext{\edtext{\edtext{\skt{पाय\-ये\-च्चा\-ग\-दां\-स्तां\-स्ता\-न्द\-धि\-क्षौ\-द्र\-घृ\-ता\-दिभिः}}{\lemma{\skt{पाय\appbrk{}ये\appbrk{}च्चा\appbrk{}ग\appbrk{}दां\appbrk{}स्तांस्}}\Afootnote{H; \skt{पायये ता गदांस्तांस्}\ A; \skt{पाययेच्चागदां तास्}\ K}}}{\lemma{\skt{तान्}}\Afootnote{K H; \skt{ता}\ A}}}{\lemma{\skt{दधि\appbrk{}क्षौ\appbrk{}द्र\appbrk{}घृ\appbrk{}ता\appbrk{}दिभिः}}\Afootnote{K H; \skt{न्क्षीर\appbrk{}क्षौ\appbrk{}द्र\appbrk{}घृ\appbrk{}ता\appbrk{}दिभिः~।}\ A}}\skt{~। }\pend
\pstart\versepar \skt{तदलाभे हिता वा }\edtext{\skt{स्यात्कृ\-ष्ण\-व\-ल्मी\-क\-मृ\-त्तिका}}{\lemma{\skt{स्यात्कृ\appbrk{}ष्ण\appbrk{}व\appbrk{}ल्मी\appbrk{}क\appbrk{}मृ\appbrk{}त्तिका}}\Afootnote{K\textsuperscript{pc} H; \skt{स्यात्कृष्णा वल्मीकमृत्तिका~॥}\ A; \skt{स्यात्कृ\appbrk{}ष्ण\appbrk{}व\appbrk{}ल्मी\appbrk{}क\appbrk{}मृ\appbrk{}त्तिकाः}\ K\textsuperscript{ac}}}\skt{~॥}\pend

\pstart\versepar \ledsidenote{\passref{5.5.18}}\edtext{\skt{कोवि\-दा\-र\-शि\-री\-षार्कं}}{\lemma{\skt{कोवि\appbrk{}दा\appbrk{}र\appbrk{}शि\appbrk{}री\appbrk{}षार्कं}}\Afootnote{K H; \skt{कोवि\appbrk{}दा\appbrk{}र\appbrk{}शि\appbrk{}री\appbrk{}षार्क}\ A}}\skt{ }\edtext{\skt{कटभीं}}{\lemma{\skt{कटभीं}}\Afootnote{H; \skt{कटभीर्}\ A K}}\skt{ }\edtext{\skt{वापि}}{\lemma{\skt{वापि}}\Afootnote{K H; \skt{वाऽपि}\ A}}\skt{ भक्षयेत्~। }\pend
\pstart\versepar \skt{न }\edtext{\skt{पिबे\-त्तै\-ल\-कौ\-लत्थं}}{\lemma{\skt{पिबे\appbrk{}त्तै\appbrk{}ल\appbrk{}कौ\appbrk{}लत्थं}}\Afootnote{\skt{पिबे\appbrk{}त्तै\appbrk{}ल\appbrk{}कौ\appbrk{}लत्थ}\ A; \skt{पिवे\appbrk{}त्तै\appbrk{}ल\appbrk{}कौ\appbrk{}लत्थं}\ K H}}\skt{ }\edtext{\skt{मद्यं}}{\lemma{\skt{मद्यं}}\Afootnote{K; \skt{मद्य}\ A H}}\skt{ }\edtext{\skt{सौवीरकं च}}{\lemma{\skt{सौवीरकं च}}\Afootnote{K; \skt{सौवीरकाणि}\ A; \skt{सौवीरकन्न}\ H}}\skt{ }\edtext{\skt{न}}{\lemma{\skt{न}}\Afootnote{K; \skt{च~॥}\ A H}}\skt{~।}\pend

\pstart\versepar \ledsidenote{\passref{5.5.19}}\edtext{\skt{द्रवमन्यत्तु}}{\lemma{\skt{अन्यत्तु}}\Afootnote{A K; \skt{अन्यन्तु}\ H}}\skt{ यत्किञ्चित्पीत्वा पीत्वा तदुद्धरेत्~। }\pend
\pstart\versepar \skt{प्रायो हि वमनेनैव }\edtext{\skt{सुखं}}{\lemma{\skt{सुखं}}\Afootnote{K H; \skt{मुखं}\ A}}\skt{ निर्ह्रियते विषम्~।}\edtext{}{\lemma{\skt{तदुद्धरेत्~। प्रायो हि}}\Afootnote{K H; \skt{तदुद्वमेत्~। प्रायो हि}\ A}}\pend

\pstart\versepar \ledsidenote{\passref{5.5.20}}\skt{फणिनां विषवेगे तु }\edtext{\skt{प्रथमं}}{\lemma{\skt{प्रथमं}}\Afootnote{K H; \skt{प्रथमे}\ A}}\skt{ शोणितं हरेत्~। }\pend
\pstart\versepar \skt{द्वितीये }\edtext{\skt{मधु\-स\-र्पि\-र्भ्या\-म\-गदं}}{\lemma{\skt{अगदं}}\Afootnote{K H; \emph{om.}\ A}}\skt{ }\edtext{\skt{सह}}{\lemma{\skt{सह}}\Afootnote{K H; \emph{om.}\ A}}\skt{ }\edtext{\skt{पाययेत्}}{\lemma{\skt{पाययेत्}}\Afootnote{K H; \skt{पाययेतागदं भिसक्~॥}\ A}}\skt{~।}\edtext{}{\lemma{\skt{हरेत्~। द्वितीये}}\Afootnote{A H; \skt{हरेत्~। द्विती}\ K}}\pend

\pstart\versepar \ledsidenote{\passref{5.5.21}}\edtext{\skt{नस्तःकर्माञ्जने}}{\lemma{\skt{नस्तःकर्माञ्जने}}\Afootnote{K; \skt{नस्यकर्माञ्जने}\ A; \skt{नस्तकर्म्माञ्जने}\ H}}\skt{ युञ्ज्यात्तृतीये }\edtext{\skt{विषनाशने}}{\lemma{\skt{विषनाशने}}\Afootnote{A; \skt{विषनाशनं~।}\ K H}}\skt{~। }\pend
\pstart\versepar \skt{}\edtext{\skt{वान्ते}}{\lemma{\skt{वान्ते}}\Afootnote{K H; \skt{वान्तं}\ A}}\skt{ चतुर्थे }\edtext{\skt{विषघ्नां}}{\lemma{\skt{विषघ्नां}}\Afootnote{K H; \emph{om.}\ A}}\skt{ }\edtext{\skt{यवागूं}}{\lemma{\skt{यवागूं}}\Afootnote{K H; \emph{om.}\ A}}\skt{ }\edtext{\skt{पाययेद्भिषक्}}{\lemma{\skt{पाययेद्भिषक्}}\Afootnote{K\textsuperscript{ac} H; \skt{पूर्वोक्तां यवागूमथ दापयेत्~॥}\ A; \skt{पाययेद्भिषक् * शीतोपचारं पुरुषं वेगयोः पञ्चषष्ठयोः पाययेत्छोधनं तीक्ष्णं यवागूञ्चापि कीर्त्तिता सप्तमे त्ववपीडेन शिरस्तीक्ष्णेन शोधयेत्पूर्वा मण्डलिनां वेगा दर्वी\appbrk{}क\appbrk{}र\appbrk{}व\appbrk{}दा\appbrk{}च\appbrk{}रे\appbrk{}त्तृ\appbrk{}तीये च विरिक्तस्य *}\ K\textsuperscript{pc}}}\skt{~॥}\pend

\pstart\versepar \ledsidenote{\passref{5.5.22}}\edtext{\skt{शीतोपचारं}}{\lemma{\skt{शीतोपचारं}}\Afootnote{A K\textsuperscript{pc} H; \emph{om.}\ K\textsuperscript{ac}}}\skt{ }\edtext{\skt{पुरुषं}}{\lemma{\skt{पुरुषं}}\Afootnote{K\textsuperscript{pc} H; \skt{कृत्वादौ}\ A; \emph{om.}\ K\textsuperscript{ac}}}\skt{ }\edtext{\skt{वेगयोः}}{\lemma{\skt{वेगयोः}}\Afootnote{K\textsuperscript{pc} H; \skt{भिषक्पञ्}\ A; \emph{om.}\ K\textsuperscript{ac}}}\skt{ }\edtext{\skt{पञ्चषष्ठयोः}}{\lemma{\skt{पञ्चषष्ठयोः}}\Afootnote{K\textsuperscript{pc} H; \skt{चमषष्ठयोः~।}\ A; \emph{om.}\ K\textsuperscript{ac}}}\skt{~। }\pend
\pstart\versepar \skt{}\edtext{\skt{पाययेच्छोधनं}}{\lemma{\skt{पाययेच्छोधनं}}\Afootnote{A H; \emph{om.}\ K\textsuperscript{ac}; \skt{पाययेत्छोधनं}\ K\textsuperscript{pc}}}\skt{ }\edtext{\skt{तीक्ष्णं}}{\lemma{\skt{तीक्ष्णं}}\Afootnote{A K\textsuperscript{pc} H; \emph{om.}\ K\textsuperscript{ac}}}\skt{ }\edtext{\edtext{\skt{यवागूञ्चापि}}{\lemma{\skt{यवागूञ्}}\Afootnote{A K\textsuperscript{pc} H; \emph{om.}\ K\textsuperscript{ac}}}}{\lemma{\skt{चापि}}\Afootnote{A K\textsuperscript{pc} H; \emph{om.}\ K\textsuperscript{ac}}}\skt{ }\edtext{\skt{कीर्त्तिताम्}}{\lemma{\skt{कीर्त्तिताम्}}\Afootnote{A H; \emph{om.}\ K\textsuperscript{ac}; \skt{कीर्त्तिता}\ K\textsuperscript{pc}}}\skt{~।}\pend

\pstart \ledsidenote{\passref{5.5.23}}\edtext{\skt{सप्तमे}}{\lemma{\skt{सप्तमे}}\Afootnote{A K\textsuperscript{pc} H; \emph{om.}\ K\textsuperscript{ac}}}\skt{ }\edtext{\skt{त्ववपीडेन}}{\lemma{\skt{त्ववपीडेन}}\Afootnote{A K\textsuperscript{pc} H; \emph{om.}\ K\textsuperscript{ac}}}\skt{ }\edtext{\edtext{\skt{शिरस्तीक्ष्णेन}}{\lemma{\skt{शिरस्}}\Afootnote{A K\textsuperscript{pc} H; \emph{om.}\ K\textsuperscript{ac}}}}{\lemma{\skt{तीक्ष्णेन}}\Afootnote{A K\textsuperscript{pc}; \emph{om.}\ K\textsuperscript{ac}; \skt{तीक्ष्णेण}\ H}}\skt{ }\edtext{\skt{शोधयेत्}}{\lemma{\skt{शोधयेत्}}\Afootnote{K\textsuperscript{pc} H; \skt{शोधयेत्~। तीक्ष्णमेवाञ्जनं दद्या\appbrk{}त्ती\appbrk{}क्ष्ण\appbrk{}श\appbrk{}स्त्रेण मूर्ध्नि च~॥}\ A; \emph{om.}\ K\textsuperscript{ac}}}\skt{~।}\pend

\pstart \ledsidenote{\passref{5.5.24}}\edtext{\skt{पूर्वो}}{\lemma{\skt{पूर्वो}}\Afootnote{K\textsuperscript{pc} H; \skt{कृत्वा काकपदं चर्म सासृग्वा पिशितं क्षिपेत्~। पूर्वो}\ A; \emph{om.}\ K\textsuperscript{ac}}}\skt{ }\edtext{\skt{मण्डलिनां}}{\lemma{\skt{मण्डलिनां}}\Afootnote{A K\textsuperscript{pc} H; \emph{om.}\ K\textsuperscript{ac}}}\skt{ }\edtext{\skt{वेगो}}{\lemma{\skt{वेगो}}\Afootnote{A K\textsuperscript{pc} H; \emph{om.}\ K\textsuperscript{ac}}}\skt{ }\edtext{\skt{दर्वी\-क\-र\-व\-दा\-च\-रेत्}}{\lemma{\skt{दर्वी\appbrk{}क\appbrk{}र\appbrk{}व\appbrk{}दा\appbrk{}च\appbrk{}रेत्}}\Afootnote{A K\textsuperscript{pc} H; \emph{om.}\ K\textsuperscript{ac}}}\skt{~।}\pend

\pstart \ledsidenote{\passref{5.5.25}}\edtext{\skt{द्वितीये}}{\lemma{\skt{द्वितीये}}\Afootnote{H; \skt{अगदं मधुसर्पिर्भ्यां द्वितीये}\ A}}\skt{ }\edtext{\skt{सर्प्पिर्म्मधुनी}}{\lemma{\skt{सर्प्पिर्म्मधुनी}}\Afootnote{H; \emph{om.}\ A}}\skt{ }\edtext{\skt{पाययित्वा च}}{\lemma{\skt{पाययित्वा च}}\Afootnote{H; \skt{पाययेत च~।}\ A}}\skt{ }\edtext{\skt{वामयेत्}}{\lemma{\skt{वामयेत्}}\Afootnote{H; \skt{वामयित्वा यवागूं च पूर्वोक्तामथ दापयेत्~॥}\ A}}\skt{~।}\pend

\pstart\versepar \ledsidenote{\passref{5.5.26}}\edtext{\skt{तृतीये}}{\lemma{\skt{तृतीये}}\Afootnote{A K\textsuperscript{pc} H; \emph{om.}\ K\textsuperscript{ac}}}\skt{ }\edtext{\skt{च विरिक्तस्य}}{\lemma{\skt{च विरिक्तस्य}}\Afootnote{K\textsuperscript{pc}; \skt{शोधितं तीक्ष्णैर्}\ A; \emph{om.}\ K\textsuperscript{ac}; \skt{सविरिक्तस्य}\ H\textsuperscript{ac}; \skt{सुविरिक्तस्य}\ H\textsuperscript{pc}}}\skt{ }\edtext{\edtext{\skt{यवा\-गू\-न्दा\-प\-ये\-द्धि\-ताम्}}{\lemma{\skt{यवागून्}}\Afootnote{A H; \skt{यवागुम्}\ K}}}{\lemma{\skt{दापयेद्धिताम्}}\Afootnote{K H; \skt{पाययेद्धिताम्~।}\ A}}\skt{~। }\pend
\pstart\versepar \skt{चतुर्थे पञ्चमे }\edtext{\skt{चापि}}{\lemma{\skt{चापि}}\Afootnote{A K; \skt{वापि}\ H}}\skt{ दर्व्वी\-क\-र\-व\-दा\-च\-रेत्~।}\pend

\pstart\versepar \ledsidenote{\passref{5.5.27}}\skt{काकोल्यादिर्हितः षष्ठे पेयश्च }\edtext{\skt{मधुरोऽगदः}}{\lemma{\skt{ऽगदः}}\Afootnote{A; \skt{गदः}\ K H}}\skt{~। }\pend
\pstart\versepar \skt{}\edtext{\skt{हितोऽवपीडे}}{\lemma{\skt{ऽवपीडे}}\Afootnote{A H; \skt{वपीडे}\ K}}\skt{ त्वगदः सप्तमे विषनाशनः~॥}\pend

\pstart \ledsidenote{\passref{5.5.28}}\edtext{\skt{अथ}}{\lemma{\skt{अथ}}\Afootnote{K H; \skt{पूर्वे}\ A}}\skt{ राजिमतां वेगे }\edtext{\skt{प्रथमे}}{\lemma{\skt{प्रथमे}}\Afootnote{K H; \skt{ऽलाबुभिः}\ A}}\skt{ शोणितं }\edtext{\skt{हरेत्}}{\lemma{\skt{हरेत्}}\Afootnote{K H; \skt{हरेत्~। अगदं मधुसर्पिर्भ्यां संयुक्तं पाययेत च~॥}\ A}}\skt{~।}\pend

\pstart\versepar \ledsidenote{\passref{5.5.29}}\skt{वान्तं द्वितीये त्वगदं पाय\-ये\-द्वि\-ष\-ना\-श\-नम्~। }\pend
\pstart\versepar \skt{तृतीयादिषु }\edtext{\skt{त्रिष्वेव}}{\lemma{\skt{एव}}\Afootnote{K H; \skt{एवं}\ A}}\skt{ विधिर्दार्वीकरो हितः~।}\pend

\pstart\versepar \ledsidenote{\passref{5.5.30}}\edtext{\edtext{\skt{षष्ठेऽञ्जनं}}{\lemma{\skt{षष्ठे}}\Afootnote{A; \skt{षष्ठे ञ्}\ K H}}}{\lemma{\skt{ऽञ्जनं}}\Afootnote{A; \skt{जनं}\ K H}}\skt{ तीक्ष्ण\-त\-म\-म\-व\-पी\-डश्च सप्तमे~। }\pend
\pstart\versepar \skt{}\edtext{\skt{गर्भि\-णी\-बा\-ल\-वृ\-द्धानां}}{\lemma{\skt{गर्भि\appbrk{}णी\appbrk{}बा\appbrk{}ल\appbrk{}वृ\appbrk{}द्धानां}}\Afootnote{A; \skt{गर्भि\appbrk{}णी\appbrk{}वा\appbrk{}ल\appbrk{}वृ\appbrk{}द्धानां}\ K; \skt{गर्व्भि\appbrk{}णी\appbrk{}बा\appbrk{}ल\appbrk{}वृ\appbrk{}द्धानां}\ H}}\skt{ }\edtext{\skt{सिरा\-वे\-ध\-वि\-व\-र्जि\-तम्}}{\lemma{\skt{सिरा\appbrk{}वे\appbrk{}ध\appbrk{}वि\appbrk{}व\appbrk{}र्जि\appbrk{}तम्}}\Afootnote{K H; \skt{सिरा\appbrk{}व्य\appbrk{}ध\appbrk{}न\appbrk{}व\appbrk{}र्जि\appbrk{}तम्~॥}\ A}}\skt{~।}\pend

\pstart\versepar \ledsidenote{\passref{5.5.31}}\edtext{\skt{विषार्त्तिषु}}{\lemma{\skt{विषार्त्तिषु}}\Afootnote{K H; \skt{विषार्तानां}\ A}}\skt{ यथोद्दिष्टं विधानं }\edtext{\skt{मृदु}}{\lemma{\skt{मृदु}}\Afootnote{K H; \emph{om.}\ A}}\skt{ शस्यते~। }\pend
\pstart\versepar \skt{}\edtext{\skt{रक्ता\-व\-से\-का\-ञ्ज\-नानि}}{\lemma{\skt{रक्ता\appbrk{}व\appbrk{}से\appbrk{}का\appbrk{}ञ्ज\appbrk{}नानि}}\Afootnote{K H; \skt{मृदु~। रक्ता\appbrk{}व\appbrk{}से\appbrk{}का\appbrk{}ञ्ज\appbrk{}नानि}\ A}}\skt{ }\edtext{\edtext{\skt{नर\-तु\-ल्या\-न्य\-जा\-विके}}{\lemma{\skt{नरतुल्यान्य्}}\Afootnote{A K\textsuperscript{pc}; \skt{नरतुल्या}\ K\textsuperscript{ac}; \skt{नरतुल्याञ्}\ H}}}{\lemma{\skt{अजाविके}}\Afootnote{A K\textsuperscript{pc}; \skt{विके~।}\ K\textsuperscript{ac}; \skt{जनाविके~॥}\ H}}\skt{~।}\pend

\pstart\versepar \ledsidenote{\passref{5.5.32}}\edtext{\edtext{\skt{गवा\-श्व\-यो\-स्त\-द्द्वि\-गुणं}}{\lemma{\skt{गवाश्वयोस्}}\Afootnote{K H; \emph{om.}\ A}}}{\lemma{\skt{तद्द्विगुणं}}\Afootnote{H; \emph{om.}\ A; \skt{त द्विगुणां}\ K}}\skt{ त्रिगुणं }\edtext{\skt{महिषोष्ट्रयोः}}{\lemma{\skt{महिषोष्ट्रयोः}}\Afootnote{K H; \skt{महिषे सोष्ट्रे गवाश्वे द्विगुणं तु}\ A}}\skt{~। }\pend
\pstart\versepar \skt{}\edtext{\skt{चतुर्गुणं तु}}{\lemma{\skt{चतुर्गुणं तु}}\Afootnote{K H; \skt{तत्~। चतुर्गुणं तु}\ A}}\skt{ नागानां केवलं सर्वपक्षिणाम्~।}\pend

\pstart\versepar \ledsidenote{\passref{5.5.34}}\edtext{\skt{देश\-प्र\-कृ\-ति\-सा\-त्म्य\-र्तु\-वि\-ष\-वे\-ग\-ब\-ला\-ब\-लम्}}{\lemma{\skt{देश\appbrk{}प्र\appbrk{}कृ\appbrk{}ति\appbrk{}सा\appbrk{}त्म्य\appbrk{}र्तु\appbrk{}वि\appbrk{}ष\appbrk{}वे\appbrk{}ग\appbrk{}ब\appbrk{}ला\appbrk{}ब\appbrk{}लम्}}\Afootnote{A; \skt{देशे प्रकृ\appbrk{}ति\appbrk{}सा\appbrk{}त्म्य\appbrk{}न्तु\appbrk{}वि\appbrk{}ष\appbrk{}वे\appbrk{}ग\appbrk{}व\appbrk{}ला\appbrk{}वलं~।}\ K; \skt{देश\appbrk{}प्र\appbrk{}कृ\appbrk{}ति\appbrk{}सा\appbrk{}त्म्य\appbrk{}न्तु\appbrk{}वि\appbrk{}ष\appbrk{}वे\appbrk{}ग\appbrk{}व\appbrk{}ला\appbrk{}वलं~।}\ H}}\skt{~। }\pend
\pstart\versepar \skt{प्रधार्य }\edtext{\skt{निपुणं}}{\lemma{\skt{निपुणं}}\Afootnote{A; \skt{निपुनं}\ K H}}\skt{ }\edtext{\skt{बुध्या}}{\lemma{\skt{बुध्या}}\Afootnote{A; \skt{वुध्या}\ K H}}\skt{ ततः कर्म समाचरेत्~।}\pend

\pstart \ledsidenote{\passref{5.5.47cd}}\skt{निःशेषं }\edtext{\skt{निर्हरेच्चैनं}}{\lemma{\skt{निर्हरेच्चैनं}}\Afootnote{K H\textsuperscript{pc}; \skt{निर्हरेच्चैवं}\ A; \skt{निर्हच्चैनम्}\ H\textsuperscript{ac}}}\skt{ विषं पर\-म\-दु\-र्ज\-यम्~।}\pend

\pstart\versepar \ledsidenote{\passref{5.5.48}}\edtext{\edtext{\skt{स्वल्प\-म\-प्य\-व\-तिष्ठं}}{\lemma{\skt{स्वल्पम्}}\Afootnote{K H; \skt{अल्पम्}\ A}}\skt{ हि}}{\lemma{\skt{अवतिष्ठं हि}}\Afootnote{K H; \skt{अवशिष्टं हि}\ A}}\skt{ भूयो वेगाय कल्पते~। }\pend
\pstart\versepar \skt{कुर्याद्वा }\edtext{\skt{साद\-वै\-व\-र्ण्ण्य\-ज्व\-र\-का\-स\-शि\-रो\-रुजः}}{\lemma{\skt{साद\appbrk{}वै\appbrk{}व\appbrk{}र्ण्ण्य\appbrk{}ज्व\appbrk{}र\appbrk{}का\appbrk{}स\appbrk{}शि\appbrk{}रो\appbrk{}रुजः}}\Afootnote{A; \skt{साद\appbrk{}वै\appbrk{}व\appbrk{}र्ण्ण्य\appbrk{}ज्व\appbrk{}र\appbrk{}का\appbrk{}स\appbrk{}सि\appbrk{}रो\appbrk{}रुजः~।}\ K; \skt{साद\appbrk{}वै\appbrk{}व\appbrk{}र्ण्ण्य\appbrk{}ज्व\appbrk{}र\appbrk{}का\appbrk{}स\appbrk{}शि\appbrk{}रो\appbrk{}रुजाः~।}\ H}}\skt{~।}\pend

\pstart\versepar \ledsidenote{\passref{5.5.49}}\edtext{\skt{शोष\-शो\-फ\-प्र\-ति\-श्या\-य\-ति\-मि\-रा\-रु\-चि\-जा\-ड्य\-ताम्}}{\lemma{\skt{शोष\appbrk{}शो\appbrk{}फ\appbrk{}प्र\appbrk{}ति\appbrk{}श्या\appbrk{}य\appbrk{}ति\appbrk{}मि\appbrk{}रा\appbrk{}रु\appbrk{}चि\appbrk{}जा\appbrk{}ड्य\appbrk{}ताम्}}\Afootnote{K H; \skt{शोफ\appbrk{}शो\appbrk{}ष\appbrk{}प्र\appbrk{}ति\appbrk{}श्या\appbrk{}य\appbrk{}ति\appbrk{}मि\appbrk{}रा\appbrk{}रु\appbrk{}चि\appbrk{}पी\appbrk{}न\appbrk{}सान्~।}\ A}}\skt{~। }\pend
\pstart\versepar \skt{}\edtext{\skt{तासु}}{\lemma{\skt{तासु}}\Afootnote{K H; \skt{तेषु}\ A}}\skt{ चापि }\edtext{\skt{यथायोगं}}{\lemma{\skt{यथायोगं}}\Afootnote{K H; \skt{यथादोषं}\ A}}\skt{ }\edtext{\skt{प्रतिकर्म}}{\lemma{\skt{प्रतिकर्म}}\Afootnote{A K H\textsuperscript{pc}; \skt{प्रति}\ H\textsuperscript{ac}}}\skt{ प्रयोजयेत्~।}\pend

\pstart\versepar \ledsidenote{\passref{5.5.50}}\skt{विषा\-र्त्तो\-प\-द्र\-वां\-श्चापि यथास्वं समुपाचरेत्~॥ }\pend
\pstart\versepar \skt{अथारिष्टां }\edtext{\skt{विमुच्याशु}}{\lemma{\skt{विमुच्याशु}}\Afootnote{A H; \skt{विमुच्याशुः}\ K}}\skt{ प्रच्छयित्वाङ्कितं तया~।}\edtext{}{\lemma{\skt{समुपाचरेत्~॥ अथारिष्टां}}\Afootnote{A K; \skt{समुपाचरेत्~॥ यथारिष्टाम्}\ H}}\pend

\pstart \ledsidenote{\passref{5.5.51}}\edtext{\skt{विद्यात्तत्र}}{\lemma{\skt{विद्यात्तत्र}}\Afootnote{K H; \skt{दह्यात्तत्र}\ A}}\skt{ विषं स्कन्नं भूयो वेगाय }\edtext{\skt{कल्पते}}{\lemma{\skt{कल्पते}}\Afootnote{K H; \skt{कल्पते~। एव\appbrk{}मौ\appbrk{}ष\appbrk{}धि\appbrk{}भि\appbrk{}र्मन्त्रैः क्रियायोगैश्च यत्नतः~॥}\ A}}\skt{~।}\pend

\pstart \ledsidenote{\passref{5.5.52.add-1}}\edtext{\skt{विषापायेऽनिलं}}{\lemma{\skt{ऽनिलं}}\Afootnote{\skt{निलं}\ K H}}\skt{ क्रुद्धं जये\-द\-नि\-ल\-वा\-रणैः~।}\pend

\pstart\versepar \ledsidenote{\passref{5.5.53}}\edtext{\skt{तैल\-म\-द्य\-कु\-ल\-त्था\-म्ल\-व\-र्जै\-र्वि\-ष\-ह\-रा\-युतैः}}{\lemma{\skt{तैल\appbrk{}म\appbrk{}द्य\appbrk{}कु\appbrk{}ल\appbrk{}त्था\appbrk{}म्ल\appbrk{}व\appbrk{}र्जैर्}}\Afootnote{\skt{तैल\appbrk{}म\appbrk{}त्स्य\appbrk{}कु\appbrk{}ल\appbrk{}त्था\appbrk{}म्ल\appbrk{}व\appbrk{}र्ज्यैर्}\ A; \skt{तैल\appbrk{}म\appbrk{}द्य\appbrk{}कु\appbrk{}ल\appbrk{}त्था\appbrk{}म्ब्ल\appbrk{}व\appbrk{}र्जैर्}\ K; \skt{तैल\appbrk{}म\appbrk{}द्य\appbrk{}कु\appbrk{}ल\appbrk{}त्था\appbrk{}म्ल\appbrk{}व\appbrk{}र्ज्यैर्}\ H}}\skt{~। }\pend
\pstart\versepar \skt{}\edtext{\skt{पित्तं}}{\lemma{\skt{पित्तं}}\Afootnote{K H; \emph{om.}\ A}}\skt{ }\edtext{\skt{पित्तज्वरहरैः}}{\lemma{\skt{पित्तज्वरहरैः}}\Afootnote{K H; \skt{पित्तज्व्रहरैः}\ A}}\skt{ }\edtext{\skt{कषा\-य\-स्ने\-ह\-रे\-चनैः}}{\lemma{\skt{कषा\appbrk{}य\appbrk{}स्ने\appbrk{}ह\appbrk{}रे\appbrk{}चनैः}}\Afootnote{K H; \skt{पित्तं कषा\appbrk{}य\appbrk{}स्ने\appbrk{}ह\appbrk{}ब\appbrk{}स्तिभिः~॥}\ A}}\skt{~।}\pend

\pstart \ledsidenote{\passref{5.5.54}}\skt{कफ\-मा\-र\-ग्व\-धा\-द्येन सक्षौद्रेण गणेन }\edtext{\skt{तु}}{\lemma{\skt{तु}}\Afootnote{K H; \skt{तु~। श्लेष्म\appbrk{}घ्नै\appbrk{}र\appbrk{}ग\appbrk{}दै\appbrk{}श्चैव तिक्तै रूक्षैश्च भोजनैः~॥}\ A}}\skt{~॥}\pend

\pstart\versepar \ledsidenote{\passref{5.5.56}}\skt{गाढं }\edtext{\edtext{\skt{बद्धेऽरिष्टया}}{\lemma{\skt{बद्धे}}\Afootnote{A; \skt{वद्धे}\ K H}}}{\lemma{\skt{ऽरिष्टया}}\Afootnote{A; \skt{रिष्टया}\ K H}}\skt{ प्रच्छिते वा~। }\pend
\pstart\versepar \skt{}\edtext{\skt{तीक्ष्णै\-र्ले\-पै\-र्वि\-ष\-शे\-षेण}}{\lemma{\skt{विषशेषेण}}\Afootnote{K H; \skt{तद्वि\appbrk{}धै\appbrk{}र्वा\appbrk{}व\appbrk{}शिष्टैः~।}\ A}}\skt{ }\edtext{\skt{वापि}}{\lemma{\skt{वापि}}\Afootnote{K H; \emph{om.}\ A}}\skt{~। }\pend
\pstart\versepar \skt{शूने गात्रे क्लिन्न\-म\-त्य\-र्थ\-पूति }\pend
\pstart\versepar \skt{}\edtext{\skt{शीर्ण्णं}}{\lemma{\skt{शीर्ण्णं}}\Afootnote{K H; \skt{ज्ञेयं}\ A}}\skt{ मांसं }\edtext{\skt{विषपूति}}{\lemma{\skt{विषपूति}}\Afootnote{K H; \skt{तद्विषात्पूति}\ A}}\skt{ }\edtext{\skt{प्रदिष्टं}}{\lemma{\skt{प्रदिष्टं}}\Afootnote{K H; \skt{कष्टम्~॥}\ A}}\skt{~॥}\pend

\pstart\versepar \ledsidenote{\passref{5.5.57}}\edtext{\skt{सद्यः}}{\lemma{\skt{सद्यः}}\Afootnote{A H; \skt{सद्य}\ K}}\skt{ }\edtext{\skt{क्षतं}}{\lemma{\skt{क्षतं}}\Afootnote{K H; \skt{विद्धं}\ A}}\skt{ }\edtext{\skt{पच्यते}}{\lemma{\skt{पच्यते}}\Afootnote{K H; \skt{निस्रवे}\ A}}\skt{ }\edtext{\skt{यस्य}}{\lemma{\skt{यस्य}}\Afootnote{K H; \emph{om.}\ A}}\skt{ }\edtext{\skt{जन्तोः}}{\lemma{\skt{जन्तोः}}\Afootnote{K; \emph{om.}\ A; \skt{जन्तो}\ H}}\skt{~। }\pend
\pstart\versepar \skt{}\edtext{\skt{कृष्णं}}{\lemma{\skt{कृष्णं}}\Afootnote{K H; \skt{त्कृष्ण}\ A}}\skt{ रक्तं }\edtext{\skt{स्रवते}}{\lemma{\skt{स्रवते}}\Afootnote{K H; \skt{पाकं याया}\ A}}\skt{ }\edtext{\skt{दह्यते}}{\lemma{\skt{दह्यते}}\Afootnote{K H; \skt{द्दह्यते चाप्}\ A}}\skt{ }\edtext{\skt{च}}{\lemma{\skt{च}}\Afootnote{K H; \skt{य्}\ A}}\skt{~। }\pend
\pstart\versepar \skt{}\edtext{\skt{श्यावीभूतं}}{\lemma{\skt{श्यावीभूतं}}\Afootnote{K H; \skt{अभीक्ष्णम्~। कृष्णीभूतं}\ A}}\skt{ }\edtext{\skt{क्लिन्न\-म\-त्य\-र्थ\-पूति}}{\lemma{\skt{अत्यर्थपूति}}\Afootnote{K H; \skt{अत्यर्थपूति शीर्णं मांसं यात्य्}\ A}}\skt{~। }\pend
\pstart\versepar \skt{}\edtext{\skt{क्षतात्मांसं}}{\lemma{\skt{क्षतात्मांसं}}\Afootnote{K H; \skt{अजस्रं क्षताच्च~॥}\ A}}\skt{ }\edtext{\skt{शीर्यते}}{\lemma{\skt{शीर्यते}}\Afootnote{K H; \emph{om.}\ A}}\skt{ }\edtext{\skt{यस्य}}{\lemma{\skt{यस्य}}\Afootnote{K H; \emph{om.}\ A}}\skt{ }\edtext{\skt{चापि}}{\lemma{\skt{चापि}}\Afootnote{K H; \emph{om.}\ A}}\skt{~।}\pend

\pstart\versepar \ledsidenote{\passref{5.5.58}}\skt{तृष्णा मूर्च्छा }\edtext{\skt{ज्वरदाहौ}}{\lemma{\skt{ज्वरदाहौ}}\Afootnote{K H; \skt{भ्रान्तिदाहौ}\ A}}\skt{ }\edtext{\skt{च}}{\lemma{\skt{च}}\Afootnote{K H; \skt{ज्वरश्च}\ A}}\skt{ }\edtext{\skt{यस्य}}{\lemma{\skt{यस्य}}\Afootnote{K H; \skt{यस्य स्युस्तं}\ A}}\skt{~। }\pend
\pstart\versepar \skt{}\edtext{\skt{दिग्धाहतं तं}}{\lemma{\skt{दिग्धाहतं तं}}\Afootnote{K H; \skt{दिग्धविद्धं}\ A}}\skt{ }\edtext{\edtext{\skt{मनु\-ज\-म्व्य\-व\-स्येत्}}{\lemma{\skt{मनुजम्}}\Afootnote{K H; \emph{om.}\ A}}}{\lemma{\skt{व्यवस्येत्}}\Afootnote{K H; \skt{व्यवस्येत्~। पूर्वोद्दिष्टं लक्षणं सर्वमेतज्जुष्टं यस्यालं विषेण व्रणाः स्युः~॥}\ A}}\skt{~॥}\pend

\pstart\versepar \ledsidenote{\passref{5.5.58.add-1}}\skt{लिङ्गान्येतान्येव वा यस्य विद्याद्~। }\pend
\pstart\versepar \skt{व्रणे विषं }\edtext{\skt{यस्य}}{\lemma{\skt{यस्य}}\Afootnote{K; \skt{तस्य}\ H}}\skt{ दत्तं प्रमादात्~। }\pend
\pstart\versepar \skt{दिग्धाहतं }\edtext{\skt{विषजुष्टं}}{\lemma{\skt{विषजुष्टं}}\Afootnote{H; \skt{विषं जुष्टं}\ K}}\skt{ व्रणञ्च~। }\pend
\pstart\versepar \skt{ये चाप्यन्ये }\edtext{\skt{विष\-पू\-ति\-व्र\-णार्त्ताः}}{\lemma{\skt{विष\appbrk{}पू\appbrk{}ति\appbrk{}व्र\appbrk{}णार्त्ताः}}\Afootnote{K\textsuperscript{pc} H; \skt{विषपू}\ K\textsuperscript{ac}}}\skt{~॥}\pend

\pstart\versepar \ledsidenote{\passref{5.5.59cd}}\edtext{\skt{तेषां}}{\lemma{\skt{तेषां}}\Afootnote{A K\textsuperscript{pc} H; \emph{om.}\ K\textsuperscript{ac}}}\skt{ }\edtext{\edtext{\skt{धीमा\-न\-धि\-मां\-सा\-न्य\-पोह्य}}{\lemma{\skt{धीमा\appbrk{}न\appbrk{}धि\appbrk{}मां\appbrk{}सान्य्}}\Afootnote{H; \skt{युक्त्या पूतिमांसान्य्}\ A; \emph{om.}\ K\textsuperscript{ac}; \skt{धीमां\appbrk{}न\appbrk{}धि\appbrk{}मां\appbrk{}सान्य्}\ K\textsuperscript{pc}}}}{\lemma{\skt{अपोह्य}}\Afootnote{A K\textsuperscript{pc} H; \emph{om.}\ K\textsuperscript{ac}}}\skt{~। }\pend
\pstart\versepar \skt{}\edtext{\skt{जलौकाभिः}}{\lemma{\skt{जलौकाभिः}}\Afootnote{H; \skt{वार्योकोभिः}\ A; \emph{om.}\ K\textsuperscript{ac}; \skt{जल्लौकाभिः}\ K\textsuperscript{pc}}}\skt{ }\edtext{\skt{शोणितं}}{\lemma{\skt{शोणितं}}\Afootnote{A K\textsuperscript{pc} H; \emph{om.}\ K\textsuperscript{ac}}}\skt{ }\edtext{\skt{चापहृत्वा}}{\lemma{\skt{चापहृत्वा}}\Afootnote{K\textsuperscript{pc} H; \skt{चापहृत्य~॥}\ A; \emph{om.}\ K\textsuperscript{ac}}}\skt{~॥}\pend

\pstart\versepar \ledsidenote{\passref{5.5.60}}\skt{हृत्वा }\edtext{\skt{दोषा\-नू\-र्ध्व\-म\-धश्च}}{\lemma{\skt{दोषा\appbrk{}नू\appbrk{}र्ध्व\appbrk{}म\appbrk{}धश्च}}\Afootnote{K; \skt{दोषान्क्षिप्रमूर्ध्वं त्वधश्च}\ A; \skt{दोषा\appbrk{}नू\appbrk{}र्द्ध्व\appbrk{}म\appbrk{}धश्च}\ H}}\skt{ सम्यक्~। }\pend
\pstart\versepar \skt{सिञ्चेच्छीतैः }\edtext{\skt{क्षीरिणां}}{\lemma{\skt{क्षीरिणां}}\Afootnote{K H; \skt{त्क्षीरिणां}\ A}}\skt{ }\edtext{\skt{त्वकषायैः}}{\lemma{\skt{त्वकषायैः}}\Afootnote{K; \skt{त्वक्कषायैः~।}\ A H}}\skt{~। }\pend
\pstart\versepar \skt{}\edtext{\skt{वस्त्रा\-न्त\-रा\-न्दा\-प\-येच्च}}{\lemma{\skt{वस्त्रान्तरान्}}\Afootnote{K H; \skt{अन्तर्वस्त्रं}\ A}}\skt{ }\edtext{\skt{प्रदेहाञ्}}{\lemma{\skt{प्रदेहाञ्}}\Afootnote{K H; \skt{प्रदेहा}\ A}}\skt{~। }\pend
\pstart\versepar \skt{}\edtext{\edtext{\edtext{\skt{च्छीतै\-र्द्र\-व्यै\-र्घृ\-त\-यु\-क्तै\-र्वि\-षघ्नैः}}{\lemma{\skt{च्छीतैर्}}\Afootnote{K H; \skt{न्शीतैर्}\ A}}}{\lemma{\skt{द्रव्यैर्}}\Afootnote{A H; \skt{द्रव्यै}\ K}}}{\lemma{\skt{घृतयुक्तैर्}}\Afootnote{K H; \skt{आज्ययुक्तैर्}\ A}}\skt{~।}\edtext{}{\lemma{\skt{सम्यक्~। सिञ्चेच्छीतैः}}\Afootnote{K H; \skt{सम्यक्सिञ्चे}\ A}}\pend

\pstart\versepar \ledsidenote{\passref{5.5.61}}\edtext{\edtext{\skt{क्षतेऽस्थनि}}{\lemma{\skt{क्षते}}\Afootnote{K H; \skt{भिन्ने त्वस्थ्ना दुष्टजा}\ A}}\skt{ स}}{\lemma{\skt{ऽस्थनि स}}\Afootnote{\skt{तेन}\ A; \skt{स्थितिः स}\ K H}}\skt{ }\edtext{\edtext{\skt{विषैरेष}}{\lemma{\skt{विषैर्}}\Afootnote{K H; \emph{om.}\ A}}}{\lemma{\skt{एष}}\Afootnote{K H; \emph{om.}\ A}}\skt{ }\edtext{\skt{एव}}{\lemma{\skt{एव}}\Afootnote{K H; \emph{om.}\ A}}\skt{~। }\pend
\pstart\versepar \skt{}\edtext{\skt{विधिः}}{\lemma{\skt{विधिः}}\Afootnote{K H; \emph{om.}\ A}}\skt{ कार्यः }\edtext{\skt{पित्तविषे}}{\lemma{\skt{पित्तविषे}}\Afootnote{K H; \skt{पूर्वो मार्गः पैत्तिके यो विषे}\ A}}\skt{ }\edtext{\skt{तथैव}}{\lemma{\skt{तथैव}}\Afootnote{K H; \skt{च~।}\ A}}\skt{~॥ }\pend
\pstart\versepar \skt{}\edtext{\skt{त्रिवृद्विशल्या}}{\lemma{\skt{त्रिवृद्विशल्या}}\Afootnote{\skt{त्रिवृद्विशल्ये}\ A; \skt{तृवृद्विशल्या}\ K H}}\skt{ मधुकं हरिद्रे~। }\pend
\pstart\versepar \skt{}\edtext{\skt{मञ्जिष्ठवक्रौ}}{\lemma{\skt{मञ्जिष्ठवक्रौ}}\Afootnote{K; \skt{रक्ता नरेन्द्रो}\ A; \skt{मञ्जिष्ठवर्ग्गो}\ H\textsuperscript{ac}; \skt{मञ्जि\appbrk{}ष्ठ\appbrk{}यु\appbrk{}क्तो\appbrk{}वर्ग्गो}\ H\textsuperscript{pc}}}\skt{ लवणश्च }\edtext{\skt{सर्वः}}{\lemma{\skt{सर्वः}}\Afootnote{K H; \skt{वर्गः~॥}\ A}}\skt{~।}\pend

\pstart\versepar \ledsidenote{\passref{5.5.62}}\skt{कटुत्रिकं }\edtext{\skt{चैव}}{\lemma{\skt{चैव}}\Afootnote{A K H\textsuperscript{pc}; \skt{चै}\ H\textsuperscript{ac}}}\skt{ }\edtext{\skt{विचूर्ण्णितानि}}{\lemma{\skt{विचूर्ण्णितानि}}\Afootnote{K H; \skt{सुचूर्णितानि}\ A}}\skt{ }\pend
\pstart\versepar \skt{शृङ्गे }\edtext{\skt{निद\-ध्या\-त्म\-धु\-सं\-यु\-तानि}}{\lemma{\skt{निद\appbrk{}ध्या\appbrk{}त्म\appbrk{}धु\appbrk{}सं\appbrk{}यु\appbrk{}तानि}}\Afootnote{K H; \skt{निद\appbrk{}ध्या\appbrk{}न्म\appbrk{}धु\appbrk{}सं\appbrk{}यु\appbrk{}तानि~।}\ A}}\skt{~। }\pend
\pstart\versepar \skt{}\edtext{\skt{एषोऽगदो}}{\lemma{\skt{ऽगदो}}\Afootnote{A H; \skt{गदो}\ K}}\skt{ हन्ति विषं प्रयुक्तः }\pend
\pstart\versepar \skt{पाना\-ञ्ज\-ना\-भ्य\-ञ्ज\-न\-न\-स्य\-योगैः~।}\pend

\pstart\versepar \ledsidenote{\passref{5.5.63}}\edtext{\skt{अवार्यवीर्यो}}{\lemma{\skt{अवार्यवीर्यो}}\Afootnote{A; \skt{अवारवीर्यो}\ K H}}\skt{ विषवेगहन्ता }\pend
\pstart\versepar \skt{महागदो नाम महाप्रभावः~॥ }\pend
\pstart\versepar \skt{}\edtext{\skt{विड\-ङ्ग\-पा\-ठा\-त्रि\-फ\-ला\-ज\-मोदा-}}{\lemma{\skt{विड\appbrk{}ङ्ग\appbrk{}पा\appbrk{}ठा\appbrk{}त्रि\appbrk{}फ\appbrk{}ला\appbrk{}ज\appbrk{}मोदा}}\Afootnote{A H; \skt{विड\appbrk{}ङ्ग\appbrk{}पा\appbrk{}ठ\appbrk{}त्रि\appbrk{}फ\appbrk{}ला\appbrk{}ज\appbrk{}मोद}\ K}}\skt{ }\pend
\pstart\versepar \skt{}\edtext{\skt{हिङ्गूनि}}{\lemma{\skt{हिङ्गूनि}}\Afootnote{A H; \skt{हिंगुनि}\ K}}\skt{ वक्रं }\edtext{\skt{त्रिकटुं}}{\lemma{\skt{त्रिकटुं}}\Afootnote{K; \skt{त्रिकटू}\ A; \skt{त्रिकटून्}\ H}}\skt{ }\edtext{\skt{तथैव}}{\lemma{\skt{तथैव}}\Afootnote{K H; \skt{नि चैव~॥}\ A}}\skt{~।}\pend

\pstart\versepar \ledsidenote{\passref{5.5.64}}\skt{सर्वश्च वर्गो लवणः }\edtext{\skt{सुसूक्ष्मः}}{\lemma{\skt{सुसूक्ष्मः}}\Afootnote{A H; \skt{ससूक्ष्मः}\ K}}\skt{ }\pend
\pstart\versepar \skt{}\edtext{\skt{सचि\-त्र\-क\-क्षौ\-द्र\-युतो}}{\lemma{\skt{सचि\appbrk{}त्र\appbrk{}क\appbrk{}क्षौ\appbrk{}द्र\appbrk{}युतो}}\Afootnote{K H; \skt{सचित्रकः क्षौद्रयुतो}\ A}}\skt{ }\edtext{\skt{निधेयः}}{\lemma{\skt{निधेयः}}\Afootnote{A K; \skt{विधेयः~।}\ H}}\skt{~। }\pend
\pstart\versepar \skt{शृङ्गे गवां शृङ्गमयेन चैव }\pend
\pstart\versepar \skt{प्रच्छादितः पक्ष\-मु\-पे\-क्षि\-तश्च~।}\pend

\pstart\versepar \ledsidenote{\passref{5.5.65}}\edtext{\skt{एषोऽगदः}}{\lemma{\skt{ऽगदः}}\Afootnote{A H; \skt{गद}\ K}}\skt{ स्थाव\-र\-ज\-ङ्ग\-मा\-नाञ् }\pend
\pstart\versepar \skt{जेता विषाणामजितो हि नाम्ना~॥ }\pend
\pstart\versepar \skt{प्रपौण्डरीकं सुरदारु }\edtext{\skt{रास्ना}}{\lemma{\skt{रास्ना}}\Afootnote{K H; \skt{मुस्ता}\ A}}\skt{ }\pend
\pstart\versepar \skt{}\edtext{\skt{कालानुसारी}}{\lemma{\skt{कालानुसारी}}\Afootnote{K H; \skt{कालानुसार्या}\ A}}\skt{ }\edtext{\skt{कटुरोहणीश्च}}{\lemma{\skt{कटुरोहणीश्}}\Afootnote{K; \skt{कटुरोहिणी}\ A; \skt{कटुरोहिणीञ्}\ H}}\skt{~।}\pend

\pstart\versepar \ledsidenote{\passref{5.5.66}}\edtext{\skt{स्थौणे\-य\-क\-ध्या\-म\-क\-प\-द्म\-कानि}}{\lemma{\skt{स्थौणे\appbrk{}य\appbrk{}क\appbrk{}ध्या\appbrk{}म\appbrk{}क\appbrk{}प\appbrk{}द्म\appbrk{}कानि}}\Afootnote{K H; \skt{स्थौणे\appbrk{}य\appbrk{}क\appbrk{}ध्या\appbrk{}म\appbrk{}क\appbrk{}गु\appbrk{}ग्गु\appbrk{}लूनि}\ A}}\skt{ }\pend
\pstart\versepar \skt{}\edtext{\skt{पुन्ना\-ग\-ता\-ली\-स\-सु\-व\-र्चि\-काश्च}}{\lemma{\skt{पुन्ना\appbrk{}ग\appbrk{}ता\appbrk{}ली\appbrk{}स\appbrk{}सु\appbrk{}व\appbrk{}र्चि\appbrk{}काश्च}}\Afootnote{K; \skt{पुन्ना\appbrk{}ग\appbrk{}ता\appbrk{}ली\appbrk{}श\appbrk{}सु\appbrk{}व\appbrk{}र्चि\appbrk{}काश्च~।}\ A; \skt{पुन्ना\appbrk{}ग\appbrk{}ता\appbrk{}ली\appbrk{}स\appbrk{}सु\appbrk{}व\appbrk{}र्च्चि\appbrk{}काञ्च~॥}\ H}}\skt{~। }\pend
\pstart\versepar \skt{}\edtext{\skt{कुट\-न्न\-टै\-ला\-सि\-त\-सि\-न्धु\-वाराः}}{\lemma{\skt{कुट\appbrk{}न्न\appbrk{}टै\appbrk{}ला\appbrk{}सि\appbrk{}त\appbrk{}सि\appbrk{}न्धु\appbrk{}वाराः}}\Afootnote{A H; \skt{कुट\appbrk{}न्न\appbrk{}टै\appbrk{}ला\appbrk{}सि\appbrk{}त\appbrk{}सि\appbrk{}न्दु\appbrk{}वाराः}\ K}}\skt{ }\pend
\pstart\versepar \skt{शैलेयकुष्ठे तगरं }\edtext{\skt{प्रियङ्गुः}}{\lemma{\skt{प्रियङ्गुः}}\Afootnote{A H; \skt{प्रियंङ्गुः~।}\ K}}\skt{~।}\pend

\pstart\versepar \ledsidenote{\passref{5.5.67}}\edtext{\skt{लोध्रं}}{\lemma{\skt{लोध्रं}}\Afootnote{K H; \skt{रोध्रं}\ A}}\skt{ }\edtext{\skt{तथा}}{\lemma{\skt{तथा}}\Afootnote{K H; \skt{जलं}\ A}}\skt{ }\edtext{\skt{गुग्गु\-ल\-गै\-रि\-कञ्च}}{\lemma{\skt{गुग्गु\appbrk{}ल\appbrk{}गै\appbrk{}रि\appbrk{}कञ्च}}\Afootnote{K; \skt{काञ्चनगैरिकं च}\ A; \skt{गुग्गु\appbrk{}लु\appbrk{}गै\appbrk{}रि\appbrk{}कञ्च}\ H}}\skt{ }\pend
\pstart\versepar \skt{}\edtext{\skt{ससैन्धवे}}{\lemma{\skt{ससैन्धवे}}\Afootnote{K H; \skt{समागधं चन्दनसैन्ध}\ A}}\skt{ }\edtext{\skt{पिप्पलिनागरे च}}{\lemma{\skt{पिप्पलिनागरे च}}\Afootnote{K H; \skt{वं च~।}\ A}}\skt{~। }\pend
\pstart\versepar \skt{}\edtext{\skt{सूक्ष्माणि}}{\lemma{\skt{सूक्ष्माणि}}\Afootnote{A K; \skt{शूक्ष्माणि}\ H}}\skt{ चूर्ण्णानि समानि कृत्वा }\pend
\pstart\versepar \skt{शृङ्गे }\edtext{\skt{निद\-ध्या\-त्म\-धु\-सं\-यु\-तानि}}{\lemma{\skt{निद\appbrk{}ध्या\appbrk{}त्म\appbrk{}धु\appbrk{}सं\appbrk{}यु\appbrk{}तानि}}\Afootnote{K H; \skt{निद\appbrk{}ध्या\appbrk{}न्म\appbrk{}धु\appbrk{}सं\appbrk{}यु\appbrk{}तानि~॥}\ A}}\skt{~।}\pend

\pstart\versepar \ledsidenote{\passref{5.5.68}}\edtext{\skt{एषोऽगदस्तार्क्ष्य}}{\lemma{\skt{ऽगदस्}}\Afootnote{A; \skt{गदस्}\ K H}}\skt{ इति प्रदिष्टो }\pend
\pstart\versepar \skt{विष\-न्नि\-ह\-न्या\-दपि तक्षकस्य~॥ }\pend
\pstart\versepar \skt{}\edtext{\skt{मांसी\-ह\-रे\-णु\-त्रि\-फ\-ला\-मु\-रुङ्गी}}{\lemma{\skt{मांसी\appbrk{}ह\appbrk{}रे\appbrk{}णु\appbrk{}त्रि\appbrk{}फ\appbrk{}ला\appbrk{}मु\appbrk{}रुङ्गी}}\Afootnote{\skt{मांसी\appbrk{}ह\appbrk{}रे\appbrk{}णु\appbrk{}त्रि\appbrk{}फ\appbrk{}ला\appbrk{}मु\appbrk{}रङ्गी}\ A; \skt{मांसी\appbrk{}ह\appbrk{}रे\appbrk{}णु\appbrk{}त्रिः\appbrk{}फ\appbrk{}ला\appbrk{}मु\appbrk{}रुङ्गी}\ K; \skt{मान्सी\appbrk{}ह\appbrk{}रे\appbrk{}णु\appbrk{}स्तृ\appbrk{}फ\appbrk{}ला\appbrk{}मु\appbrk{}रुङ्गी}\ H}}\skt{ }\pend
\pstart\versepar \skt{}\edtext{\skt{मञ्जि\-ष्ठ\-य\-ष्ट्या\-ह्व\-य\-प\-द्म\-कानि}}{\lemma{\skt{मञ्जि\appbrk{}ष्ठ\appbrk{}य\appbrk{}ष्ट्या\appbrk{}ह्व\appbrk{}य\appbrk{}प\appbrk{}द्म\appbrk{}कानि}}\Afootnote{K H; \skt{रक्ता\appbrk{}ल\appbrk{}ता\appbrk{}य\appbrk{}ष्टि\appbrk{}क\appbrk{}प\appbrk{}द्म\appbrk{}कानि~॥}\ A}}\skt{~।}\pend

\pstart\versepar \ledsidenote{\passref{5.5.69}}\edtext{\skt{विड\-ङ्ग\-ता\-ली\-स\-सु\-ग\-न्धि\-कैला}}{\lemma{\skt{विड\appbrk{}ङ्ग\appbrk{}ता\appbrk{}ली\appbrk{}स\appbrk{}सु\appbrk{}ग\appbrk{}न्धि\appbrk{}कैला}}\Afootnote{K H; \skt{विड\appbrk{}ङ्ग\appbrk{}ता\appbrk{}ली\appbrk{}श\appbrk{}सु\appbrk{}ग\appbrk{}न्धि\appbrk{}कैला}\ A}}\skt{ }\pend
\pstart\versepar \skt{}\edtext{\skt{त्वक्कुष्ठवक्राणि}}{\lemma{\skt{त्वक्कुष्ठवक्राणि}}\Afootnote{K H; \skt{त्वक्कुष्ठपत्राणि}\ A}}\skt{ सचन्दनानि~। }\pend
\pstart\versepar \skt{भार्गी }\edtext{\skt{पटोलीकिणिही}}{\lemma{\skt{पटोलीकिणिही}}\Afootnote{K H; \skt{पटोलं किणिही}\ A}}\skt{ सपाठा }\pend
\pstart\versepar \skt{मृगादनी }\edtext{\skt{क्रोष्टकमेखला}}{\lemma{\skt{क्रोष्टकमेखला}}\Afootnote{K H; \skt{कर्कटिका पुरश्}\ A}}\skt{ च~।}\pend

\pstart\versepar \ledsidenote{\passref{5.5.70}}\skt{पालिन्द्यशोकौ क्रमुकं }\edtext{\skt{सुरस्या}}{\lemma{\skt{सुरस्या}}\Afootnote{K; \skt{सुरस्याः}\ A H}}\skt{ }\pend
\pstart\versepar \skt{प्रसू\-न\-मा\-रु\-ष्क\-र\-जञ्च पुष्पम्~। }\pend
\pstart\versepar \skt{}\edtext{\edtext{\skt{चूर्ण्णान्यथैषां}}{\lemma{\skt{चूर्ण्णान्य्}}\Afootnote{K H; \skt{सूक्ष्मानि चूर्णानि}\ A}}}{\lemma{\skt{अथैषां}}\Afootnote{K H; \skt{स}\ A}}\skt{ }\edtext{\skt{निहितानि}}{\lemma{\skt{निहितानि}}\Afootnote{K H; \skt{मानि}\ A}}\skt{ शृङ्गे }\pend
\pstart\versepar \skt{}\edtext{\skt{देयानि}}{\lemma{\skt{देयानि}}\Afootnote{K H; \skt{न्यसेत्स}\ A}}\skt{ पित्तानि }\edtext{\skt{समाक्षिकानि}}{\lemma{\skt{समाक्षिकानि}}\Afootnote{K H; \skt{समाक्षिकाणि~॥}\ A}}\skt{~।}\pend

\pstart\versepar \ledsidenote{\passref{5.5.71}}\edtext{\skt{वरा\-ह\-गो\-धा\-शि\-खि\-श\-ल्य\-कानां}}{\lemma{\skt{वरा\appbrk{}ह\appbrk{}गो\appbrk{}धा\appbrk{}शि\appbrk{}खि\appbrk{}श\appbrk{}ल्य\appbrk{}कानां}}\Afootnote{K H; \skt{वरा\appbrk{}ह\appbrk{}गो\appbrk{}धा\appbrk{}शि\appbrk{}खि\appbrk{}श\appbrk{}ल्ल\appbrk{}कीनां}\ A}}\skt{ }\pend
\pstart\versepar \skt{मार्जारजं पार्षतनाकुले च~। }\pend
\pstart\versepar \skt{}\edtext{\skt{यस्यागदोऽयं}}{\lemma{\skt{ऽयं}}\Afootnote{A; \skt{यं}\ K H}}\skt{ सुकृतो }\edtext{\skt{गृहस्थो}}{\lemma{\skt{गृहस्थो}}\Afootnote{K H; \skt{गृहे स्यान्}\ A}}\skt{ }\pend
\pstart\versepar \skt{}\edtext{\skt{नाम्नार्षभो}}{\lemma{\skt{नाम्नार्षभो}}\Afootnote{K H; \skt{नाम्नर्षभो}\ A}}\skt{ नाम नरर्षभस्य~।}\pend

\pstart\versepar \ledsidenote{\passref{5.5.72}}\skt{न तत्र सर्पाः कुत एव कीटास् }\pend
\pstart\versepar \skt{त्यजन्ति वीर्याणि विषाणि चैव~। }\pend
\pstart\versepar \skt{एतेन भेर्यः पटहाश्च दिग्धाः }\pend
\pstart\versepar \skt{नानद्यमाना विषमाशु हन्युः~।}\pend

\pstart\versepar \ledsidenote{\passref{5.5.73}}\skt{दिग्धाः पताकाश्च निरीक्ष्य सद्यो }\pend
\pstart\versepar \skt{}\edtext{\skt{विषाभिभूताः}}{\lemma{\skt{विषाभिभूताः}}\Afootnote{K H; \skt{विषाभिभूता}\ A}}\skt{ }\edtext{\skt{सुखिनो}}{\lemma{\skt{सुखिनो}}\Afootnote{K H; \skt{ह्यविषा}\ A}}\skt{ भवन्ति~॥ }\pend
\pstart\versepar \skt{लाक्षा }\edtext{\skt{हरेण्वौ}}{\lemma{\skt{हरेण्वौ}}\Afootnote{K H; \skt{हरेणुर्}\ A}}\skt{ }\edtext{\skt{नलदप्रियङ्ग्वौ}}{\lemma{\skt{नलदप्रियङ्ग्वौ}}\Afootnote{H; \skt{नलदं प्रियङ्गुः}\ A; \skt{नलदप्रियग्वौ}\ K}}\skt{ }\pend
\pstart\versepar \skt{}\edtext{\skt{मञ्जि\-ष्ठ\-य\-ष्ट्या\-ह्व\-य\-पृ\-थ्वि\-काश्च}}{\lemma{\skt{मञ्जि\appbrk{}ष्ठ\appbrk{}य\appbrk{}ष्ट्या\appbrk{}ह्व\appbrk{}य\appbrk{}पृ\appbrk{}थ्वि\appbrk{}काश्}}\Afootnote{K H; \skt{शिग्रुद्वयं यष्टि\appbrk{}क\appbrk{}पृ\appbrk{}थ्वि\appbrk{}काश्}\ A}}\skt{~।}\pend

\pstart\versepar \ledsidenote{\passref{5.5.74}}\edtext{\skt{चूर्ण्णीकृतोऽयं}}{\lemma{\skt{ऽयं}}\Afootnote{A; \skt{यं}\ K H}}\skt{ }\edtext{\skt{रजनीविमिश्रो}}{\lemma{\skt{रजनीविमिश्रो}}\Afootnote{H; \skt{रजनीविमिश्रो}\ A; \skt{रजनाविमिश्रो}\ K}}\skt{ }\pend
\pstart\versepar \skt{}\edtext{\skt{वर्गो}}{\lemma{\skt{वर्गो}}\Afootnote{K H; \skt{सर्पिर्मधुभ्यां सहितो}\ A}}\skt{ निधेयो }\edtext{\skt{मधुसर्पिषाक्तः}}{\lemma{\skt{मधुसर्पिषाक्तः}}\Afootnote{K H; \emph{om.}\ A}}\skt{~। }\pend
\pstart\versepar \skt{शृङ्गे गवां पूर्ववदा }\edtext{\skt{पिधानस्}}{\lemma{\skt{पिधानस्}}\Afootnote{A H; \skt{पिधाणस्}\ K}}\skt{ }\pend
\pstart\versepar \skt{ततः }\edtext{\skt{प्रयो\-ज्योऽ\-ञ्ज\-न\-पा\-न\-नस्यैः}}{\lemma{\skt{ऽञ्ज\appbrk{}न\appbrk{}पा\appbrk{}न\appbrk{}नस्यैः}}\Afootnote{K H; \skt{ऽञ्ज\appbrk{}न\appbrk{}न\appbrk{}स्य\appbrk{}पानैः~॥}\ A}}\skt{~।}\pend

\pstart\versepar \ledsidenote{\passref{5.5.75}}\skt{सञ्जीवनो नाम गतासुकल्पान् }\pend
\pstart\versepar \skt{एषोऽगदो जीवयतीह मर्त्यान्~। }\pend
\pstart\versepar \skt{श्लेष्मा\-त\-की\-क\-ट्फ\-ल\-मा\-तु\-लुङ्ग }\pend
\pstart\versepar \skt{}\edtext{\skt{श्वेता}}{\lemma{\skt{श्वेता}}\Afootnote{A; \skt{स्वेता}\ K H}}\skt{ गिरिह्वा किणिही सिता च~॥}\edtext{}{\lemma{\skt{गता\appbrk{}सु\appbrk{}क\appbrk{}ल्पा\appbrk{}ने\appbrk{}षोऽ\appbrk{}गदो}}\Afootnote{A; \skt{गता\appbrk{}शु\appbrk{}क\appbrk{}ल्पा\appbrk{}न्गदो}\ K; \skt{गता\appbrk{}शु\appbrk{}क\appbrk{}ल्प…\appbrk{}मेषो गदो}\ H\textsuperscript{ac}; \skt{गता\appbrk{}सु\appbrk{}क\appbrk{}ल्प…\appbrk{}मेषो गदो}\ H\textsuperscript{pc}}}\edtext{}{\lemma{\skt{मर्त्यान्~। श्लेष्मा\appbrk{}त\appbrk{}की\appbrk{}क\appbrk{}ट्फ\appbrk{}ल\appbrk{}मा\appbrk{}तु\appbrk{}लुङ्ग}}\Afootnote{K; \skt{मर्त्यान्~। श्लेष्मा\appbrk{}त\appbrk{}की\appbrk{}क\appbrk{}ट्फ\appbrk{}ल\appbrk{}मा\appbrk{}तु\appbrk{}लुङ्ग्यः}\ A; \skt{मर्त्यः~॥ श्लेष्मा\appbrk{}न्त\appbrk{}की\appbrk{}क\appbrk{}ट्फ\appbrk{}ल\appbrk{}मा\appbrk{}तु\appbrk{}लुङ्गा}\ H}}\pend

\pstart\versepar \ledsidenote{\passref{5.5.76}}\edtext{\skt{सत\-ण्डु\-ली\-योऽ\-गद}}{\lemma{\skt{ऽगद}}\Afootnote{A; \skt{गद}\ K H}}\skt{ }\edtext{\skt{एष}}{\lemma{\skt{एष}}\Afootnote{A K; \skt{एव}\ H}}\skt{ मुख्यो }\pend
\pstart\versepar \skt{}\edtext{\skt{विषेषु}}{\lemma{\skt{विषेषु}}\Afootnote{A H; \skt{विशेषु}\ K}}\skt{ दर्वी\-क\-र\-रा\-जि\-ला\-नाम्~॥ }\pend
\pstart\versepar \skt{}\edtext{\skt{द्राक्षाश्वगन्धा}}{\lemma{\skt{द्राक्षाश्वगन्धा}}\Afootnote{K H; \skt{द्राक्षा सुगन्धा}\ A}}\skt{ }\edtext{\skt{गजवृत्तिका च}}{\lemma{\skt{गजवृत्तिका च}}\Afootnote{K H; \skt{नगवृत्तिका च}\ A}}\skt{ }\pend
\pstart\versepar \skt{}\edtext{\skt{श्वेता}}{\lemma{\skt{श्वेता}}\Afootnote{A; \skt{स्वेता}\ K H}}\skt{ }\edtext{\skt{च}}{\lemma{\skt{च}}\Afootnote{K; \skt{स}\ A; \skt{…}\ H}}\skt{ }\edtext{\skt{पिष्टा}}{\lemma{\skt{पिष्टा}}\Afootnote{K H; \skt{मङ्गा}\ A}}\skt{ }\edtext{\skt{समभागयुक्ताः}}{\lemma{\skt{समभागयुक्ताः}}\Afootnote{K H; \skt{समभागयुक्ता~॥}\ A}}\skt{~॥}\pend

\pstart\versepar \ledsidenote{\passref{5.5.77}}\skt{देयो द्विभागः }\edtext{\skt{सुरसच्छदस्य}}{\lemma{\skt{सुरसच्छदस्य}}\Afootnote{K H; \skt{सुरसाच्छदस्य}\ A}}\skt{ }\pend
\pstart\versepar \skt{}\edtext{\skt{कपि\-त्थ\-बि\-ल्वा\-दपि}}{\lemma{\skt{कपि\appbrk{}त्थ\appbrk{}बि\appbrk{}ल्वा\appbrk{}दपि}}\Afootnote{A; \skt{कपि\appbrk{}त्थ\appbrk{}वि\appbrk{}ल्वा\appbrk{}दपि}\ K H}}\skt{ दाडिमाच्च~। }\pend
\pstart\versepar \skt{तथा }\edtext{\skt{च}}{\lemma{\skt{च}}\Afootnote{K H; \skt{र्ध}\ A}}\skt{ }\edtext{\skt{भागो}}{\lemma{\skt{भागो}}\Afootnote{K H; \skt{भागः}\ A}}\skt{ सित\-सि\-न्दु\-वा\-राद् }\pend
\pstart\versepar \skt{अङ्को\-ल्ल\-बी\-जा\-दपि गैरिकाच्च~।}\edtext{}{\lemma{\skt{सित\appbrk{}सि\appbrk{}न्दु\appbrk{}वा\appbrk{}रा\appbrk{}द\appbrk{}ङ्को\appbrk{}ल्ल\appbrk{}बी\appbrk{}जा\appbrk{}दपि}}\Afootnote{\skt{सित\appbrk{}सि\appbrk{}न्धु\appbrk{}वा\appbrk{}रा\appbrk{}द\appbrk{}ङ्को\appbrk{}ठ\appbrk{}मू\appbrk{}ला\appbrk{}दपि}\ A; \skt{सित\appbrk{}सि\appbrk{}न्दु\appbrk{}वा\appbrk{}रा\appbrk{}द\appbrk{}ङ्को\appbrk{}ल्ल\appbrk{}वी\appbrk{}जा\appbrk{}दपि}\ K H}}\pend

\pstart\versepar \ledsidenote{\passref{5.5.78ab}}\edtext{\skt{एषोऽगदः}}{\lemma{\skt{ऽगदः}}\Afootnote{A H; \skt{गदो}\ K}}\skt{ क्षौद्रयुतो निहन्ति }\pend
\pstart\versepar \skt{}\edtext{\skt{विशेषतो}}{\lemma{\skt{विशेषतो}}\Afootnote{A H; \skt{विसेषतो}\ K}}\skt{ मण्ड\-लि\-ना\-म्वि\-षाणि~॥}\pend

\pstart \ledsidenote{\passref{5.5.84ab}}\edtext{\skt{सोम\-रा\-जी\-य\-ब\-हुला}}{\lemma{\skt{सोम\appbrk{}रा\appbrk{}जी\appbrk{}य\appbrk{}ब\appbrk{}हुला}}\Afootnote{\skt{सोमराजीफलं पुष्पं}\ A; \skt{सोम\appbrk{}रा\appbrk{}जी\appbrk{}य\appbrk{}व\appbrk{}हुला}\ K; \skt{सोम\appbrk{}रा\appbrk{}जी\appbrk{}य\appbrk{}व\appbrk{}फला}\ H}}\skt{ }\edtext{\skt{कदली}}{\lemma{\skt{कदली}}\Afootnote{K H; \skt{कटभी}\ A}}\skt{ }\edtext{\skt{सिन्दुवारकः}}{\lemma{\skt{सिन्दुवारकः}}\Afootnote{K H; \skt{सिन्धुवारकः~।}\ A}}\skt{~।}\pend

\pstart \ledsidenote{\passref{5.5.85cd}}\edtext{\skt{श्यामाम्बष्ठा}}{\lemma{\skt{श्यामाम्बष्ठा}}\Afootnote{A; \skt{श्यामाम्वष्ठा}\ K H}}\skt{ }\edtext{\skt{तालपत्री}}{\lemma{\skt{तालपत्री}}\Afootnote{K H; \skt{विडङ्गानि}\ A}}\skt{ }\edtext{\edtext{\skt{तथा\-म्रा\-श्म\-न्त\-कोऽपि}}{\lemma{\skt{तथाम्राश्मन्तको}}\Afootnote{K H; \skt{तथाम्राश्मन्तका}\ A}}}{\lemma{\skt{ऽपि}}\Afootnote{K H; \skt{नि}\ A}}\skt{ च~।}\pend

\pstart \ledsidenote{\passref{5.5.85.add-1}}\skt{मण्डूकपर्ण्णी वरुणः सप्तला स पुनर्ण्णवा~॥}\pend

\pstart \ledsidenote{\passref{5.5.84cd}}\edtext{\skt{चोरको}}{\lemma{\skt{चोरको}}\Afootnote{K H; \skt{चोरको वरुणः}\ A}}\skt{ }\edtext{\skt{नागविन्ना}}{\lemma{\skt{नागविन्ना}}\Afootnote{K H; \skt{कुष्ठं सर्पगन्धा स}\ A}}\skt{ }\edtext{\skt{च}}{\lemma{\skt{च}}\Afootnote{K H; \skt{स}\ A}}\skt{ }\edtext{\skt{तथा}}{\lemma{\skt{तथा}}\Afootnote{K H; \skt{प्तला~॥}\ A}}\skt{ }\edtext{\skt{सर्पसुगन्धिका}}{\lemma{\skt{सर्पसुगन्धिका}}\Afootnote{K H; \emph{om.}\ A}}\skt{~।}\pend

\pstart\versepar \ledsidenote{\passref{5.5.86}}\skt{भूमी }\edtext{\skt{कुरवकश्चैव}}{\lemma{\skt{कुरवकश्}}\Afootnote{K H; \skt{कुरबकश्}\ A}}\skt{ गण }\edtext{\skt{एकरसः}}{\lemma{\skt{एकरसः}}\Afootnote{H; \skt{एकसरः}\ A; \skt{एकरस}\ K}}\skt{ }\edtext{\skt{स्मृतः}}{\lemma{\skt{स्मृतः}}\Afootnote{A H; \skt{स्मृताः~॥}\ K}}\skt{~॥ }\pend
\pstart\versepar \skt{}\edtext{\skt{एकैकशो}}{\lemma{\skt{एकैकशो}}\Afootnote{K H; \skt{एकशो}\ A}}\skt{ }\edtext{\skt{द्वन्द्वशो वा}}{\lemma{\skt{द्वन्द्वशो वा}}\Afootnote{K H; \skt{द्वित्रिशो वा}\ A}}\skt{ }\edtext{\skt{प्रयोक्तव्यो}}{\lemma{\skt{प्रयोक्तव्यो}}\Afootnote{K H; \skt{ऽपि प्रयोक्तव्यो}\ A}}\skt{ }\edtext{\skt{विषापहः}}{\lemma{\skt{विषापहः}}\Afootnote{A; \skt{विषापहमिति~॥}\ K; \skt{विषापह इति~॥}\ H}}\skt{~॥}\pend

\pstart \ledsidenote{\passref{5.5.trailer}}\edtext{\skt{॥ इति}}{\lemma{\skt{॥ इति}}\Afootnote{A; \emph{om.}\ K H}}\skt{ }\edtext{\skt{कल्पे}}{\lemma{\skt{कल्पे}}\Afootnote{\skt{सुश्रु\appbrk{}त\appbrk{}सं\appbrk{}हि\appbrk{}तायां कल्पस्थाने सर्प\appbrk{}द\appbrk{}ष्ट\appbrk{}वि\appbrk{}ष\appbrk{}चि\appbrk{}कि\appbrk{}त्सितं नाम पञ्चमोऽध्यायः~॥५~॥}\ A; \skt{कल्पे ५~॥ {\symfont ❈}~॥}\ K; \skt{कल्पस्थाने पञ्चमोऽध्यायः~॥}\ H}}\skt{ ५~॥}\pend

\endnumbering
\beginnumbering
\eledsection*{Adhyāya 6}
\pstart \ledsidenote{\passref{5.7.1}}\skt{अथातो }\edtext{\skt{मूषिकाकल्पं}}{\lemma{\skt{मूषिकाकल्पं}}\Afootnote{K H; \skt{मूषिककल्पं}\ A}}\skt{ }\edtext{\skt{व्याख्यास्यामः}}{\lemma{\skt{व्याख्यास्यामः}}\Afootnote{A H; \emph{om.}\ K}}\skt{~॥}\pend

\pstart\versepar \ledsidenote{\passref{5.7.3}}\edtext{\skt{पूर्वमुक्ताः}}{\lemma{\skt{उक्ताः}}\Afootnote{K H; \emph{om.}\ A}}\skt{ शुक्रविषाः }\edtext{\skt{मूषिका ये}}{\lemma{\skt{मूषिका ये}}\Afootnote{K H; \skt{उक्ता मूषिका ये}\ A}}\skt{ समासतः~। }\pend
\pstart\versepar \skt{नाम\-ल\-क्ष\-ण\-भै\-ष\-ज्यै\-र\-ष्टा\-दश }\edtext{\skt{निबोध}}{\lemma{\skt{निबोध}}\Afootnote{A K; \skt{निवोध}\ H}}\skt{ }\edtext{\skt{ताः}}{\lemma{\skt{ताः}}\Afootnote{K; \skt{मे~॥}\ A; \skt{तां~॥}\ H}}\skt{~॥}\pend

\pstart\versepar \ledsidenote{\passref{5.7.4}}\skt{लालनः पुत्रकः }\edtext{\skt{कृष्णो}}{\lemma{\skt{कृष्णो}}\Afootnote{A H; \skt{कृष्ण}\ K}}\skt{ }\edtext{\edtext{\skt{वसि\-र\-श्चि\-क्कि\-र\-स्तथा}}{\lemma{\skt{वसिरश्}}\Afootnote{K H; \skt{हंसिरश्}\ A}}}{\lemma{\skt{चिक्किरस्}}\Afootnote{K H; \skt{चिक्वि(क्कि)रस्}\ A}}\skt{~। }\pend
\pstart\versepar \skt{}\edtext{\edtext{\skt{च्छुच्छु\-न्द\-रोऽ\-र\-ल\-श्चैव}}{\lemma{\skt{च्छुच्छुन्दरो}}\Afootnote{A K; \skt{जायते}\ H}}}{\lemma{\skt{ऽरलश्}}\Afootnote{\skt{ऽलसश्}\ A; \skt{रलश्}\ K; \skt{गरलश्}\ H}}\skt{ कषा\-य\-द\-श\-नोऽपि च~।}\pend

\pstart\versepar \ledsidenote{\passref{5.7.5}}\edtext{\edtext{\skt{कुलि\-ङ्ग\-श्चा\-जि\-त\-श्चैव}}{\lemma{\skt{कुलिङ्गश्}}\Afootnote{A K; \skt{कुलिङ्गाश्}\ H}}}{\lemma{\skt{चाजितश्}}\Afootnote{A K; \skt{चाजिवश्}\ H}}\skt{ चपलः कपिलस्तथा~। }\pend
\pstart\versepar \skt{}\edtext{\edtext{\skt{कोकि\-लोऽ\-रु\-ण\-सं\-ज्ञश्च}}{\lemma{\skt{कोकिलो}}\Afootnote{A; \skt{कोकिला}\ K H}}}{\lemma{\skt{ऽरुणसंज्ञश्च}}\Afootnote{A; \skt{रुणसंज्ञश्च}\ K; \skt{रुणसंज्ञाश्च}\ H}}\skt{ }\edtext{\skt{महा\-कृ\-ष्ण\-स्त\-थो\-न्दुरुः}}{\lemma{\skt{तथोन्दुरुः}}\Afootnote{K H; \skt{तथोन्दुरः~॥}\ A}}\skt{~।}\pend

\pstart\versepar \ledsidenote{\passref{5.7.6}}\edtext{\skt{श्वेतश्च}}{\lemma{\skt{श्वेतश्च}}\Afootnote{K H; \skt{श्वेतेन}\ A}}\skt{ महता सार्धं कपिलेनाखुना तथा~। }\pend
\pstart\versepar \skt{मूषिकश्च }\edtext{\skt{कपो\-ता\-भ\-स्त\-थै\-वा\-ष्टा\-दश}}{\lemma{\skt{तथैवाष्टादश}}\Afootnote{A H; \skt{तथैवाष्टादशः}\ K}}\skt{ स्मृताः~॥}\pend

\pstart\versepar \ledsidenote{\passref{5.7.7}}\skt{शुक्रं पतति }\edtext{\skt{यत्रैषां}}{\lemma{\skt{यत्रैषां}}\Afootnote{A K; \skt{यत्रैषा}\ H}}\skt{ शुक्रस्पृष्टैः स्पृशन्ति वा~। }\pend
\pstart\versepar \skt{}\edtext{\skt{नख\-द\-न्ता\-दि\-भि\-स्तस्मिं}}{\lemma{\skt{तस्मिं}}\Afootnote{K; \skt{तस्मि}\ A H}}\skt{ }\edtext{\skt{गात्रे}}{\lemma{\skt{गात्रे}}\Afootnote{K; \skt{न्गात्रे}\ A H}}\skt{ रक्तं प्रदुष्यति~।}\pend

\pstart\versepar \ledsidenote{\passref{5.7.8}}\edtext{\skt{जायन्ते}}{\lemma{\skt{जायन्ते}}\Afootnote{A K; \skt{जायते}\ H}}\skt{ ग्रन्थयः शोफाः }\edtext{\skt{कर्ण्णिका}}{\lemma{\skt{कर्ण्णिका}}\Afootnote{A K\textsuperscript{pc}; \skt{किकर्ण्णिका}\ K\textsuperscript{ac}; \skt{कणिका}\ H\textsuperscript{ac}; \skt{कणिर्ण्णिका}\ H\textsuperscript{pc}}}\skt{ मण्डलानि च~। }\pend
\pstart\versepar \skt{}\edtext{\edtext{\skt{पिट\-को\-प\-च\-या\-श्चोग्रा}}{\lemma{\skt{पिट\appbrk{}को\appbrk{}प\appbrk{}च\appbrk{}याश्}}\Afootnote{K H; \skt{पीड\appbrk{}को\appbrk{}प\appbrk{}च\appbrk{}यश्}\ A}}}{\lemma{\skt{चोग्रा}}\Afootnote{K H; \skt{चोग्रो}\ A}}\skt{ }\edtext{\skt{विसर्पाः}}{\lemma{\skt{विसर्पाः}}\Afootnote{A K; \skt{विसर्ग्गाः}\ H}}\skt{ किटिभानि च~।}\pend

\pstart\versepar \ledsidenote{\passref{5.7.9}}\skt{पर्वभेदो }\edtext{\skt{रुजश्चापि}}{\lemma{\skt{चापि}}\Afootnote{K H; \skt{ती}\ A}}\skt{ }\edtext{\skt{ज्वरो}}{\lemma{\skt{ज्वरो}}\Afootnote{K H; \skt{व्रा}\ A}}\skt{ मूर्च्छा }\edtext{\skt{च}}{\lemma{\skt{च}}\Afootnote{K H; \skt{ऽङ्गसदनं ज्वरः~।}\ A}}\skt{ }\edtext{\skt{दारुणाः}}{\lemma{\skt{दारुणाः}}\Afootnote{K H; \emph{om.}\ A}}\skt{~। }\pend
\pstart\versepar \skt{}\edtext{\skt{दौर्बल्यमरुचिः}}{\lemma{\skt{दौर्बल्यम्}}\Afootnote{A; \skt{दौर्वल्यम्}\ K H}}\skt{ }\edtext{\skt{सादो}}{\lemma{\skt{सादो}}\Afootnote{K H; \skt{श्वासो}\ A}}\skt{ वम\-थु\-र्लो\-म\-ह\-र्ष\-णम्~।}\pend

\pstart\versepar \ledsidenote{\passref{5.7.10}}\skt{दष्टरूपं समा\-सो\-क्त\-मे\-त\-द्व्या\-स\-मतः शृणु~॥ }\pend
\pstart\versepar \skt{लालास्रावो }\edtext{\skt{लालनेन}}{\lemma{\skt{लालनेन}}\Afootnote{A; \skt{लालना}\ K; \skt{लालना च}\ H}}\skt{ }\edtext{\skt{च्छर्दि}}{\lemma{\skt{च्छर्दि}}\Afootnote{K H; \emph{om.}\ A}}\skt{ }\edtext{\skt{हिक्का}}{\lemma{\skt{हिक्का}}\Afootnote{A K; \skt{हिक्काञ्}\ H}}\skt{ }\edtext{\skt{च}}{\lemma{\skt{च}}\Afootnote{K H; \skt{छर्दिश्च}\ A}}\skt{ जायते~।}\pend

\pstart\versepar \ledsidenote{\passref{5.7.11}}\skt{तण्डु\-ली\-य\-क\-कल्कं तु लिह्यात्तत्र समाक्षिकम्~॥ }\pend
\pstart\versepar \skt{}\edtext{\skt{पुत्र\-के\-णा\-ङ्ग\-सं\-सादः}}{\lemma{\skt{पुत्र\appbrk{}के\appbrk{}णा\appbrk{}ङ्ग\appbrk{}सं\appbrk{}सादः}}\Afootnote{K H; \skt{पुत्र\appbrk{}के\appbrk{}णा\appbrk{}ङ्ग\appbrk{}सा\appbrk{}दश्च}\ A}}\skt{ }\edtext{\skt{पाण्डुवल्गुश्च}}{\lemma{\skt{पाण्डुवल्गुश्च}}\Afootnote{K H; \skt{पाण्डुवर्णश्च}\ A}}\skt{ जायते~।}\pend

\pstart\versepar \ledsidenote{\passref{5.7.12}}\skt{चीयते ग्रन्थिभिश्चाङ्गं }\edtext{\edtext{\skt{शिशु\-र्मू\-षि\-क\-सं\-स्थितैः}}{\lemma{\skt{शिशुर्}}\Afootnote{K H; \skt{आखु}\ A}}}{\lemma{\skt{मूषिकसंस्थितैः}}\Afootnote{K H; \skt{शावकसन्निभैः~।}\ A}}\skt{~।}\skt{ }\pend
\pstart\versepar \skt{}\edtext{\skt{शिरी\-षे\-ङ्गु\-दि\-पत्रं तु}}{\lemma{\skt{शिरी\appbrk{}षे\appbrk{}ङ्गु\appbrk{}दि\appbrk{}पत्रं तु}}\Afootnote{H; \skt{शिरी\appbrk{}षे\appbrk{}ङ्गु\appbrk{}द\appbrk{}कल्कं तु}\ A; \skt{शिरी\appbrk{}शे\appbrk{}ङ्गु\appbrk{}दि\appbrk{}पत्रं तु}\ K}}\skt{ लिह्यात्तत्र समाक्षिकम्~॥}\pend

\pstart\versepar \ledsidenote{\passref{5.7.13}}\edtext{\skt{कृष्णे\-ना\-सृ\-क्छ\-र्द\-यति}}{\lemma{\skt{कृष्णे\appbrk{}ना\appbrk{}सृ\appbrk{}क्छ\appbrk{}र्द\appbrk{}यति}}\Afootnote{K H; \skt{कृष्णेन दंशे शोफोऽसृक्छर्दिः प्रायश्च}\ A}}\skt{ }\edtext{\skt{दुर्दिने तु}}{\lemma{\skt{दुर्दिने तु}}\Afootnote{K H; \skt{दुर्दिने~।}\ A}}\skt{ }\edtext{\skt{विशेषतः}}{\lemma{\skt{विशेषतः}}\Afootnote{K H; \emph{om.}\ A}}\skt{~। }\pend
\pstart\versepar \skt{}\edtext{\skt{शिरीषपत्रं}}{\lemma{\skt{शिरीषपत्रं}}\Afootnote{H; \skt{शिरीषफल}\ A; \skt{शिरीषपत्र}\ K}}\skt{ }\edtext{\skt{कुष्ठैलाः}}{\lemma{\skt{कुष्ठैलाः}}\Afootnote{\skt{कुष्ठं तु}\ A; \skt{कुष्ठैला}\ K H}}\skt{ }\edtext{\skt{पिबे\-त्किं\-शु\-क\-भ\-स्मना}}{\lemma{\skt{पिबे\appbrk{}त्किं\appbrk{}शु\appbrk{}क\appbrk{}भ\appbrk{}स्मना}}\Afootnote{A; \skt{पिवे\appbrk{}त्किं\appbrk{}शु\appbrk{}क\appbrk{}भ\appbrk{}स्मना~॥}\ K; \skt{पिवे\appbrk{}त्किं\appbrk{}शु\appbrk{}क\appbrk{}भ\appbrk{}श्मना~॥}\ H}}\skt{~॥}\pend

\pstart\versepar \ledsidenote{\passref{5.7.14}}\edtext{\skt{वसि\-रे\-णा\-न्न\-वि\-द्वेषो}}{\lemma{\skt{वसि\appbrk{}रे\appbrk{}णा\appbrk{}न्न\appbrk{}वि\appbrk{}द्वेषो}}\Afootnote{H; \skt{हंसि\appbrk{}रे\appbrk{}णा\appbrk{}न्न\appbrk{}वि\appbrk{}द्वेषो}\ A; \skt{वसि\appbrk{}रो\appbrk{}णा\appbrk{}न्न\appbrk{}वि\appbrk{}द्वेषो}\ K}}\skt{ }\edtext{\skt{जृम्भो}}{\lemma{\skt{जृम्भो}}\Afootnote{K H; \skt{जृम्भा}\ A}}\skt{ }\edtext{\skt{रोम्नाञ्च}}{\lemma{\skt{रोम्नाञ्च}}\Afootnote{K H; \skt{रोम्णां च}\ A}}\skt{ }\edtext{\skt{कुष्ठता}}{\lemma{\skt{कुष्ठता}}\Afootnote{K H; \skt{हर्षणम्~।}\ A}}\skt{~। }\pend
\pstart\versepar \skt{}\edtext{\skt{पिबे\-दा\-र\-ग्व\-धा\-दिन्तु}}{\lemma{\skt{पिबे\appbrk{}दा\appbrk{}र\appbrk{}ग्व\appbrk{}धा\appbrk{}दिन्तु}}\Afootnote{A; \skt{पिवे\appbrk{}दा\appbrk{}र\appbrk{}ग्व\appbrk{}धा\appbrk{}दिन्तु}\ K H}}\skt{ }\edtext{\edtext{\skt{वान्तस्तत्राशु}}{\lemma{\skt{वान्तस्}}\Afootnote{K H; \skt{सुवान्तस्}\ A}}}{\lemma{\skt{तत्राशु}}\Afootnote{K H; \skt{तत्र}\ A}}\skt{ मानवः~॥}\pend

\pstart\versepar \ledsidenote{\passref{5.7.15}}\edtext{\skt{चिक्किरेण}}{\lemma{\skt{चिक्किरेण}}\Afootnote{K H; \skt{चिक्विरेण (चिक्किरेण)}\ A}}\skt{ शिरोदुःखं शोफो हिक्का वमिस्तथा~। }\pend
\pstart\versepar \skt{}\edtext{\skt{सुवान्तो}}{\lemma{\skt{सुवान्तो}}\Afootnote{K; \emph{om.}\ A; \skt{वासन्तो}\ H}}\skt{ }\edtext{\skt{जालिनीक्वाथैः}}{\lemma{\skt{जालिनीक्वाथैः}}\Afootnote{K H; \skt{जालिनी}\ A}}\skt{ }\edtext{\edtext{\edtext{\skt{सार\-म\-ङ्को\-ल्ल\-ज\-म्पि\-बेत्}}{\lemma{\skt{सारम्}}\Afootnote{K H; \emph{om.}\ A}}}{\lemma{\skt{अङ्कोल्लजम्}}\Afootnote{K H; \emph{om.}\ A}}}{\lemma{\skt{पिबेत्}}\Afootnote{\skt{मद\appbrk{}ना\appbrk{}ङ्को\appbrk{}ठ\appbrk{}क\appbrk{}षा\appbrk{}यै\appbrk{}र्वा\appbrk{}म\appbrk{}येत्तु तम्~॥}\ A; \skt{पिवेत्}\ K H}}\skt{~॥}\pend

\pstart \ledsidenote{\passref{5.7.16cd}}\skt{छुच्छुन्दरेण }\edtext{\skt{विट्स\-ङ्ग\-ग्री\-वा\-स्त\-म्भ\-वि\-जृ\-म्भिकाः}}{\lemma{\skt{विट्स\appbrk{}ङ्ग\appbrk{}ग्री\appbrk{}वा\appbrk{}स्त\appbrk{}म्भ\appbrk{}वि\appbrk{}जृ\appbrk{}म्भिकाः}}\Afootnote{\skt{तट्छर्दिर्ज्वरो दौर्बल्यमेव च~॥}\ A; \skt{विट्छ\appbrk{}ङ्गः\appbrk{}ग्री\appbrk{}वा\appbrk{}स्त\appbrk{}म्भ\appbrk{}वि\appbrk{}जृ\appbrk{}म्भिकाः~।}\ K; \skt{विट्छङ्गाः ग्रीवा\appbrk{}स्त\appbrk{}म्भ\appbrk{}वि\appbrk{}जृ\appbrk{}म्भिकाः}\ H\textsuperscript{ac} H\textsuperscript{pc}}}\skt{~।}\pend

\pstart \ledsidenote{\passref{5.7.16ab}}\edtext{\skt{यव\-ना\-ला\-र्ष\-भी\-क्षारं}}{\lemma{\skt{यव\appbrk{}ना\appbrk{}ला\appbrk{}र्ष\appbrk{}भी\appbrk{}क्षारं}}\Afootnote{K H; \skt{यव\appbrk{}ना\appbrk{}ल\appbrk{}र्ष\appbrk{}भी\appbrk{}क्षारं}\ A}}\skt{ }\edtext{\skt{वृहत्यौ}}{\lemma{\skt{वृहत्यौ}}\Afootnote{K; \skt{बृहत्योश्}\ A; \skt{वृहत्योश्}\ H}}\skt{ }\edtext{\skt{चात्र}}{\lemma{\skt{चात्र}}\Afootnote{A H; \skt{…}\ K}}\skt{ }\edtext{\skt{दापयेत्}}{\lemma{\skt{दापयेत्}}\Afootnote{A H; \emph{om.}\ K}}\skt{~॥}\pend

\pstart \ledsidenote{\passref{5.7.18cd}}\edtext{\skt{ग्रीवा\-स्त\-म्भोऽ\-र\-ले\-नाथ}}{\lemma{\skt{ऽरलेनाथ}}\Afootnote{\skt{ऽलसेनोर्ध्व}\ A; \skt{रलेनाथ}\ H}}\skt{ }\edtext{\edtext{\skt{रुजा\-श्चा\-दं\-श\-म\-ण्डले}}{\lemma{\skt{रुजाश्}}\Afootnote{H; \skt{वायुर्दंशे रुजा}\ A}}}{\lemma{\skt{चादंशमण्डले}}\Afootnote{H; \skt{ज्वरः~॥}\ A}}\skt{~।}\pend

\pstart\versepar \ledsidenote{\passref{5.7.19}}\edtext{\edtext{\skt{महा\-ग\-द\-म्म\-हा\-वीर्यं}}{\lemma{\skt{महागदम्}}\Afootnote{A H; \skt{…}\ K}}}{\lemma{\skt{महावीर्यं}}\Afootnote{H; \skt{ससर्पिष्कं}\ A; \emph{om.}\ K}}\skt{ }\edtext{\skt{लिह्यात्तत्र}}{\lemma{\skt{लिह्यात्तत्र}}\Afootnote{A H; \skt{…तत्र}\ K}}\skt{ समाक्षिकम्~॥ }\pend
\pstart\versepar \skt{निद्रा कषायदन्तेन }\edtext{\skt{जायते}}{\lemma{\skt{जायते}}\Afootnote{K H; \skt{हृच्छोषः}\ A}}\skt{ कार्श्यमेव च~।}\pend

\pstart\versepar \ledsidenote{\passref{5.7.20}}\edtext{\skt{लिह्यात्तत्र}}{\lemma{\skt{लिह्यात्तत्र}}\Afootnote{K H; \skt{क्षौद्रोपेताः}\ A}}\skt{ शिरीषस्य }\edtext{\skt{मधुना}}{\lemma{\skt{मधुना}}\Afootnote{K H; \skt{लिह्यात्सा}\ A}}\skt{ सारमाषकान्~॥ }\pend
\pstart\versepar \skt{कुलिङ्गेन रुजः शोफो }\edtext{\skt{राज्य\-श्चा\-दं\-श\-म\-ण्डले}}{\lemma{\skt{चादंशमण्डले}}\Afootnote{K H; \skt{च दंशमण्डले~॥}\ A}}\skt{~।}\edtext{}{\lemma{\skt{सारमाषकान्~॥ कुलिङ्गेन}}\Afootnote{K H; \skt{रफलत्वचः~। कुलिङ्गेन}\ A}}\pend

\pstart\versepar \ledsidenote{\passref{5.7.21}}\skt{सहे }\edtext{\skt{ससिन्दुवारे च}}{\lemma{\skt{ससिन्दुवारे च}}\Afootnote{K H; \skt{ससिन्धुवारे च}\ A}}\skt{ लिह्यात्तत्र समाक्षिके~॥ }\pend
\pstart\versepar \skt{}\edtext{\skt{अजितेन}}{\lemma{\skt{अजितेन}}\Afootnote{K H; \skt{अजितेनाङ्}\ A}}\skt{ }\edtext{\skt{वमी}}{\lemma{\skt{वमी}}\Afootnote{K H; \skt{गकृष्णत्वं छर्दिर्}\ A}}\skt{ }\edtext{\skt{मूर्च्छा}}{\lemma{\skt{मूर्च्छा}}\Afootnote{K H; \skt{मूर्च्छा च}\ A}}\skt{ हृद्ग्रहः }\edtext{\skt{कृष्णगात्रता}}{\lemma{\skt{कृष्णगात्रता}}\Afootnote{H; \emph{om.}\ A; \skt{…}\ K}}\skt{~।}\pend

\pstart\versepar \ledsidenote{\passref{5.7.22}}\edtext{\skt{तत्र}}{\lemma{\skt{तत्र}}\Afootnote{H; \emph{om.}\ A K}}\skt{ }\edtext{\skt{स्नुहाक्षीरयुक्तां}}{\lemma{\skt{स्नुहाक्षीरयुक्तां}}\Afootnote{H; \skt{स्नुक्क्षीरपिष्टां पालिन्दीं}\ A; \skt{…रयुक्तां}\ K}}\skt{ मञ्जिष्ठां मधुना लिहेत्~॥ }\pend
\pstart\versepar \skt{चपलेन भवेच्छर्दिर्मूर्च्छा च सह तृष्णया~।}\pend

\pstart\versepar \ledsidenote{\passref{5.7.23}}\edtext{\skt{सभस्मकाष्ठां}}{\lemma{\skt{सभस्मकाष्ठां}}\Afootnote{H; \emph{om.}\ A; \skt{सभस्मकाष्ठा}\ K}}\skt{ }\edtext{\skt{सजटां}}{\lemma{\skt{सजटां}}\Afootnote{K H; \emph{om.}\ A}}\skt{ क्षौद्रेण }\edtext{\skt{त्रिफलां}}{\lemma{\skt{त्रिफलां}}\Afootnote{A H; \skt{त्रिफलं}\ K}}\skt{ पिबेत्~॥ }\pend
\pstart\versepar \skt{कपिलेन }\edtext{\skt{व्रणं}}{\lemma{\skt{व्रणं}}\Afootnote{K; \skt{व्रणे}\ A H}}\skt{ }\edtext{\skt{कोथं}}{\lemma{\skt{कोथं}}\Afootnote{K H; \skt{कोथो}\ A}}\skt{ ज्वरो }\edtext{\edtext{\skt{ग्रन्थ्युद्गमस्तथा}}{\lemma{\skt{ग्रन्थ्युद्गमस्}}\Afootnote{H; \skt{ग्रन्थ्युद्गमः}\ A; \skt{ग्रंन्थ्युद्गमस्}\ K}}}{\lemma{\skt{तथा}}\Afootnote{K H; \skt{सतृट्~॥}\ A}}\skt{~।}\edtext{}{\lemma{\skt{पिबेत्~॥ कपिलेन}}\Afootnote{\skt{लिह्या\appbrk{}द्भ\appbrk{}द्र\appbrk{}का\appbrk{}ष्ठ\appbrk{}ज\appbrk{}टा\appbrk{}न्वि\appbrk{}ताम्~। कपिलेन}\ A; \skt{पिवेत्~॥ कपिलेन}\ K H}}\pend

\pstart\versepar \ledsidenote{\passref{5.7.24}}\edtext{\skt{क्षौद्रेण}}{\lemma{\skt{क्षौद्रेण}}\Afootnote{K H; \emph{om.}\ A}}\skt{ }\edtext{\skt{लिह्याच्छ्वेतात्र}}{\lemma{\skt{लिह्याच्छ्वेतात्र}}\Afootnote{K H; \skt{लिह्यान्मधुयुतां श्वेतां}\ A}}\skt{ }\edtext{\skt{श्वेता}}{\lemma{\skt{श्वेता}}\Afootnote{K H; \skt{श्वेतां}\ A}}\skt{ }\edtext{\skt{वापि}}{\lemma{\skt{वापि}}\Afootnote{K H; \skt{चापि}\ A}}\skt{ }\edtext{\skt{पुनर्ण्णवा}}{\lemma{\skt{पुनर्ण्णवा}}\Afootnote{K H; \skt{पुनर्नवाम्~।}\ A}}\skt{~॥ }\pend
\pstart\versepar \skt{ग्रन्थयः }\edtext{\skt{कोकिलेनोक्ता}}{\lemma{\skt{कोकिलेनोक्ता}}\Afootnote{K H; \skt{कोकिलेनोग्रा}\ A}}\skt{ }\edtext{\skt{ज्वरो}}{\lemma{\skt{ज्वरो}}\Afootnote{A H; \skt{…}\ K}}\skt{ }\edtext{\skt{दाहश्च}}{\lemma{\skt{दाहश्च}}\Afootnote{A H; \skt{…हश्च}\ K}}\skt{ दारुणः~।}\pend

\pstart\versepar \ledsidenote{\passref{5.7.25}}\edtext{\skt{नीला\-व\-र्षा\-भु\-निः\-क्वाथैः}}{\lemma{\skt{नीला\appbrk{}व\appbrk{}र्षा\appbrk{}भु\appbrk{}निः\appbrk{}क्वाथैः}}\Afootnote{K H; \skt{वर्षा\appbrk{}भू\appbrk{}नी\appbrk{}लि\appbrk{}नी\appbrk{}क्वा\appbrk{}थ\appbrk{}कल्क}\ A}}\skt{ सिद्धं }\edtext{\skt{तत्र}}{\lemma{\skt{तत्र}}\Afootnote{K H; \skt{घृतं}\ A}}\skt{ }\edtext{\skt{पिबेद्घृतम्}}{\lemma{\skt{पिबेद्घृतम्}}\Afootnote{\skt{पिबे}\ A; \skt{पिवेद्घृतं}\ K H}}\skt{~॥ }\pend
\pstart\versepar \skt{}\edtext{\skt{अरुणेनानिलः}}{\lemma{\skt{अरुणेनानिलः}}\Afootnote{K H; \skt{त्~। अरुणेनानिलः}\ A}}\skt{ क्रुद्धो }\edtext{\skt{वातजां}}{\lemma{\skt{वातजां}}\Afootnote{K H; \skt{वातजा}\ A}}\skt{ }\edtext{\skt{कुरुते}}{\lemma{\skt{कुरुते}}\Afootnote{K H; \skt{न्कुरुते}\ A}}\skt{ }\edtext{\skt{गदान्}}{\lemma{\skt{गदान्}}\Afootnote{A; \skt{गदां~।}\ K H}}\skt{~।}\pend

\pstart\versepar \ledsidenote{\passref{5.7.26}}\skt{महाकृष्णेन पित्तञ्च }\edtext{\skt{श्वेतेन}}{\lemma{\skt{श्वेतेन}}\Afootnote{A K\textsuperscript{pc} H; \skt{श्वेते}\ K\textsuperscript{ac}}}\skt{ कफ एव च~। }\pend
\pstart\versepar \skt{महता कपि\-ले\-ना\-सृ\-क्क\-पो\-तेन चतुष्टयम्~॥}\pend

\pstart\versepar \ledsidenote{\passref{5.7.27}}\skt{भवन्ति चैषान्दंशेषु }\edtext{\skt{ग्रन्थि\-म\-ण्ड\-ल\-क\-र्ण्णिकाः}}{\lemma{\skt{ग्रन्थि\appbrk{}म\appbrk{}ण्ड\appbrk{}ल\appbrk{}क\appbrk{}र्ण्णिकाः}}\Afootnote{A H; \skt{ग्रन्थि\appbrk{}म\appbrk{}ण्ड\appbrk{}ल\appbrk{}क\appbrk{}र्ण्णिका~॥}\ K}}\skt{~॥ }\pend
\pstart\versepar \skt{}\edtext{\edtext{\skt{पिड\-को\-प\-च\-या\-श्चाङ्गे}}{\lemma{\skt{पिड\appbrk{}को\appbrk{}प\appbrk{}च\appbrk{}याश्}}\Afootnote{K H; \skt{पिड\appbrk{}को\appbrk{}प\appbrk{}च\appbrk{}यश्}\ A}}}{\lemma{\skt{चाङ्गे}}\Afootnote{K H; \skt{चोग्रः}\ A}}\skt{ }\edtext{\skt{शोफाश्च}}{\lemma{\skt{शोफाश्च}}\Afootnote{K H; \skt{शोफश्च}\ A}}\skt{ }\edtext{\skt{भृशदारुणाः}}{\lemma{\skt{भृशदारुणाः}}\Afootnote{H; \skt{भृशदारुणः~॥}\ A; \skt{भृषदारुणः~।}\ K}}\skt{~।}\pend

\pstart\versepar \ledsidenote{\passref{5.7.28}}\edtext{\skt{दधि\-क्षी\-र\-घृ\-त\-प्र\-स्था\-स्त्रयः}}{\lemma{\skt{दधि\appbrk{}क्षी\appbrk{}र\appbrk{}घृ\appbrk{}त\appbrk{}प्र\appbrk{}स्था\appbrk{}स्त्रयः}}\Afootnote{\skt{दधि\appbrk{}क्षी\appbrk{}र\appbrk{}घृ\appbrk{}त\appbrk{}प्र\appbrk{}स्था\appbrk{}स्त्रयः}\ A K; \skt{दधि\appbrk{}क्षी\appbrk{}र\appbrk{}घृ\appbrk{}त\appbrk{}म्प्र\appbrk{}स्था\appbrk{}स्त्रयः}\ H}}\skt{ }\edtext{\skt{प्रत्येकशो}}{\lemma{\skt{प्रत्येकशो}}\Afootnote{A K H\textsuperscript{pc}; \skt{प्रत्येक}\ H\textsuperscript{ac}}}\skt{ }\edtext{\skt{मिताः}}{\lemma{\skt{मिताः}}\Afootnote{K; \skt{मताः~।}\ A; \skt{सम्मिताः~॥}\ H\textsuperscript{ac}; \skt{मितासम्मिताः~॥}\ H\textsuperscript{pc}}}\skt{~। }\pend
\pstart\versepar \skt{}\edtext{\skt{करञ्जारग्वधं}}{\lemma{\skt{करञ्जारग्वधं}}\Afootnote{K H; \skt{करञ्जारग्वध}\ A}}\skt{ }\edtext{\skt{व्योषं}}{\lemma{\skt{व्योषं}}\Afootnote{H; \skt{व्योष}\ A K}}\skt{ }\edtext{\skt{बृह\-त्यं\-शु\-म\-ती\-स्थिराः}}{\lemma{\skt{बृह\appbrk{}त्यं\appbrk{}शु\appbrk{}म\appbrk{}ती\appbrk{}स्थिराः}}\Afootnote{A; \skt{वृह\appbrk{}त्यं\appbrk{}शु\appbrk{}म\appbrk{}ती\appbrk{}स्थिरा~।}\ K; \skt{वृह्य\appbrk{}त्य\appbrk{}ङ्शु\appbrk{}म\appbrk{}ती\appbrk{}स्थिराः~।}\ H}}\skt{~।}\pend

\pstart\versepar \ledsidenote{\passref{5.7.29}}\edtext{\skt{निष्क्वाथ्य}}{\lemma{\skt{निष्क्वाथ्य}}\Afootnote{A K; \skt{निःक्वाथ्य}\ H}}\skt{ }\edtext{\skt{तस्य}}{\lemma{\skt{तस्य}}\Afootnote{K H; \skt{चैषां}\ A}}\skt{ क्वाथस्य }\edtext{\skt{चतुर्थांशः}}{\lemma{\skt{चतुर्थांशः}}\Afootnote{K H; \skt{चतुर्थोऽंशः}\ A}}\skt{ पुनर्भवेत्~॥ }\pend
\pstart\versepar \skt{तृवृ\-त्ति\-ल्वा\-मृ\-ता\-व\-क्र\-स\-र्व\-ग\-न्धा\-ग\-मृ\-त्तिकाः~।}\edtext{}{\lemma{\skt{पुनर्भवेत्~॥ तृवृ\appbrk{}त्ति\appbrk{}ल्वा\appbrk{}मृ\appbrk{}ता\appbrk{}व\appbrk{}क्र\appbrk{}स\appbrk{}र्व\appbrk{}ग\appbrk{}न्धा\appbrk{}ग\appbrk{}मृ\appbrk{}त्तिकाः}}\Afootnote{\skt{पुनर्भवेत्~। त्रिवृ\appbrk{}द्गो\appbrk{}ज्य\appbrk{}मृ\appbrk{}ता\appbrk{}व\appbrk{}क्र\appbrk{}स\appbrk{}र्प\appbrk{}गन्धाः समृत्तिकाः~॥}\ A; \skt{पुनर्भवेत्~॥ तृवृ\appbrk{}त्ति\appbrk{}ल्वा\appbrk{}मृ\appbrk{}ता\appbrk{}व\appbrk{}क्र\appbrk{}स\appbrk{}र्वा\appbrk{}ग\appbrk{}न्धा\appbrk{}ग\appbrk{}मृ\appbrk{}त्तिका~।}\ K; \skt{पुनर्व्भवेत्~॥ तृवृ\appbrk{}त्ति\appbrk{}ल्वा\appbrk{}मृ\appbrk{}ता\appbrk{}व\appbrk{}क्रा\appbrk{}स\appbrk{}र्व्व\appbrk{}ग\appbrk{}न्धा\appbrk{}ग\appbrk{}मृ\appbrk{}त्तिका~॥}\ H}}\pend

\pstart\versepar \ledsidenote{\passref{5.7.30}}\skt{कपि\-त्थ\-दा\-डि\-म\-त्वक्च }\edtext{\skt{श्लक्ष्णपिष्टानि}}{\lemma{\skt{श्लक्ष्णपिष्टानि}}\Afootnote{K; \skt{श्लक्ष्णपिष्टाः प्र}\ A; \skt{श्लक्ष्णपिष्ठानि}\ H}}\skt{ दापयेत्~। }\pend
\pstart\versepar \skt{}\edtext{\skt{तत्सर्वमेकतः}}{\lemma{\skt{एकतः}}\Afootnote{A K; \skt{मेकतः}\ H}}\skt{ कृत्वा }\edtext{\edtext{\skt{शनैर्मृद्वग्निना}}{\lemma{\skt{शनैर्}}\Afootnote{A K; \skt{शनै}\ H}}}{\lemma{\skt{मृद्वग्निना}}\Afootnote{A K; \skt{र्म्मृद्वग्निना}\ H}}\skt{ पचेत्~॥}\edtext{}{\lemma{\skt{दापयेत्~। तत्सर्वम्}}\Afootnote{A K; \skt{दापयेत्~॥ तत्सर्व्व}\ H}}\pend

\pstart\versepar \ledsidenote{\passref{5.7.31}}\skt{पञ्चा\-ना\-म\-रु\-णा\-दीनां }\edtext{\edtext{\skt{विष\-मे\-त\-द्व्य\-पो\-हति}}{\lemma{\skt{विषम्}}\Afootnote{A K\textsuperscript{pc} H; \skt{विष}\ K\textsuperscript{ac}}}}{\lemma{\skt{एतद्व्यपोहति}}\Afootnote{A K\textsuperscript{pc} H; \skt{तद्व्यपोहति~।}\ K\textsuperscript{ac}}}\skt{~। }\pend
\pstart\versepar \skt{}\edtext{\skt{काका\-द\-नी\-का\-क\-माची}}{\lemma{\skt{काका\appbrk{}द\appbrk{}नी\appbrk{}का\appbrk{}क\appbrk{}माची}}\Afootnote{K H; \skt{काका\appbrk{}द\appbrk{}नी\appbrk{}का\appbrk{}क\appbrk{}माच्योः}\ A}}\skt{ }\edtext{\skt{स्वरस्वेष्वथवा}}{\lemma{\skt{स्वरस्वेष्व्}}\Afootnote{K; \skt{स्वरसेष्व्}\ A H}}\skt{ कृतम्~।}\pend

\pstart\versepar \ledsidenote{\passref{5.7.32}}\skt{सिराश्च }\edtext{\skt{व्यधयेत्प्राप्ताः}}{\lemma{\skt{व्यधयेत्प्राप्ताः}}\Afootnote{K H; \skt{स्रावयेत्प्राप्ताः}\ A}}\skt{ कुर्या\-त्सं\-शो\-ध\-नानि च~। }\pend
\pstart\versepar \skt{}\edtext{\skt{सर्व्वेषु वा}}{\lemma{\skt{सर्व्वेषु वा}}\Afootnote{H; \skt{सर्वेषां च}\ A; \emph{om.}\ K}}\skt{ }\edtext{\skt{विधिः}}{\lemma{\skt{विधिः}}\Afootnote{A H; \skt{…}\ K}}\skt{ कार्यो मूषि\-का\-णा\-म्वि\-षे\-ष्वयं~॥}\pend

\pstart\versepar \ledsidenote{\passref{5.7.33}}\skt{दग्ध्वा विस्रावयेद्दंशं }\edtext{\skt{प्रच्छितञ्च}}{\lemma{\skt{प्रच्छितञ्च}}\Afootnote{A H; \skt{प्रच्छितंञ्च}\ K}}\skt{ प्रलेपयेत्~। }\pend
\pstart\versepar \skt{शिरी\-ष\-र\-ज\-नी\-वक्रं कुंकु\-मै\-र\-मृ\-ता\-युतैः~।}\edtext{}{\lemma{\skt{प्रलेपयेत्~। शिरी\appbrk{}ष\appbrk{}र\appbrk{}ज\appbrk{}नी\appbrk{}वक्रं}}\Afootnote{K H; \skt{प्रलेपयेत्~। शिरी\appbrk{}ष\appbrk{}र\appbrk{}ज\appbrk{}नी\appbrk{}कुष्ठ}\ A}}\pend

\pstart \ledsidenote{\passref{5.7.34ab}}\skt{च्छर्दनं }\edtext{\skt{नीलिनीक्वाथैः}}{\lemma{\skt{नीलिनीक्वाथैः}}\Afootnote{K H\textsuperscript{ac}; \skt{जालिनीक्वाथैः}\ A; \skt{नीलिनीक्वाथंथैः}\ H\textsuperscript{pc}}}\skt{ }\edtext{\skt{शुका\-ख्या\-ङ्को\-ल्ल\-यो\-रपि}}{\lemma{\skt{शुका\appbrk{}ख्या\appbrk{}ङ्को\appbrk{}ल्ल\appbrk{}योर्}}\Afootnote{K H; \skt{शुका\appbrk{}ख्या\appbrk{}ङ्को\appbrk{}ठ\appbrk{}योर्}\ A}}\skt{~।}\pend

\pstart\versepar \ledsidenote{\passref{5.7.37}}\skt{विरेचने }\edtext{\skt{तृवृ\-द्द\-न्ती\-त्रि\-फ\-ला\-कल्क}}{\lemma{\skt{तृवृ\appbrk{}द्द\appbrk{}न्ती\appbrk{}त्रि\appbrk{}फ\appbrk{}ला\appbrk{}कल्क}}\Afootnote{K H; \skt{त्रिवृ\appbrk{}द्द\appbrk{}न्ती\appbrk{}त्रि\appbrk{}फ\appbrk{}ला\appbrk{}कल्क}\ A}}\skt{ इष्यते~। }\pend
\pstart\versepar \skt{शिरोविरेचने सारः }\edtext{\skt{शिरीषस्य}}{\lemma{\skt{शिरीषस्य}}\Afootnote{A H; \skt{शिरीशस्य}\ K}}\skt{ फलानि }\edtext{\skt{वा}}{\lemma{\skt{वा}}\Afootnote{K H; \skt{च~॥}\ A}}\skt{~।}\pend

\pstart\versepar \ledsidenote{\passref{5.7.38}}\edtext{\skt{कटुत्रिकाढ्यश्च}}{\lemma{\skt{कटुत्रिकाढ्यश्च}}\Afootnote{K H; \skt{हित\appbrk{}स्त्रि\appbrk{}क\appbrk{}टु\appbrk{}का\appbrk{}ढ्यश्च}\ A}}\skt{ }\edtext{\skt{हितो}}{\lemma{\skt{हितो}}\Afootnote{K H; \emph{om.}\ A}}\skt{ }\edtext{\skt{गोम\-य\-स्व\-र\-सोऽ\-ञ्जने}}{\lemma{\skt{गोमयस्वरसो}}\Afootnote{A; \skt{गोमयः स्वरसो}\ K H}}\skt{~। }\pend
\pstart\versepar \skt{कपि\-त्थ\-गो\-म\-य\-रसौ }\edtext{\skt{सक्षौद्रौ}}{\lemma{\skt{सक्षौद्रौ}}\Afootnote{K; \emph{om.}\ A; \skt{सक्षौद्रैर्}\ H}}\skt{ }\edtext{\skt{लेह}}{\lemma{\skt{लेह}}\Afootnote{K H; \emph{om.}\ A}}\skt{ }\edtext{\skt{इष्यते}}{\lemma{\skt{इष्यते}}\Afootnote{K H; \skt{लिह्या\appbrk{}न्मा\appbrk{}क्षि\appbrk{}क\appbrk{}सं\appbrk{}युतौ~॥}\ A}}\skt{~।}\pend

\pstart\versepar \ledsidenote{\passref{5.7.40}}\skt{तण्डु\-ली\-य\-क\-मू\-लेषु सर्पिः सिद्धं }\edtext{\skt{पिबेन्नरः}}{\lemma{\skt{पिबेन्}}\Afootnote{A; \skt{पिवेन्}\ K H}}\skt{~। }\pend
\pstart\versepar \skt{आस्फो\-त\-मू\-ल\-सि\-द्धम्वा }\edtext{\skt{पञ्च\-का\-पि\-त्त\-मेव}}{\lemma{\skt{पञ्चकापित्तम्}}\Afootnote{K H; \skt{पञ्चकापित्थम्}\ A}}\skt{ वा~।}\pend

\pstart\versepar \ledsidenote{\passref{5.7.41}}\edtext{\skt{मूषिकाणां}}{\lemma{\skt{मूषिकाणां}}\Afootnote{A K; \skt{मूषिकानाम्}\ H}}\skt{ विषं प्रायः कुप्यत्यभ्रेषु }\edtext{\skt{निर्हृतम्}}{\lemma{\skt{निर्हृतम्}}\Afootnote{K H; \skt{अनिर्हृतम्~।}\ A}}\skt{~। }\pend
\pstart\versepar \skt{तत्राप्येष विधिः कार्यः यश्च दूषीविषापहः~॥}\pend

\pstart\versepar \ledsidenote{\passref{5.7.42}}\edtext{\skt{स्थिरां }\edtext{\skt{मन्दरुजाश्चापि}}{\lemma{\skt{चापि}}\Afootnote{\skt{वाऽपि}\ A; \emph{om.}\ K; \skt{चा}\ H}}}{\lemma{\skt{स्थिरां मन्दरुजाश्}}\Afootnote{\skt{स्थिराणां रुजतां}\ A; \skt{स्थिरांमन्दरु}\ K; \skt{स्थिरा\appbrk{}न्म\appbrk{}न्द\appbrk{}रु\appbrk{}जाञ्}\ H}}\skt{ }\edtext{\skt{कर्ण्णिकां}}{\lemma{\skt{कर्ण्णिकां}}\Afootnote{\skt{व्रणानां कर्णिकां}\ A; \skt{जश्चापि कर्ण्णिकां}\ K; \skt{पि कर्ण्णिकाम्}\ H}}\skt{ प्रच्छयेद्भिषक्~। }\pend
\pstart\versepar \skt{सर्वस्मिन्नेव }\edtext{\skt{तु}}{\lemma{\skt{तु}}\Afootnote{K H; \skt{था}\ A}}\skt{ }\edtext{\skt{विषे}}{\lemma{\skt{विषे}}\Afootnote{K H; \skt{दोषं}\ A}}\skt{ }\edtext{\skt{व्रण\-व\-च्चा\-च\-रे\-त्क्रि\-याम्}}{\lemma{\skt{व्रण\appbrk{}व\appbrk{}च्चा\appbrk{}च\appbrk{}रे\appbrk{}त्क्रि\appbrk{}याम्}}\Afootnote{K H; \skt{व्रणवच्चापि शोधयेत्~॥}\ A}}\skt{~॥}\edtext{}{\lemma{\skt{प्रच्छयेद्भिषक्~। सर्वस्मिन्नेव}}\Afootnote{K H; \skt{भिषक्~। पाटयित्वा य}\ A}}\pend

\pstart\versepar \ledsidenote{\passref{5.7.43}}\edtext{\edtext{\skt{श्वशृ\-गा\-ल\-वृ\-क\-व्या\-घ्र\-त\-र\-क्ष्वा\-दे\-र्विषं}}{\lemma{\skt{श्वशृ\appbrk{}गा\appbrk{}ल\appbrk{}वृ\appbrk{}क\appbrk{}व्या\appbrk{}घ्र\appbrk{}त\appbrk{}र\appbrk{}क्ष्वा\appbrk{}देर्}}\Afootnote{K H; \skt{श्वशृ\appbrk{}गा\appbrk{}ल\appbrk{}त\appbrk{}र\appbrk{}क्ष्वृक्ष}\ A}}}{\lemma{\skt{विषं}}\Afootnote{K H; \skt{व्याघ्रादीनां}\ A}}\skt{ यदा~। }\pend
\pstart\versepar \skt{}\edtext{\skt{श्लेष्मा}}{\lemma{\skt{श्लेष्मा}}\Afootnote{K H; \skt{ऽनिलः~। श्लेष्म}\ A}}\skt{ }\edtext{\skt{प्रदुष्टो}}{\lemma{\skt{प्रदुष्टो}}\Afootnote{A K; \skt{प्रदुष्टा}\ H}}\skt{ }\edtext{\skt{मुष्णाति}}{\lemma{\skt{मुष्णाति}}\Afootnote{A K; \skt{पुष्णाति}\ H}}\skt{ }\edtext{\skt{संज्ञां}}{\lemma{\skt{संज्ञां}}\Afootnote{A H; \skt{संज्ञं}\ K}}\skt{ }\edtext{\skt{संज्ञावहाश्रितः}}{\lemma{\skt{संज्ञावहाश्रितः}}\Afootnote{A H; \skt{स्रोत्रोवलाश्रितः~॥}\ K}}\skt{~॥}\pend

\pstart\versepar \ledsidenote{\passref{5.7.44}}\skt{तदा प्रस्र\-स्त\-ला\-ङ्गू\-ल\-ह\-नु\-स्क\-न्धोऽ\-भि\-ला\-लि\-मान्~। }\pend
\pstart\versepar \skt{}\edtext{\skt{अव्य\-क्त\-व\-धि\-रोऽ\-न्धश्च }\edtext{\skt{सोऽन्यो\-न्य\-म\-भि\-धा\-वति}}{\lemma{\skt{ऽन्योन्यम्}}\Afootnote{A H; \skt{न्योन्यम्}\ K}}}{\lemma{\skt{ऽन्धश्च सो}}\Afootnote{A K; \skt{न्धश्च सो}\ H}}\skt{~।}\edtext{}{\lemma{\skt{ऽभिलालिमान्~। अव्यक्तवधिरो}}\Afootnote{\skt{ऽतिलालवान्~। अत्यर्थबधिरो}\ A; \skt{भिलालिमान्~। अव्यक्तवधिरो}\ K; \skt{ऽकिलालिमान्~। अव्यक्तवधिरो}\ H}}\pend

\pstart \ledsidenote{\passref{5.7.45}}\edtext{\skt{तेन}}{\lemma{\skt{तेन}}\Afootnote{K H; \skt{तेनोन्मत्तेन}\ A}}\skt{ दष्टस्य }\edtext{\skt{चाङ्गे स्युः}}{\lemma{\skt{चाङ्गे स्युः}}\Afootnote{K H; \skt{दंष्ट्रिणा सविषेण तु~।}\ A}}\skt{ }\edtext{\skt{सुप्तः}}{\lemma{\skt{सुप्तः}}\Afootnote{K H; \skt{सुप्तता जायते दंशे}\ A}}\skt{ }\edtext{\skt{कृष्णं}}{\lemma{\skt{कृष्णं}}\Afootnote{K H; \skt{कृष्णं चाति}\ A}}\skt{ }\edtext{\skt{क्षरत्यसृक्}}{\lemma{\skt{क्षरत्य्}}\Afootnote{K H; \skt{स्रवत्य्}\ A}}\skt{~॥}\pend

\pstart\versepar \ledsidenote{\passref{5.7.46}}\skt{दिग्धविद्धस्य लिङ्गेन }\edtext{\skt{प्राय\-श\-श्चा\-भि\-लि\-ङ्गितः}}{\lemma{\skt{चाभिलिङ्गितः}}\Afootnote{K H\textsuperscript{ac}; \skt{चोपलक्षितः~।}\ A; \skt{चाप्य\appbrk{}श्चा\appbrk{}भि\appbrk{}लि\appbrk{}ङ्गितः~।}\ H\textsuperscript{pc}}}\skt{~। }\pend
\pstart\versepar \skt{येन चापि भवेद्दष्टस्तस्य }\edtext{\skt{चेष्टारुतन्नरः}}{\lemma{\skt{चेष्टारुतन्}}\Afootnote{K H; \skt{चेष्टां रुतं}\ A}}\skt{~।}\pend

\pstart\versepar \ledsidenote{\passref{5.7.47}}\edtext{\skt{बहुशः}}{\lemma{\skt{बहुशः}}\Afootnote{A; \skt{वहुशः}\ K H}}\skt{ प्रतिकुर्वाणः क्रियाहीनो }\edtext{\skt{विनश्यति}}{\lemma{\skt{विनश्यति}}\Afootnote{A H; \skt{विश्यति~॥}\ K}}\skt{~॥ }\pend
\pstart\versepar \skt{दंष्ट्रिणा येन }\edtext{\skt{दष्टस्तु}}{\lemma{\skt{दष्टस्तु}}\Afootnote{K H; \skt{दष्टश्च}\ A}}\skt{ }\edtext{\skt{तं}}{\lemma{\skt{तं}}\Afootnote{K H; \skt{त}\ A}}\skt{ }\edtext{\skt{दष्टो}}{\lemma{\skt{दष्टो}}\Afootnote{K H; \skt{द्रूपं}\ A}}\skt{ }\edtext{\skt{यदि}}{\lemma{\skt{यदि}}\Afootnote{K H; \skt{यस्तु}\ A}}\skt{ पश्यति~।}\pend

\pstart\versepar \ledsidenote{\passref{5.7.48}}\skt{अप्सु वा यदि }\edtext{\edtext{\skt{वादर्शेऽरिष्टं}}{\lemma{\skt{वादर्शे}}\Afootnote{K H; \skt{वादर्शो}\ A}}}{\lemma{\skt{ऽरिष्टं}}\Afootnote{A; \skt{रिष्टं}\ K H}}\skt{ तस्य विनिर्दिशेत्~। }\pend
\pstart\versepar \skt{यदि }\edtext{\edtext{\skt{त्रस्यत्यदष्टोऽपि}}{\lemma{\skt{त्रस्यत्य्}}\Afootnote{H; \skt{त्~। त्रस्यत्य्}\ A; \skt{त्रस्य}\ K}}}{\lemma{\skt{अदष्टो}}\Afootnote{H; \skt{अकस्माद्योऽभीक्ष्णं दृष्ट्वा स्पृष्ट्वा}\ A; \skt{दष्टो}\ K\textsuperscript{ac}; \skt{त्यदष्टो}\ K\textsuperscript{pc}}}\skt{ }\edtext{\skt{शब्द\-स्प\-र्श\-न\-द\-र्शनैः}}{\lemma{\skt{शब्द\appbrk{}स्प\appbrk{}र्श\appbrk{}न\appbrk{}द\appbrk{}र्शनैः}}\Afootnote{\skt{वा जलम्~॥}\ A; \skt{शब्द\appbrk{}स्य\appbrk{}र्श\appbrk{}न\appbrk{}द\appbrk{}र्शनैः~।}\ K; \skt{शब्द\appbrk{}स्य\appbrk{}र्श\appbrk{}न\appbrk{}द\appbrk{}र्शमैः~॥}\ H}}\skt{~।}\edtext{}{\lemma{\skt{विनिर्दिशेत्~। यदि}}\Afootnote{K H; \skt{विनिर्दिशे}\ A}}\pend

\pstart \ledsidenote{\passref{5.7.49ab}}\edtext{\skt{जलत्रासं तु तं}}{\lemma{\skt{जलत्रासं तु तं}}\Afootnote{K H; \skt{जलत्रासं तु}\ A}}\skt{ }\edtext{\skt{विद्यादृष्टं}}{\lemma{\skt{विद्यादृष्टं}}\Afootnote{\skt{विद्यात्तं रिष्टं}\ A; \skt{विद्याद्दष्टं}\ K H}}\skt{ तदपि कीर्त्तितम्~।}\pend

\pstart \ledsidenote{\passref{5.7.50cd}}\edtext{\skt{विस्राव्य}}{\lemma{\skt{विस्राव्य}}\Afootnote{K H\textsuperscript{pc}; \skt{दंशं विस्राव्य}\ A; \skt{विस्रा}\ H\textsuperscript{ac}}}\skt{ }\edtext{\skt{दंशं तं}}{\lemma{\skt{दंशं तं}}\Afootnote{K H; \skt{तैर्}\ A}}\skt{ दष्टे सर्पिषा परिदाहितम्~।}\pend

\pstart\versepar \ledsidenote{\passref{5.7.51}}\skt{प्रदिह्यादगदैः सर्पिः पुराणं }\edtext{\skt{चापि}}{\lemma{\skt{चापि}}\Afootnote{K H; \emph{om.}\ A}}\skt{ पाययेत्~। }\pend
\pstart\versepar \skt{}\edtext{\skt{अर्क\-क्षी\-र\-यु\-त\-ञ्चापि}}{\lemma{\skt{चापि}}\Afootnote{K H; \skt{ह्य्}\ A}}\skt{ }\edtext{\edtext{\skt{शीघ्र\-न्द\-द्या\-द्वि\-रे\-च\-नम्}}{\lemma{\skt{शीघ्रन्}}\Afootnote{K H; \skt{अस्य}\ A}}}{\lemma{\skt{दद्या\appbrk{}द्वि\appbrk{}रे\appbrk{}च\appbrk{}नम्}}\Afootnote{K H; \skt{दद्याच्चापि विशोधनम्~॥}\ A}}\skt{~।}\edtext{}{\lemma{\skt{पाययेत्~। अर्क\appbrk{}क्षी\appbrk{}र\appbrk{}यु\appbrk{}तञ्}}\Afootnote{K H; \skt{पाययेत च~। अर्कक्षीरयुतं}\ A}}\pend

\pstart \ledsidenote{\passref{5.7.52ab}}\skt{श्वेतां }\edtext{\edtext{\skt{पुनर्ण्णवाञ्चास्यै}}{\lemma{\skt{पुनर्ण्णवाञ्}}\Afootnote{K H; \skt{पुनर्नवां}\ A}}}{\lemma{\skt{चास्यै}}\Afootnote{K; \skt{चास्य}\ A; \skt{चास्ये}\ H}}\skt{ }\edtext{\skt{दद्या\-द्धु\-त्तू\-र\-का\-यु\-ताम्}}{\lemma{\skt{दद्या\appbrk{}द्धु\appbrk{}त्तू\appbrk{}र\appbrk{}का\appbrk{}यु\appbrk{}ताम्}}\Afootnote{K H; \skt{दद्या\appbrk{}द्ध\appbrk{}त्तू\appbrk{}र\appbrk{}का\appbrk{}यु\appbrk{}ताम्~।}\ A}}\skt{~।}\pend

\pstart \ledsidenote{\passref{5.7.60ab}}\skt{स्नापयेत्तं नदीतीरे समन्त्रैर्वा चतुष्पथे~।}\pend

\pstart \ledsidenote{\passref{5.7.60.add-1}}\edtext{\skt{बीज\-र\-त्नौ\-ष\-धी\-गर्भैः}}{\lemma{\skt{बीज\appbrk{}र\appbrk{}त्नौ\appbrk{}ष\appbrk{}धी\appbrk{}गर्भैः}}\Afootnote{\skt{वीज\appbrk{}र\appbrk{}त्नौ\appbrk{}ष\appbrk{}धी\appbrk{}गर्भैः}\ K; \skt{वीज\appbrk{}र\appbrk{}त्नौ\appbrk{}ष\appbrk{}धी\appbrk{}गव्भैः}\ H}}\skt{ कुम्भैः }\edtext{\skt{शीताम्बुपूरितैः}}{\lemma{\skt{शीताम्बुपूरितैः}}\Afootnote{\skt{शीताम्वुपूरितैः}\ K H}}\skt{~॥}\pend

\pstart \ledsidenote{\passref{5.7.61cd}}\edtext{\skt{अलर्काधिपते}}{\lemma{\skt{अलर्काधिपते}}\Afootnote{K; \skt{अलकाधिपते}\ A H}}\skt{ यक्ष सार\-मे\-य\-ग\-णा\-धिप~।}\pend

\pstart\versepar \ledsidenote{\passref{5.7.62}}\skt{अल\-र्क\-जु\-ष्ट\-मे\-तन्मे निर्विषं कुरु मा }\edtext{\skt{चिरात्~॥ स्वाहा}}{\lemma{\skt{चिरात्~॥ स्वाहा}}\Afootnote{K H; \skt{चिरा}\ A}}\skt{~॥ }\pend
\pstart\versepar \skt{}\edtext{\edtext{\skt{दद्या\-त्सं\-शो\-ध\-न\-न्ती\-क्ष्ण\-मेवं}}{\lemma{\skt{दद्या\appbrk{}त्सं\appbrk{}शो\appbrk{}ध\appbrk{}नन्}}\Afootnote{H; \skt{त्~। दद्यात्संशोधनं}\ A; \skt{दद्या\appbrk{}त्स\appbrk{}शो\appbrk{}ध\appbrk{}नन्}\ K}}}{\lemma{\skt{एवं}}\Afootnote{A; \skt{अव}\ K H}}\skt{ }\edtext{\skt{स्नातस्य}}{\lemma{\skt{स्नातस्य}}\Afootnote{A; \skt{स्यांन्तस्यं}\ K; \skt{स्यन्तस्य}\ H}}\skt{ देहिनः~।}\pend

\pstart \ledsidenote{\passref{5.7.63ab}}\skt{अशुद्धस्य }\edtext{\skt{सुरूढेऽपि}}{\lemma{\skt{ऽपि}}\Afootnote{A; \skt{पि}\ K H}}\skt{ व्रणे कुप्यति तद्विषम्~।}\pend

\pstart\versepar \ledsidenote{\passref{5.7.63.add-1}}\skt{प्रसुप्तो वोत्थितो वापि स्वस्थः }\edtext{\skt{त्रस्तो न}}{\lemma{\skt{त्रस्तो न}}\Afootnote{K; \skt{स्वस्थो न}\ H}}\skt{ सिध्यति~। }\pend
\pstart\versepar \skt{जलत्रासी च यो }\edtext{\skt{मर्त्यो}}{\lemma{\skt{मर्त्यो}}\Afootnote{K H\textsuperscript{pc}; \skt{मर्त्यो}\ H\textsuperscript{ac}}}\skt{ }\edtext{\skt{दष्टे}}{\lemma{\skt{दष्टे}}\Afootnote{K; \emph{om.}\ H\textsuperscript{ac}; \skt{दष्टो यश्च}\ H\textsuperscript{pc}}}\skt{ }\edtext{\skt{यश्च}}{\lemma{\skt{यश्च}}\Afootnote{K; \emph{om.}\ H\textsuperscript{ac}; \skt{दर्शयश्च}\ H\textsuperscript{pc}}}\skt{ }\edtext{\skt{प्रकुप्यतीति}}{\lemma{\skt{प्रकुप्यतीति}}\Afootnote{K H\textsuperscript{pc}; \emph{om.}\ H\textsuperscript{ac}}}\skt{~॥}\pend

\pstart \ledsidenote{\passref{5.7.trailer}}\edtext{\skt{कल्पे}}{\lemma{\skt{कल्पे}}\Afootnote{\skt{इति सुश्रु\appbrk{}त\appbrk{}सं\appbrk{}हि\appbrk{}तायां कल्पस्थाने मूषिककल्पो}\ A; \skt{कल्पे ६~॥ {\symfont ❈}~॥}\ K; \skt{॥ कल्पस्थाने}\ H}}\skt{ ६~॥ }\edtext{\skt{ओ}}{\lemma{\skt{ओ}}\Afootnote{\skt{नाम सप्तमोऽध्यायः~॥७~॥}\ A; \emph{om.}\ K; \skt{षष्ठोऽध्यायः~॥}\ H}}\skt{~॥}\pend

\endnumbering
\beginnumbering
\eledsection*{Adhyāya 7}
\pstart \ledsidenote{\passref{5.6.1}}\skt{अथातो दुन्दुभिस्वनीयं }\edtext{\skt{कल्पं}}{\lemma{\skt{कल्पं}}\Afootnote{A H; \skt{कंल्प व्याख्यास्यामः~॥}\ K}}\skt{ }\edtext{\skt{व्याख्यास्यामः}}{\lemma{\skt{व्याख्यास्यामः}}\Afootnote{A H; \emph{om.}\ K}}\skt{~॥ }\pend

\pstart \ledsidenote{\passref{5.6.3}}\edtext{\skt{धवा\-श्व\-क\-र्ण्ण\-ति\-नि\-श\-पि\-चु\-म\-र्द\-पा\-ट\-ली\-पा\-रि\-भ\-द्र\-को\-दु\-म्ब\-र\-क\-र\-घा\-ट\-का\-र्जु\-न\-स\-र्ज्ज\-क\-पी\-त\-न\-श्ले\-ष्मा\-त\-का\-ङ्को\-ठ\-कु\-ट\-ज\-श\-मी\-क\-पि\-त्था\-श्म\-न्त\-का\-र्क\-चि\-रि\-बि\-ल्व\-म\-हा\-वृ\-क्षा\-र\-ल\-म\-धु\-क\-म\-धु\-क\-शि\-ग्रु\-शा\-क\-गो\-जी\-भू\-ज\-ति\-ल्व\-के\-क्षु\-र\-क\-गो\-प\-घो\-ण्टा\-रि\-मे\-दानां}}{\lemma{\skt{धवा\appbrk{}श्व\appbrk{}क\appbrk{}र्ण्ण\appbrk{}ति\appbrk{}नि\appbrk{}श\appbrk{}पि\appbrk{}चु\appbrk{}म\appbrk{}र्द\appbrk{}पा\appbrk{}ट\appbrk{}ली\appbrk{}पा\appbrk{}रि\appbrk{}भ\appbrk{}द्र\appbrk{}को\appbrk{}दु\appbrk{}म्ब\appbrk{}र\appbrk{}क\appbrk{}र\appbrk{}घा\appbrk{}ट\appbrk{}का\appbrk{}र्जु\appbrk{}न\appbrk{}स\appbrk{}र्ज्ज\appbrk{}क\appbrk{}पी\appbrk{}त\appbrk{}न\appbrk{}श्ले\appbrk{}ष्मा\appbrk{}त\appbrk{}का\appbrk{}ङ्को\appbrk{}ठ\appbrk{}कु\appbrk{}ट\appbrk{}ज\appbrk{}श\appbrk{}मी\appbrk{}क\appbrk{}पि\appbrk{}त्था\appbrk{}श्म\appbrk{}न्त\appbrk{}का\appbrk{}र्क\appbrk{}चि\appbrk{}रि\appbrk{}बि\appbrk{}ल्व\appbrk{}म\appbrk{}हा\appbrk{}वृ\appbrk{}क्षा\appbrk{}र\appbrk{}ल\appbrk{}म\appbrk{}धु\appbrk{}क\appbrk{}म\appbrk{}धु\appbrk{}क\appbrk{}शि\appbrk{}ग्रु\appbrk{}शा\appbrk{}क\appbrk{}गो\appbrk{}जी\appbrk{}भू\appbrk{}ज\appbrk{}ति\appbrk{}ल्व\appbrk{}के\appbrk{}क्षु\appbrk{}र\appbrk{}क\appbrk{}गो\appbrk{}प\appbrk{}घो\appbrk{}ण्टा\appbrk{}रि\appbrk{}मे\appbrk{}दानां}}\Afootnote{\skt{धवा\appbrk{}श्व\appbrk{}क\appbrk{}र्ण\appbrk{}शि\appbrk{}री\appbrk{}ष\appbrk{}ति\appbrk{}नि\appbrk{}श\appbrk{}प\appbrk{}ला\appbrk{}श\appbrk{}पि\appbrk{}चु\appbrk{}म\appbrk{}र्द\appbrk{}पा\appbrk{}ट\appbrk{}लि\appbrk{}पा\appbrk{}रि\appbrk{}भ\appbrk{}द्र\appbrk{}का\appbrk{}म्रो\appbrk{}दु\appbrk{}म्ब\appbrk{}र\appbrk{}क\appbrk{}र\appbrk{}हा\appbrk{}ट\appbrk{}का\appbrk{}र्जु\appbrk{}न\appbrk{}क\appbrk{}कु\appbrk{}भ\appbrk{}स\appbrk{}र्ज\appbrk{}क\appbrk{}पी\appbrk{}त\appbrk{}न\appbrk{}श्ले\appbrk{}ष्मा\appbrk{}त\appbrk{}का\appbrk{}ङ्को\appbrk{}ठा\appbrk{}म\appbrk{}ल\appbrk{}क\appbrk{}प्र\appbrk{}ग्र\appbrk{}ह\appbrk{}कु\appbrk{}ट\appbrk{}ज\appbrk{}श\appbrk{}मी\appbrk{}क\appbrk{}पि\appbrk{}त्था\appbrk{}श्म\appbrk{}न्त\appbrk{}का\appbrk{}र्क\appbrk{}चि\appbrk{}र\appbrk{}बि\appbrk{}ल्व\appbrk{}म\appbrk{}हा\appbrk{}वृ\appbrk{}क्षा\appbrk{}रु\appbrk{}ष्क\appbrk{}रा\appbrk{}र\appbrk{}लु\appbrk{}म\appbrk{}धु\appbrk{}क\appbrk{}म\appbrk{}धु\appbrk{}शि\appbrk{}ग्रु\appbrk{}शा\appbrk{}क\appbrk{}गो\appbrk{}जी\appbrk{}मू\appbrk{}र्वा\appbrk{}भू\appbrk{}र्ज\appbrk{}ति\appbrk{}ल्व\appbrk{}के\appbrk{}क्षु\appbrk{}र\appbrk{}क\appbrk{}गो\appbrk{}प\appbrk{}घो\appbrk{}ण्टा\appbrk{}रि\appbrk{}मे\appbrk{}दानां}\ A; \skt{धवा\appbrk{}श्व\appbrk{}क\appbrk{}र्ण्ण\appbrk{}ति\appbrk{}नि\appbrk{}श\appbrk{}पि\appbrk{}चु\appbrk{}म\appbrk{}र्द\appbrk{}पा\appbrk{}ट\appbrk{}ली\appbrk{}पा\appbrk{}रि\appbrk{}भ\appbrk{}द्र\appbrk{}को\appbrk{}दु\appbrk{}म्व\appbrk{}र\appbrk{}क\appbrk{}र\appbrk{}घा\appbrk{}त\appbrk{}का\appbrk{}र्जु\appbrk{}न\appbrk{}स\appbrk{}र्ज्ज\appbrk{}क\appbrk{}पी\appbrk{}त\appbrk{}न\appbrk{}श्ले\appbrk{}ष्मा\appbrk{}त\appbrk{}का\appbrk{}ङ्को\appbrk{}ठ\appbrk{}कु\appbrk{}ट\appbrk{}ज\appbrk{}श\appbrk{}मी\appbrk{}क\appbrk{}पि\appbrk{}त्था\appbrk{}श्म\appbrk{}न्त\appbrk{}का\appbrk{}र्क\appbrk{}चि\appbrk{}रि\appbrk{}वि\appbrk{}ल्व\appbrk{}म\appbrk{}हा\appbrk{}वृ\appbrk{}क्षा\appbrk{}र\appbrk{}ल\appbrk{}म\appbrk{}धु\appbrk{}क\appbrk{}म\appbrk{}धु\appbrk{}कं\appbrk{}श्रि\appbrk{}ग्रु\appbrk{}शा\appbrk{}क\appbrk{}गो\appbrk{}जी\appbrk{}भू\appbrk{}ज\appbrk{}ति\appbrk{}ल्व\appbrk{}के\appbrk{}क्षु\appbrk{}र\appbrk{}क\appbrk{}गो\appbrk{}प\appbrk{}घो\appbrk{}ण्टा\appbrk{}रि\appbrk{}मे\appbrk{}दानां}\ K; \skt{धवा\appbrk{}श्व\appbrk{}क\appbrk{}र्ण्ण\appbrk{}ति\appbrk{}नि\appbrk{}श\appbrk{}पि\appbrk{}चु\appbrk{}म\appbrk{}र्द्द\appbrk{}पि\appbrk{}चु\appbrk{}म\appbrk{}र्द्द\appbrk{}पा\appbrk{}ट\appbrk{}ली\appbrk{}पा\appbrk{}रि\appbrk{}भ\appbrk{}द्र\appbrk{}को\appbrk{}दु\appbrk{}म्व\appbrk{}र\appbrk{}क\appbrk{}र\appbrk{}घा\appbrk{}टा\appbrk{}र्ज्जु\appbrk{}न\appbrk{}स\appbrk{}र्ज्ज\appbrk{}क\appbrk{}पी\appbrk{}त\appbrk{}न\appbrk{}श्ले\appbrk{}ष्मा\appbrk{}न्त\appbrk{}कां\appbrk{}को\appbrk{}ठ\appbrk{}कु\appbrk{}ट\appbrk{}ज\appbrk{}श\appbrk{}मी\appbrk{}क\appbrk{}पि\appbrk{}त्था\appbrk{}श्म\appbrk{}न्त\appbrk{}का\appbrk{}र्क्क\appbrk{}चि\appbrk{}रि\appbrk{}वि\appbrk{}ल्व\appbrk{}म\appbrk{}हा\appbrk{}वृ\appbrk{}क्षा\appbrk{}र\appbrk{}लु\appbrk{}म\appbrk{}धु\appbrk{}क\appbrk{}म\appbrk{}धु\appbrk{}शि\appbrk{}ग्रु\appbrk{}शा\appbrk{}क\appbrk{}गो\appbrk{}जी\appbrk{}भू\appbrk{}ज\appbrk{}ति\appbrk{}ल्व\appbrk{}के\appbrk{}क्षु\appbrk{}र\appbrk{}क\appbrk{}गो\appbrk{}प\appbrk{}घो\appbrk{}ण्टा\appbrk{}रि\appbrk{}मे\appbrk{}दा\appbrk{}नाम्}\ H}}\skt{ }\edtext{\skt{भस्मान्याहृत्य}}{\lemma{\skt{भस्मान्य्}}\Afootnote{A K; \skt{भश्मान्य्}\ H}}\skt{ गवां मूत्रेण क्षारकल्पेन परिस्राव्य विपचेत्~। दद्याच्चात्र }\edtext{\skt{पिप्प\-ली\-पि\-प्प\-ली\-मू\-ल\-त\-ण्डु\-ली\-य\-क\-व\-रा\-ङ्ग\-चो\-र\-क\-म\-ञ्जि\-ष्ठा\-क\-र\-ञ्जि\-का\-ह\-स्ति\-पि\-प्प\-ली\-वि\-ड\-ङ्गा\-गृ\-ह\-धू\-मा\-न\-न्त\-सो\-म\-स\-र\-ल\-बा\-ह्ली\-क\-कु\-शा\-म्र\-स\-र्ष\-प\-व\-रु\-ण\-प्ल\-क्ष\-नि\-चु\-ल\-व\-र्ध\-मा\-न\-व\-ञ्जु\-ल\-पु\-त्र\-श्रे\-णी\-स\-प्त\-प\-र्ण्ण\-टु\-ण्टु\-कै\-ल\-वा\-लु\-क\-ना\-ग\-द\-न्त्य\-ति\-वि\-षा\-भ\-द्र\-दा\-रु\-म\-रि\-च\-कु\-ष्ठ\-व\-चा\-चू\-र्ण्णानि}}{\lemma{\skt{पिप्प\appbrk{}ली\appbrk{}पि\appbrk{}प्प\appbrk{}ली\appbrk{}मू\appbrk{}ल\appbrk{}त\appbrk{}ण्डु\appbrk{}ली\appbrk{}य\appbrk{}क\appbrk{}व\appbrk{}रा\appbrk{}ङ्ग\appbrk{}चो\appbrk{}र\appbrk{}क\appbrk{}म\appbrk{}ञ्जि\appbrk{}ष्ठा\appbrk{}क\appbrk{}र\appbrk{}ञ्जि\appbrk{}का\appbrk{}ह\appbrk{}स्ति\appbrk{}पि\appbrk{}प्प\appbrk{}ली\appbrk{}वि\appbrk{}ड\appbrk{}ङ्गा\appbrk{}गृ\appbrk{}ह\appbrk{}धू\appbrk{}मा\appbrk{}न\appbrk{}न्त\appbrk{}सो\appbrk{}म\appbrk{}स\appbrk{}र\appbrk{}ल\appbrk{}बा\appbrk{}ह्ली\appbrk{}क\appbrk{}कु\appbrk{}शा\appbrk{}म्र\appbrk{}स\appbrk{}र्ष\appbrk{}प\appbrk{}व\appbrk{}रु\appbrk{}ण\appbrk{}प्ल\appbrk{}क्ष\appbrk{}नि\appbrk{}चु\appbrk{}ल\appbrk{}व\appbrk{}र्ध\appbrk{}मा\appbrk{}न\appbrk{}व\appbrk{}ञ्जु\appbrk{}ल\appbrk{}पु\appbrk{}त्र\appbrk{}श्रे\appbrk{}णी\appbrk{}स\appbrk{}प्त\appbrk{}प\appbrk{}र्ण्ण\appbrk{}टु\appbrk{}ण्टु\appbrk{}कै\appbrk{}ल\appbrk{}वा\appbrk{}लु\appbrk{}क\appbrk{}ना\appbrk{}ग\appbrk{}द\appbrk{}न्त्य\appbrk{}ति\appbrk{}वि\appbrk{}षा\appbrk{}भ\appbrk{}द्र\appbrk{}दा\appbrk{}रु\appbrk{}म\appbrk{}रि\appbrk{}च\appbrk{}कु\appbrk{}ष्ठ\appbrk{}व\appbrk{}चा\appbrk{}चू\appbrk{}र्ण्णानि}}\Afootnote{\skt{पिप्प\appbrk{}ली\appbrk{}मू\appbrk{}ल\appbrk{}त\appbrk{}ण्डु\appbrk{}ली\appbrk{}य\appbrk{}क\appbrk{}व\appbrk{}रा\appbrk{}ङ्ग\appbrk{}चो\appbrk{}च\appbrk{}म\appbrk{}ञ्जि\appbrk{}ष्ठा\appbrk{}क\appbrk{}र\appbrk{}ञ्जि\appbrk{}का\appbrk{}ह\appbrk{}स्ति\appbrk{}पि\appbrk{}प्प\appbrk{}ली\appbrk{}म\appbrk{}रि\appbrk{}च\appbrk{}वि\appbrk{}ड\appbrk{}ङ्ग\appbrk{}गृ\appbrk{}ह\appbrk{}धू\appbrk{}मा\appbrk{}न\appbrk{}न्ता\appbrk{}सो\appbrk{}म\appbrk{}स\appbrk{}र\appbrk{}ला\appbrk{}बा\appbrk{}ह्ली\appbrk{}क\appbrk{}गु\appbrk{}हा\appbrk{}को\appbrk{}शा\appbrk{}म्र\appbrk{}श्वे\appbrk{}त\appbrk{}स\appbrk{}र्ष\appbrk{}प\appbrk{}व\appbrk{}रु\appbrk{}ण\appbrk{}ल\appbrk{}व\appbrk{}ण\appbrk{}प्ल\appbrk{}क्ष\appbrk{}नि\appbrk{}चु\appbrk{}ल\appbrk{}क\appbrk{}व\appbrk{}ञ्जु\appbrk{}ल\appbrk{}व\appbrk{}क्रा\appbrk{}ल\appbrk{}व\appbrk{}र्ध\appbrk{}मा\appbrk{}न\appbrk{}पु\appbrk{}त्र\appbrk{}श्रे\appbrk{}णी\appbrk{}स\appbrk{}प्त\appbrk{}प\appbrk{}र्ण\appbrk{}टु\appbrk{}ण्टु\appbrk{}कै\appbrk{}ल\appbrk{}वा\appbrk{}लु\appbrk{}क\appbrk{}ना\appbrk{}ग\appbrk{}द\appbrk{}न्त्य\appbrk{}ति\appbrk{}वि\appbrk{}षा\appbrk{}भ\appbrk{}या\appbrk{}भ\appbrk{}द्र\appbrk{}दा\appbrk{}रु\appbrk{}कु\appbrk{}ष्ठ\appbrk{}ह\appbrk{}रि\appbrk{}द्रा\appbrk{}व\appbrk{}चा\appbrk{}चू\appbrk{}र्णानि}\ A; \skt{पिप्प\appbrk{}ली\appbrk{}पि\appbrk{}प्प\appbrk{}ली\appbrk{}मू\appbrk{}ल\appbrk{}त\appbrk{}ण्डु\appbrk{}ली\appbrk{}य\appbrk{}क\appbrk{}व\appbrk{}रां\appbrk{}ग\appbrk{}चो\appbrk{}र\appbrk{}क\appbrk{}म\appbrk{}ञ्जि\appbrk{}ष्ठा\appbrk{}क\appbrk{}र\appbrk{}ञ्जि\appbrk{}का\appbrk{}ह\appbrk{}स्ति\appbrk{}पि\appbrk{}प्प\appbrk{}ली\appbrk{}वि\appbrk{}ड\appbrk{}ङ्गा\appbrk{}गृ\appbrk{}ह\appbrk{}धू\appbrk{}मा\appbrk{}न\appbrk{}न्त\appbrk{}सो\appbrk{}म\appbrk{}स\appbrk{}र\appbrk{}ल\appbrk{}वा\appbrk{}वा\appbrk{}ल्ही\appbrk{}क\appbrk{}कु\appbrk{}शा\appbrk{}म्र\appbrk{}स\appbrk{}र्ष\appbrk{}प\appbrk{}व\appbrk{}रु\appbrk{}ण\appbrk{}प्ल\appbrk{}क्ष\appbrk{}नि\appbrk{}चु\appbrk{}ल\appbrk{}व\appbrk{}र्ध\appbrk{}मा\appbrk{}न\appbrk{}व\appbrk{}ञ्जु\appbrk{}ल\appbrk{}पु\appbrk{}त्र\appbrk{}श्रे\appbrk{}णी\appbrk{}स\appbrk{}प्त\appbrk{}प\appbrk{}र्ण्ण\appbrk{}टु\appbrk{}ण्टु\appbrk{}कै\appbrk{}ल\appbrk{}वा\appbrk{}लु\appbrk{}क\appbrk{}ना\appbrk{}ग\appbrk{}द\appbrk{}न्त्य\appbrk{}ति\appbrk{}वि\appbrk{}षा\appbrk{}भ\appbrk{}द्र\appbrk{}दा\appbrk{}रु\appbrk{}म\appbrk{}रि\appbrk{}च\appbrk{}कु\appbrk{}ष्ठ\appbrk{}व\appbrk{}चा\appbrk{}चू\appbrk{}र्ण्णानि}\ K; \skt{पिप्प\appbrk{}ली\appbrk{}मू\appbrk{}ल\appbrk{}त\appbrk{}ण्डु\appbrk{}ली\appbrk{}य\appbrk{}क\appbrk{}व\appbrk{}रा\appbrk{}ङ्ग\appbrk{}चो\appbrk{}च\appbrk{}क\appbrk{}म\appbrk{}ञ्जि\appbrk{}ष्ठा\appbrk{}क\appbrk{}र\appbrk{}ञ्जि\appbrk{}का\appbrk{}ह\appbrk{}स्ति\appbrk{}पि\appbrk{}प्प\appbrk{}ली\appbrk{}वि\appbrk{}ड\appbrk{}ङ्गा\appbrk{}गृ\appbrk{}ह\appbrk{}धू\appbrk{}मा\appbrk{}न\appbrk{}न्त\appbrk{}सो\appbrk{}म\appbrk{}स\appbrk{}र\appbrk{}ल\appbrk{}वा\appbrk{}ह्ली\appbrk{}क\appbrk{}कु\appbrk{}शा\appbrk{}म्र\appbrk{}स\appbrk{}र्ष\appbrk{}प\appbrk{}व\appbrk{}रु\appbrk{}ण\appbrk{}प्ल\appbrk{}क्ष\appbrk{}नि\appbrk{}चु\appbrk{}ल\appbrk{}व\appbrk{}र्द्ध\appbrk{}मा\appbrk{}न\appbrk{}व\appbrk{}ञ्जु\appbrk{}ल\appbrk{}पु\appbrk{}त्र\appbrk{}श्रे\appbrk{}णी\appbrk{}स\appbrk{}प्त\appbrk{}प\appbrk{}र्ण्ण\appbrk{}टु\appbrk{}ण्टु\appbrk{}कै\appbrk{}ल\appbrk{}वा\appbrk{}लु\appbrk{}ना\appbrk{}ग\appbrk{}द\appbrk{}न्त्य\appbrk{}ति\appbrk{}वि\appbrk{}षा\appbrk{}भ\appbrk{}द्र\appbrk{}दा\appbrk{}रु\appbrk{}म\appbrk{}रि\appbrk{}च\appbrk{}कु\appbrk{}ष्ठ\appbrk{}व\appbrk{}चा\appbrk{}चू\appbrk{}र्ण्णानि}\ H}}\skt{ }\edtext{\skt{लोहानां}}{\lemma{\skt{लोहानां}}\Afootnote{K H; \skt{लोहानां च}\ A}}\skt{ समभागानि ततः क्षार\-व\-दा\-ग\-त\-पा\-क\-म\-व\-तार्य लोहकुम्भे निदध्यात्~।}\pend

\pstart\versepar \ledsidenote{\passref{5.6.4}}\edtext{\skt{एतेन}}{\lemma{\skt{एतेन}}\Afootnote{K H; \skt{अनेन}\ A}}\skt{ दुन्दुभिं }\edtext{\skt{लिम्पे\-त्प\-ता\-का\-स्त\-र\-णानि च}}{\lemma{\skt{लिम्पे\appbrk{}त्प\appbrk{}ता\appbrk{}का\appbrk{}स्त\appbrk{}र\appbrk{}णानि च}}\Afootnote{K H; \skt{लिम्पेत्पताकां तोरणानि च~।}\ A}}\skt{~॥}\skt{ }\pend
\pstart\versepar \skt{}\edtext{\skt{दर्श\-ना\-च्छ्र\-व\-णा\-च्चापि}}{\lemma{\skt{दर्श\appbrk{}ना\appbrk{}च्छ्र\appbrk{}व\appbrk{}णा\appbrk{}च्चापि}}\Afootnote{K H; \skt{श्रव\appbrk{}णा\appbrk{}द्द\appbrk{}र्श\appbrk{}ना\appbrk{}त्स्पर्शा}\ A}}\skt{ }\edtext{\skt{विषा\-न्स\-र्वा\-न्प्र\-मु\-च्यते}}{\lemma{\skt{विषा\appbrk{}न्स\appbrk{}र्वा\appbrk{}न्प्र\appbrk{}मु\appbrk{}च्यते}}\Afootnote{K H; \skt{त्विषा\appbrk{}त्सं\appbrk{}प्र\appbrk{}ति\appbrk{}मु\appbrk{}च्यते~॥}\ A}}\skt{~।}\pend

\pstart\versepar \ledsidenote{\passref{5.6.5}}\skt{एष क्षारागदो नाम शर्क\-रा\-स्व\-श्म\-रीषु च~। }\pend
\pstart\versepar \skt{अर्शस्सु वातगुल्मेषु कास\-शू\-लो\-द\-रेषु च~।}\pend

\pstart\versepar \ledsidenote{\passref{5.6.6}}\skt{अजीर्ण्णे }\edtext{\skt{ग्रहणीदोषे}}{\lemma{\skt{ग्रहणीदोषे}}\Afootnote{A H; \skt{ग्रहणेदोषे}\ K}}\skt{ भक्तद्वेषे च दारुणे~। }\pend
\pstart\versepar \skt{शोफे सर्वसरे चापि देयः }\edtext{\skt{श्वासे च}}{\lemma{\skt{श्वासे च}}\Afootnote{A K H\textsuperscript{pc}; \skt{श्वासे}\ H\textsuperscript{ac}}}\skt{ }\edtext{\skt{दुस्तरे}}{\lemma{\skt{दुस्तरे}}\Afootnote{K H; \skt{दारुणे~॥}\ A}}\skt{~।}\pend

\pstart\versepar \ledsidenote{\passref{5.6.7}}\edtext{\skt{एष}}{\lemma{\skt{एष}}\Afootnote{K H; \skt{सदा}\ A}}\skt{ सर्वविषार्त्तानां सर्व\-थै\-वो\-प\-यु\-ज्यते~। }\pend
\pstart\versepar \skt{}\edtext{\skt{तथा}}{\lemma{\skt{तथा}}\Afootnote{K H; \skt{एष}\ A}}\skt{ तक्ष\-क\-मु\-ख्या\-ना\-मपि }\edtext{\edtext{\skt{सर्पांकुशोऽगदः}}{\lemma{\skt{सर्पांकुशो}}\Afootnote{K H; \skt{दर्पाङ्कुशो}\ A}}}{\lemma{\skt{ऽगदः}}\Afootnote{A; \skt{गदः~॥}\ K H}}\skt{~॥}\pend

\pstart\versepar \ledsidenote{\passref{5.6.12}}\skt{अपामार्गस्य }\edtext{\skt{बीजानि}}{\lemma{\skt{बीजानि}}\Afootnote{A; \skt{वीजानि}\ K H\textsuperscript{pc}; \skt{च वीजानि}\ H\textsuperscript{ac}}}\skt{ शिरीषस्य च }\edtext{\skt{माषकाम्}}{\lemma{\skt{माषकाम्}}\Afootnote{K H; \skt{माषका}\ A}}\skt{~। }\pend
\pstart\versepar \skt{}\edtext{\skt{श्वेते द्वे}}{\lemma{\skt{श्वेते द्वे}}\Afootnote{K H; \skt{न्~। श्वेते द्वे}\ A}}\skt{ काकमाचीञ्च गवाम्मूत्रेण }\edtext{\skt{पीषयेत्}}{\lemma{\skt{पीषयेत्}}\Afootnote{K H; \skt{पेषयेत्~॥}\ A}}\skt{~।}\pend

\pstart\versepar \ledsidenote{\passref{5.6.13}}\edtext{\skt{सर्पिरेतेषु}}{\lemma{\skt{एतेषु}}\Afootnote{K H; \skt{एतैस्तु}\ A}}\skt{ संसिद्धं विषसंशमनं परम्~। }\pend
\pstart\versepar \skt{अमृतं नाम विख्यातमपि }\edtext{\skt{सञ्जी\-व\-ये\-त्मृ\-तम्}}{\lemma{\skt{सञ्जी\appbrk{}व\appbrk{}ये\appbrk{}त्मृ\appbrk{}तम्}}\Afootnote{K H; \skt{संजी\appbrk{}व\appbrk{}ये\appbrk{}न्मृ\appbrk{}तम्~॥}\ A}}\skt{~॥}\pend

\pstart\versepar \ledsidenote{\passref{5.6.14}}\edtext{\skt{चन्दनागुरुणी}}{\lemma{\skt{चन्दनागुरुणी}}\Afootnote{A K; \skt{चन्दनागरुणी}\ H}}\skt{ कुष्ठं तगरं }\edtext{\skt{तैलपर्ण्णिकम्}}{\lemma{\skt{तैलपर्ण्णिकम्}}\Afootnote{K; \skt{तिलपर्णिकम्~।}\ A; \skt{तैलपर्ण्णिकां~।}\ H}}\skt{~। }\pend
\pstart\versepar \skt{प्रपौ\-ण्ड\-री\-क\-न्न\-लदं सरलं देवदारु च~।}\pend

\pstart\versepar \ledsidenote{\passref{5.6.15}}\skt{भद्रश्रियं यवफलां भार्गीन्नीलीं सुगन्धिकाम्~। }\pend
\pstart\versepar \skt{कालेयकं पद्मकञ्च मधुकं }\edtext{\skt{सनखां}}{\lemma{\skt{सनखां}}\Afootnote{H; \skt{नागरं}\ A; \skt{सनखं}\ K}}\skt{ जटाम्~।}\pend

\pstart\versepar \ledsidenote{\passref{5.6.16}}\skt{पुन्ना\-गै\-लै\-ल\-वा\-लूनि गैरिकं ध्यामकं }\edtext{\skt{तथा}}{\lemma{\skt{तथा}}\Afootnote{K H; \skt{बलाम्~।}\ A}}\skt{~। }\pend
\pstart\versepar \skt{तोयं सर्जरसं मांसीं शतपुष्पां हरेणुकाम्~।}\pend

\pstart\versepar \ledsidenote{\passref{5.6.17}}\edtext{\skt{तालीसपत्रं}}{\lemma{\skt{तालीसपत्रं}}\Afootnote{K H; \skt{तालीशपत्रं}\ A}}\skt{ क्षुद्रैलां }\edtext{\skt{प्रियङ्गू}}{\lemma{\skt{प्रियङ्गू}}\Afootnote{K; \skt{प्रियङ्गुं}\ A; \skt{प्रियंगूं}\ H}}\skt{ }\edtext{\skt{सकुटन्नटाम्}}{\lemma{\skt{सकुटन्नटाम्}}\Afootnote{H; \skt{सकुटन्नटम्~।}\ A; \skt{सकुटंन्नटां~।}\ K}}\skt{~। }\pend
\pstart\versepar \skt{}\edtext{\skt{तिलपुष्पं}}{\lemma{\skt{तिलपुष्पं}}\Afootnote{K H; \skt{शिलापुष्पं}\ A}}\skt{ सशैलेयम्पत्रं }\edtext{\skt{काला\-नु\-सा\-रि\-वाम्}}{\lemma{\skt{काला\appbrk{}नु\appbrk{}सा\appbrk{}रि\appbrk{}वाम्}}\Afootnote{A K; \skt{कालानुशारिवां~॥}\ H}}\skt{~।}\pend

\pstart\versepar \ledsidenote{\passref{5.6.18}}\skt{कटुत्रिकं शीतशिवं काश्मर्यं कटुरोहिणीम्~। }\pend
\pstart\versepar \skt{}\edtext{\edtext{\skt{सोम\-रा\-जी\-म\-ति\-विषां}}{\lemma{\skt{सोमराजीम्}}\Afootnote{H; \skt{सोमराजी}\ A K}}}{\lemma{\skt{अतिविषां}}\Afootnote{\skt{मतिविषां}\ A K; \skt{अतिविषा}\ H}}\skt{ }\edtext{\edtext{\skt{पृथ्वी\-का\-मि\-न्द्र\-वा\-रु\-णीम्}}{\lemma{\skt{पृथ्वीकाम्}}\Afootnote{H; \skt{पृथ्विकाम्}\ A; \skt{पृथ्वीका}\ K}}}{\lemma{\skt{इन्द्रवारुणीम्}}\Afootnote{A H; \skt{मिन्द्रवारुणीं~।}\ K}}\skt{~।}\pend

\pstart\versepar \ledsidenote{\passref{5.6.19}}\edtext{\skt{उशीरे द्वे}}{\lemma{\skt{उशीरे द्वे}}\Afootnote{K H; \skt{उशीरं}\ A}}\skt{ }\edtext{\skt{वरुणकं}}{\lemma{\skt{वरुणकं}}\Afootnote{K H; \skt{वरुणं मुस्तं}\ A}}\skt{ }\edtext{\skt{कुस्तुम्बुर्यो}}{\lemma{\skt{कुस्तुम्बुर्यो}}\Afootnote{\skt{कुस्तुम्बुरु}\ A; \skt{कुस्तुम्वुर्यो}\ K H}}\skt{ }\edtext{\skt{नखानि च}}{\lemma{\skt{नखानि च}}\Afootnote{K H; \skt{न}\ A}}\skt{~। }\pend
\pstart\versepar \skt{}\edtext{\skt{त्वचं}}{\lemma{\skt{त्वचं}}\Afootnote{K H; \emph{om.}\ A}}\skt{ }\edtext{\skt{तस्करसाह्वञ्च}}{\lemma{\skt{तस्करसाह्वञ्च}}\Afootnote{K H; \emph{om.}\ A}}\skt{ }\edtext{\skt{ग्रन्थिलां}}{\lemma{\skt{ग्रन्थिलां}}\Afootnote{K H; \emph{om.}\ A}}\skt{ }\edtext{\skt{सहरीतकीम्}}{\lemma{\skt{सहरीतकीम्}}\Afootnote{K H; \skt{खं तथा~।}\ A}}\skt{~। }\pend
\pstart\versepar \skt{श्वेते हरिद्रे स्थौणेयं लाक्षाञ्च }\edtext{\skt{लवणानि}}{\lemma{\skt{लवणानि}}\Afootnote{A H; \skt{लवनानि}\ K}}\skt{ च~।}\pend

\pstart\versepar \ledsidenote{\passref{5.6.20}}\skt{कुमु\-दो\-त्प\-ल\-प\-द्मानि पुष्पञ्चापि }\edtext{\skt{तथार्जकम्}}{\lemma{\skt{तथार्जकम्}}\Afootnote{K H; \skt{तथाऽर्कजम्~।}\ A}}\skt{~। }\pend
\pstart\versepar \skt{}\edtext{\skt{चम्प\-का\-शो\-क\-सु\-म\-ना\-ति\-ल\-क\-प्र\-स\-वानि}}{\lemma{\skt{चम्प\appbrk{}का\appbrk{}शो\appbrk{}क\appbrk{}सु\appbrk{}म\appbrk{}ना\appbrk{}ति\appbrk{}ल\appbrk{}क\appbrk{}प्र\appbrk{}स\appbrk{}वानि}}\Afootnote{K; \skt{चम्प\appbrk{}का\appbrk{}शो\appbrk{}क\appbrk{}सु\appbrk{}म\appbrk{}न\appbrk{}स्ति\appbrk{}ल्व\appbrk{}क\appbrk{}प्र\appbrk{}स\appbrk{}वानि}\ A; \skt{चम्प\appbrk{}का\appbrk{}शो\appbrk{}क\appbrk{}सु\appbrk{}म\appbrk{}ना\appbrk{}स्ति\appbrk{}ल\appbrk{}क\appbrk{}प्र\appbrk{}स\appbrk{}वानि}\ H}}\skt{ च~।}\pend

\pstart\versepar \ledsidenote{\passref{5.6.21}}\edtext{\skt{पाट\-ली\-शा\-ल्म\-ली\-शे\-लू\-शि\-री\-षा\-णा\-न्त\-थैव}}{\lemma{\skt{पाट\appbrk{}ली\appbrk{}शा\appbrk{}ल्म\appbrk{}ली\appbrk{}शे\appbrk{}लू\appbrk{}शि\appbrk{}री\appbrk{}षा\appbrk{}णान्}}\Afootnote{H; \skt{पाट\appbrk{}ली\appbrk{}शा\appbrk{}ल्म\appbrk{}ली\appbrk{}शै\appbrk{}लु\appbrk{}शि\appbrk{}री\appbrk{}षाणां}\ A; \skt{पाट\appbrk{}ली\appbrk{}शा\appbrk{}ल्म\appbrk{}ली\appbrk{}शे\appbrk{}लू\appbrk{}शि\appbrk{}री\appbrk{}शा\appbrk{}णान्}\ K}}\skt{ च~। }\pend
\pstart\versepar \skt{}\edtext{\edtext{\skt{सुर\-स्या\-स्तृ\-ण\-मू\-ल्यस्य}}{\lemma{\skt{सुरस्यास्}}\Afootnote{K H; \skt{कुसुमं}\ A}}}{\lemma{\skt{तृणमूल्यस्य}}\Afootnote{K; \skt{तृणमूल्याश्च सुर}\ A; \skt{तृणशूल्यस्य}\ H}}\skt{ }\edtext{\skt{सिन्दुवारस्य}}{\lemma{\skt{सिन्दुवारस्य}}\Afootnote{K H; \skt{भीसि\appbrk{}न्धु\appbrk{}वा\appbrk{}र\appbrk{}जम्~॥}\ A}}\skt{ }\edtext{\skt{यानि}}{\lemma{\skt{यानि}}\Afootnote{K H; \emph{om.}\ A}}\skt{ }\edtext{\skt{च}}{\lemma{\skt{च}}\Afootnote{K H; \emph{om.}\ A}}\skt{~।}\pend

\pstart \ledsidenote{\passref{5.6.22}}\edtext{\edtext{\skt{धवा\-श्व\-क\-र्ण्ण\-यो\-श्चापि}}{\lemma{\skt{धवा\appbrk{}श्व\appbrk{}क\appbrk{}र्ण्ण\appbrk{}योश्}}\Afootnote{K H; \skt{धवाश्वकर्णपा}\ A}}}{\lemma{\skt{चापि}}\Afootnote{K H; \skt{र्थानां}\ A}}\skt{ }\edtext{\skt{पुष्पाणि}}{\lemma{\skt{पुष्पाणि}}\Afootnote{A H; \skt{पुष्पानि}\ K}}\skt{ तिनिशस्य }\edtext{\skt{च}}{\lemma{\skt{च}}\Afootnote{K H; \skt{च~। गुग्गुलुं कुङ्कुमं बिम्बीं सर्पाक्षीं गन्धनाकुलीम्~॥}\ A}}\skt{~।}\pend

\pstart\versepar \ledsidenote{\passref{5.6.23}}\edtext{\skt{एतत्सम्भृत्य}}{\lemma{\skt{एतत्सम्भृत्य}}\Afootnote{K H; \skt{एतत्संभृत्य}\ A}}\skt{ }\edtext{\skt{सम्भारं}}{\lemma{\skt{सम्भारं}}\Afootnote{K H; \skt{संभारं}\ A}}\skt{ }\edtext{\skt{सूक्ष्मं}}{\lemma{\skt{सूक्ष्मं}}\Afootnote{K; \skt{सूक्ष्म}\ A; \skt{शूक्ष्म}\ H}}\skt{ }\edtext{\skt{चूर्ण्णं तु}}{\lemma{\skt{चूर्ण्णं तु}}\Afootnote{K H; \skt{चूर्णानि}\ A}}\skt{ कारयेत्~। }\pend
\pstart\versepar \skt{गोपि\-त्त\-म\-धु\-स\-र्पि\-र्भि\-र्युक्तं शृङ्गे निधापयेत्~।}\pend

\pstart\versepar \ledsidenote{\passref{5.6.24}}\edtext{\skt{भग्न\-स्क\-न्ध\-वि\-वृ\-त्ताक्षं}}{\lemma{\skt{भग्न\appbrk{}स्क\appbrk{}न्ध\appbrk{}वि\appbrk{}वृ\appbrk{}त्ताक्षं}}\Afootnote{K; \skt{भग्नस्कन्धं विवृताक्षं}\ A H}}\skt{ मृत्योर्दंष्ट्रान्तरं गतम्~। }\pend
\pstart\versepar \skt{अने\-ना\-ग\-द\-मु\-ख्येन मनुष्यं }\edtext{\skt{पुनरानयेत्}}{\lemma{\skt{आनयेत्}}\Afootnote{K H; \skt{आहरेत्~॥}\ A}}\skt{~।}\pend

\pstart\versepar \ledsidenote{\passref{5.6.25}}\skt{एषोऽग्निकल्पं दुर्वारं क्रुद्ध\-स्या\-मि\-त\-ते\-जसः~। }\pend
\pstart\versepar \skt{}\edtext{\edtext{\skt{सर्व\-ना\-ग\-प\-ते\-र्ह\-न्या\-दपि}}{\lemma{\skt{सर्व\appbrk{}ना\appbrk{}ग\appbrk{}प\appbrk{}तेर्}}\Afootnote{\skt{विषं नागपतेर्}\ A; \skt{सर्व\appbrk{}ना\appbrk{}ग\appbrk{}ग\appbrk{}तेर्}\ K H}}\skt{ वा}}{\lemma{\skt{हन्यादपि वा}}\Afootnote{K H; \skt{हन्यात्प्रसभं}\ A}}\skt{ }\edtext{\skt{वासुकेर्विषम्}}{\lemma{\skt{विषम्}}\Afootnote{K H; \skt{अपि~॥}\ A}}\skt{~।}\pend

\pstart\versepar \ledsidenote{\passref{5.6.26}}\edtext{\skt{महासुगन्धो}}{\lemma{\skt{महासुगन्धो}}\Afootnote{K H; \skt{महासुगन्धि}\ A}}\skt{ }\edtext{\skt{नाम्नायं}}{\lemma{\skt{नाम्नायं}}\Afootnote{K H; \skt{नामाऽयं}\ A}}\skt{ }\edtext{\skt{पञ्चा\-शी\-त्य\-ङ्ग\-सं\-भृतः}}{\lemma{\skt{पञ्चा\appbrk{}शी\appbrk{}त्य\appbrk{}ङ्ग\appbrk{}सं\appbrk{}भृतः}}\Afootnote{K; \skt{पञ्चा\appbrk{}शी\appbrk{}त्य\appbrk{}ङ्ग\appbrk{}सं\appbrk{}युतः~।}\ A; \skt{पञ्चा\appbrk{}शी\appbrk{}त्य\appbrk{}ङ्ग\appbrk{}स\appbrk{}म्भृतः~।}\ H}}\skt{~। }\pend
\pstart\versepar \skt{}\edtext{\skt{राजागदानां}}{\lemma{\skt{राजागदानां}}\Afootnote{K H; \skt{राजाऽगदानां}\ A}}\skt{ सर्वेषां राज्ञो हस्ते भवेत्सदा~।}\pend

\pstart\versepar \ledsidenote{\passref{5.6.27}}\edtext{\skt{तेनानुलिप्तश्च}}{\lemma{\skt{तेनानुलिप्तश्च}}\Afootnote{K H; \skt{स्नातानुलिप्तस्तु}\ A}}\skt{ नृपो भवे\-त्स\-र्व\-ज\-न\-प्रियः~। }\pend
\pstart\versepar \skt{भ्राजिष्णुताञ्च लभते शत्रु\-म\-ध्य\-ग\-तोऽपि }\edtext{\skt{सः}}{\lemma{\skt{सः}}\Afootnote{K H; \skt{सन्~॥}\ A}}\skt{~।}\pend

\pstart\versepar \ledsidenote{\passref{5.6.28}}\skt{उष्णवर्ज्यो विधिः कार्यो विषार्त्तानां विजानता~। }\pend
\pstart\versepar \skt{}\edtext{\skt{त्यक्त्वा}}{\lemma{\skt{त्यक्त्वा}}\Afootnote{K H; \skt{मुक्त्वा}\ A}}\skt{ कीटविषं तद्धि शीते\-ना\-भि\-प्र\-व\-र्धते~।}\pend

\pstart\versepar \ledsidenote{\passref{5.6.31}}\skt{दिवास्वप्नं व्यवायञ्च व्यायामं }\edtext{\edtext{\skt{क्रोधमातपम्}}{\lemma{\skt{क्रोधम्}}\Afootnote{K H; \skt{क्रोध}\ A}}}{\lemma{\skt{आतपम्}}\Afootnote{K H; \skt{मातपम्~।}\ A}}\skt{~। }\pend
\pstart\versepar \skt{}\edtext{\skt{सुरा\-ति\-ल\-कु\-ल\-त्थांश्च}}{\lemma{\skt{सुरा\appbrk{}ति\appbrk{}ल\appbrk{}कु\appbrk{}ल\appbrk{}त्थांश्च}}\Afootnote{A K; \skt{सुरा\appbrk{}ति\appbrk{}ल\appbrk{}कु\appbrk{}ल\appbrk{}त्थाश्च}\ H}}\skt{ }\edtext{\skt{वर्जयीत}}{\lemma{\skt{वर्जयीत}}\Afootnote{K H; \skt{वर्जयेद्धि}\ A}}\skt{ विषातुरः~॥}\pend

\pstart\versepar \ledsidenote{\passref{5.6.32}}\skt{प्रसन्नदोषं }\edtext{\skt{प्रकृ\-ति\-स्थ\-धा\-तु\-म\-न्ना\-भि\-कामं}}{\lemma{\skt{अन्नाभिकामं}}\Afootnote{K H; \skt{अन्नाभिकाङ्क्षं}\ A}}\skt{ }\edtext{\skt{सम\-मू\-त्र\-वि\-ट्कम्}}{\lemma{\skt{सम\appbrk{}मू\appbrk{}त्र\appbrk{}वि\appbrk{}ट्कम्}}\Afootnote{K H; \skt{सम\appbrk{}सू\appbrk{}त्र\appbrk{}जि\appbrk{}ह्वम्~।}\ A}}\skt{~। }\pend
\pstart\versepar \skt{}\edtext{\skt{प्रस\-न्न\-स\-र्वे\-न्द्रि\-य\-चि\-त्त\-चेष्टं}}{\lemma{\skt{प्रस\appbrk{}न्न\appbrk{}स\appbrk{}र्वे\appbrk{}न्द्रि\appbrk{}य\appbrk{}चि\appbrk{}त्त\appbrk{}चेष्टं}}\Afootnote{K H; \skt{प्रस\appbrk{}न्न\appbrk{}व\appbrk{}र्णे\appbrk{}न्द्रि\appbrk{}य\appbrk{}चि\appbrk{}त्त\appbrk{}चेष्टं}\ A}}\skt{ }\edtext{\skt{वैद्योऽ\-व\-ग\-च्छे\-द\-विषं}}{\lemma{\skt{ऽव\appbrk{}ग\appbrk{}च्छे\appbrk{}द\appbrk{}विषं}}\Afootnote{A H; \skt{वगच्छेदविषं}\ K}}\skt{ मनुष्यम्~॥}\pend

\pstart \ledsidenote{\passref{5.6.trailer}}\edtext{\skt{इति}}{\lemma{\skt{इति}}\Afootnote{A K; \emph{om.}\ H}}\skt{ }\edtext{\skt{कल्पे}}{\lemma{\skt{कल्पे}}\Afootnote{K; \skt{सुश्रु\appbrk{}त\appbrk{}सं\appbrk{}हि\appbrk{}तायां कल्पस्थाने दुन्दु\appbrk{}भि\appbrk{}स्व\appbrk{}नी\appbrk{}य\appbrk{}कल्पो नाम षष्ठोऽध्यायः~॥६~॥}\ A; \skt{इति~॥ कल्पस्थाने सप्तमोऽध्यायः~॥}\ H}}\skt{ ७}\pend

\endnumbering
\beginnumbering
\eledsection*{Adhyāya 8}
\pstart \ledsidenote{\passref{5.8.1}}\skt{अथातः कीटकल्पं }\edtext{\skt{व्याख्यास्यामः}}{\lemma{\skt{व्याख्यास्यामः}}\Afootnote{A H; \emph{om.}\ K}}\skt{~॥}\pend

\pstart\versepar \ledsidenote{\passref{5.8.3}}\skt{सर्पाणां }\edtext{\skt{शुक्र\-वि\-ण्मू\-त्र\-श\-व\-पू\-त्य\-ण्ड\-स\-म्भवाः}}{\lemma{\skt{शुक्र\appbrk{}वि\appbrk{}ण्मू\appbrk{}त्र\appbrk{}श\appbrk{}व\appbrk{}पू\appbrk{}त्य\appbrk{}ण्ड\appbrk{}स\appbrk{}म्भवाः}}\Afootnote{K H; \skt{शुक्र\appbrk{}वि\appbrk{}ण्मू\appbrk{}त्र\appbrk{}श\appbrk{}व\appbrk{}पू\appbrk{}त्य\appbrk{}ण्ड\appbrk{}सं\appbrk{}भवाः~।}\ A}}\skt{~। }\pend
\pstart\versepar \skt{}\edtext{\skt{वाय्व\-ग्न्य\-म्बु\-प्र\-कृ\-तयः}}{\lemma{\skt{वाय्व\appbrk{}ग्न्य\appbrk{}म्बु\appbrk{}प्र\appbrk{}कृ\appbrk{}तयः}}\Afootnote{A; \skt{वाय्व\appbrk{}ग्न्य\appbrk{}म्वु\appbrk{}प्र\appbrk{}कृ\appbrk{}तयः}\ K H}}\skt{ कीटास्तु }\edtext{\skt{त्रिविधाः}}{\lemma{\skt{त्रिविधाः}}\Afootnote{K; \skt{विविधाः}\ A H}}\skt{ स्मृताः~।}\pend

\pstart\versepar \ledsidenote{\passref{5.8.4}}\edtext{\edtext{\edtext{\skt{सर्व\-दो\-ष\-प्र\-कृ\-ति\-भि\-र्यु\-क्ता\-श्चा\-प्य\-परे}}{\lemma{\skt{युक्ताश्}}\Afootnote{A K; \skt{युकाश्}\ H}}}{\lemma{\skt{चाप्य्}}\Afootnote{K H; \skt{ते}\ A}}}{\lemma{\skt{अपरे}}\Afootnote{K H; \skt{परि}\ A}}\skt{ }\edtext{\skt{मताः}}{\lemma{\skt{मताः}}\Afootnote{K H; \skt{णामतः~।}\ A}}\skt{~। }\pend
\pstart\versepar \skt{}\edtext{\edtext{\skt{कीटास्तेऽपि}}{\lemma{\skt{कीटास्ते}}\Afootnote{K H; \skt{कीटत्वे}\ A}}}{\lemma{\skt{ऽपि}}\Afootnote{A H; \skt{पि}\ K}}\skt{ }\edtext{\skt{सुघोरास्ते}}{\lemma{\skt{सुघोरास्ते}}\Afootnote{K; \skt{सुघोराः स्युः}\ A H}}\skt{ सर्व एव }\edtext{\skt{चतुर्विधाः}}{\lemma{\skt{चतुर्विधाः}}\Afootnote{A H; \skt{चतुर्विधा~।}\ K}}\skt{~।}\pend

\pstart\versepar \ledsidenote{\passref{5.8.5}}\edtext{\skt{उण्डु\-ना\-भ\-स्तु\-ण्डि\-केरी}}{\lemma{\skt{उण्डुनाभस्}}\Afootnote{K H; \skt{कुम्भीनसस्}\ A}}\skt{ शृङ्गी }\edtext{\skt{शतकुलिम्भकाः}}{\lemma{\skt{शतकुलिम्भकाः}}\Afootnote{K H; \skt{शतकुलीरकः~।}\ A}}\skt{~। }\pend
\pstart\versepar \skt{}\edtext{\skt{उच्चि\-टि\-ङ्गा\-ग्न्य\-ल्प\-वाचः}}{\lemma{\skt{उच्चि\appbrk{}टि\appbrk{}ङ्गा\appbrk{}ग्न्य\appbrk{}ल्प\appbrk{}वाचः}}\Afootnote{\skt{उच्चि\appbrk{}टि\appbrk{}ङ्गोऽ\appbrk{}ग्नि\appbrk{}नामा च}\ A; \skt{उच्चि\appbrk{}टि\appbrk{}ङ्ग\appbrk{}स्त्य\appbrk{}ल्प\appbrk{}वाचः}\ K; \skt{उच्चि\appbrk{}टि\appbrk{}ङ्गा\appbrk{}स्त्य\appbrk{}ल्प\appbrk{}वाचो}\ H}}\skt{ }\edtext{\skt{विचि\-टि\-ङ्ग\-म\-सू\-रिकाः}}{\lemma{\skt{विचि\appbrk{}टि\appbrk{}ङ्ग\appbrk{}म\appbrk{}सू\appbrk{}रिकाः}}\Afootnote{K; \skt{चिच्चिटिङ्गो मयूरिका~॥}\ A; \skt{विचिटिङ्गा मसूरिकाः~॥}\ H}}\skt{~।}\pend

\pstart\versepar \ledsidenote{\passref{5.8.6}}\skt{आव\-र्त्त\-क\-स्त\-थो\-रभ्रः }\edtext{\skt{शारि\-का\-मु\-ख\-वै\-दलौ}}{\lemma{\skt{शारि\appbrk{}का\appbrk{}मु\appbrk{}ख\appbrk{}वै\appbrk{}दलौ}}\Afootnote{K H; \skt{सारि\appbrk{}का\appbrk{}मु\appbrk{}ख\appbrk{}वै\appbrk{}दलौ~।}\ A}}\skt{~। }\pend
\pstart\versepar \skt{}\edtext{\edtext{\skt{शत\-कु\-र्दोऽ\-भि\-रा\-जीव}}{\lemma{\skt{शतकुर्दो}}\Afootnote{K H; \skt{शरावकुर्दो}\ A}}}{\lemma{\skt{ऽभिराजीव}}\Afootnote{H; \skt{ऽभीराजिः}\ A; \skt{ऽभिराजी च}\ K}}\skt{ परु\-ष\-श्चि\-त्र\-शी\-र्षकः~।}\pend

\pstart \ledsidenote{\passref{5.8.7cd}}\edtext{\skt{अष्टादशैते}}{\lemma{\skt{अष्टादशैते}}\Afootnote{K H; \skt{अष्टादशेति}\ A}}\skt{ वायव्याः }\edtext{\skt{कीटाः}}{\lemma{\skt{कीटाः}}\Afootnote{K H; \skt{कीटाः}\ A}}\skt{ }\edtext{\skt{वातप्रकोपनाः}}{\lemma{\skt{वातप्रकोपनाः}}\Afootnote{K H; \skt{पवनकोपनाः~॥}\ A}}\skt{~।}\pend

\pstart\versepar \ledsidenote{\passref{5.8.8}}\skt{तैर्भवन्तीह दष्टानां रोगा वातनिमित्तजाः~॥ }\pend
\pstart\versepar \skt{}\edtext{\skt{कौण्डिन्यः}}{\lemma{\skt{कौण्डिन्यः}}\Afootnote{K H; \skt{कौण्डिन्य}\ A}}\skt{ }\edtext{\skt{कणभः}}{\lemma{\skt{कणभः}}\Afootnote{K H; \skt{कः कण}\ A}}\skt{ }\edtext{\skt{स्वर्गो}}{\lemma{\skt{स्वर्गो}}\Afootnote{K H; \skt{भको}\ A}}\skt{ }\edtext{\skt{वारणी}}{\lemma{\skt{वारणी}}\Afootnote{K H; \skt{वरटी}\ A}}\skt{ पत्रवृश्चिकः~।}\pend

\pstart\versepar \ledsidenote{\passref{5.8.9}}\skt{विनासिका }\edtext{\skt{ब्रह्मणीका}}{\lemma{\skt{ब्रह्मणीका}}\Afootnote{\skt{ब्राह्मणिका}\ A; \skt{व्रह्मणीका}\ K H}}\skt{ }\edtext{\skt{बिन्दुलो}}{\lemma{\skt{बिन्दुलो}}\Afootnote{A; \skt{विन्दुलो}\ K H}}\skt{ भ्रमरस्तथा~। }\pend
\pstart\versepar \skt{}\edtext{\skt{बाह्यकः}}{\lemma{\skt{बाह्यकः}}\Afootnote{\skt{बाह्यकी}\ A; \skt{वाह्यकः}\ K H}}\skt{ }\edtext{\skt{पिच्चटाः}}{\lemma{\skt{पिच्चटाः}}\Afootnote{K H; \skt{पिच्चिटः}\ A}}\skt{ कुम्भीवर्चः }\edtext{\edtext{\skt{कीरोऽरिमेदकः}}{\lemma{\skt{कीरो}}\Afootnote{K H; \skt{कीटो}\ A}}}{\lemma{\skt{ऽरिमेदकः}}\Afootnote{A H; \skt{रिमेदकः~।}\ K}}\skt{~।}\pend

\pstart\versepar \ledsidenote{\passref{5.8.10}}\skt{पद्मकीटो }\edtext{\skt{दुन्दुभको}}{\lemma{\skt{दुन्दुभको}}\Afootnote{K; \skt{दुन्दुभिको}\ A; \skt{दुन्दुभेको}\ H}}\skt{ }\edtext{\skt{मशकः}}{\lemma{\skt{मशकः}}\Afootnote{K H; \skt{मकरः}\ A}}\skt{ शतपादकः~। }\pend
\pstart\versepar \skt{पञ्चालकः }\edtext{\skt{पाकमत्स्यः}}{\lemma{\skt{पाकमत्स्यः}}\Afootnote{A H; \skt{पाकमत्स्य}\ K}}\skt{ }\edtext{\edtext{\skt{कृष्णतुण्डोऽथ}}{\lemma{\skt{कृष्णतुण्डो}}\Afootnote{A H; \skt{कृष्णतुण्डो}\ K}}}{\lemma{\skt{ऽथ}}\Afootnote{A H; \skt{थ}\ K}}\skt{ }\edtext{\skt{गर्द्दभी}}{\lemma{\skt{गर्द्दभी}}\Afootnote{A H; \skt{गर्भभी~।}\ K}}\skt{~।}\pend

\pstart\versepar \ledsidenote{\passref{5.8.11}}\edtext{\skt{कीटाः}}{\lemma{\skt{कीटाः}}\Afootnote{K H; \skt{क्लीतः}\ A}}\skt{ }\edtext{\skt{क्रिमिसरावी च }\edtext{\skt{यश्चान्यः}}{\lemma{\skt{चान्यः}}\Afootnote{K H; \skt{चाप्यु}\ A}}}{\lemma{\skt{क्रिमिसरावी च यश्}}\Afootnote{K H; \skt{कृमिसरारी च यश्}\ A}}\skt{ }\edtext{\skt{श्लेष्मकः}}{\lemma{\skt{श्लेष्मकः}}\Afootnote{K H; \skt{त्क्लेशकस्}\ A}}\skt{ }\edtext{\skt{स्मृतः}}{\lemma{\skt{स्मृतः}}\Afootnote{K H; \skt{तथा~।}\ A}}\skt{~। }\pend
\pstart\versepar \skt{एते }\edtext{\skt{ह्यग्नि\-प्र\-कृ\-त\-य\-श्च\-तु\-र्विं\-श\-ति\-री\-रिताः}}{\lemma{\skt{ईरिताः}}\Afootnote{K H; \skt{एव च~॥}\ A}}\skt{~।}\pend

\pstart\versepar \ledsidenote{\passref{5.8.12}}\skt{तैर्भवन्तीह दष्टानां }\edtext{\skt{रोगाः}}{\lemma{\skt{रोगाः}}\Afootnote{A; \skt{वेगा}\ K; \skt{वेगाः}\ H}}\skt{ पित्तनिमित्तजाः~।}\skt{ }\pend
\pstart\versepar \skt{}\edtext{\skt{वैश्वम्भरः}}{\lemma{\skt{वैश्वम्भरः}}\Afootnote{K H; \skt{विश्वम्भरः}\ A}}\skt{ पञ्चशुक्लः }\edtext{\edtext{\skt{पञ्चकृष्णोऽथ}}{\lemma{\skt{पञ्चकृष्णो}}\Afootnote{A; \skt{पञ्चकृष्णो}\ K H}}}{\lemma{\skt{ऽथ}}\Afootnote{A; \skt{थ}\ K H}}\skt{ कोकिलः~।}\pend

\pstart\versepar \ledsidenote{\passref{5.8.13}}\edtext{\skt{शैर्यकः}}{\lemma{\skt{शैर्यकः}}\Afootnote{K H; \skt{सैरेयकः}\ A}}\skt{ }\edtext{\edtext{\skt{प्रवलाकश्च}}{\lemma{\skt{प्रवलाकश्}}\Afootnote{K; \skt{प्रचलको}\ A; \skt{प्रचलाकश्}\ H}}}{\lemma{\skt{च}}\Afootnote{K H; \skt{व}\ A}}\skt{ }\edtext{\skt{भटाभः}}{\lemma{\skt{भटाभः}}\Afootnote{K H; \skt{लभः}\ A}}\skt{ }\edtext{\edtext{\skt{किटिभोऽटकी}}{\lemma{\skt{किटिभो}}\Afootnote{A K; \skt{किटिभो}\ H}}}{\lemma{\skt{ऽटकी}}\Afootnote{\skt{तथा~।}\ A; \skt{टजी~।}\ K; \skt{टजी~॥}\ H\textsuperscript{ac}; \skt{टकी~॥}\ H\textsuperscript{pc}}}\skt{~। }\pend
\pstart\versepar \skt{सूचीमुखः }\edtext{\skt{कृष्णगोधा}}{\lemma{\skt{कृष्णगोधा}}\Afootnote{A K; \skt{कृष्णगोधाः}\ H}}\skt{ }\edtext{\skt{कुष्टः}}{\lemma{\skt{कुष्टः}}\Afootnote{K; \skt{यश्च}\ A; \skt{कुकुः}\ H}}\skt{ काषायवासिकः~।}\pend

\pstart \ledsidenote{\passref{5.8.14}}\edtext{\skt{त्रयोदशैते}}{\lemma{\skt{त्रयोदशैते}}\Afootnote{K H; \skt{कीटो गर्दभकश्चैव तथा त्रोटक एव च~। त्रयोदशैते}\ A}}\skt{ }\edtext{\skt{सौम्यास्तु}}{\lemma{\skt{सौम्यास्तु}}\Afootnote{K H; \skt{सौम्याः स्युः}\ A}}\skt{ कीटाः }\edtext{\skt{श्लेष्मप्रकोपनाः}}{\lemma{\skt{श्लेष्मप्रकोपनाः}}\Afootnote{K; \skt{श्लेष्मप्रकोपणाः~॥}\ A; \skt{श्लेष्मप्रकोपजाः~॥}\ H}}\skt{~।}\pend

\pstart\versepar \ledsidenote{\passref{5.8.15}}\skt{तैर्भवन्तीह दष्टानां रोगाः }\edtext{\skt{श्लेष्मनिमित्तजाः}}{\lemma{\skt{श्लेष्मनिमित्तजाः}}\Afootnote{K H; \skt{कफनिमित्तजाः~।}\ A}}\skt{~॥ }\pend
\pstart\versepar \skt{}\edtext{\skt{तुङ्गनासो}}{\lemma{\skt{तुङ्गनासो}}\Afootnote{K H; \skt{तुङ्गीनासो}\ A}}\skt{ }\edtext{\skt{वलभिकः}}{\lemma{\skt{वलभिकः}}\Afootnote{K H; \skt{विचिलक}\ A}}\skt{ }\edtext{\skt{तोलको}}{\lemma{\skt{तोलको}}\Afootnote{K H; \skt{स्तालको}\ A}}\skt{ }\edtext{\skt{नाहनस्तथा}}{\lemma{\skt{नाहनस्}}\Afootnote{K H; \skt{वाहकस्}\ A}}\skt{~।}\pend

\pstart\versepar \ledsidenote{\passref{5.8.16}}\edtext{\skt{कोण्टागीरी}}{\lemma{\skt{कोण्टागीरी}}\Afootnote{K; \skt{कोष्ठागारी}\ A; \skt{कोण्टागारी}\ H\textsuperscript{ac}; \skt{कोण्टाञ्चागारी}\ H\textsuperscript{pc}}}\skt{ क्रिमिकरो यश्च }\edtext{\skt{मण्डलपुष्पकः}}{\lemma{\skt{मण्डलपुष्पकः}}\Afootnote{K H; \skt{मण्डलपुच्छकः~।}\ A}}\skt{~। }\pend
\pstart\versepar \skt{}\edtext{\skt{तुण्डवक्त्रः}}{\lemma{\skt{तुण्डवक्त्रः}}\Afootnote{K H; \skt{तुण्डनाभः (तुङ्गनाभः)}\ A}}\skt{ }\edtext{\skt{सर्षपकः}}{\lemma{\skt{सर्षपकः}}\Afootnote{H; \skt{सर्षपिको}\ A; \skt{सर्षपक}\ K}}\skt{ }\edtext{\skt{स्फोटकः}}{\lemma{\skt{स्फोटकः}}\Afootnote{K; \skt{वल्गुलिः}\ A; \skt{स्फोटाकः}\ H}}\skt{ }\edtext{\edtext{\skt{शम्बुकश्च}}{\lemma{\skt{शम्बुकश्}}\Afootnote{A; \skt{शम्वुकश्}\ K H}}}{\lemma{\skt{च}}\Afootnote{K H; \skt{त}\ A}}\skt{ }\edtext{\skt{यः}}{\lemma{\skt{यः}}\Afootnote{K H; \skt{था~॥}\ A}}\skt{~।}\pend

\pstart\versepar \ledsidenote{\passref{5.8.17}}\edtext{\skt{अग्निकीटाश्च}}{\lemma{\skt{अग्निकीटाश्च}}\Afootnote{K H; \skt{अग्निकीटश्च}\ A}}\skt{ }\edtext{\skt{घोराः}}{\lemma{\skt{घोराः}}\Afootnote{H; \emph{om.}\ A; \skt{घो}\ K}}\skt{ }\edtext{\edtext{\skt{स्युर्द्वादशैते}}{\lemma{\skt{स्युर्}}\Afootnote{H\textsuperscript{ac}; \skt{विज्ञेया}\ A; \skt{रा स्यु}\ K; \skt{कीटा स्युर्}\ H\textsuperscript{pc}}}}{\lemma{\skt{द्वादशैते}}\Afootnote{K H; \skt{द्वादश प्रा}\ A}}\skt{ त्रिदोषजाः~। }\pend
\pstart\versepar \skt{}\edtext{\skt{तैर्भवन्ति ह}}{\lemma{\skt{भवन्ति ह}}\Afootnote{K; \skt{भवन्तीह}\ A H}}\skt{ दष्टानां वेगज्ञानानि सर्पवत्~।}\edtext{}{\lemma{\skt{त्रिदोषजाः~। तैर्}}\Afootnote{K H; \skt{णनाशनाः~। तैर्}\ A}}\pend

\pstart \ledsidenote{\passref{5.8.20cd}}\edtext{\skt{पिटकोपचयः}}{\lemma{\skt{पिटकोपचयः}}\Afootnote{K H; \skt{पिडकोपचयः}\ A}}\skt{ }\edtext{\skt{शोफाः}}{\lemma{\skt{शोफाः}}\Afootnote{K H; \skt{शोफो}\ A}}\skt{ ग्रन्थयो मण्डलानि च~।}\pend

\pstart \ledsidenote{\passref{5.8.21ab}}\edtext{\skt{दर्द्रुश्च}}{\lemma{\skt{दर्द्रुश्च}}\Afootnote{K H; \skt{दद्रवः}\ A}}\skt{ कर्ण्णिकाश्चैव विसर्पाः किटिभानि च~।}\pend

\pstart \ledsidenote{\passref{5.8.21.add-1}}\skt{भवन्ति दंशपर्यन्ते देहे वापि विषाकुले~।}\pend

\pstart \ledsidenote{\passref{5.8.25cd}}\skt{एक\-जा\-ती\-न\-त\-स्तूर्ध्वं }\edtext{\skt{कीटान्भेदेन}}{\lemma{\skt{कीटान्भेदेन}}\Afootnote{K H; \skt{कीटा}\ A}}\skt{ }\edtext{\skt{वक्ष्यते}}{\lemma{\skt{वक्ष्यते}}\Afootnote{K H; \skt{न्वक्ष्यामि भेदतः~॥}\ A}}\skt{~।}\pend

\pstart\versepar \ledsidenote{\passref{5.8.26}}\skt{सामान्यतो दष्टलिङ्गैः }\edtext{\skt{साध्या\-सा\-ध्य\-क्र\-मेण च}}{\lemma{\skt{साध्या\appbrk{}सा\appbrk{}ध्य\appbrk{}क्र\appbrk{}मेण च}}\Afootnote{A H; \skt{साध्या\appbrk{}सा\appbrk{}ध्य\appbrk{}क्र\appbrk{}मेन च~।}\ K}}\skt{~। }\pend
\pstart\versepar \skt{}\edtext{\skt{त्रिकण्टकः}}{\lemma{\skt{त्रिकण्टकः}}\Afootnote{K H; \skt{त्रिकण्टः क}\ A}}\skt{ }\edtext{\skt{कुनी}}{\lemma{\skt{कुनी}}\Afootnote{K; \skt{रिणी}\ A; \skt{कुणी}\ H}}\skt{ चापि }\edtext{\skt{हस्ति\-क\-क्ष्योऽ\-प\-रा\-जितः}}{\lemma{\skt{हस्तिकक्ष्यो}}\Afootnote{K; \skt{हस्तिकक्षो}\ A H}}\skt{~। }\pend
\pstart\versepar \skt{चत्वार एते }\edtext{\skt{कणभाः}}{\lemma{\skt{कणभाः}}\Afootnote{A K; \skt{कणभाः}\ H}}\skt{ व्याख्या\-ता\-स्ती\-व्र\-वे\-दनाः~।}\pend

\pstart\versepar \ledsidenote{\passref{5.8.27-verse}}\edtext{\skt{एभिर्दष्टस्य}}{\lemma{\skt{दष्टस्य}}\Afootnote{\skt{दष्टेति}\ K H}}\skt{ गुरुता }\edtext{\skt{गात्रा\-णा\-म\-ङ्ग\-वे\-दना}}{\lemma{\skt{अङ्गवेदना}}\Afootnote{K H\textsuperscript{ac}; \skt{दंशवेदना~॥}\ H\textsuperscript{pc}}}\skt{~।}\skt{ }\pend
\pstart\versepar \skt{लालास्रावश्च भवति गात्रभेदश्च दारुणः~।}\pend

\pstart\versepar \ledsidenote{\passref{5.8.28-verse-1}}\skt{प्रतिसूर्यः पिङ्गभासो }\edtext{\skt{बहुवर्ण्णो}}{\lemma{\skt{बहुवर्ण्णो}}\Afootnote{\skt{वहुवर्ण्णो}\ K H}}\skt{ महाशिराः~। }\pend
\pstart\versepar \skt{तथा निरुपमश्चापि पञ्च गोधेरकाः स्मृताः~।}\pend

\pstart\versepar \ledsidenote{\passref{5.8.28-verse-2}}\skt{तैर्भवन्तीह दष्टानां }\edtext{\skt{वेगज्ञानानि}}{\lemma{\skt{वेगज्ञानानि}}\Afootnote{K H\textsuperscript{pc}; \skt{वेगज्ञानि}\ H\textsuperscript{ac}}}\skt{ सर्पवत्~। }\pend
\pstart\versepar \skt{रुजश्च विविधाकारा ग्रन्थयश्च सुदारुणाः~।}\pend

\pstart\versepar \ledsidenote{\passref{5.8.29-verse-1}}\edtext{\skt{श्वेता}}{\lemma{\skt{श्वेता}}\Afootnote{\skt{स्वेता}\ K H}}\skt{ कृष्णा कृष्णराजी रक्ता रक्तैश्च मण्डलैः~। }\pend
\pstart\versepar \skt{}\edtext{\skt{सर्वश्वेता}}{\lemma{\skt{सर्वश्वेता}}\Afootnote{\skt{सर्वस्वेता}\ K H}}\skt{ सर्षपिका षडेता }\edtext{\skt{गृहगोलिकाः}}{\lemma{\skt{गृहगोलिकाः}}\Afootnote{K; \skt{गृहगोधिकाः~।}\ H}}\skt{~।}\pend

\pstart\versepar \ledsidenote{\passref{5.8.29-verse-2}}\edtext{\skt{ताभिर्दष्टे}}{\lemma{\skt{दष्टे}}\Afootnote{H; \skt{दष्ट}\ K}}\skt{ दंशतोदो हृत्पीडा दाह एव च~। }\pend
\pstart\versepar \skt{दंशशोफश्च भवति ग्रन्थिजन्म च }\edtext{\skt{दारुणम्}}{\lemma{\skt{दारुणम्}}\Afootnote{H; \skt{दारुणंः~।}\ K}}\skt{~।}\pend

\pstart\versepar \ledsidenote{\passref{5.8.30-verse-1}}\skt{परुषा कृष्णचित्रे च कपिला पीतिका तथा~॥ }\pend
\pstart\versepar \skt{रक्ता }\edtext{\skt{श्वेताग्निवर्ण्णा च}}{\lemma{\skt{श्वेताग्निवर्ण्णा च}}\Afootnote{\skt{स्वेताग्निवर्ण्णा च}\ K H}}\skt{ }\edtext{\skt{शतपादोऽष्टधा}}{\lemma{\skt{शतपादो}}\Afootnote{H; \skt{शतपाद्यो}\ K}}\skt{ स्मृताः~।}\pend

\pstart\versepar \ledsidenote{\passref{5.8.30-verse-2}}\edtext{\skt{ताभिर्दष्टे}}{\lemma{\skt{दष्टे}}\Afootnote{K; \skt{द्दष्ट}\ H}}\skt{ रुजास्तीव्रा दंशशोफश्च दारुणः~। }\pend
\pstart\versepar \skt{दंशे च }\edtext{\edtext{\skt{पिट\-को\-त्प\-त्ति\-र्मूर्च्छां}}{\lemma{\skt{पिटकोत्पत्तिर्}}\Afootnote{K; \skt{पिटकोत्पाटिम्}\ H}}}{\lemma{\skt{मूर्च्छां}}\Afootnote{K; \skt{मूर्च्छा}\ H}}\skt{ चापि }\edtext{\skt{सुदारुणाः}}{\lemma{\skt{सुदारुणाः}}\Afootnote{K; \skt{सुदारुणा~॥}\ H}}\skt{~।}\pend

\pstart\versepar \ledsidenote{\passref{5.8.31-verse-1}}\skt{श्वेतश्च कृष्णवर्ण्णश्च }\edtext{\skt{शर\-व\-र्ण्णोऽ\-य\-म\-प्रभः}}{\lemma{\skt{ऽयम्}}\Afootnote{\skt{यम्}\ K H\textsuperscript{ac}; \skt{ग्निम्यम्}\ H\textsuperscript{pc}}}\skt{ }\pend
\pstart\versepar \skt{कुहरो हरितश्चापि भृकुटी कोटिकश्च यः~।}\pend

\pstart\versepar \ledsidenote{\passref{5.8.31-verse-2}}\skt{अष्टावेते }\edtext{\skt{कीटासञ्ज्ञा}}{\lemma{\skt{कीटासञ्ज्ञा}}\Afootnote{K; \skt{कीटसंज्ञा}\ H}}\skt{ दर्दुराः परिकीर्त्तिताः~। }\pend
\pstart\versepar \skt{तैर्दष्टः }\edtext{\skt{कण्डुसंयुक्तो}}{\lemma{\skt{कण्डुसंयुक्तो}}\Afootnote{K; \skt{कण्डूसंयुक्तो}\ H}}\skt{ हरितं मूर्छितो वमेत्~।}\pend

\pstart\versepar \ledsidenote{\passref{5.8.31.add}}\skt{जलौकाः षट्समाख्याताः सल\-क्ष\-ण\-चि\-कि\-त्सिताः~। }\pend
\pstart\versepar \skt{अहिकुत्थुः कुत्थुकश्च वृत्त\-शू\-क\-स्त\-थैव च~॥}\pend

\pstart\versepar \ledsidenote{\passref{5.8.32-verse}}\skt{त्रयो विश्वम्भराः प्रोक्ताः दाह\-ज्व\-र\-रु\-जा\-वहाः~। }\pend
\pstart\versepar \skt{तैर्दष्टमात्रे श्वयथुरादंशे कण्डुरेव च~।}\edtext{}{\lemma{\skt{दाह\appbrk{}ज्व\appbrk{}र\appbrk{}रु\appbrk{}जा\appbrk{}वहाः~। तैर्}}\Afootnote{\skt{दाह\appbrk{}ज्व\appbrk{}र\appbrk{}रु\appbrk{}जा\appbrk{}वहाः~। तै}\ K; \skt{दाह\appbrk{}ज्व\appbrk{}र\appbrk{}रु\appbrk{}जा\appbrk{}पहाः~॥ तैर्}\ H}}\pend

\pstart\versepar \ledsidenote{\passref{5.8.34-verse-1}}\skt{फेना\-ग\-मोऽ\-ति\-सा\-रश्च }\edtext{\skt{कोठजन्मं च}}{\lemma{\skt{कोठजन्मं च}}\Afootnote{\skt{कोठजन्मे च}\ K; \skt{कोठजन्मेव}\ H}}\skt{ दारुणम्~।}\skt{ }\pend
\pstart\versepar \skt{}\edtext{\skt{सम्वाहिका}}{\lemma{\skt{सम्वाहिका}}\Afootnote{\skt{गम्वाहिका}\ K H}}\skt{ स्थूलशीर्षा }\edtext{\skt{ब्राह्मण्यङ्गुलिका}}{\lemma{\skt{ब्राह्मण्य्}}\Afootnote{\skt{व्राह्मण्य्}\ K H}}\skt{ तथा~।}\pend

\pstart\versepar \ledsidenote{\passref{5.8.34-verse-2}}\edtext{\skt{विवर्ण्णा}}{\lemma{\skt{विवर्ण्णा}}\Afootnote{K H\textsuperscript{ac}; \skt{विवर्ण्णा कर्ण्णिला ग्रन्थान्तरे}\ H\textsuperscript{pc}}}\skt{ }\edtext{\skt{कपिला}}{\lemma{\skt{कपिला}}\Afootnote{K\textsuperscript{pc} H; \skt{कपिका}\ K\textsuperscript{ac}}}\skt{ }\edtext{\skt{चापि}}{\lemma{\skt{चापि}}\Afootnote{K\textsuperscript{pc} H; \skt{ला चा}\ K\textsuperscript{ac}}}\skt{ }\edtext{\skt{षट्प्रोक्तास्तु}}{\lemma{\skt{षट्प्रोक्तास्तु}}\Afootnote{K\textsuperscript{pc} H; \emph{om.}\ K\textsuperscript{ac}}}\skt{ }\edtext{\skt{पिपीलिकाः}}{\lemma{\skt{पिपीलिकाः}}\Afootnote{K\textsuperscript{pc} H; \skt{पि~।}\ K\textsuperscript{ac}}}\skt{~। }\pend
\pstart\versepar \skt{ताभिर्दष्टे }\edtext{\skt{रुजा}}{\lemma{\skt{रुजा}}\Afootnote{K; \skt{रु}\ H}}\skt{ }\edtext{\skt{दाहः}}{\lemma{\skt{दाहः}}\Afootnote{K; \emph{om.}\ H}}\skt{ }\edtext{\edtext{\skt{कण्डु\-श्व\-य\-थु\-रेव}}{\lemma{\skt{कण्डुश्वयथुर्}}\Afootnote{K; \emph{om.}\ H}}\skt{ च}}{\lemma{\skt{एव च}}\Afootnote{K; \emph{om.}\ H}}\skt{~। }\pend
\pstart\versepar \skt{}\edtext{\skt{विशेषेण}}{\lemma{\skt{विशेषेण}}\Afootnote{K; \emph{om.}\ H}}\skt{ }\edtext{\edtext{\skt{दंशत्येताः}}{\lemma{\skt{दंशत्य्}}\Afootnote{K; \emph{om.}\ H}}}{\lemma{\skt{एताः}}\Afootnote{K; \emph{om.}\ H}}\skt{ }\edtext{\edtext{\skt{नेत्र\-यो\-र्ने\-त्र\-व\-ल्लभाः}}{\lemma{\skt{नेत्रयोर्}}\Afootnote{K; \emph{om.}\ H}}}{\lemma{\skt{नेत्रवल्लभाः}}\Afootnote{K; \skt{जास्तीव्रा शोफ\appbrk{}श्चा\appbrk{}द\appbrk{}ञ्श\appbrk{}म\appbrk{}ण्डले~।}\ H}}\skt{~।}\pend

\pstart \ledsidenote{\passref{5.8.36-verse-1}}\skt{मण्डलः }\edtext{\skt{पार्वतश्चैव}}{\lemma{\skt{पार्वतश्}}\Afootnote{K; \skt{पर्व्वतश्}\ H}}\skt{ कृष्णः }\edtext{\skt{सामुद्र}}{\lemma{\skt{सामुद्र}}\Afootnote{K; \skt{समुद्र}\ H}}\skt{ एव च~।}\pend

\pstart\versepar \ledsidenote{\passref{5.8.36-verse-2}}\skt{मशको हस्तिनामा च मशकाः पञ्चकीर्त्तिताः~। }\pend
\pstart\versepar \skt{तैर्दष्टे }\edtext{\skt{रोषसंयुक्तं}}{\lemma{\skt{रोषसंयुक्तं}}\Afootnote{H; \skt{रोषसंयुक्तां}\ K}}\skt{ शून\-मा\-दं\-श\-म\-ण्ड\-लम्~।}\pend

\pstart \ledsidenote{\passref{5.8.36-verse-3}}\skt{वेदना }\edtext{\skt{रागबहुलं}}{\lemma{\skt{रागबहुलं}}\Afootnote{\skt{रागवहुलं}\ K H}}\skt{ कण्डूयुक्तं क्षरत्यसृक्~।}\pend

\pstart\versepar \ledsidenote{\passref{5.8.38}}\edtext{\skt{गोधेरकः}}{\lemma{\skt{गोधेरकः}}\Afootnote{K H; \skt{भवन्ति चात्र गोधेरकः}\ A}}\skt{ }\edtext{\skt{स्थालका च ये च}}{\lemma{\skt{स्थालका च ये च}}\Afootnote{K H; \skt{स्थालिका च ये च}\ A}}\skt{ श्वेताग्निसप्रभे~। }\pend
\pstart\versepar \skt{भृकुटी }\edtext{\skt{कोटिकश्चैव न}}{\lemma{\skt{चैव न}}\Afootnote{A K; \skt{चापि न}\ H}}\skt{ सिध्य\-न्त्ये\-क\-जा\-तिषु~।}\pend

\pstart\versepar \ledsidenote{\passref{5.8.42}}\edtext{\skt{कीटै\-र्द\-ष्टा\-नु\-ग्र\-विषैः}}{\lemma{\skt{कीटैर्}}\Afootnote{K H; \emph{om.}\ A}}\skt{ सर्प\-व\-त्स\-मु\-पा\-च\-रेत्~। }\pend
\pstart\versepar \skt{त्रिविधानान्तु }\edtext{\skt{शेषाणां}}{\lemma{\skt{शेषाणां}}\Afootnote{K H; \skt{पूर्वेषां}\ A}}\skt{ }\edtext{\skt{त्रैविध्यं}}{\lemma{\skt{त्रैविध्यं}}\Afootnote{K H; \skt{त्रैविध्ये}\ A}}\skt{ }\edtext{\skt{भवति}}{\lemma{\skt{भवति}}\Afootnote{K H; \skt{न}\ A}}\skt{ }\edtext{\skt{क्रिया}}{\lemma{\skt{क्रिया}}\Afootnote{K H; \skt{क्रिया हिताः~॥}\ A}}\skt{~।}\edtext{}{\lemma{\skt{सर्प\appbrk{}व\appbrk{}त्स\appbrk{}मु\appbrk{}पा\appbrk{}च\appbrk{}रेत्~। त्रिविधानान्तु}}\Afootnote{K H; \skt{कीटैः सर्प\appbrk{}व\appbrk{}त्स\appbrk{}मु\appbrk{}पा\appbrk{}च\appbrk{}रेत्~। त्रिविधानां तु}\ A}}\pend

\pstart \ledsidenote{\passref{5.8.43}}\edtext{\skt{स्वेदां}}{\lemma{\skt{स्वेदां}}\Afootnote{K; \skt{स्वेदमालेपनं सेकं}\ A; \skt{स्वेदा\appbrk{}न्व\appbrk{}हु\appbrk{}प्र\appbrk{}का\appbrk{}राश्}\ H}}\skt{ }\edtext{\edtext{\skt{बहुप्रकारांश्च}}{\lemma{\skt{बहुप्रकारांश्}}\Afootnote{\skt{चोष्णमत्राव}\ A; \skt{वहुप्रकारांश्}\ K; \emph{om.}\ H}}}{\lemma{\skt{च}}\Afootnote{K H; \skt{चा}\ A}}\skt{ }\edtext{\skt{युञ्ज्यादन्यत्र}}{\lemma{\skt{युञ्ज्यादन्यत्र}}\Afootnote{K; \skt{रयेत्~। अन्यत्र}\ A; \skt{युञ्ज्यावन्यत्र}\ H}}\skt{ }\edtext{\skt{मूर्च्छितात्}}{\lemma{\skt{मूर्च्छितात्}}\Afootnote{K; \skt{मूर्च्छि\appbrk{}ता\appbrk{}द्दं\appbrk{}शा\appbrk{}त्पा\appbrk{}क\appbrk{}को\appbrk{}थ\appbrk{}प्र\appbrk{}पी\appbrk{}डि\appbrk{}तात्~॥}\ A; \skt{मूर्च्छितान्~।}\ H}}\skt{~।}\pend

\pstart \ledsidenote{\passref{5.8.44ab}}\skt{विषघ्नञ्च विधिं }\edtext{\skt{कुर्या\-त्कु\-र्या\-त्सं\-शो\-ध\-नानि}}{\lemma{\skt{कुर्या\appbrk{}त्कु\appbrk{}र्या\appbrk{}त्सं\appbrk{}शो\appbrk{}ध\appbrk{}नानि}}\Afootnote{H; \skt{सर्वं बहुशः शोधनानि}\ A; \skt{कुर्या\appbrk{}त्सं\appbrk{}शो\appbrk{}ध\appbrk{}नानि}\ K}}\skt{ च~।}\pend

\pstart \ledsidenote{\passref{5.8.56ef}}\skt{त्रिविधा वृश्चिकाः प्रोक्ता मन्द\-म\-ध्य\-म\-हा\-विषाः~।}\pend

\pstart \ledsidenote{\passref{5.8.57cd}}\skt{सर्प\-को\-थो\-द्भ\-वा\-स्तीक्ष्णा }\edtext{\skt{दिग्धदष्टं}}{\lemma{\skt{दिग्धदष्टं}}\Afootnote{K; \emph{om.}\ A; \skt{दिग्धदंष्ट्र}\ H}}\skt{ }\edtext{\edtext{\skt{विषैर्हते}}{\lemma{\skt{विषैर्}}\Afootnote{K H; \emph{om.}\ A}}}{\lemma{\skt{हते}}\Afootnote{K H; \skt{ये चान्ये विषसंभवाः~॥}\ A}}\skt{~।}\pend

\pstart\versepar \ledsidenote{\passref{5.8.58}}\edtext{\skt{कोथे}}{\lemma{\skt{कोथे}}\Afootnote{K H; \skt{मन्दा द्वादश}\ A}}\skt{ }\edtext{\skt{मध्या}}{\lemma{\skt{मध्या}}\Afootnote{\skt{मध्यास्}\ A; \skt{मध्ये}\ K H}}\skt{ }\edtext{\skt{गवादीनां}}{\lemma{\skt{गवादीनां}}\Afootnote{K H; \skt{तु त्रयः पञ्}\ A}}\skt{ }\edtext{\skt{शकृत्कोथे}}{\lemma{\skt{शकृत्कोथे}}\Afootnote{K H; \skt{चदशोत्तमाः~।}\ A}}\skt{ }\edtext{\skt{वराः}}{\lemma{\skt{वराः}}\Afootnote{K H; \emph{om.}\ A}}\skt{ }\edtext{\skt{स्मृताः}}{\lemma{\skt{स्मृताः}}\Afootnote{K H; \emph{om.}\ A}}\skt{~।}\skt{ }\pend
\pstart\versepar \skt{}\edtext{\edtext{\skt{सप्त\-विं\-श\-ति\-रे\-वैते}}{\lemma{\skt{सप्तविंशतिर्}}\Afootnote{H; \skt{दश विंशतिर्}\ A; \skt{सप्तविंशति}\ K}}}{\lemma{\skt{एवैते}}\Afootnote{H; \skt{इत्येते}\ A; \skt{रेवैते}\ K}}\skt{ }\edtext{\skt{सङ्ख्यया}}{\lemma{\skt{सङ्ख्यया}}\Afootnote{K; \skt{संख्यया}\ A; \skt{संख्यायाः}\ H}}\skt{ परिकीर्त्तिताः~।}\pend

\pstart\versepar \ledsidenote{\passref{5.8.59}}\edtext{\skt{कृष्णः}}{\lemma{\skt{कृष्णः}}\Afootnote{A H; \skt{कृष्ण}\ K}}\skt{ }\edtext{\skt{श्यावः}}{\lemma{\skt{श्यावः}}\Afootnote{A H; \skt{स्यावः}\ K}}\skt{ }\edtext{\skt{कर्बुरो}}{\lemma{\skt{कर्बुरो}}\Afootnote{A; \skt{कर्वुरो}\ K H}}\skt{ }\edtext{\skt{रोमशश्च}}{\lemma{\skt{रोमशश्च}}\Afootnote{K H; \skt{पाण्डुवर्णो}\ A}}\skt{ }\pend
\pstart\versepar \skt{गोमूत्राभः }\edtext{\skt{परुषो}}{\lemma{\skt{परुषो}}\Afootnote{K H; \skt{कर्कशो}\ A}}\skt{ }\edtext{\skt{मेचकश्च}}{\lemma{\skt{मेचकश्च}}\Afootnote{A; \skt{दकश्च~।}\ K; \skt{मोदकश्च~।}\ H}}\skt{~।}\skt{ }\pend
\pstart\versepar \skt{}\edtext{\skt{श्वेतो}}{\lemma{\skt{श्वेतो}}\Afootnote{K H; \skt{पीतो}\ A}}\skt{ }\edtext{\skt{रक्तो}}{\lemma{\skt{रक्तो}}\Afootnote{K H; \skt{धूम्रो}\ A}}\skt{ }\edtext{\skt{रोम\-शी\-र्षो\-ग्र\-धूमः}}{\lemma{\skt{रोम\appbrk{}शी\appbrk{}र्षो\appbrk{}ग्र\appbrk{}धूमः}}\Afootnote{\skt{रोमशः शाड्वलाभो}\ A; \skt{रोम\appbrk{}शी\appbrk{}र्षो\appbrk{}ग्र\appbrk{}धूम्रः}\ K H}}\skt{ }\pend
\pstart\versepar \skt{}\edtext{\edtext{\edtext{\skt{सर्वेऽप्येते}}{\lemma{\skt{सर्वे}}\Afootnote{K H; \skt{रक्तः श्वेते}\ A}}}{\lemma{\skt{ऽप्य्}}\Afootnote{\skt{नोदरे}\ A; \skt{प्य्}\ K H}}}{\lemma{\skt{एते}}\Afootnote{K H; \skt{णेति}\ A}}\skt{ }\edtext{\skt{मन्दविषा}}{\lemma{\skt{मन्दविषा}}\Afootnote{K H; \skt{मन्दाः~॥}\ A}}\skt{ }\edtext{\edtext{\skt{मतास्तु}}{\lemma{\skt{मतास्}}\Afootnote{K H; \emph{om.}\ A}}}{\lemma{\skt{तु}}\Afootnote{K H; \emph{om.}\ A}}\skt{~।}\pend

\pstart\versepar \ledsidenote{\passref{5.8.60cd}}\skt{एभिर्दष्टे वेदना }\edtext{\skt{वेपथुश्च}}{\lemma{\skt{वेपथुश्च}}\Afootnote{A H; \skt{वेथुश्च}\ K}}\skt{ }\pend
\pstart\versepar \skt{}\edtext{\skt{गात्रस्तब्धः}}{\lemma{\skt{गात्रस्तब्धः}}\Afootnote{\skt{गात्रस्तम्भः}\ A; \skt{गात्रस्तव्धः}\ K H}}\skt{ कृष्ण\-र\-क्ता\-ग\-मश्च~।}\pend

\pstart\versepar \ledsidenote{\passref{5.8.61}}\edtext{\skt{शाखाविद्धे}}{\lemma{\skt{शाखाविद्धे}}\Afootnote{K H; \skt{शाखादष्टे}\ A}}\skt{ }\edtext{\edtext{\skt{वेद\-ना\-ञ्चो\-र्ध्व\-मेति}}{\lemma{\skt{वेदनाञ्}}\Afootnote{K H; \skt{वेदना}\ A}}}{\lemma{\skt{चोर्ध्वम्}}\Afootnote{A K; \skt{चोर्द्धम्}\ H}}\skt{ }\pend
\pstart\versepar \skt{}\edtext{\skt{दंशस्वेदो}}{\lemma{\skt{दंशस्वेदो}}\Afootnote{K; \skt{दाहस्वेदौ}\ A; \skt{दंशस्वेदो}\ H\textsuperscript{ac}; \skt{दंशस्तीव्रास्वेदो}\ H\textsuperscript{pc}}}\skt{ }\edtext{\edtext{\skt{मुखशोफश्च}}{\lemma{\skt{मुखशोफश्}}\Afootnote{K H\textsuperscript{pc}; \skt{दंशशोफो}\ A; \skt{मुखशोफ}\ H\textsuperscript{ac}}}\skt{ तीव्रः}}{\lemma{\skt{च तीव्रः}}\Afootnote{K H\textsuperscript{pc}; \skt{ज्वरश्च~।}\ A; \skt{तीव्रः~॥}\ H\textsuperscript{ac}}}\skt{~। }\pend
\pstart\versepar \skt{}\edtext{\skt{रक्तं}}{\lemma{\skt{रक्तं}}\Afootnote{K H; \skt{रक्तः}\ A}}\skt{ }\edtext{\skt{पीतं}}{\lemma{\skt{पीतं}}\Afootnote{K H; \skt{पीतः}\ A}}\skt{ }\edtext{\skt{कपिलं}}{\lemma{\skt{कपिलं}}\Afootnote{K H; \skt{कापिले}\ A}}\skt{ }\edtext{\skt{चोदरं तु}}{\lemma{\skt{चोदरं तु}}\Afootnote{K; \skt{नोदरेण}\ A; \skt{चोदरञ्}\ H}}\skt{ }\pend
\pstart\versepar \skt{}\edtext{\skt{धूम्रो}}{\lemma{\skt{धूम्रो}}\Afootnote{K; \skt{सर्वे धूम्राः}\ A; \skt{च धूम्रो}\ H}}\skt{ }\edtext{\edtext{\skt{वर्ण्णस्तत्र}}{\lemma{\skt{वर्ण्णस्}}\Afootnote{K H; \skt{पर्वभिश्}\ A}}\skt{ यो}}{\lemma{\skt{तत्र यो}}\Afootnote{K H; \skt{च त्रिभिः}\ A}}\skt{ }\edtext{\skt{मध्यवीर्याः}}{\lemma{\skt{मध्यवीर्याः}}\Afootnote{K H; \skt{स्युः~॥}\ A}}\skt{~।}\pend

\pstart\versepar \ledsidenote{\passref{5.8.63}}\skt{जिह्वाशोफो }\edtext{\skt{रस\-न\-स्यो\-प\-घातो}}{\lemma{\skt{रस\appbrk{}न\appbrk{}स्यो\appbrk{}प\appbrk{}घातो}}\Afootnote{K; \skt{भोज\appbrk{}न\appbrk{}स्या\appbrk{}व\appbrk{}रोधो}\ A; \skt{रस\appbrk{}न\appbrk{}स्या\appbrk{}प\appbrk{}घातो}\ H}}\skt{ }\pend
\pstart\versepar \skt{मूर्छा चोग्रा }\edtext{\skt{मध्य\-वि\-षा\-भि\-दष्टे}}{\lemma{\skt{मध्य\appbrk{}वि\appbrk{}षा\appbrk{}भि\appbrk{}दष्टे}}\Afootnote{K H; \skt{मध्य\appbrk{}वी\appbrk{}र्या\appbrk{}भि\appbrk{}दष्टे~।}\ A}}\skt{~। }\pend
\pstart\versepar \skt{श्वेतश्चित्रः }\edtext{\skt{शबलो}}{\lemma{\skt{शबलो}}\Afootnote{\skt{श्यामलो}\ A; \skt{शवलो}\ K H}}\skt{ लोहिताभः }\pend
\pstart\versepar \skt{}\edtext{\skt{कृष्णः}}{\lemma{\skt{कृष्णः}}\Afootnote{K H; \skt{रक्तः}\ A}}\skt{ }\edtext{\skt{श्यावः}}{\lemma{\skt{श्यावः}}\Afootnote{K H; \skt{श्वेतो}\ A}}\skt{ }\edtext{\skt{श्वेतनीलोदरौ}}{\lemma{\skt{श्वेतनीलोदरौ}}\Afootnote{K H; \skt{रक्तनीलोदरौ}\ A}}\skt{ च~।}\pend

\pstart\versepar \ledsidenote{\passref{5.8.64cd}}\skt{रक्तो }\edtext{\skt{बभ्रुः}}{\lemma{\skt{बभ्रुः}}\Afootnote{A; \skt{वभ्रु}\ K; \skt{वभ्रुः}\ H}}\skt{ }\edtext{\skt{पूर्व\-व\-दे\-क\-पर्वा}}{\lemma{\skt{पूर्व\appbrk{}व\appbrk{}दे\appbrk{}क\appbrk{}पर्वा}}\Afootnote{K H; \skt{पूर्व\appbrk{}व\appbrk{}च्चै\appbrk{}क\appbrk{}पर्वा यश्}\ A}}\skt{~। }\pend
\pstart\versepar \skt{}\edtext{\skt{पूर्वा}}{\lemma{\skt{पूर्वा}}\Afootnote{K H; \skt{चापर्वा}\ A}}\skt{ }\edtext{\skt{चापि}}{\lemma{\skt{चापि}}\Afootnote{K H; \emph{om.}\ A}}\skt{ पर्वणी द्वे च यस्य~।}\pend

\pstart\versepar \ledsidenote{\passref{5.8.65ab}}\edtext{\skt{नानावर्ण्णा}}{\lemma{\skt{नानावर्ण्णा}}\Afootnote{K H; \skt{नानारूपा}\ A}}\skt{ }\edtext{\skt{रूपतश्चापि}}{\lemma{\skt{रूपतश्}}\Afootnote{K H; \skt{वर्णतश्}\ A}}\skt{ घोराः~। }\pend
\pstart\versepar \skt{}\edtext{\skt{ज्ञेया}}{\lemma{\skt{ज्ञेया}}\Afootnote{K H; \skt{ज्ञेयाश्}\ A}}\skt{ }\edtext{\edtext{\skt{ह्येते}}{\lemma{\skt{ह्य्}}\Afootnote{K H; \skt{चै}\ A}}}{\lemma{\skt{एते}}\Afootnote{K H; \skt{ते}\ A}}\skt{ वृश्चिकाः }\edtext{\skt{प्राणनाशाः}}{\lemma{\skt{प्राणनाशाः}}\Afootnote{K H; \skt{प्राणचौराः~।}\ A}}\skt{~॥}\pend

\pstart\versepar \ledsidenote{\passref{5.8.66}}\skt{एभिर्दष्टे }\edtext{\skt{विष\-वे\-ग\-प्र\-वृत्तिः}}{\lemma{\skt{विष\appbrk{}वे\appbrk{}ग\appbrk{}प्र\appbrk{}वृत्तिः}}\Afootnote{K H; \skt{सर्प\appbrk{}वे\appbrk{}ग\appbrk{}प्र\appbrk{}वृत्तिः}\ A}}\skt{ }\pend
\pstart\versepar \skt{}\edtext{\skt{स्फोटो\-त्प\-त्ति\-र्ज्व\-र\-दाहौ}}{\lemma{\skt{ज्वरदाहौ}}\Afootnote{K H; \skt{भ्रान्तिदाहौ}\ A}}\skt{ }\edtext{\skt{भ्रमश्च}}{\lemma{\skt{भ्रमश्च}}\Afootnote{K H; \skt{ज्वरश्च~।}\ A}}\skt{~। }\pend
\pstart\versepar \skt{खेभ्यः कृष्णं शोणितं }\edtext{\skt{चातितीव्रं}}{\lemma{\skt{चातितीव्रं}}\Afootnote{H; \skt{याति तीव्रं}\ A; \skt{चतितीव्रं}\ K}}\skt{ }\pend
\pstart\versepar \skt{}\edtext{\skt{ततः}}{\lemma{\skt{ततः}}\Afootnote{K H\textsuperscript{pc}; \skt{तस्मा}\ A; \skt{स्रावः}\ H\textsuperscript{ac}}}\skt{ }\edtext{\edtext{\skt{प्राणैस्त्याज्यते}}{\lemma{\skt{प्राणैस्}}\Afootnote{K H; \skt{त्प्राणैस्}\ A}}}{\lemma{\skt{त्याज्यते}}\Afootnote{K H; \skt{त्यज्यते}\ A}}\skt{ }\edtext{\skt{क्षिप्रमेव}}{\lemma{\skt{क्षिप्रम्}}\Afootnote{K H; \skt{शीघ्रम्}\ A}}\skt{~।}\pend

\pstart \ledsidenote{\passref{5.8.67ab}}\edtext{\skt{उग्र\-म\-ध्य\-वि\-षै\-र्द\-ष्टां\-श्चि\-कि\-त्से\-त्स\-र्प\-द\-ष्ट\-वत्}}{\lemma{\skt{दष्टांश्}}\Afootnote{K H; \skt{दष्टं}\ A}}\skt{~।}\pend

\pstart\versepar \ledsidenote{\passref{5.8.70}}\edtext{\skt{दंश\-म\-न्द\-वि\-षा\-णान्तु}}{\lemma{\skt{दंश\appbrk{}म\appbrk{}न्द\appbrk{}वि\appbrk{}षा\appbrk{}णान्तु}}\Afootnote{K; \skt{दंशं मन्दविषाणां तु}\ A; \skt{दंशो मन्दविषाणान्तु}\ H\textsuperscript{ac}; \skt{दंशो मध्यविषाणान्तु}\ H\textsuperscript{pc}}}\skt{ चक्रतैलेन सेचयेत्~। }\pend
\pstart\versepar \skt{विदा\-रि\-ग\-न्धा\-दि\-तै\-लेन सुखोष्णेनाथवा पुनः~।}\edtext{}{\lemma{\skt{सेचयेत्~। विदा\appbrk{}रि\appbrk{}ग\appbrk{}न्धा\appbrk{}दि\appbrk{}तै\appbrk{}लेन}}\Afootnote{K; \skt{सेचयेत्~। विदा\appbrk{}री\appbrk{}ग\appbrk{}ण\appbrk{}सि\appbrk{}द्धेन}\ A; \skt{सेचयेत्~॥ विदा\appbrk{}रि\appbrk{}गं\appbrk{}धा\appbrk{}दि\appbrk{}तै\appbrk{}लेन}\ H}}\pend

\pstart \ledsidenote{\passref{5.8.67cd}}\edtext{\skt{आदंशं}}{\lemma{\skt{आदंशं}}\Afootnote{A K\textsuperscript{pc} H; \skt{आदं}\ K\textsuperscript{ac}}}\skt{ स्वेदितं }\edtext{\skt{चूर्ण्णैः}}{\lemma{\skt{चूर्ण्णैः}}\Afootnote{A H; \skt{चूर्ण्णैंः}\ K}}\skt{ प्रच्छि\-त\-म्प्र\-ति\-सा\-र\-येत्~।}\pend

\pstart\versepar \ledsidenote{\passref{5.8.68}}\skt{रज\-नी\-सै\-न्ध\-व\-व्यो\-ष\-शि\-री\-ष\-फ\-ल\-पु\-ष्पजैः~। }\pend
\pstart\versepar \skt{मातु\-लुं\-गा\-म्ल\-गो\-मू\-त्र\-पि\-ष्टञ्च सुरसाग्रजं~।}\pend

\pstart\versepar \ledsidenote{\passref{5.8.69}}\edtext{\skt{लेपे}}{\lemma{\skt{लेपे}}\Afootnote{K H; \skt{लेपे स्वेदे}\ A}}\skt{ सुखोष्णञ्च }\edtext{\skt{तथा}}{\lemma{\skt{तथा}}\Afootnote{K H; \emph{om.}\ A}}\skt{ गोमयं }\edtext{\skt{हितमुच्यते}}{\lemma{\skt{उच्यते}}\Afootnote{K H; \skt{इष्यते~।}\ A}}\skt{~॥ }\pend
\pstart\versepar \skt{}\edtext{\skt{सर्पिः}}{\lemma{\skt{सर्पिः}}\Afootnote{K H; \skt{पाने}\ A}}\skt{ क्षौद्रयुतं }\edtext{\skt{पाने}}{\lemma{\skt{पाने}}\Afootnote{K H; \skt{सर्पिः}\ A}}\skt{ क्षीरम्वा }\edtext{\skt{वहुशर्करं}}{\lemma{\skt{वहुशर्करं}}\Afootnote{K H; \skt{बहुशर्करम्~॥}\ A}}\skt{~।}\pend

\pstart \ledsidenote{\passref{5.8.71cd}}\skt{गुडोदकं वा }\edtext{\skt{सुहितं}}{\lemma{\skt{सुहितं}}\Afootnote{K H; \skt{सुहिमं}\ A}}\skt{ }\edtext{\skt{चतु\-र्जा\-त\-क\-वा\-सितं}}{\lemma{\skt{चतु\appbrk{}र्जा\appbrk{}त\appbrk{}क\appbrk{}वा\appbrk{}सितं}}\Afootnote{H; \skt{चातु\appbrk{}र्जा\appbrk{}त\appbrk{}क\appbrk{}सं\appbrk{}यु\appbrk{}तम्~॥}\ A; \skt{चतु\appbrk{}र्जा\appbrk{}त\appbrk{}वा\appbrk{}सितं~।}\ K}}\skt{~।}\pend

\pstart \ledsidenote{\passref{5.8.72cd}}\edtext{\skt{शिखि\-कु\-क्कु\-ट\-ब\-र्हाणि}}{\lemma{\skt{शिखि\appbrk{}कु\appbrk{}क्कु\appbrk{}ट\appbrk{}ब\appbrk{}र्हाणि}}\Afootnote{A; \skt{शिखि\appbrk{}कु\appbrk{}क्कु\appbrk{}ट\appbrk{}व\appbrk{}र्हाणि}\ K H}}\skt{ सैन्धवं }\edtext{\edtext{\skt{तैलमेव}}{\lemma{\skt{तैलम्}}\Afootnote{K H; \skt{तै}\ A}}}{\lemma{\skt{एव}}\Afootnote{K H; \skt{ल}\ A}}\skt{ }\edtext{\skt{च}}{\lemma{\skt{च}}\Afootnote{K H; \skt{सर्पिषी~॥}\ A}}\skt{~।}\pend

\pstart \ledsidenote{\passref{5.8.71ab}}\skt{कुर्या\-च्चो\-त्का\-रि\-का\-स्वेदं }\edtext{\skt{विष\-घ्नै\-रु\-प\-ना\-हनैः}}{\lemma{\skt{उपनाहनैः}}\Afootnote{K H; \skt{उपनाहयेत्~।}\ A}}\skt{~।}\pend

\pstart\versepar \ledsidenote{\passref{5.8.73}}\edtext{\skt{धूपो}}{\lemma{\skt{धूपो}}\Afootnote{K H\textsuperscript{ac}; \skt{धूमो}\ A; \skt{धूपोमो}\ H\textsuperscript{pc}}}\skt{ हन्ति }\edtext{\skt{प्रयुक्तोऽयं}}{\lemma{\skt{ऽयं}}\Afootnote{\skt{तु}\ A; \skt{यं}\ K H}}\skt{ शीघ्रं वृश्चिकजं विषम्~।}\skt{ }\pend
\pstart\versepar \skt{}\edtext{\skt{कुसुंभपुष्पं}}{\lemma{\skt{कुसुंभपुष्पं}}\Afootnote{K; \skt{कुसुम्भपुष्पं}\ A; \skt{कुरुम्भपुष्पं}\ H}}\skt{ रजनी }\edtext{\skt{निष्या वा}}{\lemma{\skt{निष्या वा}}\Afootnote{K H; \skt{निशा वा}\ A}}\skt{ }\edtext{\skt{क्षौद्रकं}}{\lemma{\skt{क्षौद्रकं}}\Afootnote{K H; \skt{कोद्रवं}\ A}}\skt{ तृणम्~।}\pend

\pstart\versepar \ledsidenote{\passref{5.8.74}}\skt{एभि\-र्घृ\-ता\-क्तै\-र्धू\-पस्तु }\edtext{\skt{पायु\-दे\-श\-प्र\-यो\-जितः}}{\lemma{\skt{पायु\appbrk{}दे\appbrk{}श\appbrk{}प्र\appbrk{}यो\appbrk{}जितः}}\Afootnote{K H; \skt{पायुदेशे प्रयोजितः~।}\ A}}\skt{~॥ }\pend
\pstart\versepar \skt{नाशयेदाशु कीटोत्थं }\edtext{\skt{वृश्चिकस्य च}}{\lemma{\skt{वृश्चिकस्य च}}\Afootnote{A; \skt{वृश्चिकस्या च}\ K H}}\skt{ यद्विषम्~।}\pend

\pstart\versepar \ledsidenote{\passref{5.8.75}}\skt{लूताविषं घोरतमं दुर्वि\-ज्ञे\-य\-त\-मञ्च यत्~। }\pend
\pstart\versepar \skt{दुश्चि\-कि\-त्स्य\-त\-म\-ञ्चापि }\edtext{\edtext{\skt{भिष\-ग्भि\-र्म\-न्द\-बु\-द्धिभिः}}{\lemma{\skt{भिषग्भिर्}}\Afootnote{A K; \skt{भिषग्भिम्}\ H}}}{\lemma{\skt{मन्दबुद्धिभिः}}\Afootnote{A; \skt{मन्दवुद्धिभिः}\ K H}}\skt{~।}\edtext{}{\lemma{\skt{यत्~। दुश्चि\appbrk{}कि\appbrk{}त्स्य\appbrk{}त\appbrk{}मञ्}}\Afootnote{K H; \skt{तत्~। दुश्चिकित्स्यतमं}\ A}}\pend

\pstart\versepar \ledsidenote{\passref{5.8.76}}\skt{सविषं }\edtext{\edtext{\skt{निर्वि\-ष\-ञ्चे\-द\-मि\-त्येवं}}{\lemma{\skt{चेदम्}}\Afootnote{K H; \skt{चैत}\ A}}}{\lemma{\skt{इत्य्}}\Afootnote{K H; \skt{दित्य्}\ A}}\skt{ }\edtext{\skt{सविशंकिते}}{\lemma{\skt{सविशंकिते}}\Afootnote{K H; \skt{परिशङ्किते~।}\ A}}\skt{~। }\pend
\pstart\versepar \skt{विषघ्नमेव कर्त्त\-व्य\-म\-वि\-रोधि यदौषधम्~।}\pend

\pstart\versepar \ledsidenote{\passref{5.8.77}}\skt{अगदानां हि संयोगो }\edtext{\skt{विषदुष्टस्य}}{\lemma{\skt{विषदुष्टस्य}}\Afootnote{K; \skt{विषजुष्टस्य}\ A; \skt{विषयुक्तस्य}\ H}}\skt{ }\edtext{\skt{युज्यते}}{\lemma{\skt{युज्यते}}\Afootnote{A K; \skt{युंज्यते~॥}\ H}}\skt{~। }\pend
\pstart\versepar \skt{निर्विषे मानवे }\edtext{\skt{युक्तोऽगदः}}{\lemma{\skt{ऽगदः}}\Afootnote{A; \skt{गदः}\ K H}}\skt{ संपद्यते }\edtext{\skt{गदः}}{\lemma{\skt{गदः}}\Afootnote{K; \skt{ऽसुखम्~॥}\ A; \skt{गतः~।}\ H}}\skt{~।}\pend

\pstart\versepar \ledsidenote{\passref{5.8.78}}\skt{तस्मा\-त्स\-र्व\-प्र\-य\-त्नेन ज्ञातव्यो विषनिश्चयः~॥ }\pend
\pstart\versepar \skt{अज्ञात्वा }\edtext{\skt{विषसद्भावं}}{\lemma{\skt{विषसद्भावं}}\Afootnote{A; \skt{द्विषसद्भावं}\ K H}}\skt{ भिष\-ग्व्या\-पा\-द\-ये\-न्न\-रम्~॥}\pend

\pstart\versepar \ledsidenote{\passref{5.8.79}}\edtext{\skt{यद्वत्प्रसूतेन}}{\lemma{\skt{यद्वत्प्रसूतेन}}\Afootnote{K H; \skt{प्रोद्भिद्यमानस्तु य}\ A}}\skt{ नवांकुरेण }\pend
\pstart\versepar \skt{न }\edtext{\skt{व्यक्तजाति}}{\lemma{\skt{व्यक्तजाति}}\Afootnote{K H; \skt{व्यक्तजातिः}\ A}}\skt{ }\edtext{\skt{प्रतिभाति}}{\lemma{\skt{प्रतिभाति}}\Afootnote{K H; \skt{प्रविभाति}\ A}}\skt{ वृक्षः~।}\skt{ }\pend
\pstart\versepar \skt{तद्व\-द्दु\-रा\-ल\-क्ष्य\-तमं हि तासां }\pend
\pstart\versepar \skt{विषं }\edtext{\skt{शरीरे}}{\lemma{\skt{शरीरे}}\Afootnote{A H; \skt{र्शरीरे}\ K}}\skt{ प्रविकीर्ण्णमात्रं~।}\edtext{}{\lemma{\skt{नवांकुरेण न}}\Afootnote{K H; \skt{थाङ्कुरेण न}\ A}}\pend

\pstart\versepar \ledsidenote{\passref{5.8.80}}\edtext{\skt{ईषत्सकण्डू}}{\lemma{\skt{ईषत्सकण्डू}}\Afootnote{K H; \skt{ईषत्सकण्डु}\ A}}\skt{ }\edtext{\skt{प्रचलञ्च}}{\lemma{\skt{च}}\Afootnote{K H; \skt{स}\ A}}\skt{ कोठम् }\pend
\pstart\versepar \skt{अव्यक्तवर्ण्णं }\edtext{\skt{प्रथमेऽहनि}}{\lemma{\skt{ऽहनि}}\Afootnote{A H; \skt{हनि}\ K}}\skt{ स्यात्~। }\pend
\pstart\versepar \skt{अन्तेषु शूनं }\edtext{\skt{परि\-नि\-म्न\-म\-ध्यम्}}{\lemma{\skt{परि\appbrk{}नि\appbrk{}म्न\appbrk{}म\appbrk{}ध्यम्}}\Afootnote{A K; \skt{परिनिमध्यम्}\ H}}\skt{ }\pend
\pstart\versepar \skt{}\edtext{\skt{प्रव्यक्तवर्ण्णञ्च}}{\lemma{\skt{प्रव्यक्तवर्ण्णञ्च}}\Afootnote{K; \skt{प्रव्यक्तरूपं च}\ A; \skt{अव्यक्तवर्ण्णञ्च}\ H}}\skt{ दिने द्वितीये~।}\pend

\pstart\versepar \ledsidenote{\passref{5.8.81}}\skt{त्र्यहेण तद्दर्शयतीह }\edtext{\skt{दंशं}}{\lemma{\skt{दंशं}}\Afootnote{K H; \skt{रूपं}\ A}}\skt{ }\pend
\pstart\versepar \skt{विषं }\edtext{\skt{चतुर्थेऽहनि}}{\lemma{\skt{ऽहनि}}\Afootnote{A H; \skt{हनि}\ K}}\skt{ }\edtext{\edtext{\skt{कोपमेति}}{\lemma{\skt{कोपम्}}\Afootnote{A H; \skt{कोप}\ K}}}{\lemma{\skt{एति}}\Afootnote{A H; \skt{मेति~।}\ K}}\skt{~। }\pend
\pstart\versepar \skt{}\edtext{\edtext{\skt{अतोऽधिकेऽह्नि}}{\lemma{\skt{ऽधिके}}\Afootnote{A H; \skt{धिके}\ K}}}{\lemma{\skt{ऽह्नि}}\Afootnote{A; \skt{ह्नि}\ K H}}\skt{ प्रकरोति }\edtext{\skt{जन्तो}}{\lemma{\skt{जन्तो}}\Afootnote{K; \skt{जन्तोर्}\ A H}}\skt{ }\pend
\pstart\versepar \skt{विष\-प्र\-को\-प\-प्र\-भ\-वा\-न्वि\-का\-रान्~।}\pend

\pstart\versepar \ledsidenote{\passref{5.8.82}}\skt{षष्ठे दिने विप्रसृतन्तु सर्वान् }\pend
\pstart\versepar \skt{मर्म\-प्र\-दे\-शा\-न्भृ\-श\-मा\-वृ\-णोति~। }\pend
\pstart\versepar \skt{}\edtext{\edtext{\skt{तत्स\-प्त\-मेऽ\-त्य\-र्थ\-प\-री\-त\-गात्रं}}{\lemma{\skt{तत्सप्तमे}}\Afootnote{A H; \skt{तत्समे}\ K}}}{\lemma{\skt{ऽत्य\appbrk{}र्थ\appbrk{}प\appbrk{}री\appbrk{}त\appbrk{}गात्रं}}\Afootnote{A; \skt{त्यर्थ\appbrk{}प\appbrk{}री\appbrk{}त\appbrk{}गात्रं}\ K H}}\skt{ }\pend
\pstart\versepar \skt{}\edtext{\skt{व्यापा\-द\-ये\-त्म\-र्त्य\-म\-ति\-प्र\-वृद्धं}}{\lemma{\skt{व्यापा\appbrk{}द\appbrk{}ये\appbrk{}त्म\appbrk{}र्त्यम्}}\Afootnote{K H; \skt{व्यापा\appbrk{}द\appbrk{}ये\appbrk{}न्म\appbrk{}र्त्यम्}\ A}}\skt{~।}\pend

\pstart\versepar \ledsidenote{\passref{5.8.83}}\edtext{\edtext{\skt{यास्ती\-क्ष्ण\-च\-ण्डो\-ग्र\-विषा}}{\lemma{\skt{यास्}}\Afootnote{A; \skt{या}\ K H}}\skt{ हि}}{\lemma{\skt{तीक्ष्ण\appbrk{}च\appbrk{}ण्डो\appbrk{}ग्र\appbrk{}विषा हि}}\Afootnote{A; \skt{स्तीक्ष्ण\appbrk{}च\appbrk{}ण्डो\appbrk{}ग्र\appbrk{}विषा हि}\ K H}}\skt{ लूतास् }\pend
\pstart\versepar \skt{ताः सप्तरात्रेण नरं निहन्युः~। }\pend
\pstart\versepar \skt{}\edtext{\skt{अतोऽधिकेनापि}}{\lemma{\skt{ऽधिकेनापि}}\Afootnote{A H; \skt{धिकेनापि}\ K}}\skt{ निहन्युरन्या }\pend
\pstart\versepar \skt{}\edtext{\skt{येषां}}{\lemma{\skt{येषां}}\Afootnote{K H; \skt{यासां}\ A}}\skt{ विषं }\edtext{\edtext{\skt{मध्य\-म\-वी\-र्य\-मुक्तं}}{\lemma{\skt{मध्यमवीर्यम्}}\Afootnote{A; \skt{मध्यमवीर्य}\ K; \skt{मध्य\appbrk{}म\appbrk{}वी\appbrk{}र्य\appbrk{}र्यम्}\ H}}}{\lemma{\skt{उक्तं}}\Afootnote{A H; \skt{मुक्तं~।}\ K}}\skt{~।}\edtext{}{\lemma{\skt{लूतास्ताः}}\Afootnote{A; \skt{लूतास्ताः}\ K H}}\pend

\pstart\versepar \ledsidenote{\passref{5.8.84}}\skt{यासां कनीयो विषवीर्यमुक्तं }\pend
\pstart\versepar \skt{ताः पक्षमात्रेण विनाशयन्ति~। }\pend
\pstart\versepar \skt{तस्मात्प्रयत्नं भिषगत्र कुर्याद् }\pend
\pstart\versepar \skt{आ }\edtext{\skt{दंश\-पा\-ता\-द्वि\-ष\-घा\-त\-वेगैः}}{\lemma{\skt{दंश\appbrk{}पा\appbrk{}ता\appbrk{}द्वि\appbrk{}ष\appbrk{}घा\appbrk{}त\appbrk{}वेगैः}}\Afootnote{K; \skt{दंश\appbrk{}पा\appbrk{}ता\appbrk{}द्वि\appbrk{}ष\appbrk{}घा\appbrk{}ति\appbrk{}योगैः~॥}\ A; \skt{दंश\appbrk{}पा\appbrk{}ता\appbrk{}द्वि\appbrk{}ष\appbrk{}घा\appbrk{}त\appbrk{}वेर्गैः~॥}\ H\textsuperscript{ac}; \skt{दंश\appbrk{}पा\appbrk{}ता\appbrk{}द्वि\appbrk{}ष\appbrk{}घा\appbrk{}त\appbrk{}योर्गैः~॥}\ H\textsuperscript{pc}}}\skt{~।}\edtext{}{\lemma{\skt{विषवीर्यमुक्तं ताः}}\Afootnote{A; \skt{विषवीर्यमुक्तं ताः}\ K; \skt{विषवीर्ययुक्तं ताः}\ H}}\pend

\pstart\versepar \ledsidenote{\passref{5.8.85}}\edtext{\skt{विषन्तु}}{\lemma{\skt{विषन्तु}}\Afootnote{A H; \skt{विषंन्तु}\ K}}\skt{ लाला\-न\-ख\-मू\-त्र\-दंष्ट्रा }\pend
\pstart\versepar \skt{रजः\-पु\-री\-षै\-रथ चेन्द्रियेण~। }\pend
\pstart\versepar \skt{सप्तप्रकारं विसृजन्ति लूतास् }\pend
\pstart\versepar \skt{}\edtext{\edtext{\skt{तदु\-ग्र\-म\-ध्या\-व\-र\-वी\-र्य\-मुक्तं}}{\lemma{\skt{तदु\appbrk{}ग्र\appbrk{}म\appbrk{}ध्या\appbrk{}व\appbrk{}र\appbrk{}वी\appbrk{}र्यम्}}\Afootnote{K H; \skt{तदु\appbrk{}ग्र\appbrk{}म\appbrk{}ध्या\appbrk{}व\appbrk{}र\appbrk{}वीर्य}\ A}}}{\lemma{\skt{उक्तं}}\Afootnote{K H; \skt{युक्तम्~॥}\ A}}\skt{~॥}\pend

\pstart\versepar \ledsidenote{\passref{5.8.86}}\edtext{\skt{कोठं सक\-ण्डू\-स्थि\-र\-म\-ल्प\-मूलं}}{\lemma{\skt{कोठं सकण्डूस्थिरम्}}\Afootnote{\skt{सकण्डुकोठं स्थिरम्}\ A; \skt{कोठं सगण्डस्थिरम्}\ K; \skt{कोठ\appbrk{}न्स\appbrk{}क\appbrk{}ण्डू\appbrk{}स्थि\appbrk{}रम्}\ H}}\skt{ }\pend
\pstart\versepar \skt{लालाकृतं मन्दरुजं वदन्ति~। }\pend
\pstart\versepar \skt{}\edtext{\skt{चोषश्च}}{\lemma{\skt{चोषश्च}}\Afootnote{K H\textsuperscript{ac}; \skt{शोफश्च}\ A; \skt{श्वासश्च}\ H\textsuperscript{pc}}}\skt{ कण्डूश्च }\edtext{\skt{पुलायिकाश्च}}{\lemma{\skt{पुलायिकाश्च}}\Afootnote{K H; \skt{पुलालिका च}\ A}}\skt{ }\pend
\pstart\versepar \skt{धूमायनं चैव नखाग्रदंशे~।}\pend

\pstart\versepar \ledsidenote{\passref{5.8.87}}\edtext{\skt{दंशे तु}}{\lemma{\skt{दंशे तु}}\Afootnote{K H; \skt{दंशं तु}\ A}}\skt{ मूत्रेण सकृष्णमध्यम्~। }\pend
\pstart\versepar \skt{}\edtext{\skt{सर\-क्त\-प\-र्य\-न्त\-म\-वैहि}}{\lemma{\skt{अवैहि}}\Afootnote{K H; \skt{अवेहि}\ A}}\skt{ दीर्ण्णम्~। }\pend
\pstart\versepar \skt{दंष्ट्राभिरुग्रं कठिनं विवर्ण्णम्~। }\pend
\pstart\versepar \skt{}\edtext{\skt{जानीष्व}}{\lemma{\skt{जानीष्व}}\Afootnote{K H; \skt{जानीहि}\ A}}\skt{ दंशं स्थिरमण्डलञ्च~।}\pend

\pstart\versepar \ledsidenote{\passref{5.8.88}}\edtext{\skt{रजः\-पु\-री\-षे\-न्द्रि\-य\-जञ्च}}{\lemma{\skt{रजः\appbrk{}पु\appbrk{}री\appbrk{}षे\appbrk{}न्द्रि\appbrk{}य\appbrk{}जञ्च}}\Afootnote{K H; \skt{रजः\appbrk{}पु\appbrk{}री\appbrk{}षे\appbrk{}न्द्रि\appbrk{}यजं हि}\ A}}\skt{ विद्धि }\pend
\pstart\versepar \skt{स्फोटं }\edtext{\skt{प्रप\-क्वा\-म\-ल\-पी\-लु\-पा\-ण्डुम्}}{\lemma{\skt{प्रप\appbrk{}क्वा\appbrk{}म\appbrk{}ल\appbrk{}पी\appbrk{}लु\appbrk{}पा\appbrk{}ण्डुम्}}\Afootnote{K; \skt{विप\appbrk{}क्वा\appbrk{}म\appbrk{}ल\appbrk{}पी\appbrk{}लु\appbrk{}पा\appbrk{}ण्डुम्~।}\ A; \skt{प्रपक्वां मलपीलुपाण्डुं~॥}\ H}}\skt{~। }\pend
\pstart\versepar \skt{}\edtext{\skt{एता\-व\-दे\-त\-त्स\-मु\-दा\-हृतं ते}}{\lemma{\skt{एता\appbrk{}व\appbrk{}दे\appbrk{}त\appbrk{}त्स\appbrk{}मु\appbrk{}दा\appbrk{}हृतं ते}}\Afootnote{K H; \skt{एता\appbrk{}व\appbrk{}दे\appbrk{}त\appbrk{}त्स\appbrk{}मु\appbrk{}दा\appbrk{}हृतं तु}\ A}}\skt{ }\pend
\pstart\versepar \skt{वक्ष्यामि लूताप्रभवं }\edtext{\skt{प्रमाणम्}}{\lemma{\skt{प्रमाणम्}}\Afootnote{H; \skt{पुराणम्~॥}\ A; \skt{प्रमानं~।}\ K}}\skt{~।}\pend

\pstart\versepar \ledsidenote{\passref{5.8.89}}\skt{सामान्यतो दष्ट\-म\-सा\-ध्य\-साध्यं }\pend
\pstart\versepar \skt{चिकित्सितं चापि }\edtext{\edtext{\skt{विशेषणञ्च}}{\lemma{\skt{विशेषणञ्}}\Afootnote{K H; \skt{यथाविशेषम्~।}\ A}}}{\lemma{\skt{च}}\Afootnote{K H; \emph{om.}\ A}}\skt{~॥}\pend

\pstart\versepar \ledsidenote{\passref{5.8.90}}\skt{विश्वामित्रो नृपवरः }\edtext{\skt{कदा\-चि\-दृ\-षि\-स\-त्तमं}}{\lemma{\skt{कदा\appbrk{}चि\appbrk{}दृ\appbrk{}षि\appbrk{}स\appbrk{}त्तमं}}\Afootnote{A K\textsuperscript{pc} H; \skt{कदा\appbrk{}चि\appbrk{}दृ\appbrk{}षि\appbrk{}सत्त}\ K\textsuperscript{ac}}}\skt{~। }\pend
\pstart\versepar \skt{वशिष्ठं कोपयामास गत्वाश्रमपदं किल~।}\pend

\pstart\versepar \ledsidenote{\passref{5.8.91}}\skt{कुपितस्य मुनेस्तस्य }\edtext{\skt{लला\-टा\-त्स्वे\-द\-बि\-न्दवः}}{\lemma{\skt{लला\appbrk{}टा\appbrk{}त्स्वे\appbrk{}द\appbrk{}बि\appbrk{}न्दवः}}\Afootnote{A; \skt{लला\appbrk{}टा\appbrk{}त्स्वे\appbrk{}द\appbrk{}वि\appbrk{}न्दवः~।}\ K H}}\skt{~। }\pend
\pstart\versepar \skt{}\edtext{\skt{निये\-तु\-र्द\-र्श\-ना\-देव}}{\lemma{\skt{नियेतुर्}}\Afootnote{K H; \skt{अपतन्}\ A}}\skt{ }\edtext{\edtext{\skt{रवे\-स्त\-त्स\-म\-व\-र्चसः}}{\lemma{\skt{रवेस्}}\Afootnote{K H; \skt{स्वेस्}\ A}}}{\lemma{\skt{तत्समवर्चसः}}\Afootnote{K H; \skt{तत्समतेजसः~॥}\ A}}\skt{~।}\pend

\pstart \ledsidenote{\passref{5.8.92ab}}\edtext{\skt{लूने}}{\lemma{\skt{लूने}}\Afootnote{K H; \emph{om.}\ A}}\skt{ तृणे महर्षिणा }\edtext{\skt{धेन्वर्थे}}{\lemma{\skt{धेन्वर्थे}}\Afootnote{\skt{लूने धेन्वर्थं}\ A; \skt{वेन्धेन्वर्थे}\ K; \skt{वेन्धेण्वर्थे}\ H}}\skt{ }\edtext{\edtext{\skt{सम्भृतेऽपि}}{\lemma{\skt{सम्भृते}}\Afootnote{K H; \skt{संभृते}\ A}}}{\lemma{\skt{ऽपि}}\Afootnote{A; \skt{पि}\ K H}}\skt{ }\edtext{\skt{च}}{\lemma{\skt{च}}\Afootnote{A H; \skt{चः~।}\ K}}\skt{~।}\pend

\pstart \ledsidenote{\passref{5.8.92ef}}\skt{अपकाराय वर्त्तन्ते }\edtext{\skt{नृप\-शा\-स\-न\-वा\-हने}}{\lemma{\skt{नृप\appbrk{}शा\appbrk{}स\appbrk{}न\appbrk{}वा\appbrk{}हने}}\Afootnote{K; \skt{नृप\appbrk{}सा\appbrk{}ध\appbrk{}न\appbrk{}वा\appbrk{}हने~॥}\ A; \skt{नृप\appbrk{}शा\appbrk{}स\appbrk{}न\appbrk{}वा\appbrk{}हनैः~॥}\ H}}\skt{~।}\pend

\pstart\versepar \ledsidenote{\passref{5.8.93}}\edtext{\skt{यस्माल्लूनं}}{\lemma{\skt{यस्माल्लूनं}}\Afootnote{A; \skt{यस्माः लूनं}\ K; \skt{यस्माल्लूतन्}\ H}}\skt{ तृणं प्राप्ता }\edtext{\edtext{\skt{मुनेस्ते}}{\lemma{\skt{मुनेस्}}\Afootnote{A; \skt{न्मुनेस्}\ K; \skt{त्मुनेस्}\ H}}}{\lemma{\skt{ते}}\Afootnote{K H; \skt{प्र}\ A}}\skt{ }\edtext{\skt{स्वेदबिन्दवः}}{\lemma{\skt{स्वेदबिन्दवः}}\Afootnote{A; \skt{स्वेदविन्दवः~।}\ K H}}\skt{~। }\pend
\pstart\versepar \skt{}\edtext{\skt{तस्माल्लूता}}{\lemma{\skt{तस्माल्लूता}}\Afootnote{K H; \skt{तस्माल्लूते}\ A}}\skt{ }\edtext{\skt{विभाव्यन्ते}}{\lemma{\skt{विभाव्यन्ते}}\Afootnote{K H; \skt{ति भाष्यन्ते}\ A}}\skt{ }\edtext{\skt{संख्यया ताश्च}}{\lemma{\skt{संख्यया ताश्च}}\Afootnote{K H; \skt{सङ्ख्यया ताश्च}\ A}}\skt{ }\edtext{\skt{षोडशः}}{\lemma{\skt{षोडशः}}\Afootnote{K; \skt{षोडश~॥}\ A H}}\skt{~।}\pend

\pstart\versepar \ledsidenote{\passref{5.8.94}}\skt{कृच्छ्र\-सा\-ध्या\-स्त\-था\-साध्या लूतास्तु द्विविधाः स्मृताः~। }\pend
\pstart\versepar \skt{तासामष्टौ कृच्छ्रसाध्या वर्ज्यास्तावत्य एव तु~॥}\pend

\pstart\versepar \ledsidenote{\passref{5.8.95}}\edtext{\skt{त्रिमण्डला}}{\lemma{\skt{त्रिमण्डला}}\Afootnote{A; \skt{त्रिमण्डल}\ K; \skt{त्रिमण्डलस्}\ H}}\skt{ तथा }\edtext{\skt{स्वेता}}{\lemma{\skt{स्वेता}}\Afootnote{K H; \skt{श्वेता}\ A}}\skt{ कपिला पीतिका तथा~। }\pend
\pstart\versepar \skt{}\edtext{\skt{अलमूत्रविषे}}{\lemma{\skt{अलमूत्रविषे}}\Afootnote{K; \skt{आलमूत्रविषा}\ A; \skt{मलमूत्रविषे}\ H}}\skt{ रक्ता कसना चाष्टमी }\edtext{\skt{स्मृताः}}{\lemma{\skt{स्मृताः}}\Afootnote{K H; \skt{स्मृता~॥}\ A}}\skt{~॥}\pend

\pstart\versepar \ledsidenote{\passref{5.8.96}}\edtext{\skt{ताभिर्दष्टे}}{\lemma{\skt{दष्टे}}\Afootnote{A K; \skt{द्द्ष्टे}\ H}}\skt{ }\edtext{\skt{शिरो\-दुः\-ख\-मा\-दंशे}}{\lemma{\skt{आदंशे}}\Afootnote{K H; \emph{om.}\ A}}\skt{ }\edtext{\edtext{\skt{कण्डुरेव}}{\lemma{\skt{कण्डुर्}}\Afootnote{K H; \skt{कण्डूर्}\ A}}}{\lemma{\skt{एव}}\Afootnote{K H; \skt{दंशे च वेदना}\ A}}\skt{ }\edtext{\skt{च}}{\lemma{\skt{च}}\Afootnote{K H; \skt{।}\ A}}\skt{~। }\pend
\pstart\versepar \skt{भवन्ति च विशेषेण गदाः }\edtext{\skt{श्लेष्मिकवातिकाः}}{\lemma{\skt{श्लेष्मिकवातिकाः}}\Afootnote{K H; \skt{श्लैष्मिकवातिकाः~॥}\ A}}\skt{~।}\pend

\pstart\versepar \ledsidenote{\passref{5.8.97}}\skt{सौवर्ण्णिका लाजवर्ण्णा जालिन्येणीपदी तथा~। }\pend
\pstart\versepar \skt{}\edtext{\skt{कृष्णाग्निमुख्यौ}}{\lemma{\skt{कृष्णाग्निमुख्यौ}}\Afootnote{K; \skt{कृष्णाग्निवर्णा}\ A; \skt{कृष्णाग्निमुख्या}\ H}}\skt{ काकाण्डा }\edtext{\skt{मालागुण्यष्टमी}}{\lemma{\skt{मालागुण्य्}}\Afootnote{K H; \skt{मालागुणा}\ A}}\skt{ }\edtext{\skt{स्मृताः}}{\lemma{\skt{स्मृताः}}\Afootnote{K H; \skt{तथा~॥}\ A}}\skt{~।}\pend

\pstart\versepar \ledsidenote{\passref{5.8.98}}\skt{ताभिर्दष्टे दंशकोथः प्रवृत्तिः क्षतजस्य च~। }\pend
\pstart\versepar \skt{ज्वरो दाहोऽतिसारश्च गदाः स्युश्च त्रिदोषजाः~।}\pend

\pstart\versepar \ledsidenote{\passref{5.8.99}}\edtext{\skt{पिटका}}{\lemma{\skt{पिटका}}\Afootnote{K; \skt{पिडका}\ A H}}\skt{ विविधाकारा मण्डलानि महान्ति च~॥ }\pend
\pstart\versepar \skt{}\edtext{\skt{शोफा}}{\lemma{\skt{शोफा}}\Afootnote{K H; \emph{om.}\ A}}\skt{ }\edtext{\skt{महान्तो}}{\lemma{\skt{महान्तो}}\Afootnote{A K; \skt{महातो}\ H}}\skt{ }\edtext{\skt{मृदवो}}{\lemma{\skt{मृदवो}}\Afootnote{A H; \skt{मृदुवो}\ K}}\skt{ }\edtext{\skt{रक्ताः}}{\lemma{\skt{रक्ताः}}\Afootnote{K; \skt{शोफा रक्ताः}\ A; \skt{रक्ता}\ H}}\skt{ श्यावाश्चलास्तथा~।}\pend

\pstart\versepar \ledsidenote{\passref{5.8.100}}\skt{सामान्यं }\edtext{\skt{सर्व\-लू\-ता\-ना\-मे\-त\-दा\-दं\-श\-ल\-क्ष\-णम्}}{\lemma{\skt{सर्वलूतानाम्}}\Afootnote{A K\textsuperscript{pc} H; \skt{सर्वलूताम्}\ K\textsuperscript{ac}}}\skt{~। }\pend
\pstart\versepar \skt{विशेषलक्षणं तासां वक्ष्यामि सचिकित्सितम्~।}\pend

\pstart\versepar \ledsidenote{\passref{5.8.101}}\skt{त्रिमण्डलाया }\edtext{\skt{बहलं}}{\lemma{\skt{बहलं}}\Afootnote{\skt{दं}\ A; \skt{वहलं}\ K H}}\skt{ }\edtext{\skt{दंशः}}{\lemma{\skt{दंशः}}\Afootnote{K H; \skt{शेऽ}\ A}}\skt{ }\edtext{\skt{कृष्णं}}{\lemma{\skt{कृष्णं}}\Afootnote{K; \skt{सृक्कृष्णं}\ A; \skt{कृष्णः}\ H}}\skt{ }\edtext{\skt{क्षरत्यसृक्}}{\lemma{\skt{क्षरत्य्}}\Afootnote{K H; \skt{स्रवति}\ A}}\skt{~। }\pend
\pstart\versepar \skt{बाधिर्यं कलुषा दृष्टिस्तथा दाहश्च नेत्रयोः~।}\edtext{}{\lemma{\skt{असृक्~। बाधिर्यं}}\Afootnote{\skt{दीर्यते~। बाधिर्यं}\ A; \skt{असृक्~। वाधिर्यं}\ K H}}\pend

\pstart\versepar \ledsidenote{\passref{5.8.102}}\skt{तत्रार्कमूलं रजनी नाकुली पृश्निपर्ण्णिका~॥ }\pend
\pstart\versepar \skt{}\edtext{\skt{नस्तः}}{\lemma{\skt{नस्तः}}\Afootnote{K H; \skt{पान}\ A}}\skt{ कर्मणि }\edtext{\skt{शस्यन्ते}}{\lemma{\skt{शस्यन्ते}}\Afootnote{A K; \emph{om.}\ H}}\skt{ }\edtext{\skt{पादा\-भ्य\-ङ्गा\-ञ्ज\-नेषु}}{\lemma{\skt{पादा\appbrk{}भ्य\appbrk{}ङ्गा\appbrk{}ञ्ज\appbrk{}नेषु}}\Afootnote{K H; \skt{नस्या\appbrk{}ले\appbrk{}पा\appbrk{}ञ्ज\appbrk{}नेषु}\ A}}\skt{ च~।}\pend

\pstart\versepar \ledsidenote{\passref{5.8.103}}\edtext{\skt{श्वेतायाः}}{\lemma{\skt{श्वेतायाः}}\Afootnote{A K; \skt{श्वेतया}\ H}}\skt{ पिडका दंशे श्वेता कण्डूमती भवेत्~। }\pend
\pstart\versepar \skt{दाह\-मू\-र्च्छा\-ज्व\-र\-वती विस\-र्प\-क्ले\-द\-रु\-क्करी~॥}\pend

\pstart\versepar \ledsidenote{\passref{5.8.104}}\skt{तत्र }\edtext{\skt{चन्द\-न\-रा\-स्नै\-ल\-ह\-रे\-णु\-न\-ल\-व\-ञ्जुलाः}}{\lemma{\skt{चन्द\appbrk{}न\appbrk{}रा\appbrk{}स्नै\appbrk{}ल\appbrk{}ह\appbrk{}रे\appbrk{}णु\appbrk{}न\appbrk{}ल\appbrk{}व\appbrk{}ञ्जुलाः}}\Afootnote{K; \skt{चन्द\appbrk{}न\appbrk{}रा\appbrk{}स्नै\appbrk{}ला\appbrk{}ह\appbrk{}रे\appbrk{}णु\appbrk{}न\appbrk{}ल\appbrk{}व\appbrk{}ञ्जुलाः~।}\ A H}}\skt{~। }\pend
\pstart\versepar \skt{}\edtext{\skt{कुष्ठलामज्जकं}}{\lemma{\skt{कुष्ठलामज्जकं}}\Afootnote{K; \skt{कुष्ठं लामज्जकं}\ A H}}\skt{ वक्रं नलदं चागदो हितः~।}\pend

\pstart\versepar \ledsidenote{\passref{5.8.105}}\skt{आदंशे }\edtext{\skt{पिटका}}{\lemma{\skt{पिटका}}\Afootnote{\skt{पिडका}\ A; \skt{पिटकास्}\ K; \skt{पिडकास्}\ H}}\skt{ ताम्रा }\edtext{\skt{स्थिरा}}{\lemma{\skt{स्थिरा}}\Afootnote{K H; \emph{om.}\ A}}\skt{ }\edtext{\skt{कपिलया}}{\lemma{\skt{कपिलया}}\Afootnote{K H; \skt{कपिलायाः स्थिरा}\ A}}\skt{ भवेत्~। }\pend
\pstart\versepar \skt{शिरसो गौरवं दाहो }\edtext{\skt{भवेज्जन्तोश्च}}{\lemma{\skt{भवेज्जन्तोश्च}}\Afootnote{K H; \skt{तिमिरं भ्रम एव च~॥}\ A}}\skt{ }\edtext{\skt{नेत्रयोः}}{\lemma{\skt{नेत्रयोः}}\Afootnote{K H; \emph{om.}\ A}}\skt{~॥}\pend

\pstart\versepar \ledsidenote{\passref{5.8.106}}\edtext{\skt{पद्म\-प\-द्म\-क\-कु\-ष्ठै\-ल\-क\-र\-ञ्ज\-क\-कु\-भ\-त्वचः}}{\lemma{\skt{पद्म\appbrk{}प\appbrk{}द्म\appbrk{}क\appbrk{}कु\appbrk{}ष्ठै\appbrk{}ल\appbrk{}क\appbrk{}र\appbrk{}ञ्ज\appbrk{}क\appbrk{}कु\appbrk{}भ\appbrk{}त्वचः}}\Afootnote{K; \skt{तत्र पद्म\appbrk{}क\appbrk{}कु\appbrk{}ष्ठै\appbrk{}ला\appbrk{}क\appbrk{}र\appbrk{}ञ्ज\appbrk{}क\appbrk{}कु\appbrk{}भ\appbrk{}त्वचः~।}\ A; \skt{पद्म\appbrk{}प\appbrk{}द्म\appbrk{}क\appbrk{}कु\appbrk{}ष्ठै\appbrk{}ला\appbrk{}क\appbrk{}र\appbrk{}ञ्ज\appbrk{}क\appbrk{}कु\appbrk{}भ\appbrk{}त्वचः~।}\ H}}\skt{~। }\pend
\pstart\versepar \skt{}\edtext{\skt{स्थिरा\-क\-म्पि\-ल्या\-पा\-मा\-र्ग\-दू\-र्वा\-ब्राह्म्यौ}}{\lemma{\skt{स्थिरा\appbrk{}क\appbrk{}म्पि\appbrk{}ल्या\appbrk{}पा\appbrk{}मा\appbrk{}र्ग\appbrk{}दू\appbrk{}र्वा\appbrk{}ब्राह्म्यौ}}\Afootnote{\skt{स्थिरा\appbrk{}र्क\appbrk{}प\appbrk{}र्ण्य\appbrk{}पा\appbrk{}मा\appbrk{}र्ग\appbrk{}दू\appbrk{}र्वा\appbrk{}ब्राह्म्यो}\ A; \skt{स्थिरा\appbrk{}क\appbrk{}म्पि\appbrk{}ल्य\appbrk{}पा\appbrk{}मा\appbrk{}र्ग\appbrk{}दू\appbrk{}र्वा\appbrk{}व्राह्म्यौ}\ K; \skt{स्थिरा\appbrk{}कं\appbrk{}पि\appbrk{}ल्या\appbrk{}पा\appbrk{}मा\appbrk{}र्ग्ग\appbrk{}दु\appbrk{}र्व्वा\appbrk{}व्रह्म्यो}\ H}}\skt{ }\edtext{\skt{विषापहाः}}{\lemma{\skt{विषापहाः}}\Afootnote{A K; \skt{विषापहा~॥}\ H}}\skt{~।}\pend

\pstart\versepar \ledsidenote{\passref{5.8.107}}\skt{आदंशे }\edtext{\skt{पिटका}}{\lemma{\skt{पिटका}}\Afootnote{K H; \skt{पीतिका}\ A}}\skt{ }\edtext{\skt{पीता}}{\lemma{\skt{पीता}}\Afootnote{K H; \skt{यास्तु}\ A}}\skt{ }\edtext{\skt{पीतया}}{\lemma{\skt{पीतया}}\Afootnote{K H; \skt{पिडका}\ A}}\skt{ }\edtext{\skt{जायते}}{\lemma{\skt{जायते}}\Afootnote{K H; \skt{पीतिका}\ A}}\skt{ स्थिरा~। }\pend
\pstart\versepar \skt{}\edtext{\skt{तथा}}{\lemma{\skt{तथा}}\Afootnote{K H; \skt{भवे}\ A}}\skt{ }\edtext{\skt{च्छर्दिज्वरः}}{\lemma{\skt{च्छर्दिज्वरः}}\Afootnote{K; \skt{च्छर्दिर्ज्वरः}\ A; \skt{च्छर्द्दिर्ज्वरः}\ H}}\skt{ }\edtext{\skt{शूलो}}{\lemma{\skt{शूलो}}\Afootnote{K H; \skt{शूलं}\ A}}\skt{ }\edtext{\skt{रक्ते}}{\lemma{\skt{रक्ते}}\Afootnote{K H; \skt{मूर्ध्नि रक्ते}\ A}}\skt{ }\edtext{\skt{स्याताञ्च}}{\lemma{\skt{स्याताञ्च}}\Afootnote{K H; \emph{om.}\ A}}\skt{ }\edtext{\skt{लोचने}}{\lemma{\skt{लोचने}}\Afootnote{K H; \skt{तथाऽक्षिणी~॥}\ A}}\skt{~।}\pend

\pstart\versepar \ledsidenote{\passref{5.8.108}}\edtext{\skt{तत्रेष्ठाः}}{\lemma{\skt{तत्रेष्ठाः}}\Afootnote{K; \skt{तत्रेष्टाः}\ A; \skt{तत्रेष्ठा}\ H}}\skt{ }\edtext{\skt{ककु\-भो\-शी\-र\-मु\-ञ्जा\-ब\-ल्व\-ज\-व\-ञ्जलाः}}{\lemma{\skt{ककु\appbrk{}भो\appbrk{}शी\appbrk{}र\appbrk{}मु\appbrk{}ञ्जा\appbrk{}ब\appbrk{}ल्व\appbrk{}ज\appbrk{}व\appbrk{}ञ्जलाः}}\Afootnote{\skt{कुट\appbrk{}जो\appbrk{}शी\appbrk{}र\appbrk{}तु\appbrk{}ङ्ग\appbrk{}प\appbrk{}द्म\appbrk{}क\appbrk{}व\appbrk{}ञ्जुलाः~।}\ A; \skt{ककु\appbrk{}भो\appbrk{}शी\appbrk{}र\appbrk{}मु\appbrk{}ञ्जा\appbrk{}व\appbrk{}ल्व\appbrk{}ज\appbrk{}व\appbrk{}ञ्जलाः~।}\ K H}}\skt{~। }\pend
\pstart\versepar \skt{}\edtext{\skt{कुश\-का\-श\-वं\-श\-कि\-णि\-ही\-शि\-री\-ष\-क\-कु\-भ\-त्वचः}}{\lemma{\skt{कुश\appbrk{}का\appbrk{}श\appbrk{}वं\appbrk{}श\appbrk{}कि\appbrk{}णि\appbrk{}ही\appbrk{}शि\appbrk{}री\appbrk{}ष\appbrk{}क\appbrk{}कु\appbrk{}भ\appbrk{}त्वचः}}\Afootnote{K; \skt{शिरी\appbrk{}ष\appbrk{}कि\appbrk{}णि\appbrk{}ही\appbrk{}शे\appbrk{}लु\appbrk{}क\appbrk{}द\appbrk{}म्ब\appbrk{}क\appbrk{}कु\appbrk{}भ\appbrk{}त्वचः~॥}\ A; \skt{कुश\appbrk{}का\appbrk{}श\appbrk{}व\appbrk{}श\appbrk{}कि\appbrk{}णि\appbrk{}ही\appbrk{}शि\appbrk{}री\appbrk{}ष\appbrk{}क\appbrk{}कु\appbrk{}भ\appbrk{}त्वचः~॥}\ H}}\skt{~।}\pend

\pstart\versepar \ledsidenote{\passref{5.8.109}}\edtext{\skt{रक्त\-म\-ण्ड\-ल\-व\-द्दंशे}}{\lemma{\skt{रक्त\appbrk{}म\appbrk{}ण्ड\appbrk{}ल\appbrk{}व\appbrk{}द्दंशे}}\Afootnote{K H; \skt{रक्तमण्डनिभे दंशे}\ A}}\skt{ }\edtext{\skt{पिटकाः}}{\lemma{\skt{पिटकाः}}\Afootnote{K H; \skt{पिडकाः}\ A}}\skt{ }\edtext{\skt{सर्षपा}}{\lemma{\skt{सर्षपा}}\Afootnote{A K; \skt{सर्ष}\ H}}\skt{ }\edtext{\skt{इव}}{\lemma{\skt{इव}}\Afootnote{A K; \emph{om.}\ H}}\skt{~। }\pend
\pstart\versepar \skt{}\edtext{\skt{दूयते}}{\lemma{\skt{दूयते}}\Afootnote{K; \skt{जायन्ते}\ A; \emph{om.}\ H}}\skt{ }\edtext{\skt{तालुशोषश्च}}{\lemma{\skt{तालुशोषश्च}}\Afootnote{A K; \emph{om.}\ H}}\skt{ }\edtext{\edtext{\skt{दाह\-श्चा\-ल\-वि\-षा\-न्विते}}{\lemma{\skt{दाहश्}}\Afootnote{A K; \emph{om.}\ H}}}{\lemma{\skt{चालविषान्विते}}\Afootnote{\skt{चालविषार्दिते~॥}\ A; \skt{चलविषान्विते~।}\ K; \skt{पान्विते~।}\ H}}\skt{~।}\pend

\pstart\versepar \ledsidenote{\passref{5.8.110}}\skt{तत्र }\edtext{\skt{प्रिय\-ङ्गु\-ह्री\-वे\-र\-कु\-ष्ठ\-ला\-म\-ज्ज\-कानि वा}}{\lemma{\skt{प्रिय\appbrk{}ङ्गु\appbrk{}ह्री\appbrk{}वे\appbrk{}र\appbrk{}कु\appbrk{}ष्ठ\appbrk{}ला\appbrk{}म\appbrk{}ज्ज\appbrk{}कानि वा}}\Afootnote{K H; \skt{प्रिय\appbrk{}ङ्गु\appbrk{}ह्री\appbrk{}बे\appbrk{}र\appbrk{}कु\appbrk{}ष्ठ\appbrk{}ला\appbrk{}म\appbrk{}ज्ज\appbrk{}व\appbrk{}ञ्जुलाः~।}\ A}}\skt{~। }\pend
\pstart\versepar \skt{}\edtext{\skt{अगदः}}{\lemma{\skt{अगदः}}\Afootnote{A K; \skt{अगुरुः}\ H}}\skt{ शतपुष्पा च सपि\-प्प\-ल\-व\-टा\-ङ्कुराः~।}\pend

\pstart\versepar \ledsidenote{\passref{5.8.111}}\edtext{\skt{पूति\-मू\-त्र\-वि\-षा\-दंशो}}{\lemma{\skt{पूति\appbrk{}मू\appbrk{}त्र\appbrk{}वि\appbrk{}षा\appbrk{}दंशो}}\Afootnote{K H; \skt{पूति\appbrk{}र्मू\appbrk{}त्र\appbrk{}वि\appbrk{}षा\appbrk{}दंशो}\ A}}\skt{ विसर्पी कृष्णशोणितः~। }\pend
\pstart\versepar \skt{कास\-श्वा\-स\-व\-मी\-मू\-र्च्छा\-ज्व\-र\-दा\-ह\-स\-म\-न्वितः~।}\pend

\pstart\versepar \ledsidenote{\passref{5.8.112}}\edtext{\skt{मनः\-शि\-ला\-ल\-म\-धु\-क\-कु\-ष्ठ\-प\-द्म\-क\-च\-न्दनैः}}{\lemma{\skt{मनः\appbrk{}शि\appbrk{}ला\appbrk{}ल\appbrk{}म\appbrk{}धु\appbrk{}क\appbrk{}कु\appbrk{}ष्ठ\appbrk{}प\appbrk{}द्म\appbrk{}क\appbrk{}च\appbrk{}न्दनैः}}\Afootnote{K H; \skt{मनः\appbrk{}शि\appbrk{}ला\appbrk{}ल\appbrk{}म\appbrk{}धु\appbrk{}क\appbrk{}कु\appbrk{}ष्ठ\appbrk{}च\appbrk{}न्द\appbrk{}न\appbrk{}प\appbrk{}द्मकैः~।}\ A}}\skt{~। }\pend
\pstart\versepar \skt{}\edtext{\skt{लाम\-ज्ज\-क\-यु\-तै\-स्तत्र}}{\lemma{\skt{लाम\appbrk{}ज्ज\appbrk{}क\appbrk{}यु\appbrk{}तैस्}}\Afootnote{K H; \skt{मधुमिश्रैः सला\appbrk{}म\appbrk{}ज्जै\appbrk{}र\appbrk{}ग\appbrk{}दस्}\ A}}\skt{ }\edtext{\skt{विषनाशः}}{\lemma{\skt{विषनाशः}}\Afootnote{K H; \emph{om.}\ A}}\skt{ }\edtext{\skt{प्रकीर्त्तितः}}{\lemma{\skt{प्रकीर्त्तितः}}\Afootnote{K H; \skt{कीर्तितः~॥}\ A}}\skt{~।}\pend

\pstart\versepar \ledsidenote{\passref{5.8.113}}\edtext{\skt{दंशः}}{\lemma{\skt{दंशः}}\Afootnote{K H; \emph{om.}\ A}}\skt{ }\edtext{\skt{सपाण्डुपिटको}}{\lemma{\skt{सपाण्डुपिटको}}\Afootnote{K H; \skt{आपाण्डुपिडको}\ A}}\skt{ }\edtext{\skt{दाह\-क्ले\-द\-स\-म\-न्वितः}}{\lemma{\skt{दाह\appbrk{}क्ले\appbrk{}द\appbrk{}स\appbrk{}म\appbrk{}न्वितः}}\Afootnote{K; \skt{दंशो दाह\appbrk{}क्ले\appbrk{}द\appbrk{}स\appbrk{}म\appbrk{}न्वितः~।}\ A; \skt{दाहः क्लेदसमन्वितः~॥}\ H}}\skt{~। }\pend
\pstart\versepar \skt{}\edtext{\skt{रक्तया}}{\lemma{\skt{रक्तया}}\Afootnote{K H; \skt{रक्ताया}\ A}}\skt{ रक्तपर्यन्तो }\edtext{\skt{विज्ञे\-य\-श्चो\-ष\-सं\-युतः}}{\lemma{\skt{चोषसंयुतः}}\Afootnote{K H; \skt{रक्तसंयुतः~॥}\ A}}\skt{~।}\pend

\pstart\versepar \ledsidenote{\passref{5.8.114}}\edtext{\skt{चिकित्सा}}{\lemma{\skt{चिकित्सा}}\Afootnote{K; \skt{कार्यस्}\ A; \skt{चिकित्सान्}\ H}}\skt{ }\edtext{\skt{तत्र}}{\lemma{\skt{तत्र}}\Afootnote{K H; \skt{तत्रा}\ A}}\skt{ }\edtext{\skt{ह्रीवे\-र\-च\-न्द\-नो\-शी\-र\-प\-द्मकैः}}{\lemma{\skt{ह्रीवे\appbrk{}र\appbrk{}च\appbrk{}न्द\appbrk{}नो\appbrk{}शी\appbrk{}र\appbrk{}प\appbrk{}द्मकैः}}\Afootnote{K; \skt{गद\appbrk{}स्तो\appbrk{}य\appbrk{}च\appbrk{}न्द\appbrk{}नो\appbrk{}शी\appbrk{}र\appbrk{}प\appbrk{}द्मकैः~।}\ A; \skt{ह्रीवे\appbrk{}र\appbrk{}श्च\appbrk{}न्द\appbrk{}नो\appbrk{}शी\appbrk{}र\appbrk{}प\appbrk{}द्मकैः~॥}\ H}}\skt{~। }\pend
\pstart\versepar \skt{}\edtext{\skt{कर्त्त\-व्या\-र्जु\-न\-शे\-लुभ्यां}}{\lemma{\skt{कर्त्त\appbrk{}व्या\appbrk{}र्जु\appbrk{}न\appbrk{}शे\appbrk{}लुभ्यां}}\Afootnote{K H; \skt{तथै\appbrk{}वा\appbrk{}र्जु\appbrk{}न\appbrk{}शे\appbrk{}लुभ्यां}\ A}}\skt{ }\edtext{\skt{त्वग्भि\-रा\-म्रा\-त\-कस्य}}{\lemma{\skt{आम्रातकस्य}}\Afootnote{A H; \skt{आम्रातकस्य}\ K}}\skt{ च~।}\pend

\pstart\versepar \ledsidenote{\passref{5.8.115}}\skt{पिच्छिलं }\edtext{\skt{कसनादंशो}}{\lemma{\skt{कसनादंशो}}\Afootnote{\skt{कसनादंशा}\ A K; \skt{कसनादं}\ H}}\skt{ }\edtext{\skt{रुधिरं}}{\lemma{\skt{रुधिरं}}\Afootnote{K; \skt{द्रुधिरं}\ A; \skt{शो रुधिरं}\ H}}\skt{ शीतलं स्रवेत्~। }\pend
\pstart\versepar \skt{श्वासकासौ च }\edtext{\skt{तत्रोक्तं}}{\lemma{\skt{तत्रोक्तं}}\Afootnote{A H; \emph{om.}\ K}}\skt{ रक्त\-लू\-ता\-चि\-कि\-त्सि\-तम्~।}\edtext{}{\lemma{\skt{स्रवेत्~। श्वासकासौ च}}\Afootnote{K H; \skt{स्रवेत्~। कासश्वासौ च}\ A}}\pend

\pstart\versepar \ledsidenote{\passref{5.8.120}}\edtext{\skt{सर्वेषामेव}}{\lemma{\skt{सर्वेषाम्}}\Afootnote{K H; \skt{सर्वासाम्}\ A}}\skt{ युञ्जीत }\edtext{\skt{विषे}}{\lemma{\skt{विषे}}\Afootnote{A H; \skt{विशे}\ K}}\skt{ }\edtext{\skt{श्लेष्मा\-त\-क\-त्व\-चम्}}{\lemma{\skt{श्लेष्मा\appbrk{}त\appbrk{}क\appbrk{}त्व\appbrk{}चम्}}\Afootnote{A K; \skt{श्लेष्मान्तकत्वचं~॥}\ H}}\skt{~। }\pend
\pstart\versepar \skt{}\edtext{\skt{धीरः}}{\lemma{\skt{धीरः}}\Afootnote{K H; \skt{भिष}\ A}}\skt{ }\edtext{\skt{सर्वविकारेषु}}{\lemma{\skt{सर्वविकारेषु}}\Afootnote{K H; \skt{क्सर्वप्रकारेण}\ A}}\skt{ तथा चाक्षी\-व\-पि\-प्प\-लम्~।}\pend

\pstart\versepar \ledsidenote{\passref{5.8.121}}\edtext{\skt{कृच्छ्रसाध्या}}{\lemma{\skt{कृच्छ्रसाध्या}}\Afootnote{K H; \skt{कृच्छ्रसाध्य}\ A}}\skt{ विषा ह्यष्टौ }\edtext{\skt{लूताः}}{\lemma{\skt{लूताः}}\Afootnote{K; \emph{om.}\ A; \skt{लूता}\ H}}\skt{ प्रोक्ता }\edtext{\skt{यथागमं}}{\lemma{\skt{यथागमं}}\Afootnote{K; \skt{द्वे च यदृ}\ A; \emph{om.}\ H}}\skt{~। }\pend
\pstart\versepar \skt{}\edtext{\skt{अवा\-र्य\-वि\-ष\-वी\-र्याणां}}{\lemma{\skt{अवा\appbrk{}र्य\appbrk{}वि\appbrk{}ष\appbrk{}वी\appbrk{}र्याणां}}\Afootnote{\skt{च्छया~। अवा\appbrk{}र्य\appbrk{}वि\appbrk{}ष\appbrk{}वी\appbrk{}र्याणां}\ A; \skt{अवा\appbrk{}र्य\appbrk{}वि\appbrk{}ष\appbrk{}वी\appbrk{}र्याणां}\ K; \skt{यथाक्रमं~॥ …र\appbrk{}वि\appbrk{}ष\appbrk{}वी\appbrk{}र्याणां}\ H}}\skt{ लक्षणानि }\edtext{\skt{निबोध}}{\lemma{\skt{निबोध}}\Afootnote{A; \skt{निवोध}\ K H}}\skt{ मे~।}\pend

\pstart\versepar \ledsidenote{\passref{5.8.122}}\edtext{\skt{ध्यामः}}{\lemma{\skt{ध्यामः}}\Afootnote{A; \skt{ध्यातः}\ K H}}\skt{ }\edtext{\skt{सौवर्ण्णिकादंशः}}{\lemma{\skt{सौवर्ण्णिकादंशः}}\Afootnote{A K; \skt{सौपर्ण्णिकादंशः}\ H}}\skt{ सफेनो }\edtext{\skt{मत्स्यगन्धिकः}}{\lemma{\skt{मत्स्यगन्धिकः}}\Afootnote{K H; \skt{मत्स्यगन्धकः~।}\ A}}\skt{~।}\skt{ }\pend
\pstart\versepar \skt{}\edtext{\skt{कासश्वासो}}{\lemma{\skt{कासश्वासो}}\Afootnote{K H; \skt{श्वासः कासो}\ A}}\skt{ ज्वरस्तृष्णा मूर्च्छा चात्र }\edtext{\skt{सुदारुणाः}}{\lemma{\skt{सुदारुणाः}}\Afootnote{K H; \skt{सुदारुणा~॥}\ A}}\skt{~।}\pend

\pstart\versepar \ledsidenote{\passref{5.8.123}}\edtext{\skt{ध्यामं}}{\lemma{\skt{ध्यामं}}\Afootnote{\emph{om.}\ A; \skt{ध्यामः}\ K H}}\skt{ }\edtext{\skt{पूति}}{\lemma{\skt{पूति}}\Afootnote{K; \emph{om.}\ A; \skt{पूतिः}\ H}}\skt{ }\edtext{\edtext{\skt{स्रवे\-द्र\-क्त\-मा\-दांशे}}{\lemma{\skt{स्रवेद्रक्तम्}}\Afootnote{K H; \emph{om.}\ A}}}{\lemma{\skt{आदांशे}}\Afootnote{K; \skt{आदंशे}\ A H}}\skt{ }\edtext{\skt{लाजवर्ण्णया}}{\lemma{\skt{लाजवर्ण्णया}}\Afootnote{K H; \skt{लाजवर्णाया}\ A}}\skt{~।}\skt{ }\pend
\pstart\versepar \skt{}\edtext{\skt{दाहो}}{\lemma{\skt{दाहो}}\Afootnote{K H; \skt{ध्यामं पूति स्रवेदसृक्~। दाहो}\ A}}\skt{ }\edtext{\skt{मूर्च्छातिसारौ च}}{\lemma{\skt{मूर्च्छातिसारौ च}}\Afootnote{K; \skt{मूर्च्छाऽतिसारश्च}\ A; \skt{मूर्च्छा च सारौ च}\ H}}\skt{ }\edtext{\skt{शिरोदुःख च}}{\lemma{\skt{शिरोदुःख च}}\Afootnote{K; \skt{शिरोदुःखं च}\ A H}}\skt{ जायते~।}\pend

\pstart\versepar \ledsidenote{\passref{5.8.124}}\skt{घोरो }\edtext{\skt{दंशश्च}}{\lemma{\skt{दंशश्च}}\Afootnote{K H; \skt{दंशस्तु}\ A}}\skt{ जालिन्या }\edtext{\skt{राजि\-मा\-न\-व\-दी\-र्यते}}{\lemma{\skt{राजि\appbrk{}मा\appbrk{}न\appbrk{}व\appbrk{}दी\appbrk{}र्यते}}\Afootnote{H; \skt{राजि\appbrk{}मा\appbrk{}न\appbrk{}व\appbrk{}की\appbrk{}र्यते~।}\ A; \skt{राजि\appbrk{}मा\appbrk{}ण\appbrk{}व\appbrk{}दी\appbrk{}र्यते~।}\ K}}\skt{~। }\pend
\pstart\versepar \skt{स्तम्भः श्वास\-स्त\-मो\-वृ\-द्धि\-स्ता\-लु\-शो\-षश्च }\edtext{\skt{तत्कृतः}}{\lemma{\skt{तत्कृतः}}\Afootnote{K H; \skt{जायते~॥}\ A}}\skt{~।}\pend

\pstart\versepar \ledsidenote{\passref{5.8.125}}\edtext{\skt{एणीपद्या}}{\lemma{\skt{एणीपद्या}}\Afootnote{\skt{एणीपद्यास्}\ A; \skt{एणीपदा}\ K H}}\skt{ }\edtext{\skt{महादाहो}}{\lemma{\skt{महादाहो}}\Afootnote{K H; \skt{तथा दंशो}\ A}}\skt{ }\edtext{\skt{दंशः}}{\lemma{\skt{दंशः}}\Afootnote{K H; \skt{भवे}\ A}}\skt{ }\edtext{\skt{कृष्णतिलाकृतिः}}{\lemma{\skt{कृष्णतिलाकृतिः}}\Afootnote{K H; \skt{त्कृष्णतिलाकृतिः~।}\ A}}\skt{~।}\skt{ }\pend
\pstart\versepar \skt{}\edtext{\skt{तृष्णा\-मू\-र्च्छा\-ज्व\-र\-च्छर्दिः}}{\lemma{\skt{तृष्णा\appbrk{}मू\appbrk{}र्च्छा\appbrk{}ज्व\appbrk{}र\appbrk{}च्छर्दिः}}\Afootnote{K; \skt{तृष्णा मूर्च्छा ज्वरच्छर्दि}\ A; \skt{तृष्णा\appbrk{}मू\appbrk{}र्च्छा\appbrk{}ज्व\appbrk{}र\appbrk{}श्च्छर्द्दिः~।}\ H}}\skt{ }\edtext{\skt{श्वास\-का\-स\-स\-म\-न्वितः}}{\lemma{\skt{श्वास\appbrk{}का\appbrk{}स\appbrk{}स\appbrk{}म\appbrk{}न्वितः}}\Afootnote{K; \skt{कास\appbrk{}श्वा\appbrk{}स\appbrk{}स\appbrk{}म\appbrk{}न्वितः~॥}\ A; \skt{श्वस\appbrk{}का\appbrk{}स\appbrk{}स\appbrk{}म\appbrk{}न्वितः~।}\ H}}\skt{~।}\pend

\pstart\versepar \ledsidenote{\passref{5.8.125.add-1}}\skt{कृष्णया }\edtext{\skt{कृष्णपर्यन्तो}}{\lemma{\skt{कृष्णपर्यन्तो}}\Afootnote{K; \skt{हृष्णपर्यन्तो}\ H}}\skt{ निम्न\-म\-ध्योऽ\-ति\-चो\-ष\-वान्~। }\pend
\pstart\versepar \skt{पाण्डु\-मू\-र्च्छा\-व\-मी\-दाहः }\edtext{\skt{श्वास\-का\-स\-स\-म\-न्वितः}}{\lemma{\skt{श्वास\appbrk{}का\appbrk{}स\appbrk{}स\appbrk{}म\appbrk{}न्वितः}}\Afootnote{K; \skt{श्वस\appbrk{}का\appbrk{}स\appbrk{}स\appbrk{}म\appbrk{}न्वितः~।}\ H}}\skt{~।}\pend

\pstart\versepar \ledsidenote{\passref{5.8.125.add-2}}\edtext{\skt{दंशोऽग्निमुख्यो}}{\lemma{\skt{ऽग्निमुख्यो}}\Afootnote{H; \skt{ग्निमुख्यो}\ K}}\skt{ विज्ञेयो दग्धः स्फोटाः सवेदनः~। }\pend
\pstart\versepar \skt{चोष\-क\-ण्डू\-रो\-म\-हर्षो }\edtext{\skt{दाह\-ज्व\-र\-नि\-पी\-डतः}}{\lemma{\skt{दाह\appbrk{}ज्व\appbrk{}र\appbrk{}नि\appbrk{}पी\appbrk{}डतः}}\Afootnote{K; \skt{दाह\appbrk{}ज्व\appbrk{}र\appbrk{}स\appbrk{}म\appbrk{}न्वितः~।}\ H}}\skt{~।}\pend

\pstart\versepar \ledsidenote{\passref{5.8.126}}\skt{दंशः काकाण्डिकादष्टे पाण्डु\-र\-क्तोऽ\-ति\-वे\-दनः~। }\pend
\pstart\versepar \skt{}\edtext{\edtext{\skt{हिक्का\-का\-स\-स्तृ\-षा\-मू\-र्च्छा\-नि\-द्रा\-हृ\-द्रो\-ग\-पी\-डितः}}{\lemma{\skt{हिक्काकासस्}}\Afootnote{K; \emph{om.}\ A; \skt{हिक्वाकास}\ H}}}{\lemma{\skt{तृषा\appbrk{}मू\appbrk{}र्च्छा\appbrk{}नि\appbrk{}द्रा\appbrk{}हृ\appbrk{}द्रो\appbrk{}ग\appbrk{}पी\appbrk{}डितः}}\Afootnote{K H; \skt{तृण्मू\appbrk{}र्च्छा\appbrk{}श्वा\appbrk{}स\appbrk{}हृ\appbrk{}द्रो\appbrk{}ग\appbrk{}हि\appbrk{}क्का\appbrk{}कासाः स्युरुच्छ्रिताः~॥}\ A}}\skt{~।}\pend

\pstart\versepar \ledsidenote{\passref{5.8.127}}\edtext{\skt{रक्तो}}{\lemma{\skt{रक्तो}}\Afootnote{A K; \emph{om.}\ H}}\skt{ }\edtext{\skt{दंशो}}{\lemma{\skt{दंशो}}\Afootnote{K; \skt{मालागुणादंशो}\ A; \skt{रक्तो दंशो}\ H}}\skt{ }\edtext{\edtext{\skt{धूम\-ग\-न्धि\-र्मा\-ला\-गु\-ण्या\-ति\-वे\-दनः}}{\lemma{\skt{धूमगन्धिर्}}\Afootnote{H; \skt{धूमगन्धो}\ A; \skt{मधूगन्धिर्}\ K}}}{\lemma{\skt{माला\appbrk{}गु\appbrk{}ण्या\appbrk{}ति\appbrk{}वे\appbrk{}दनः}}\Afootnote{K; \skt{ऽतिवेदनः~।}\ A; \skt{म्माला\appbrk{}गु\appbrk{}ण्या\appbrk{}दि\appbrk{}वे\appbrk{}दनः~॥}\ H}}\skt{~। }\skt{ }\pend
\pstart\versepar \skt{}\edtext{\skt{विदीर्यते च}}{\lemma{\skt{विदीर्यते च}}\Afootnote{K H; \emph{om.}\ A}}\skt{ }\edtext{\skt{बहुधा}}{\lemma{\skt{बहुधा}}\Afootnote{\skt{बहुधा च}\ A; \skt{वहुधा}\ K H}}\skt{ }\edtext{\skt{दाह\-मू\-र्च्छा\-ज्व\-रा\-न्वितः}}{\lemma{\skt{दाह\appbrk{}मू\appbrk{}र्च्छा\appbrk{}ज्व\appbrk{}रा\appbrk{}न्वितः}}\Afootnote{K H; \skt{विशीर्येत दाह\appbrk{}मू\appbrk{}र्च्छा\appbrk{}ज्व\appbrk{}रा\appbrk{}न्वितः~॥}\ A}}\skt{~।}\pend

\pstart\versepar \ledsidenote{\passref{5.8.128}}\edtext{\skt{असाध्यानामपि}}{\lemma{\skt{असाध्यानाम्}}\Afootnote{K H; \skt{असाध्यास्व्}\ A}}\skt{ }\edtext{\skt{भिषक्प्रयुञ्जीत}}{\lemma{\skt{भिषक्प्रयुञ्जीत}}\Afootnote{K; \skt{अभिहितं प्र}\ A; \skt{भिषष्प्रयुञ्जीत}\ H}}\skt{ }\edtext{\skt{चिकित्सितम्}}{\lemma{\skt{चिकित्सितम्}}\Afootnote{K H; \skt{त्याख्यायाशु योज}\ A}}\skt{~॥ }\pend
\pstart\versepar \skt{}\edtext{\skt{दोषो\-च्छ्रा\-य\-वि\-शे\-षेण}}{\lemma{\skt{दोषो\appbrk{}च्छ्रा\appbrk{}य\appbrk{}वि\appbrk{}शे\appbrk{}षेण}}\Afootnote{H; \skt{येत्~। दोषो\appbrk{}च्छ्रा\appbrk{}य\appbrk{}वि\appbrk{}शे\appbrk{}षेण}\ A; \skt{दोषोच्छ्रायो विशेषेण}\ K}}\skt{ }\edtext{\skt{च्छेद\-क\-र्म\-वि\-व\-र्जि\-तम्}}{\lemma{\skt{च्छेद\appbrk{}क\appbrk{}र्म\appbrk{}वि\appbrk{}व\appbrk{}र्जि\appbrk{}तम्}}\Afootnote{K H; \skt{दाह\appbrk{}च्छे\appbrk{}द\appbrk{}वि\appbrk{}व\appbrk{}र्जि\appbrk{}तम्~॥}\ A}}\skt{~।}\pend

\pstart\versepar \ledsidenote{\passref{5.8.129}}\edtext{\skt{साध्याभिरथ}}{\lemma{\skt{अथ}}\Afootnote{K H; \skt{आभिर्}\ A}}\skt{ }\edtext{\skt{लूता\-भि\-र्द\-ष्ट\-मा\-त्रस्य}}{\lemma{\skt{दष्टमात्रस्य}}\Afootnote{A K; \skt{द्द्ष्टमात्रस्य}\ H}}\skt{ देहिनः~। }\pend
\pstart\versepar \skt{वृद्धिपत्रेण }\edtext{\skt{मतिमां}}{\lemma{\skt{मतिमां}}\Afootnote{K; \skt{मतिमा}\ A H}}\skt{ }\edtext{\skt{सम्य\-गा\-दं\-श\-मु\-द्ध\-रेत्}}{\lemma{\skt{सम्यगादंशम्}}\Afootnote{K; \skt{न्सम्यगादंशम्}\ A H}}\skt{~।}\pend

\pstart \ledsidenote{\passref{5.8.129.add-1}}\edtext{\skt{जम्बोष्ठेन}}{\lemma{\skt{जम्बोष्ठेन}}\Afootnote{\skt{जम्वोष्ठेन}\ K H}}\skt{ सुतप्तेन दहे\-दा\-क\-र\-वा\-र\-णात्~।}\pend

\pstart \ledsidenote{\passref{5.8.131ab}}\edtext{\skt{मधु\-सै\-न्ध\-व\-सं\-यु\-क्तै\-र\-ग\-दै\-र्ले\-प\-ये\-त्ततः}}{\lemma{\skt{अगदैर्}}\Afootnote{A H; \skt{अगर्दै}\ K}}\skt{~।}\pend

\pstart \ledsidenote{\passref{5.8.133ab}}\skt{क्षीरिणां त्वक्कषायेण }\edtext{\skt{कुसु\-म्भ\-म\-धु\-सै\-न्धवैः}}{\lemma{\skt{कुसु\appbrk{}म्भ\appbrk{}म\appbrk{}धु\appbrk{}सै\appbrk{}न्धवैः}}\Afootnote{K H; \skt{सुशीतेन च सेचयेत्~।}\ A}}\skt{~।}\pend

\pstart\versepar \ledsidenote{\passref{5.8.137}}\edtext{\edtext{\skt{किण्व\-गु\-ग्गु\-ल\-गो\-द\-न्त\-पा\-रा\-व\-त\-म\-लै\-रपि}}{\lemma{\skt{किण्व\appbrk{}गु\appbrk{}ग्गु\appbrk{}ल\appbrk{}गो\appbrk{}द\appbrk{}न्त\appbrk{}पा\appbrk{}रा\appbrk{}व\appbrk{}त\appbrk{}म\appbrk{}लैर्}}\Afootnote{K; \skt{गुग्गुलुः सैन्धवं किण्वं वर्चः पारावत}\ A; \skt{किण्व\appbrk{}गु\appbrk{}ग्गु\appbrk{}लु\appbrk{}गो\appbrk{}द\appbrk{}न्त\appbrk{}पा\appbrk{}रा\appbrk{}व\appbrk{}त\appbrk{}म\appbrk{}लैर्}\ H}}}{\lemma{\skt{अपि}}\Afootnote{K H; \skt{स्य च~।}\ A}}\skt{~। }\pend
\pstart\versepar \skt{विष\-वृ\-द्धि\-क\-र\-ञ्चान्नं हित्वा }\edtext{\skt{संभोजनं}}{\lemma{\skt{संभोजनं}}\Afootnote{A K; \skt{सभोजनं}\ H}}\skt{ हितम्~।}\pend

\pstart\versepar \ledsidenote{\passref{5.8.138}}\edtext{\skt{विषेभ्यः}}{\lemma{\skt{विषेभ्यः}}\Afootnote{A K; \skt{विसेभ्यः}\ H}}\skt{ खलु सर्वेभ्यः कर्ण्णिकामरुजां स्थिराम्~। }\pend
\pstart\versepar \skt{प्रच्छयित्वा }\edtext{\skt{मधुयुतैः}}{\lemma{\skt{मधुयुतैः}}\Afootnote{K H; \skt{मधून्मिश्रैः}\ A}}\skt{ शोध\-नी\-यै\-रु\-पा\-च\-रेत्~।}\pend

\pstart \ledsidenote{\passref{5.8.138.add}}\skt{सदा\-ह\-पा\-का\-न्य\-न्नेन }\edtext{\skt{चिकि\-त्से\-द्दृ\-ष्ट\-वा\-न्भि\-षक्}}{\lemma{\skt{चिकि\appbrk{}त्से\appbrk{}द्दृ\appbrk{}ष्ट\appbrk{}वा\appbrk{}न्भि\appbrk{}षक्}}\Afootnote{H; \skt{चिकित्सेद्दृष्टवा भिषक्~।}\ K}}\skt{~।}\pend

\pstart\versepar \ledsidenote{\passref{5.8.139}}\skt{सप्तषष्ठस्य कीटानां शत\-स्यै\-त\-द्वि\-भा\-गशः~। }\pend
\pstart\versepar \skt{दष्ट\-ल\-क्ष\-ण\-मा\-ख्यातं }\edtext{\edtext{\skt{चिकित्साञ्चाप्यतः}}{\lemma{\skt{चिकित्साञ्}}\Afootnote{K H; \skt{चिकित्सा}\ A}}}{\lemma{\skt{अतः}}\Afootnote{K; \skt{अनन्}\ A H}}\skt{ }\edtext{\skt{परम्}}{\lemma{\skt{परम्}}\Afootnote{K; \skt{तरम्}\ A H}}\skt{~।}\pend

\pstart\versepar \ledsidenote{\passref{5.8.140}}\skt{सविं\-श\-म\-ध्या\-य\-श\-त\-मे\-त\-दुक्तं विभागशः~। }\pend
\pstart\versepar \skt{}\edtext{\skt{इहो\-द्दि\-ष्टा\-न\-नि\-र्दि\-ष्टा\-न\-र्था\-न्व\-क्ष्यामि}}{\lemma{\skt{इहो\appbrk{}द्दि\appbrk{}ष्टा\appbrk{}न\appbrk{}नि\appbrk{}र्दि\appbrk{}ष्टा\appbrk{}न\appbrk{}र्था\appbrk{}न्व\appbrk{}क्ष्यामि}}\Afootnote{A H; \skt{इहो\appbrk{}द्दि\appbrk{}ष्टा\appbrk{}न\appbrk{}नि\appbrk{}र्दि\appbrk{}ष्टा\appbrk{}नर्था वक्ष्यामि}\ K}}\skt{ }\edtext{\skt{चोत्तरे}}{\lemma{\skt{चोत्तरे}}\Afootnote{K H; \skt{अथोत्तरे~॥}\ A}}\skt{~॥}\pend

\pstart\versepar \ledsidenote{\passref{5.8.140.add-1}}\skt{शास्त्रं }\edtext{\skt{शास्त्रसमुत्पत्तिं}}{\lemma{\skt{शास्त्रसमुत्पत्तिं}}\Afootnote{K; \skt{शास्त्रसमुत्पत्ति}\ H}}\skt{ }\edtext{\skt{व्याधि\-का\-र्य\-बा\-ला\-बलं}}{\lemma{\skt{व्याधि\appbrk{}का\appbrk{}र्य\appbrk{}बा\appbrk{}ला\appbrk{}बलं}}\Afootnote{\skt{व्याधि\appbrk{}का\appbrk{}र्य\appbrk{}वा\appbrk{}ला\appbrk{}वलं}\ K; \skt{व्याधि\appbrk{}का\appbrk{}र्य\appbrk{}व\appbrk{}ला\appbrk{}वलं~॥}\ H}}\skt{ }\pend
\pstart\versepar \skt{सूत्रभूतं समासेन श्लोकस्थानं }\edtext{\skt{प्रचक्षते}}{\lemma{\skt{प्रचक्षते}}\Afootnote{K; \skt{प्रवक्ष्यते~।}\ H}}\skt{~॥}\pend

\pstart\versepar \ledsidenote{\passref{5.8.140.add-2}}\skt{दोषा\-हा\-रा\-प\-चा\-रैश्च साग\-न्तु\-व्या\-धि\-ल\-क्ष\-णम्~। }\pend
\pstart\versepar \skt{अव\-स्था\-ल\-क्ष\-ण\-ञ्चैव निदानं स्थानमुच्यते~॥}\pend

\pstart\versepar \ledsidenote{\passref{5.8.140.add-3}}\skt{सम्भवश्चैव देहस्य धातु\-रि\-न्द्रि\-य\-म\-र्मसु~।}\skt{ }\pend
\pstart\versepar \skt{सिरादीनाञ्च सर्वेषां शारीरे कथितम्मया~।}\pend

\pstart\versepar \ledsidenote{\passref{5.8.140.add-4}}\skt{यथा\-स्था\-नो\-प\-दि\-ष्टानां }\edtext{\skt{विदितानाञ्च}}{\lemma{\skt{विदितानाञ्च}}\Afootnote{K; \skt{विदिता…}\ H}}\skt{ }\edtext{\skt{लक्षणैः}}{\lemma{\skt{लक्षणैः}}\Afootnote{K; \skt{…क्षणैः~॥}\ H}}\skt{~।}\skt{ }\pend
\pstart\versepar \skt{}\edtext{\skt{व्याधीनां}}{\lemma{\skt{व्याधीनां}}\Afootnote{H; \skt{व्याधीना}\ K}}\skt{ साधनं शास्त्रे चिकित्सितमिति }\edtext{\skt{स्मृतम्}}{\lemma{\skt{स्मृतम्}}\Afootnote{H; \skt{स्मृतां~॥}\ K}}\skt{~॥}\pend

\pstart\versepar \ledsidenote{\passref{5.8.140.add-5}}\skt{स्थावरे जङ्गमे चैव }\edtext{\skt{विषे}}{\lemma{\skt{विषे}}\Afootnote{H; \skt{विशे}\ K}}\skt{ हित\-वि\-क\-ल्प\-नम्~। }\pend
\pstart\versepar \skt{साधनं चैव }\edtext{\skt{कार्त्स्न्ये}}{\lemma{\skt{कार्त्स्न्ये}}\Afootnote{K; \skt{कार्त्स्न्येन}\ H}}\skt{ कल्पस्थानं तदुच्यते~॥}\pend

\pstart\versepar \ledsidenote{\passref{5.8.140.add-6}}\skt{साहस्रो विस्तरः पूर्वं प्रजा\-प\-ति\-मु\-खो\-द्भवः~। }\pend
\pstart\versepar \skt{}\edtext{\skt{सविं\-श\-द\-ध्या\-य\-शतं}}{\lemma{\skt{सविं\appbrk{}श\appbrk{}द\appbrk{}ध्या\appbrk{}य\appbrk{}शतं}}\Afootnote{K; \skt{सवि\appbrk{}ङ्श\appbrk{}म\appbrk{}ध्या\appbrk{}य\appbrk{}श\appbrk{}तम्}\ H}}\skt{ मया }\edtext{\skt{वत्स}}{\lemma{\skt{वत्स}}\Afootnote{K\textsuperscript{pc} H; \skt{व}\ K\textsuperscript{ac}}}\skt{ प्रकीर्त्तितम्~॥}\pend

\pstart\versepar \ledsidenote{\passref{5.8.141}}\skt{सना\-त\-न\-त्वा\-त्वे\-दा\-ना\-म\-क्ष\-र\-त्वा\-त्त\-थैव च~। }\pend
\pstart\versepar \skt{}\edtext{\skt{दृष्टा\-दृ\-ष्ट\-फ\-ल\-त्वाच्च}}{\lemma{\skt{दृष्टा\appbrk{}दृ\appbrk{}ष्ट\appbrk{}फ\appbrk{}ल\appbrk{}त्वाच्च}}\Afootnote{K H; \skt{तथा दृष्टफलत्वाच्च}\ A}}\skt{ }\edtext{\skt{हितत्वाच्चापि}}{\lemma{\skt{हितत्वाच्चापि}}\Afootnote{K H; \skt{हितत्वादपि}\ A}}\skt{ देहिनाम्~।}\pend

\pstart\versepar \ledsidenote{\passref{5.8.142}}\skt{वाक्स\-मू\-हा\-र्थ\-वि\-स्ता\-रा\-त्पू\-जि\-त\-त्वाच्च }\edtext{\skt{देहिषु}}{\lemma{\skt{देहिषु}}\Afootnote{K; \skt{देहिभिः~।}\ A; \skt{देहिनः~।}\ H}}\skt{~। }\pend
\pstart\versepar \skt{}\edtext{\skt{चिकि\-त्सि\-ता\-त्पु\-ण्य\-तमं न}}{\lemma{\skt{चिकि\appbrk{}त्सि\appbrk{}ता\appbrk{}त्पु\appbrk{}ण्य\appbrk{}तमं न}}\Afootnote{A K; \skt{चिकि\appbrk{}त्सि\appbrk{}ता\appbrk{}त्पू\appbrk{}ण्य\appbrk{}तमं न}\ H}}\skt{ }\edtext{\skt{किञ्चिदपि}}{\lemma{\skt{किञ्चिदपि}}\Afootnote{A H; \skt{किञ्चिदपि}\ K}}\skt{ }\edtext{\skt{सुश्रुत}}{\lemma{\skt{सुश्रुत}}\Afootnote{\skt{शुश्रुमः~॥}\ A H; \skt{सुश्रुतः~।}\ K}}\skt{~।}\pend

\pstart\versepar \ledsidenote{\passref{5.8.143}}\edtext{\skt{ऋषे\-रि\-न्द्र\-प्र\-भा\-वस्य}}{\lemma{\skt{इन्द्रप्रभावस्य}}\Afootnote{K H; \skt{इन्द्रप्रभाव}\ A}}\skt{ }\edtext{\skt{तस्मा\-द\-मृ\-त\-ज\-न्मनः}}{\lemma{\skt{तस्मा\appbrk{}द\appbrk{}मृ\appbrk{}त\appbrk{}ज\appbrk{}न्मनः}}\Afootnote{K H; \skt{स्यामृ\appbrk{}त\appbrk{}यो\appbrk{}ने\appbrk{}र्भि\appbrk{}ष\appbrk{}ग्गुरोः~।}\ A}}\skt{~। }\pend
\pstart\versepar \skt{}\edtext{\edtext{\skt{धार\-यि\-त्वे\-द\-म\-मलं}}{\lemma{\skt{धारयित्वेदम्}}\Afootnote{K H; \skt{धारयित्वा तु}\ A}}}{\lemma{\skt{अमलं}}\Afootnote{K H; \skt{विमलं}\ A}}\skt{ मत\-म्प\-र\-म\-स\-म्म\-तम्~। }\pend
\pstart\versepar \skt{}\edtext{\skt{उक्ता\-चा\-र\-स\-मा\-चारः}}{\lemma{\skt{उक्ता\appbrk{}चा\appbrk{}र\appbrk{}स\appbrk{}मा\appbrk{}चारः}}\Afootnote{K; \skt{उक्ता\appbrk{}हा\appbrk{}र\appbrk{}स\appbrk{}मा\appbrk{}चार}\ A; \skt{उक्ता\appbrk{}चा\appbrk{}र\appbrk{}स\appbrk{}मा\appbrk{}चाराः}\ H}}\skt{ }\edtext{\skt{प्रेत्य}}{\lemma{\skt{प्रेत्य}}\Afootnote{K; \skt{इह प्रेत्य}\ A; \skt{प्रत्य}\ H}}\skt{ }\edtext{\skt{चेह च}}{\lemma{\skt{चेह च}}\Afootnote{K; \skt{च}\ A; \skt{वेह च}\ H}}\skt{ }\edtext{\skt{नन्दति}}{\lemma{\skt{नन्दति}}\Afootnote{K H; \skt{मोदते~॥}\ A}}\skt{~।}\pend

\pstart \ledsidenote{\passref{5.8.143.add-1}}\skt{शेषाणामपि तन्त्राणां युक्तिज्ञो }\edtext{\skt{लोकबान्धवः}}{\lemma{\skt{लोकबान्धवः}}\Afootnote{\skt{लोकवान्धवः}\ K H}}\skt{~।}\pend

\pstart\versepar \ledsidenote{\passref{5.8.143.add-2}}\edtext{\skt{यत्किं\-चि\-दा\-बा\-ध\-क\-र\-न्त\-द्य\-स्मा\-च्छ\-ल्य\-सं\-ज्ञितं}}{\lemma{\skt{यत्किं\appbrk{}चि\appbrk{}दा\appbrk{}बा\appbrk{}ध\appbrk{}क\appbrk{}रन्}}\Afootnote{K NE-PUL\_818; \skt{यत्कि\appbrk{}ञ्चि\appbrk{}दा\appbrk{}वा\appbrk{}ध\appbrk{}क\appbrk{}रन्}\ H}}\skt{~। }\pend
\pstart\versepar \skt{}\edtext{\skt{व्याप्ता\-न्य\-ङ्गा\-न्य\-त\-सतो}}{\lemma{\skt{अतस्}}\Afootnote{K H; \skt{अतो}\ NE-PUL\_818}}\skt{ तेन }\edtext{\skt{शल्यज्ञानेन सूरिणा}}{\lemma{\skt{सूरिणा}}\Afootnote{K H; \skt{भूरिणा~।}\ NE-PUL\_818}}\skt{ भूरिणा~।}\pend

\pstart\versepar \ledsidenote{\passref{5.8.143.add-3}}\skt{अतश्चास्य विशेषेण गतिर्न प्रतिषिध्यते~। }\pend
\pstart\versepar \skt{यथा स्वविषयस्थस्य राज्ञो }\edtext{\skt{बलवतो}}{\lemma{\skt{बलवतो}}\Afootnote{\skt{वलवतो}\ K H}}\skt{ गतिः~।}\pend

\pstart\versepar \ledsidenote{\passref{5.8.143.add-4}}\edtext{\skt{उपद्रवाणां}}{\lemma{\skt{उपद्रवाणां}}\Afootnote{H; \skt{उपद्रवानां}\ K}}\skt{ निर्देशो निदानं }\edtext{\skt{व्यञ्जनानि च}}{\lemma{\skt{व्यञ्जनानि च}}\Afootnote{K; \skt{व्यजनानि च~।}\ H}}\skt{~। }\pend
\pstart\versepar \skt{ज्वरा\-दी\-ना\-ञ्चि\-कि\-त्सा\-र्थ\-मु\-त्त\-र\-न्त\-न्त्र\-मु\-च्यते~॥}\pend

\pstart\versepar \ledsidenote{\passref{5.8.143.add-5}}\edtext{\skt{भवति}}{\lemma{\skt{भवति}}\Afootnote{K NE-PUL\_818; \skt{भवन्ति}\ H}}\skt{ चात्र~॥ }\pend
\pstart\versepar \skt{इदन्तु यः पञ्चसु सन्निवेशितम्~। }\pend
\pstart\versepar \skt{सविं\-श\-द\-ध्या\-य\-शतं सहोत्तरम्~। }\pend
\pstart\versepar \skt{पठेत्स }\edtext{\skt{राज्ञोऽर्हति}}{\lemma{\skt{ऽर्हति}}\Afootnote{H NE-PUL\_818; \skt{र्हति}\ K}}\skt{ वैद्यपूजितः~। }\pend
\pstart\versepar \skt{क्रियाम्प्रयोक्तुं }\edtext{\skt{भिष\-गा\-ग\-त\-क्रमः}}{\lemma{\skt{भिष\appbrk{}गा\appbrk{}ग\appbrk{}त\appbrk{}क्रमः}}\Afootnote{K; \skt{भिषभा गतक्रम इति~॥~॥}\ H\textsuperscript{ac}; \skt{भिष गागतक्रम इति~॥~॥}\ H\textsuperscript{pc}; \skt{भिषगागमज्ञः~॥}\ NE-PUL\_818}}\skt{ आगमज्ञः~॥}\pend

\pstart\versepar \ledsidenote{\passref{5.8.143.add-6}}\skt{अन्नरक्षा स्थावरविषं जङ्गमञ्च विषं तथा~। }\pend
\pstart\versepar \skt{}\edtext{\skt{सर्प\-द\-ष्ट\-वि\-ष\-ज्ञानं}}{\lemma{\skt{सर्प\appbrk{}द\appbrk{}ष्ट\appbrk{}वि\appbrk{}ष\appbrk{}ज्ञानं}}\Afootnote{H; \skt{सर्प\appbrk{}द\appbrk{}ष्ट\appbrk{}वि\appbrk{}षा\appbrk{}ज्ञनं}\ K}}\skt{ सर्प\-द\-ष्ट\-चि\-कि\-त्सि\-तम्~॥ }\pend
\pstart\versepar \skt{मूषिका दुन्दुभिश्चैव कीटाकल्पेन चाष्टमः~॥}\pend

\pstart \ledsidenote{\passref{5.8.trailer1}}\edtext{\skt{सौश्रुते शल्यतन्त्रे}}{\lemma{\skt{सौश्रुते शल्यतन्त्रे}}\Afootnote{\skt{इति सुश्रु\appbrk{}त\appbrk{}सं\appbrk{}हि\appbrk{}तायां}\ A; \emph{om.}\ H}}\skt{ }\edtext{\skt{कल्पस्थानं समाप्तम्~॥ {\symfont ❈}}}{\lemma{\skt{कल्पस्थानं समाप्तम्~॥ {\symfont ❈}}}\Afootnote{\skt{कल्पस्थाने कीटकल्पो नामाष्टमोऽध्यायः~॥८~॥}\ A; \skt{कल्पेषु कीट\appbrk{}चि\appbrk{}कि\appbrk{}त्सि\appbrk{}त\appbrk{}म\appbrk{}ष्टमं~॥ ओ~॥}\ H}}\skt{~॥}\pend

\pstart \ledsidenote{\passref{5.8.trailer2}}\edtext{\skt{सौश्रुते}}{\lemma{\skt{सौश्रुते}}\Afootnote{K H; \skt{इति भगवता श्रीध\appbrk{}न्व\appbrk{}न्त\appbrk{}रि\appbrk{}णो\appbrk{}प\appbrk{}दि\appbrk{}ष्टायां तच्छिष्येण महर्षिणा सुश्रुतेन विरचितायां सुश्रुत}\ A}}\skt{ }\edtext{\skt{शल्यतन्त्रे}}{\lemma{\skt{शल्यतन्त्रे}}\Afootnote{K H; \skt{संहितायां पञ्चमं}\ A}}\skt{ कल्पस्थानं }\edtext{\skt{समाप्तं}}{\lemma{\skt{समाप्तं}}\Afootnote{A H; \skt{समाप्तः~॥ {\symfont ❈}~॥}\ K}}\skt{~॥}\pend

\endnumbering
